% 农业上市公司画像：生猪养殖、肉鸡养殖、其他养殖（共 22 家，按流通市值降序）
% 由 scripts/make_company_profiles.py 自动生成，请勿手工修改

\subsubsection{牧原股份（002714）}
\label{co:002714}

\pfl{公司基本情况} 公司前身可追溯至 2014 年，自繁自养一体化模式的全球最大生猪养殖企业，2026 年 H 股上市后形成 A+H 双平台；2026 年上半年猪价深跌录得大额亏损，公司主动调减能繁母猪产能并压降养殖成本。 公司于 2014-01-28 在深交所主板首发上市，注册资本 57.73 亿元，总股本 54.65 亿股。 截至 2026-09-24收盘价 40.48 元，流通市值 1,327.9 亿元，近一年-25.0\%，流通市值在三级行业“生猪养殖”的 11 家公司中排第 \zhnumber{1} 位，近一年涨跌幅低于全样本中位数（-10.1\%）14.9 个百分点。 按 2026 年半年报披露的产品构成，收入占比前三位为生猪 88.3\%；屠宰、肉食产品 37.1\%；饲料原料 7.3\%。 发展过程中的关键节点包括：上市公告书载明上市时间为 2014 年 1 月 28 日，股票简称“牧原股份”，法定代表……（2014）；公司定向增发 49.77 亿元 A 股，用以新增生猪产能；同年设立牧原肉食、正阳肉食、商……（2019）；2 月 6 日 H 股在香港联交所挂牌上市，最终发行价 39.0 港元/股，基础发行 2……（2026）。\footnote{\texttt{http://vip.stock.finance.sina.com.cn/corp/view/vISSUE\_MarketBulletinDetail.php}}\footnote{\texttt{https://pdf.dfcfw.com/pdf/H3\_AP202204141559239147\_1.pdf}}\footnote{\texttt{https://www.thepaper.cn/newsDetail\_forward\_32553701}}

\pfl{历史股价} 以 2021 年 1 月为起点、2026 年 9 月为终点，股价由 88.90 元变为 40.48 元，累计 -54.5\%，区间最高 113.90 元（2021 年 2 月）、最低 34.99 元（2026 年 6 月）；当前价相当于区间最高价的 35.5\%，为区间最低价的 1.16 倍。 月度收益的年化波动率 38.7\%，高于全样本中位数 37.4\%，在三级行业“生猪养殖”的 11 家公司中波动率排第 \zhnumber{5} 位。 从事件看，2024 年 1 月 至 2025 年 8 月的上涨（+55.9\%）与公司在越南胡志明市设立全资子公司“越南牧原有限公司”，推进国际化布局（2025-03）在时间上重合——股价的拐点多数能在公司自身的资本运作与风险事件上找到对应。\footnote{\texttt{https://www.stcn.com/article/detail/3235472.html}}

\begin{center}
\begin{tikzpicture}
\begin{axis}[
  width=12.6cm, height=3.9cm, scale only axis,
  xmin=2020.922, xmax=2026.888, ymin=32.54, ymax=121.87,
  xtick={2021,2022,2023,2024,2025,2026},
  yearticks,
  yticklabel style={font=\scriptsize},
  ylabel={元/股}, ylabel style={font=\scriptsize},
  grid=major, grid style={LightGray, line width=0.3pt},
  axis line style={InkGray, line width=0.5pt},
  every axis plot/.append style={line width=0.9pt},
]
\addplot[AgriGreen] table[x=x,y=y]{data/clean/profiles/px_002714.dat};
\addplot[SoftRed, dashed, line width=0.7pt] table[x=x,y=y]{data/clean/profiles/pxzz_002714.dat};
\addplot[only marks, mark=*, mark size=1.1pt, SoftRed] table[x=x,y=y]{data/clean/profiles/pxzz_002714.dat};
\end{axis}
\end{tikzpicture}

\captionof{figure}[牧原股份（002714）股价与周期骨架]{牧原股份（002714）月度收盘价（元/股）与 ZigZag 周期骨架（阈值 25\%，圆点为周期转折点）}
\end{center}

\pfl{过去周期分析} 以 25\% 的 ZigZag 阈值划分，样本区间内股价形成 2 个主升段与 3 个主跌段。 最大主升段出现在 2024 年 1 月 至 2025 年 8 月，19 个月+55.9\%；最大主跌段为 2021 年 1 月 至 2021 年 7 月，6 个月-52.5\%。 最近一次转折出现在 2026 年 6 月（下跌段自 2025 年 8 月起，10 个月-36.3\%），当前处于该段的延续之中。 公司股价与申万农林牧渔指数几乎同步见底，个股并未走出独立行情。 生猪养殖的定价主线是产能去化与成本曲线的赛跑（第 \ref{sec:hog} 章）：2026 年三季度行业仍处在“产能去化未完成、价格磨底”的阶段，公司 2026 年 6 月末资产负债率 54.2\%、半年 ROE -7.5\%（未年化），在本轮下行中属于随行业同步调整的一类。

\pfl{盈利情况} 2026 年上半年营业收入 594.10 亿元（同比 -22.3\%），归母净利润 -60.78 亿元，由上年同期的盈利转为亏损。 以年报为趋势对照，2023---2025 年营业收入由 1,108.61 亿元变为 1,441.45 亿元（两年年均复合增速 +14.0\%），归母净利润由 -42.63 亿元变为 154.87 亿元。 按半年报口径，2026 年上半年的 ROE 为 -7.5\%（未年化）、销售净利率 -10.2\%、毛利率 -1.6\%，ROE 在“生猪养殖”11 家公司中排第 \zhnumber{3} 位。 同期全样本半年 ROE 中位数 0.75\%，公司与之相差 8.23 个百分点（弱于中位数公司）。 结合公司披露，2026 年上半年净利润 -60.78 亿元，主要系生猪销售价格同比大幅下降所致；响应国家生猪产能调控政策主动压减产能，能繁母猪存栏较上年阶段性高点下降 60 万头左右（计划 9 月底降……。\footnote{\texttt{https://www.stcn.com/article/detail/4105451.html}}

\begin{center}\small\setlength{\tabcolsep}{4pt}
\captionof{table}[牧原股份（002714）关键财务指标]{牧原股份（002714）关键财务指标（合并报表口径；两个半年报期均未年化；现金短债比 $=$ 货币资金 $\div$ 短期债务）}
\begin{tabular}{@{}lrrrrr@{}}
\toprule
指标 & 2023 年 & 2024 年 & 2025 年 & 2025 年半年报 & 2026 年半年报 \\
\midrule
营业收入（亿元） & 1,108.61 & 1,379.47 & 1,441.45 & 764.63 & 594.10 \\
归母净利润（亿元） & -42.63 & 178.81 & 154.87 & 105.30 & -60.78 \\
毛利率（\%） & 3.1 & 19.1 & 17.8 & 19.0 & -1.6 \\
ROE（\%） & -6.4 & 25.1 & 20.6 & 13.8 & -7.5 \\
资产负债率（\%） & 62.1 & 58.7 & 54.2 & 56.1 & 54.2 \\
经营活动现金流净额（亿元） & 98.93 & 375.43 & 300.56 & 173.51 & -22.24 \\
货币资金（亿元） & 194.29 & 169.52 & 138.62 & 202.85 & 165.47 \\
有息负债合计（亿元） & 787.62 & 760.80 & 702.49 & 773.01 & 748.02 \\
现金短债比（倍） & 0.35 & 0.32 & 0.28 & 0.37 & 0.29 \\
\bottomrule
\end{tabular}
\end{center}

\begin{center}
\begin{tikzpicture}
\begin{groupplot}[
  group style={group size=2 by 1, horizontal sep=0.95cm},
  width=7.2cm, height=4.6cm,
  tick label style={font=\tiny},
  title style={font=\scriptsize\heiti},
  yticklabel style={font=\tiny},
  grid=major, grid style={LightGray, line width=0.3pt},
  axis line style={InkGray, line width=0.5pt},
  enlarge x limits={abs=0.8},
  legend style={font=\scriptsize, at={(0.5,-0.22)}, anchor=north,
                legend columns=2, draw=none, fill=none, column sep=6pt},
]
\nextgroupplot[
  ybar, bar width=4pt,
  ymin=-211.00, ymax=1621.72,
  xtick={0,1,2,3,4,5.7},
  xticklabels={2021,2022,2023,2024,2025,2026H1},
  ylabel={亿元},
]
\addplot[fill=AgriGreen!75, draw=AgriGreen, bar shift=-2.6pt] table[x=i, y=rev]{data/clean/profiles/fin_002714.dat};
\addplot[fill=SoftRed!60, draw=SoftRed, bar shift=2.6pt] table[x=i, y=ni]{data/clean/profiles/fin_002714.dat};
\addplot[fill=AgriGreen!30, draw=AgriGreen, pattern=north east lines, bar shift=-2.6pt] table[x=x, y=rev]{data/clean/profiles/finh_002714.dat};
\addplot[fill=SoftRed!20, draw=SoftRed, pattern=north east lines, bar shift=2.6pt] table[x=x, y=ni]{data/clean/profiles/finh_002714.dat};
\legend{营业收入（年/半年）, 归母净利润（年/半年）}
\nextgroupplot[
  ymin=-13.0, ymax=69.1,
  xtick={0,1,2,3,4,5.7},
  xticklabels={2021,2022,2023,2024,2025,2026H1},
  ylabel={\%},
]
\addplot[AgriGold, mark=*, mark size=1.3pt] table[x=i, y=roe]{data/clean/profiles/fin_002714.dat};
\addplot[OilBlue, mark=square*, mark size=1.3pt] table[x=i, y=debt]{data/clean/profiles/fin_002714.dat};
\addplot[only marks, AgriGold, mark=*, mark size=2.0pt] table[x=x, y=roe]{data/clean/profiles/finh_002714.dat};
\addplot[only marks, OilBlue, mark=square*, mark size=2.0pt] table[x=x, y=debt]{data/clean/profiles/finh_002714.dat};
\legend{ROE, 资产负债率}
\end{groupplot}
\end{tikzpicture}
\begin{tikzpicture}
\begin{axis}[
  width=15.2cm, height=3.1cm, scale only axis,
  ymin=-103.22, ymax=884.80,
  xtick={0,1,2,3,4,5.7},
  xticklabels={2021,2022,2023,2024,2025,2026H1},
  tick label style={font=\tiny},
  yticklabel style={font=\tiny},
  ylabel={亿元}, ylabel style={font=\scriptsize},
  ybar, bar width=4pt,
  grid=major, grid style={LightGray, line width=0.3pt},
  axis line style={InkGray, line width=0.5pt},
  enlarge x limits={abs=0.8},
  legend style={font=\scriptsize, at={(0.5,-0.30)}, anchor=north,
                legend columns=3, draw=none, fill=none, column sep=10pt},
]
\addplot[fill=AgriGreen!55, draw=AgriGreen, bar shift=-3.0pt] table[x=i, y=cash]{data/clean/profiles/fin_002714.dat};
\addplot[fill=SoftRed!45, draw=SoftRed, bar shift=3.0pt] table[x=i, y=ibd]{data/clean/profiles/fin_002714.dat};
\addplot[fill=AgriGreen!25, draw=AgriGreen, pattern=north east lines, bar shift=-3.0pt] table[x=x, y=cash]{data/clean/profiles/finh_002714.dat};
\addplot[fill=SoftRed!20, draw=SoftRed, pattern=north east lines, bar shift=3.0pt] table[x=x, y=ibd]{data/clean/profiles/finh_002714.dat};
\addplot[OilBlue, mark=*, mark size=1.1pt, line width=0.8pt] table[x=i, y=ocf]{data/clean/profiles/fin_002714.dat};
\addplot[only marks, OilBlue, mark=*, mark size=1.8pt] table[x=x, y=ocf]{data/clean/profiles/finh_002714.dat};
\legend{货币资金, 有息负债, 经营活动现金流净额}
\end{axis}
\end{tikzpicture}

\captionof{figure}[牧原股份（002714）近年主要财务指标]{牧原股份（002714）近年主要财务指标（左上：营业收入与归母净利润；右上：ROE 与资产负债率；下：货币资金、有息负债与经营活动现金流净额。
2021---2025 年为年报口径，2026H1 为 2026 年半年报/半年末，与其前一年报之间留出间隔、以斜纹或空心点区分；半年报数据未年化）}
\end{center}

\pfl{财务分析} 2026 年 6 月末资产负债率 54.2\%，较 2025 年末 上升 0.0 个百分点（样本中位数 52.2\%，所处三级行业内排第 \zhnumber{5} 位）。 2026 年 6 月末货币资金 165.47 亿元、短期债务（短期借款与一年内到期非流动负债）568.32 亿元、长期债务 179.70 亿元，有息负债合计 748.02 亿元，现金短债比 0.29 倍——缺口明显。 净有息负债 582.56 亿元，相当于其 2026 年上半年营业收入的 98\%。 流动比率 0.79、速动比率 0.28，短期偿付压力较大。 2026 年上半年经营活动现金流净额 -22.24 亿元，经营现金流与利润同时为负，现金消耗较快。 2026 年 6 月末在建工程 23.80 亿元，较上年末减少 1.09 亿元，在建投入在收缩。 公司披露的财务与合规风险包括：商票逾期事件发酵后，中诚信国际将公司评级展望下调至负面，认为公司商票管理等内控制度有待提……；上海票据交易所公示：2021 年 8 月 1 日至 11 月 30 日公司 32 家子公……。\footnote{\texttt{http://finance.ce.cn/stock/gsgdbd/202112/16/t20211216\_37176818.shtml}}\footnote{\texttt{https://m.bjnews.com.cn/detail/1654056736169495.html}}

\pfl{重大战略转型} 近三年公司的战略动作集中在以下几个方面：2024-09-27 披露股份回购方案，回购总金额不低于 30 亿元且不超过 40 亿元；截至 2025 年 8 月 26 日累计回购……；2025-03 在越南胡志明市设立全资子公司“越南牧原有限公司”，推进国际化布局；2026 响应国家生猪产能调控政策主动压减产能，能繁母猪存栏较上年阶段性高点下降 60 万头左右（计划 9 月底降至 30……；2026-02 H 股上市募资用于拓展全球商业版图及研发投入；公司明确表示不会用于国内新增产能建设。 公告口径上，近三年（2023 年 9 月---2026 年 9 月）公司披露重大重组与并购类公告 0 条、对外投资与转型类公告 1 条、回购与股权激励类公告 34 条，合计公告 615 条。 第一大产品仍是“生猪”，收入占比 88.3\%，与 2021 年年报相比没有方向性变化。 综合判断，公司在报告期内有明确的外延扩张或业务结构调整动作，转型方向与融资安排之间存在直接关联。\footnote{\texttt{https://pdf.dfcfw.com/pdf/H3\_AP202204141559239147\_1.pdf}}\footnote{\texttt{https://www.cnfin.com/announ/detail/index.html}}\footnote{\texttt{https://www.stcn.com/article/detail/3235472.html}}

\begin{center}\small\setlength{\tabcolsep}{4pt}
\captionof{table}[牧原股份（002714）历史融资与资本运作]{牧原股份（002714）历史融资与资本运作（据公司公告与公开报道整理；金额为披露值，含首次公开发行、再融资、债券、银行借款与重整投资等）}
\begin{tabular}{@{}p{1.5cm}p{3.2cm}p{2.4cm}p{6.3cm}@{}}
\toprule
时间 & 方式 & 规模 & 用途与说明 \\
\midrule
2014-01 & IPO（深交所中小板） & 发行 6,050 万股，发行价 24.07 元/股，募集资金总额 7.221 亿元、净额 6.678 亿元 & 生猪养殖产能及补充流动资金；其中新股 3,000 万股、老股转让 3,050 万股…… \\
2015-12 & 定向增发 & 3,287.31 万股，发行价 30.42 元/股，实际募集资金净额 9.9269 亿元 & 生猪养殖项目建设；2015 年 12 月 22 日发行…… \\
2017-04 & 定向增发 & 124,717,239 股，发行价 24.67 元/股，募集资金总额 30.7677 亿元 & 生猪养殖产能扩张；2017 年 4 月 14 日验资 \\
2018 & 优先股 & 24.6 亿元 & 募集资金用于生猪产能扩张等项目；票面利率 6.8\%，向商业银行和信托公司募集…… \\
2019 & 公司债券（扶贫专项） & 8 亿元 & 债券简称“19 牧原 01”，代码 112849，票面利率 7.00\%…… \\
2019-09 & 定向增发 & 76,663,600 股，发行价 65.22 元/股，募集资金总额 50.00 亿元、净额 49.77 亿元 & 新增约 475 万头生猪产能（2018 年 12 月启动预案）…… \\
2021-08 & 可转换公司债券（牧原转债，127045） & 95.5 亿元（9,550 万张，面值 100 元） & 51 亿元投生猪养殖项目、19 亿元投生猪屠宰项目…… \\
2022-12 & 定向增发（2021 年度非公开发行） & 150,112,584 股，发行价 39.97 元/股，募集资金总额 60.00 亿元、净额 59.87 亿元 & 全部用于补充流动资金；2021 年 10 月 18 日董事会通过预案…… \\
2024-03 & 中期票据（银行间市场） & 注册总额度 30 亿元，2024 年度第一期发行金额上限 5 亿元、期限 3 年 & 补充流动资金及债务结构调整；在中国银行间市场交易商协会注册 \\
2026-02 & H 股 IPO（香港联交所主板） & 基础发行 273,951,400 股，发行价 39.0 港元/股；募资总额约 106.84 亿港元、净额约 104.70 亿港元（未行使超额配售权） & 拓展全球商业版图及研发投入，不用于国内新增产能建设…… \\
\bottomrule
\end{tabular}
\end{center}

\pfl{历史融投资情况} 从现金流量表看，2025 年公司取得借款收到的现金 619.05 亿元、偿还债务支付的现金 679.52 亿元、吸收权益性投资 10.53 亿元，筹资活动现金流净额 -228.78 亿元（2023 年为取得借款 716.49 亿元、偿还债务 533.66 亿元）。 2026 年上半年筹资活动现金流净额为 111.60 亿元（取得借款 471.88 亿元、偿还债务 419.89 亿元，未年化），仍处于净融入状态。 2025 年分配股利、利润或偿付利息支付的现金合计 103.45 亿元。 公司披露的股本与债务融资脉络包括：H 股 IPO（香港联交所主板）、IPO（深交所中小板）、中期票据（银行间市场）、优先股、公司债券（扶贫专项）、可转换公司债券（牧原转债，127045）。 其中规模较大的两笔为：2026-02H 股 IPO（香港联交所主板）（基础发行 273,951,400 股，发行价 39.0 港元……）；2022-12定向增发（2021 年度非公开发行）（150,112,584 股，发行价 39.97 元/股，募集……）。 股本方面，近三年披露股本变动 34 次（可转债转股 28 次、激励股份解禁 3 次、股份回购 2 次、限售股份上市 1 次）；总股本由 2021 年末的 52.62 亿股变为 54.65 亿股（+3.86\%）。 分红方面，近三年现金分红 5 次、累计每股派现 3.50 元（最近一次 2025-12-31）——在利润承压的阶段仍维持了分红。 近三年“再融资”类公告 23 条，涉及定向增发、可转债。\footnote{\texttt{http://vip.stock.finance.sina.com.cn/corp/go.php/vISSUE\_NewStock/stockid/002714.phtml}}\footnote{\texttt{https://pdf.dfcfw.com/pdf/H2\_AN202005281380416085\_1.pdf}}\footnote{\texttt{http://vip.stock.finance.sina.com.cn/corp/view/vISSUE\_MarketBulletinDetail.php}}

\pfl{未来潜在融资需求分析} 2026 年 6 月末货币资金 165.47 亿元，而短期债务 568.32 亿元（现金短债比 0.29 倍），2026 年上半年经营现金流净额 -22.24 亿元。按“短期债务减去账面货币资金”的滚动偿付口径，静态缺口约 402.86 亿元；若再计入上半年亏损的年化额（121.55 亿元），维持一年正常经营所需的外部资金约 402.86---524.41 亿元。 公司 2026 年上半年归母净利润为负（-60.78 亿元），但 6 月末资产负债率 54.2\% 尚未进入第一档（亏损且负债率 $>$65\%）的区间，当前融资需求以补充流动资金、支撑低谷期经营性支出为主。 公开披露的下一步融资线索包括：2026 2025 年报披露：2025 年年底负债总额较年初下降 171 亿元，资产负债率较年初下降 4.53 个百分点，财务结构……；2026-02 H 股发行已落地（净募资约 104.70 亿港元），为海外扩张与研发提供资金；可转债“牧原转债”余额仍在存续期内（202……。 资金需求还要叠加在建工程：期末在建工程 23.80 亿元、2026 年上半年购建固定资产等支出 65.21 亿元（未年化），这部分支出在周期底部通常会被主动放缓。\footnote{\texttt{https://money.finance.sina.com.cn/corp/view/vCB\_AllBulletinDetail.php}}\footnote{\texttt{https://www.thepaper.cn/newsDetail\_forward\_32553701}}


\subsubsection{温氏股份（300498）}
\label{co:300498}

\pfl{公司基本情况} 以“公司+农户”模式经营的畜禽养殖龙头，肉鸡规模居全国首位、肉猪规模居行业前列，2015 年换股吸收合并大华农实现整体上市；2026 年上半年猪价深跌导致大额亏损，肉禽板块盈利部分对冲。 公司于 2015-11-02 在深交所创业板首发上市，注册资本 66.54 亿元，总股本 66.53 亿股。 截至 2026-09-24收盘价 14.51 元，流通市值 862.7 亿元，近一年-23.6\%，流通市值在三级行业“生猪养殖”的 11 家公司中排第 \zhnumber{2} 位，近一年涨跌幅低于全样本中位数（-10.1\%）13.5 个百分点。 按 2026 年半年报披露的行业构成，收入占比前三位为养殖行业 131.1\%；工程及设备制造业 2.5\%；动保行业 2.1\%。 发展过程中的关键节点包括：上市公告书披露换股吸收合并大华农事项经证监会核准后实施，股票在深交所创业板上市（2015）；3 月 29 日发行 92.97 亿元可转换公司债券（温氏转债，123107），4 月 ……（2021）。\footnote{\texttt{http://money.finance.sina.com.cn/corp/view/vISSUE\_MarketBulletinDetail.php}}\footnote{\texttt{https://www.wens.com.cn/uploadfiles/2023/09/202309111130185504.pdf}}

\pfl{历史股价} 本报告样本区间（2021 年 1 月 至 2026 年 9 月）内，股价由 16.96 元变为 14.51 元，累计 -14.4\%，区间最高 24.04 元（2022 年 7 月）、最低 12.32 元（2026 年 6 月）；当前价相当于区间最高价的 60.4\%，为区间最低价的 1.18 倍。 月度收益的年化波动率 29.8\%，低于全样本中位数 37.4\%，在三级行业“生猪养殖”的 11 家公司中波动率排第 \zhnumber{10} 位。 从事件看，2024 年 5 月 至 2026 年 6 月的下跌（-42.5\%）与公司2025 年 4 月 23 日董事会、2025 年 5 月 14 日 202……（2025-05）在时间上重合——股价的拐点多数能在公司自身的资本运作与风险事件上找到对应。\footnote{\texttt{http://reportdocs.static.szse.cn/UpFiles/GsptFiles/mjsms/5043B8533F794E9E93457F533E05614}}

\begin{center}
\begin{tikzpicture}
\begin{axis}[
  width=12.6cm, height=3.9cm, scale only axis,
  xmin=2020.922, xmax=2026.888, ymin=11.46, ymax=25.72,
  xtick={2021,2022,2023,2024,2025,2026},
  yearticks,
  yticklabel style={font=\scriptsize},
  ylabel={元/股}, ylabel style={font=\scriptsize},
  grid=major, grid style={LightGray, line width=0.3pt},
  axis line style={InkGray, line width=0.5pt},
  every axis plot/.append style={line width=0.9pt},
]
\addplot[AgriGreen] table[x=x,y=y]{data/clean/profiles/px_300498.dat};
\addplot[SoftRed, dashed, line width=0.7pt] table[x=x,y=y]{data/clean/profiles/pxzz_300498.dat};
\addplot[only marks, mark=*, mark size=1.1pt, SoftRed] table[x=x,y=y]{data/clean/profiles/pxzz_300498.dat};
\end{axis}
\end{tikzpicture}

\captionof{figure}[温氏股份（300498）股价与周期骨架]{温氏股份（300498）月度收盘价（元/股）与 ZigZag 周期骨架（阈值 25\%，圆点为周期转折点）}
\end{center}

\pfl{过去周期分析} 以 25\% 的 ZigZag 阈值划分，样本区间内股价形成 2 个主升段与 3 个主跌段。 最大主升段出现在 2021 年 7 月 至 2022 年 7 月，12 个月+91.1\%；最大主跌段为 2024 年 5 月 至 2026 年 6 月，25 个月-42.5\%。 最近一次转折出现在 2026 年 6 月（下跌段自 2024 年 5 月起，25 个月-42.5\%），当前处于该段的延续之中。 公司股价与申万农林牧渔指数几乎同步见底，个股并未走出独立行情。 生猪养殖的定价主线是产能去化与成本曲线的赛跑（第 \ref{sec:hog} 章）：2026 年三季度行业仍处在“产能去化未完成、价格磨底”的阶段，公司 2026 年 6 月末资产负债率 58.9\%、半年 ROE -11.1\%（未年化），在本轮下行中属于随行业同步调整的一类。

\pfl{盈利情况} 2026 年上半年营业收入 467.65 亿元（同比 -6.2\%），归母净利润 -43.66 亿元，由上年同期的盈利转为亏损。 以年报为趋势对照，2023---2025 年营业收入由 899.21 亿元变为 1,038.62 亿元（两年年均复合增速 +7.5\%），归母净利润由 -63.90 亿元变为 52.66 亿元。 按半年报口径，2026 年上半年的 ROE 为 -11.1\%（未年化）、销售净利率 -9.4\%、毛利率 -2.9\%，ROE 在“生猪养殖”11 家公司中排第 \zhnumber{4} 位。 行业同口径半年 ROE 的中位数为 0.75\%，该公司低于中位数 11.88 个百分点。 结合公司披露，禽类业务量价齐增：上半年肉鸡销量同比提升，鲜品与熟食占比提高，形成对养猪亏损的对冲；董事会通过回购公司股份方案，回购资金总额不低于 9 亿元且不超过 18 亿元，回购价格不超过 27.01……。\footnote{\texttt{https://finance.sina.com.cn/jjxw/2026-08-27/doc-iniptnuz5843346.shtml}}\footnote{\texttt{https://file.finance.qq.com/finance/hs/pdf/2025/08/01/1224366923.PDF}}

\begin{center}\small\setlength{\tabcolsep}{4pt}
\captionof{table}[温氏股份（300498）关键财务指标]{温氏股份（300498）关键财务指标（合并报表口径；两个半年报期均未年化；现金短债比 $=$ 货币资金 $\div$ 短期债务）}
\begin{tabular}{@{}lrrrrr@{}}
\toprule
指标 & 2023 年 & 2024 年 & 2025 年 & 2025 年半年报 & 2026 年半年报 \\
\midrule
营业收入（亿元） & 899.21 & 1,056.45 & 1,038.62 & 498.75 & 467.65 \\
归母净利润（亿元） & -63.90 & 92.80 & 52.66 & 34.75 & -43.66 \\
毛利率（\%） & 0.8 & 17.3 & 11.2 & 13.8 & -2.9 \\
ROE（\%） & -17.6 & 24.7 & 12.6 & 8.3 & -11.1 \\
资产负债率（\%） & 61.4 & 53.1 & 49.8 & 50.6 & 58.9 \\
经营活动现金流净额（亿元） & 75.94 & 196.07 & 118.07 & 64.65 & 4.22 \\
货币资金（亿元） & 42.66 & 47.66 & 10.13 & 24.97 & 36.75 \\
有息负债合计（亿元） & 326.61 & 259.72 & 193.85 & 215.37 & 286.24 \\
现金短债比（倍） & 0.47 & 0.67 & 0.25 & 0.41 & 0.25 \\
\bottomrule
\end{tabular}
\end{center}

\begin{center}
\begin{tikzpicture}
\begin{groupplot}[
  group style={group size=2 by 1, horizontal sep=0.95cm},
  width=7.2cm, height=4.6cm,
  tick label style={font=\tiny},
  title style={font=\scriptsize\heiti},
  yticklabel style={font=\tiny},
  grid=major, grid style={LightGray, line width=0.3pt},
  axis line style={InkGray, line width=0.5pt},
  enlarge x limits={abs=0.8},
  legend style={font=\scriptsize, at={(0.5,-0.22)}, anchor=north,
                legend columns=2, draw=none, fill=none, column sep=6pt},
]
\nextgroupplot[
  ybar, bar width=4pt,
  ymin=-253.09, ymax=1199.31,
  xtick={0,1,2,3,4,5.7},
  xticklabels={2021,2022,2023,2024,2025,2026H1},
  ylabel={亿元},
]
\addplot[fill=AgriGreen!75, draw=AgriGreen, bar shift=-2.6pt] table[x=i, y=rev]{data/clean/profiles/fin_300498.dat};
\addplot[fill=SoftRed!60, draw=SoftRed, bar shift=2.6pt] table[x=i, y=ni]{data/clean/profiles/fin_300498.dat};
\addplot[fill=AgriGreen!30, draw=AgriGreen, pattern=north east lines, bar shift=-2.6pt] table[x=x, y=rev]{data/clean/profiles/finh_300498.dat};
\addplot[fill=SoftRed!20, draw=SoftRed, pattern=north east lines, bar shift=2.6pt] table[x=x, y=ni]{data/clean/profiles/finh_300498.dat};
\legend{营业收入（年/半年）, 归母净利润（年/半年）}
\nextgroupplot[
  ymin=-42.0, ymax=73.9,
  xtick={0,1,2,3,4,5.7},
  xticklabels={2021,2022,2023,2024,2025,2026H1},
  ylabel={\%},
]
\addplot[AgriGold, mark=*, mark size=1.3pt] table[x=i, y=roe]{data/clean/profiles/fin_300498.dat};
\addplot[OilBlue, mark=square*, mark size=1.3pt] table[x=i, y=debt]{data/clean/profiles/fin_300498.dat};
\addplot[only marks, AgriGold, mark=*, mark size=2.0pt] table[x=x, y=roe]{data/clean/profiles/finh_300498.dat};
\addplot[only marks, OilBlue, mark=square*, mark size=2.0pt] table[x=x, y=debt]{data/clean/profiles/finh_300498.dat};
\legend{ROE, 资产负债率}
\end{groupplot}
\end{tikzpicture}
\begin{tikzpicture}
\begin{axis}[
  width=15.2cm, height=3.1cm, scale only axis,
  ymin=-44.96, ymax=503.57,
  xtick={0,1,2,3,4,5.7},
  xticklabels={2021,2022,2023,2024,2025,2026H1},
  tick label style={font=\tiny},
  yticklabel style={font=\tiny},
  ylabel={亿元}, ylabel style={font=\scriptsize},
  ybar, bar width=4pt,
  grid=major, grid style={LightGray, line width=0.3pt},
  axis line style={InkGray, line width=0.5pt},
  enlarge x limits={abs=0.8},
  legend style={font=\scriptsize, at={(0.5,-0.30)}, anchor=north,
                legend columns=3, draw=none, fill=none, column sep=10pt},
]
\addplot[fill=AgriGreen!55, draw=AgriGreen, bar shift=-3.0pt] table[x=i, y=cash]{data/clean/profiles/fin_300498.dat};
\addplot[fill=SoftRed!45, draw=SoftRed, bar shift=3.0pt] table[x=i, y=ibd]{data/clean/profiles/fin_300498.dat};
\addplot[fill=AgriGreen!25, draw=AgriGreen, pattern=north east lines, bar shift=-3.0pt] table[x=x, y=cash]{data/clean/profiles/finh_300498.dat};
\addplot[fill=SoftRed!20, draw=SoftRed, pattern=north east lines, bar shift=3.0pt] table[x=x, y=ibd]{data/clean/profiles/finh_300498.dat};
\addplot[OilBlue, mark=*, mark size=1.1pt, line width=0.8pt] table[x=i, y=ocf]{data/clean/profiles/fin_300498.dat};
\addplot[only marks, OilBlue, mark=*, mark size=1.8pt] table[x=x, y=ocf]{data/clean/profiles/finh_300498.dat};
\legend{货币资金, 有息负债, 经营活动现金流净额}
\end{axis}
\end{tikzpicture}

\captionof{figure}[温氏股份（300498）近年主要财务指标]{温氏股份（300498）近年主要财务指标（左上：营业收入与归母净利润；右上：ROE 与资产负债率；下：货币资金、有息负债与经营活动现金流净额。
2021---2025 年为年报口径，2026H1 为 2026 年半年报/半年末，与其前一年报之间留出间隔、以斜纹或空心点区分；半年报数据未年化）}
\end{center}

\pfl{财务分析} 2026 年 6 月末资产负债率 58.9\%，较 2025 年末 上升 9.1 个百分点（样本中位数 52.2\%，所处三级行业内排第 \zhnumber{6} 位）。 2026 年 6 月末货币资金 36.75 亿元、短期债务（短期借款与一年内到期非流动负债）146.20 亿元、长期债务 140.04 亿元，有息负债合计 286.24 亿元，现金短债比 0.25 倍——缺口明显。 净有息负债 249.48 亿元，相当于其 2026 年上半年营业收入的 53\%。 流动比率 0.83、速动比率 0.29，短期偿付压力较大。 2026 年上半年经营活动现金流净额 4.22 亿元，净利润为负而经营现金流为正，差额主要来自折旧摊销与营运资本变动，说明亏损尚未完全转化为现金流出。 2026 年 6 月末在建工程 28.48 亿元，较上年末增加 8.95 亿元，仍有在建投入。 2026 年 6 月末商誉 1.68 亿元，占归母净资产的 0.4\%，是净资产中最脆弱的一块。

\pfl{重大战略转型} 近三年公司的战略动作集中在以下几个方面：2024-09-18 董事会通过回购公司股份方案，回购资金总额不低于 9 亿元且不超过 18 亿元，回购价格不超过 27.01 元/股……；2025-04 2024 年度利润分配：向全体股东每 10 股派发现金红利 2 元（含税），不送股、不转增；2026-04 2025 年度归母净利润 52.66 亿元；2025 年度利润分配预案为每 10 股派发现金红利 2 元（含税）；2026-08 禽类业务量价齐增：上半年肉鸡销量同比提升，鲜品与熟食占比提高，形成对养猪亏损的对冲。 公告口径上，近三年（2023 年 9 月---2026 年 9 月）公司披露重大重组与并购类公告 3 条、对外投资与转型类公告 1 条、回购与股权激励类公告 82 条，合计公告 686 条。 第一大行业“养殖行业”的收入占比由 2021 年年报的 96.3\% 变为 131.1\%（34.8 个百分点），收入结构出现再平衡。 综合判断，公司在报告期内有明确的外延扩张或业务结构调整动作，转型方向与融资安排之间存在直接关联。\footnote{\texttt{https://file.finance.qq.com/finance/hs/pdf/2025/08/01/1224366923.PDF}}\footnote{\texttt{https://paper.cnstock.com/html/2025-04/24/content\_2054137.htm}}\footnote{\texttt{https://pdf.dfcfw.com/pdf/H2\_AN202604211821387744\_1.pdf}}

\begin{center}\small\setlength{\tabcolsep}{4pt}
\captionof{table}[温氏股份（300498）历史融资与资本运作]{温氏股份（300498）历史融资与资本运作（据公司公告与公开报道整理；金额为披露值，含首次公开发行、再融资、债券、银行借款与重整投资等）}
\begin{tabular}{@{}p{1.5cm}p{3.2cm}p{2.4cm}p{6.3cm}@{}}
\toprule
时间 & 方式 & 规模 & 用途与说明 \\
\midrule
2017-03-20 & 公司债券（17 温氏 01，112506） & 5 亿元 & 公司经营资金；票面利率 4.60\%，期限 5 年（附第 3 年末利率调整及回售选择权）…… \\
2017-07-04 & 公司债券（17 温氏 02，112539） & 20 亿元 & 公司经营资金；票面利率 4.87\%，期限 5 年…… \\
2020-07-01 & 公司债券（20 温氏 01） & 23.00 亿元 & 公司经营资金；起息日 2020 年 7 月 1 日，2025 年 7 月 1 日到期…… \\
2020-08-11 & 中期票据（20 温氏食品 MTN001） & 9.00 亿元 & 补充流动资金等；到期兑付日 2023 年 8 月 11 日 \\
2021-03-29 & 可转换公司债券（温氏转债，123107） & 92.97 亿元（92,970,000 张，每张 100 元） & 生猪养殖项目及补充流动资金；截至 2025 年 3 月底募投项目已使用 36.16 亿元…… \\
2025-05 & 公司债券（注册） & 拟发行面值总额不超过 50 亿元、分期发行 & 包括一般公司债券、绿色公司债券、乡村振兴公司债券、科技创新公司债券等品种…… \\
\bottomrule
\end{tabular}
\end{center}

\pfl{历史融投资情况} 从现金流量表看，2025 年公司取得借款收到的现金 7.06 亿元、偿还债务支付的现金 77.76 亿元、吸收权益性投资 0.14 亿元，筹资活动现金流净额 -134.76 亿元（2023 年为取得借款 100.67 亿元、偿还债务 127.01 亿元）。 2026 年上半年筹资活动现金流净额为 62.81 亿元（取得借款 128.27 亿元、偿还债务 36.99 亿元，未年化），仍处于净融入状态。 2025 年分配股利、利润或偿付利息支付的现金合计 40.71 亿元。 公开信息中与公司直接相关的融资与资本运作包括：中期票据（20 温氏食品 MTN001）、公司债券（17 温氏 01，112506）、公司债券（17 温氏 02，112539）、公司债券（20 温氏 01）、公司债券（注册）、可转换公司债券（温氏转债，123107）。 其中规模较大的两笔为：2025-05公司债券（注册）（拟发行面值总额不超过 50 亿元、分期发行）；2021-03-29可转换公司债券（温氏转债，123107）（92.97 亿元（92,970,000 张，每张 100 元……）。 股本方面，近三年披露股本变动 17 次（可转债转股 13 次、股权激励 3 次、增发新股上市 2 次、股份回购 1 次）；总股本由 2021 年末的 63.52 亿股变为 66.53 亿股（+4.74\%）。 分红方面，近三年现金分红 7 次、累计每股派现 1.25 元（最近一次 2025-12-31）——在利润承压的阶段仍维持了分红。 近三年“再融资”类公告 58 条，涉及可转债。\footnote{\texttt{https://www.wens.com.cn/uploadfiles/2019/04/20190410095823852.pdf}}\footnote{\texttt{https://pdf.dfcfw.com/pdf/H2\_AN202007231393627444\_1.pdf}}\footnote{\texttt{https://www.lhratings.com/reports/B0389-P51259-2020-GZ2023.pdf}}

\pfl{未来潜在融资需求分析} 2026 年 6 月末货币资金 36.75 亿元，而短期债务 146.20 亿元（现金短债比 0.25 倍），2026 年上半年经营现金流净额 4.22 亿元。按“短期债务减去账面货币资金”的滚动偿付口径，静态缺口约 109.44 亿元；若再计入上半年亏损的年化额（87.32 亿元），维持一年正常经营所需的外部资金约 109.44---196.76 亿元。 公司 2026 年上半年归母净利润为负（-43.66 亿元），但 6 月末资产负债率 58.9\% 尚未进入第一档（亏损且负债率 $>$65\%）的区间，当前融资需求以补充流动资金、支撑低谷期经营性支出为主。 公开披露的下一步融资线索包括：2025-2026 拟公开发行不超过 50 亿元公司债券（分期发行）已进入注册程序；温氏转债存续至 2027 年 3 月，2025 年 3 ……。 资金需求还要叠加在建工程：期末在建工程 28.48 亿元、2026 年上半年购建固定资产等支出 56.80 亿元（未年化），这部分支出在周期底部通常会被主动放缓。\footnote{\texttt{https://www.lhratings.com/reports/B0389-P51259-2020-GZ2025.pdf}}


\subsubsection{正邦科技（002157）}
\label{co:002157}

\pfl{公司基本情况} 2007 年设立至今，正邦科技曾经的江西养猪龙头企业，2022—2023 年陷入债务危机并被法院裁定重整，由双胞胎集团牵头联合体接盘；2024 年 6 月撤销退市风险警示，现依托双胞胎协同恢复生猪与饲料业务，控股股东承诺 4 年内完成资产整体上市。 公司于 2007-08-17 在深交所主板首发上市，注册资本 92.50 亿元，总股本 92.98 亿股。 截至 2026-09-24收盘价 2.92 元，流通市值 227.4 亿元，近一年-3.0\%，流通市值在三级行业“生猪养殖”的 11 家公司中排第 \zhnumber{3} 位，近一年涨跌幅高于全样本中位数（-10.1\%）7.1 个百分点。 按 2026 年半年报披露的行业构成，收入占比前三位为养殖 53.0\%；饲料 44.3\%；其他(补充) 1.8\%。 发展过程中的关键节点包括：经证监发行字[2007]197 号文核准发行，2007 年 8 月 15 日获深证上[2……（2007）；报告期内完成总额 16 亿元的可转债（正邦转债，128114）发行，并披露总额 75 亿……（2020）；2020 年度非公开发行新增股份 569,908,811 股于 2021 年 2 月 1……（2021）。\footnote{\texttt{http://static.cninfo.com.cn/finalpage/2007-12-25/36215528.PDF}}\footnote{\texttt{http://file.finance.qq.com/finance/hs/pdf/2020/08/28/1208279405.PDF}}\footnote{\texttt{http://file.finance.sina.com.cn/211.154.219.97:9494/MRGG/CNSESZ\_STOCK/2021/2021-8/2021-0}}

\pfl{历史股价} 自 2021 年 1 月至 2026 年 9 月，股价由 17.34 元变为 2.92 元，累计 -83.2\%，区间最高 17.34 元（2021 年 1 月）、最低 2.39 元（2023 年 5 月）；当前价相当于区间最高价的 16.8\%，为区间最低价的 1.22 倍。 月度收益的年化波动率 36.3\%，低于全样本中位数 37.4\%，在三级行业“生猪养殖”的 11 家公司中波动率排第 \zhnumber{7} 位。 从事件看，2021 年 1 月 至 2023 年 5 月的下跌（-86.2\%）与公司因 2022 年度经审计净资产为负值，公司股票被深交所实施退市风险警示，简称……（2023-05-05）在时间上重合；2023 年 5 月 至 2026 年 4 月的上涨（+45.2\%）与公司公司公告双胞胎已全面启动资产注入前期筹备（资产评估、财务审计、同业业务拆分、……（2025-12-26）在时间上重合——股价的拐点多数能在公司自身的资本运作与风险事件上找到对应。\footnote{\texttt{https://m.bjnews.com.cn/detail/1718168037169755.html}}\footnote{\texttt{https://caifuhao.eastmoney.com/news/20260706123702388586080}}

\begin{center}
\begin{tikzpicture}
\begin{axis}[
  width=12.6cm, height=3.9cm, scale only axis,
  xmin=2020.922, xmax=2026.888, ymin=2.22, ymax=18.55,
  xtick={2021,2022,2023,2024,2025,2026},
  yearticks,
  yticklabel style={font=\scriptsize},
  ylabel={元/股}, ylabel style={font=\scriptsize},
  grid=major, grid style={LightGray, line width=0.3pt},
  axis line style={InkGray, line width=0.5pt},
  every axis plot/.append style={line width=0.9pt},
]
\addplot[AgriGreen] table[x=x,y=y]{data/clean/profiles/px_002157.dat};
\addplot[SoftRed, dashed, line width=0.7pt] table[x=x,y=y]{data/clean/profiles/pxzz_002157.dat};
\addplot[only marks, mark=*, mark size=1.1pt, SoftRed] table[x=x,y=y]{data/clean/profiles/pxzz_002157.dat};
\end{axis}
\end{tikzpicture}

\captionof{figure}[正邦科技（002157）股价与周期骨架]{正邦科技（002157）月度收盘价（元/股）与 ZigZag 周期骨架（阈值 25\%，圆点为周期转折点）}
\end{center}

\pfl{过去周期分析} 以 25\% 的 ZigZag 阈值划分，样本区间内股价形成 1 个主升段与 1 个主跌段。 最大主跌段为 2021 年 1 月 至 2023 年 5 月，28 个月-86.2\%；最大主升段出现在 2023 年 5 月 至 2026 年 4 月，35 个月+45.2\%。 最近一次转折出现在 2026 年 4 月（上涨段自 2023 年 5 月起，35 个月+45.2\%），当前处于该段之后的震荡之中。 公司股价较申万农林牧渔指数提前约 36 个月见底，是板块中较早定价周期反转的公司之一。 生猪养殖的定价主线是产能去化与成本曲线的赛跑（第 \ref{sec:hog} 章）：2026 年三季度行业仍处在“产能去化未完成、价格磨底”的阶段，公司 2026 年 6 月末资产负债率 53.5\%、半年 ROE -7.0\%（未年化），在本轮下行中属于随行业同步调整的一类。

\pfl{盈利情况} 2026 年上半年营业收入 77.20 亿元（同比 +12.9\%），归母净利润 -7.36 亿元，由上年同期的盈利转为亏损。 以年报为趋势对照，2023---2025 年营业收入由 69.92 亿元变为 147.95 亿元（两年年均复合增速 +45.5\%），归母净利润由 85.29 亿元变为 -5.46 亿元。 按半年报口径，2026 年上半年的 ROE 为 -7.0\%（未年化）、销售净利率 -12.7\%、毛利率 -0.2\%，ROE 在“生猪养殖”11 家公司中排第 \zhnumber{1} 位。 同期全样本半年 ROE 中位数 0.75\%，公司与之相差 7.74 个百分点（弱于中位数公司）。 结合公司披露，公司公告双胞胎已全面启动资产注入前期筹备（资产评估、财务审计、同业业务拆分、合规整改）；原定 2025 ……；2024 年末已逾期未偿还短期借款 16,840.95 万元（重整后遗留部分子公司借款逾期事项）。\footnote{\texttt{https://caifuhao.eastmoney.com/news/20260706123702388586080}}\footnote{\texttt{https://money.finance.sina.com.cn/corp/view/vCB\_AllBulletinDetail.php}}

\begin{center}\small\setlength{\tabcolsep}{4pt}
\captionof{table}[正邦科技（002157）关键财务指标]{正邦科技（002157）关键财务指标（合并报表口径；两个半年报期均未年化；现金短债比 $=$ 货币资金 $\div$ 短期债务）}
\begin{tabular}{@{}lrrrrr@{}}
\toprule
指标 & 2023 年 & 2024 年 & 2025 年 & 2025 年半年报 & 2026 年半年报 \\
\midrule
营业收入（亿元） & 69.92 & 88.70 & 147.95 & 68.38 & 77.20 \\
归母净利润（亿元） & 85.29 & 2.16 & -5.46 & 2.02 & -7.36 \\
毛利率（\%） & -14.3 & 14.8 & 8.0 & 14.8 & -0.2 \\
ROE（\%） & 202.0 & 2.0 & -4.9 & 1.7 & -7.0 \\
资产负债率（\%） & 54.0 & 46.4 & 50.1 & 45.7 & 53.5 \\
经营活动现金流净额（亿元） & -13.83 & 1.74 & 8.30 & 0.20 & -4.07 \\
货币资金（亿元） & 24.42 & 33.21 & 31.20 & 34.18 & 16.97 \\
有息负债合计（亿元） & 23.19 & 18.28 & 18.17 & 17.17 & 19.10 \\
现金短债比（倍） & 3.42 & 6.05 & 5.29 & 6.80 & 2.43 \\
\bottomrule
\end{tabular}
\end{center}

\begin{center}
\begin{tikzpicture}
\begin{groupplot}[
  group style={group size=2 by 1, horizontal sep=0.95cm},
  width=7.2cm, height=4.6cm,
  tick label style={font=\tiny},
  title style={font=\scriptsize\heiti},
  yticklabel style={font=\tiny},
  grid=major, grid style={LightGray, line width=0.3pt},
  axis line style={InkGray, line width=0.5pt},
  enlarge x limits={abs=0.8},
  legend style={font=\scriptsize, at={(0.5,-0.22)}, anchor=north,
                legend columns=2, draw=none, fill=none, column sep=6pt},
]
\nextgroupplot[
  ybar, bar width=4pt,
  ymin=-254.68, ymax=556.49,
  xtick={0,1,2,3,4,5.7},
  xticklabels={2021,2022,2023,2024,2025,2026H1},
  ylabel={亿元},
]
\addplot[fill=AgriGreen!75, draw=AgriGreen, bar shift=-2.6pt] table[x=i, y=rev]{data/clean/profiles/fin_002157.dat};
\addplot[fill=SoftRed!60, draw=SoftRed, bar shift=2.6pt] table[x=i, y=ni]{data/clean/profiles/fin_002157.dat};
\addplot[fill=AgriGreen!30, draw=AgriGreen, pattern=north east lines, bar shift=-2.6pt] table[x=x, y=rev]{data/clean/profiles/finh_002157.dat};
\addplot[fill=SoftRed!20, draw=SoftRed, pattern=north east lines, bar shift=2.6pt] table[x=x, y=ni]{data/clean/profiles/finh_002157.dat};
\legend{营业收入（年/半年）, 归母净利润（年/半年）}
\nextgroupplot[
  ymin=-4409.1, ymax=628.9,
  xtick={0,1,2,3,4,5.7},
  xticklabels={2021,2022,2023,2024,2025,2026H1},
  ylabel={\%},
]
\addplot[AgriGold, mark=*, mark size=1.3pt] table[x=i, y=roe]{data/clean/profiles/fin_002157.dat};
\addplot[OilBlue, mark=square*, mark size=1.3pt] table[x=i, y=debt]{data/clean/profiles/fin_002157.dat};
\addplot[only marks, AgriGold, mark=*, mark size=2.0pt] table[x=x, y=roe]{data/clean/profiles/finh_002157.dat};
\addplot[only marks, OilBlue, mark=square*, mark size=2.0pt] table[x=x, y=debt]{data/clean/profiles/finh_002157.dat};
\legend{ROE, 资产负债率}
\end{groupplot}
\end{tikzpicture}
\begin{tikzpicture}
\begin{axis}[
  width=15.2cm, height=3.1cm, scale only axis,
  ymin=-88.87, ymax=328.33,
  xtick={0,1,2,3,4,5.7},
  xticklabels={2021,2022,2023,2024,2025,2026H1},
  tick label style={font=\tiny},
  yticklabel style={font=\tiny},
  ylabel={亿元}, ylabel style={font=\scriptsize},
  ybar, bar width=4pt,
  grid=major, grid style={LightGray, line width=0.3pt},
  axis line style={InkGray, line width=0.5pt},
  enlarge x limits={abs=0.8},
  legend style={font=\scriptsize, at={(0.5,-0.30)}, anchor=north,
                legend columns=3, draw=none, fill=none, column sep=10pt},
]
\addplot[fill=AgriGreen!55, draw=AgriGreen, bar shift=-3.0pt] table[x=i, y=cash]{data/clean/profiles/fin_002157.dat};
\addplot[fill=SoftRed!45, draw=SoftRed, bar shift=3.0pt] table[x=i, y=ibd]{data/clean/profiles/fin_002157.dat};
\addplot[fill=AgriGreen!25, draw=AgriGreen, pattern=north east lines, bar shift=-3.0pt] table[x=x, y=cash]{data/clean/profiles/finh_002157.dat};
\addplot[fill=SoftRed!20, draw=SoftRed, pattern=north east lines, bar shift=3.0pt] table[x=x, y=ibd]{data/clean/profiles/finh_002157.dat};
\addplot[OilBlue, mark=*, mark size=1.1pt, line width=0.8pt] table[x=i, y=ocf]{data/clean/profiles/fin_002157.dat};
\addplot[only marks, OilBlue, mark=*, mark size=1.8pt] table[x=x, y=ocf]{data/clean/profiles/finh_002157.dat};
\legend{货币资金, 有息负债, 经营活动现金流净额}
\end{axis}
\end{tikzpicture}

\captionof{figure}[正邦科技（002157）近年主要财务指标]{正邦科技（002157）近年主要财务指标（左上：营业收入与归母净利润；右上：ROE 与资产负债率；下：货币资金、有息负债与经营活动现金流净额。
2021---2025 年为年报口径，2026H1 为 2026 年半年报/半年末，与其前一年报之间留出间隔、以斜纹或空心点区分；半年报数据未年化）}
\end{center}

\pfl{财务分析} 2026 年 6 月末资产负债率 53.5\%，较 2025 年末 上升 3.4 个百分点（样本中位数 52.2\%，所处三级行业内排第 \zhnumber{4} 位）。 2026 年 6 月末货币资金 16.97 亿元、短期债务（短期借款与一年内到期非流动负债）6.97 亿元、长期债务 12.12 亿元，有息负债合计 19.10 亿元，现金短债比 2.43 倍——覆盖充分。 净有息负债 2.13 亿元，相当于其 2026 年上半年营业收入的 3\%。 流动比率 0.74、速动比率 0.44，短期偿付压力较大。 2026 年上半年经营活动现金流净额 -4.07 亿元，经营现金流与利润同时为负，日常运营对债务与股东投入的依赖度较高。 2026 年 6 月末在建工程 26.34 亿元，较上年末减少 2.56 亿元，在建投入在收缩。 公司披露的财务与合规风险包括：因 2022 年度经审计净资产为负值，公司股票被深交所实施退市风险警示，简称变更为“*S……；2024 年末已逾期未偿还短期借款 16,840.95 万元（重整后遗留部分子公司借款逾……。\footnote{\texttt{https://m.bjnews.com.cn/detail/1718168037169755.html}}\footnote{\texttt{https://money.finance.sina.com.cn/corp/view/vCB\_AllBulletinDetail.php}}

\pfl{重大战略转型} 近三年公司的战略动作集中在以下几个方面：2023-11 双胞胎信达联合体（牵头投资人：江西双胞胎农业有限公司）被确定为中选重整投资人，通过新设合伙企业以现金及抵债资产出……；2023-12-27 双胞胎集团作为产业投资人联合 16 家财务投资人（含央企、国企、民企）正式接管正邦科技，重整投资人合计提供资金 ……；2024-05 双胞胎集团承诺：重整完成后 24 个月内逐步启动自身生猪养殖、饲料等业务资产置入上市公司的重组程序，并在重整完成……；2025-12-26 公司公告双胞胎已全面启动资产注入前期筹备（资产评估、财务审计、同业业务拆分、合规整改）；原定 2025 年 12……。 公告口径上，近三年（2023 年 9 月---2026 年 9 月）公司披露重大重组与并购类公告 0 条、对外投资与转型类公告 6 条、回购与股权激励类公告 14 条，合计公告 473 条。 第一大行业仍是“养殖”，收入占比 53.0\%，与 2021 年年报相比没有方向性变化。 综合判断，公司当前的战略重心不是扩张而是化解债务与完成重整，业务层面的调整让位于司法重整程序。\footnote{\texttt{https://file.finance.qq.com/finance/hs/pdf/2023/11/28/1218454349.PDF}}\footnote{\texttt{https://app.xinhuanet.com/news/article.html}}\footnote{\texttt{https://pdf.dfcfw.com/pdf/H2\_AN202405101632828391\_1.pdf}}

\begin{center}\small\setlength{\tabcolsep}{4pt}
\captionof{table}[正邦科技（002157）历史融资与资本运作]{正邦科技（002157）历史融资与资本运作（据公司公告与公开报道整理；金额为披露值，含首次公开发行、再融资、债券、银行借款与重整投资等）}
\begin{tabular}{@{}p{1.5cm}p{3.2cm}p{2.4cm}p{6.3cm}@{}}
\toprule
时间 & 方式 & 规模 & 用途与说明 \\
\midrule
2007-08 & IPO（深交所） & 1,900 万股，发行价 11.09 元/股（约 2.11 亿元） & 饲料及养殖项目建设；主承销商国信证券；发行后总股本 7,529.68 万股 \\
2020 & 可转换公司债券（正邦转债，128114） & 16 亿元 & 养殖产能扩张；重整期间债券持有人会议决议：正邦转债最后交易日为 2023 年 8 月 4…… \\
2021-02-01 & 定向增发（2020 年度非公开发行） & 569,908,811 股，发行价 13.16 元/股，募集资金总额 75.00 亿元、净额 74.78 亿元 & 发展生猪养殖和补充流动资金；2020 年 6 月 29 日董事会通过预案…… \\
2023-12 & 破产重整投资款（转增股票引入重整投资人） & 重整投资人合计提供资金 43.40 亿元 & 支付重整费用、共益债务、清偿各类债务及补充上市公司流动资金…… \\
2026 & 关联方借款额度 & 任一时点借款余额不超过 15 亿元 & 满足短期内、临时性补充流动资金需求；向控股股东江西双胞胎农业有限公司同一控制下的关联方双…… \\
\bottomrule
\end{tabular}
\end{center}

\pfl{历史融投资情况} 从现金流量表看，2025 年公司取得借款收到的现金 — 亿元、偿还债务支付的现金 0.01 亿元、吸收权益性投资 0.08 亿元，筹资活动现金流净额 -2.90 亿元（2023 年为取得借款 7.64 亿元、偿还债务 1.91 亿元）。 2026 年上半年筹资活动现金流净额为 1.89 亿元（取得借款 — 亿元、偿还债务 — 亿元，未年化），仍处于净融入状态。 2025 年分配股利、利润或偿付利息支付的现金合计 0.26 亿元。 公司披露的股本与债务融资脉络包括：IPO（深交所）、关联方借款额度、可转换公司债券（正邦转债，128114）、定向增发（2020 年度非公开发行）、破产重整投资款（转增股票引入重整投资人）。 其中规模较大的两笔为：2026关联方借款额度（任一时点借款余额不超过 15 亿元）；2023-12破产重整投资款（转增股票引入重整投资人）（重整投资人合计提供资金 43.40 亿元）。 股本方面，近三年披露股本变动 10 次（可转债转股 6 次、转增 1 次、限售股份上市 1 次）；总股本由 2021 年末的 31.46 亿股变为 92.98 亿股（+195.51\%）。 分红方面，近三年未实施现金分红，与重整期间的资金约束一致。 近三年“再融资”类公告 12 条，涉及定向增发、可转债。\footnote{\texttt{http://money.finance.sina.com.cn/corp/go.php/vISSUE\_NewStock/stockid/002157.phtml}}\footnote{\texttt{https://vip.stock.finance.sina.com.cn/corp/view/vCB\_AllBulletinDetail.php}}\footnote{\texttt{https://vip.stock.finance.sina.com.cn/corp/view/vISSUE\_MarketBulletinDetail.php}}

\pfl{未来潜在融资需求分析} 2026 年 6 月末货币资金 16.97 亿元，而短期债务 6.97 亿元（现金短债比 2.43 倍），2026 年上半年经营现金流净额 -4.07 亿元。账面货币资金可以覆盖短期债务，缺口不体现在静态偿付层面。 公司 2026 年上半年归母净利润为负（-7.36 亿元），但 6 月末资产负债率 53.5\% 尚未进入第一档（亏损且负债率 $>$65\%）的区间，当前融资需求以补充流动资金、支撑低谷期经营性支出为主。 公开披露的下一步融资线索包括：2026 2026 年度担保额度安排：资产负债率 70\% 以上的子公司可享担保总额度 20 亿元，70\% 以下的子公司可享 10 ……；2026-08 2026 年半年报披露：控股股东已启动资产注入前期准备工作，并将在重整完成后 4 年内完成相关业务及资产的整体上市；公司……。 资金需求还要叠加在建工程：期末在建工程 26.34 亿元、2026 年上半年购建固定资产等支出 6.94 亿元（未年化），这部分支出在周期底部通常会被主动放缓。\footnote{\texttt{https://paper.cnstock.com/html/2026-08/25/content\_2258766.htm}}\footnote{\texttt{https://stock.stockstar.com/notice/SN2026082400022781.shtml}}


\subsubsection{圣农发展（002299）}
\label{co:002299}

\pfl{公司基本情况} 圣农发展成立于 2009 年，国内白羽肉鸡全产业链龙头，覆盖自主育种、饲料、养殖、屠宰与深加工，产能规模居行业前列；2025 年完成对原嘉吉中国白羽鸡项目（太阳谷）的全资控股，2026 年上半年营收突破百亿元、扣非净利同比增长近四成。 公司于 2009-10-21 在深交所主板首发上市，注册资本 12.43 亿元，总股本 12.43 亿股。 截至 2026-09-24收盘价 15.70 元，流通市值 191.3 亿元，近一年-10.1\%，流通市值在三级行业“肉鸡养殖”的 8 家公司中排第 \zhnumber{1} 位，近一年涨跌幅低于全样本中位数（-10.1\%）0.0 个百分点。 按 2026 年半年报披露的行业构成，收入占比前三位为家禽饲养加工行业 49.8\%；食品加工行业 45.5\%；其他(补充) 4.7\%。 发展过程中的关键节点包括：2 月启动发行股份购买圣农食品 100\% 股权，10 月完成资产过户，11 月新增股份在……（2017）；5 月 21 日完成对安徽太阳谷食品科技（集团）有限公司 54\% 股权的收购，持股比例由……（2025）。\footnote{\texttt{http://pdf.dfcfw.com/pdf/H2\_AN201804221127691596\_1.pdf}}\footnote{\texttt{http://notice.10jqka.com.cn/api/pdf/1defcca22c1d799b.pdf}}

\pfl{历史股价} 自 2021 年 1 月至 2026 年 9 月，股价由 28.65 元变为 15.70 元，累计 -45.2\%，区间最高 28.80 元（2021 年 2 月）、最低 12.17 元（2024 年 8 月）；当前价相当于区间最高价的 54.5\%，为区间最低价的 1.29 倍。 月度收益的年化波动率 31.4\%，低于全样本中位数 37.4\%，在三级行业“肉鸡养殖”的 8 家公司中波动率排第 \zhnumber{8} 位。 从事件看，2024 年 8 月 至 2026 年 2 月的上涨（+58.6\%）与公司完成太阳谷剩余 54\% 股权收购，实现 100\% 控股，公司产能向 10 亿……（2025-05-21）在时间上重合——股价的拐点多数能在公司自身的资本运作与风险事件上找到对应。\footnote{\texttt{https://epaper.stcn.com/pic/202504/07/f02db5a1e1e433f8e1c9a3480c8ee32c.pdf}}

\begin{center}
\begin{tikzpicture}
\begin{axis}[
  width=12.6cm, height=3.9cm, scale only axis,
  xmin=2020.922, xmax=2026.888, ymin=11.32, ymax=30.82,
  xtick={2021,2022,2023,2024,2025,2026},
  yearticks,
  yticklabel style={font=\scriptsize},
  ylabel={元/股}, ylabel style={font=\scriptsize},
  grid=major, grid style={LightGray, line width=0.3pt},
  axis line style={InkGray, line width=0.5pt},
  every axis plot/.append style={line width=0.9pt},
]
\addplot[AgriGreen] table[x=x,y=y]{data/clean/profiles/px_002299.dat};
\addplot[SoftRed, dashed, line width=0.7pt] table[x=x,y=y]{data/clean/profiles/pxzz_002299.dat};
\addplot[only marks, mark=*, mark size=1.1pt, SoftRed] table[x=x,y=y]{data/clean/profiles/pxzz_002299.dat};
\end{axis}
\end{tikzpicture}

\captionof{figure}[圣农发展（002299）股价与周期骨架]{圣农发展（002299）月度收盘价（元/股）与 ZigZag 周期骨架（阈值 25\%，圆点为周期转折点）}
\end{center}

\pfl{过去周期分析} 以 25\% 的 ZigZag 阈值划分，样本区间内股价形成 3 个主升段与 3 个主跌段。 最大主升段出现在 2024 年 8 月 至 2026 年 2 月，18 个月+58.6\%；最大主跌段为 2022 年 11 月 至 2024 年 8 月，21 个月-51.4\%。 最近一次转折出现在 2026 年 2 月（上涨段自 2024 年 8 月起，18 个月+58.6\%），当前处于该段之后的震荡之中。 公司股价较申万农林牧渔指数提前约 21 个月见底，是板块中较早定价周期反转的公司之一。 禽链的产能周期只有 6---12 个月、反转快于生猪（第 \ref{sec:other-livestock} 节），2026 年前四个月毛鸡养殖利润同比大幅回升，链条利润的修复先于生猪，公司 2026 年上半年实现盈利、半年 ROE 4.4\%（未年化），好于全样本中位数（0.75\%），下行周期对其报表的冲击小于同业。

\pfl{盈利情况} 2026 年上半年营业收入 107.81 亿元（同比 +21.7\%），归母净利润 5.02 亿元，同比 -44.8\%。 以年报为趋势对照，2023---2025 年营业收入由 184.87 亿元变为 200.94 亿元（两年年均复合增速 +4.3\%），归母净利润由 6.64 亿元变为 13.80 亿元。 按半年报口径，2026 年上半年的 ROE 为 4.4\%（未年化）、销售净利率 4.7\%、毛利率 13.1\%，ROE 在“肉鸡养殖”8 家公司中排第 \zhnumber{2} 位。 行业同口径半年 ROE 的中位数为 0.75\%，该公司高于中位数 3.65 个百分点。

\begin{center}\small\setlength{\tabcolsep}{4pt}
\captionof{table}[圣农发展（002299）关键财务指标]{圣农发展（002299）关键财务指标（合并报表口径；两个半年报期均未年化；现金短债比 $=$ 货币资金 $\div$ 短期债务）}
\begin{tabular}{@{}lrrrrr@{}}
\toprule
指标 & 2023 年 & 2024 年 & 2025 年 & 2025 年半年报 & 2026 年半年报 \\
\midrule
营业收入（亿元） & 184.87 & 185.86 & 200.94 & 88.56 & 107.81 \\
归母净利润（亿元） & 6.64 & 7.24 & 13.80 & 9.10 & 5.02 \\
毛利率（\%） & 10.6 & 11.1 & 12.6 & 12.0 & 13.1 \\
ROE（\%） & 6.6 & 7.0 & 12.6 & 8.4 & 4.4 \\
资产负债率（\%） & 54.1 & 50.0 & 51.7 & 55.2 & 51.9 \\
经营活动现金流净额（亿元） & 22.91 & 30.49 & 40.32 & 14.25 & 9.14 \\
货币资金（亿元） & 7.38 & 7.54 & 14.47 & 27.55 & 11.73 \\
有息负债合计（亿元） & 81.38 & 59.17 & 67.99 & 79.71 & 76.86 \\
现金短债比（倍） & 0.11 & 0.15 & 0.24 & 0.38 & 0.17 \\
\bottomrule
\end{tabular}
\end{center}

\begin{center}
\begin{tikzpicture}
\begin{groupplot}[
  group style={group size=2 by 1, horizontal sep=0.95cm},
  width=7.2cm, height=4.6cm,
  tick label style={font=\tiny},
  title style={font=\scriptsize\heiti},
  yticklabel style={font=\tiny},
  grid=major, grid style={LightGray, line width=0.3pt},
  axis line style={InkGray, line width=0.5pt},
  enlarge x limits={abs=0.8},
  legend style={font=\scriptsize, at={(0.5,-0.22)}, anchor=north,
                legend columns=2, draw=none, fill=none, column sep=6pt},
]
\nextgroupplot[
  ybar, bar width=4pt,
  ymin=-20.09, ymax=225.06,
  xtick={0,1,2,3,4,5.7},
  xticklabels={2021,2022,2023,2024,2025,2026H1},
  ylabel={亿元},
]
\addplot[fill=AgriGreen!75, draw=AgriGreen, bar shift=-2.6pt] table[x=i, y=rev]{data/clean/profiles/fin_002299.dat};
\addplot[fill=SoftRed!60, draw=SoftRed, bar shift=2.6pt] table[x=i, y=ni]{data/clean/profiles/fin_002299.dat};
\addplot[fill=AgriGreen!30, draw=AgriGreen, pattern=north east lines, bar shift=-2.6pt] table[x=x, y=rev]{data/clean/profiles/finh_002299.dat};
\addplot[fill=SoftRed!20, draw=SoftRed, pattern=north east lines, bar shift=2.6pt] table[x=x, y=ni]{data/clean/profiles/finh_002299.dat};
\legend{营业收入（年/半年）, 归母净利润（年/半年）}
\nextgroupplot[
  ymin=-4.3, ymax=59.6,
  xtick={0,1,2,3,4,5.7},
  xticklabels={2021,2022,2023,2024,2025,2026H1},
  ylabel={\%},
]
\addplot[AgriGold, mark=*, mark size=1.3pt] table[x=i, y=roe]{data/clean/profiles/fin_002299.dat};
\addplot[OilBlue, mark=square*, mark size=1.3pt] table[x=i, y=debt]{data/clean/profiles/fin_002299.dat};
\addplot[only marks, AgriGold, mark=*, mark size=2.0pt] table[x=x, y=roe]{data/clean/profiles/finh_002299.dat};
\addplot[only marks, OilBlue, mark=square*, mark size=2.0pt] table[x=x, y=debt]{data/clean/profiles/finh_002299.dat};
\legend{ROE, 资产负债率}
\end{groupplot}
\end{tikzpicture}
\begin{tikzpicture}
\begin{axis}[
  width=15.2cm, height=3.1cm, scale only axis,
  ymin=-8.14, ymax=91.14,
  xtick={0,1,2,3,4,5.7},
  xticklabels={2021,2022,2023,2024,2025,2026H1},
  tick label style={font=\tiny},
  yticklabel style={font=\tiny},
  ylabel={亿元}, ylabel style={font=\scriptsize},
  ybar, bar width=4pt,
  grid=major, grid style={LightGray, line width=0.3pt},
  axis line style={InkGray, line width=0.5pt},
  enlarge x limits={abs=0.8},
  legend style={font=\scriptsize, at={(0.5,-0.30)}, anchor=north,
                legend columns=3, draw=none, fill=none, column sep=10pt},
]
\addplot[fill=AgriGreen!55, draw=AgriGreen, bar shift=-3.0pt] table[x=i, y=cash]{data/clean/profiles/fin_002299.dat};
\addplot[fill=SoftRed!45, draw=SoftRed, bar shift=3.0pt] table[x=i, y=ibd]{data/clean/profiles/fin_002299.dat};
\addplot[fill=AgriGreen!25, draw=AgriGreen, pattern=north east lines, bar shift=-3.0pt] table[x=x, y=cash]{data/clean/profiles/finh_002299.dat};
\addplot[fill=SoftRed!20, draw=SoftRed, pattern=north east lines, bar shift=3.0pt] table[x=x, y=ibd]{data/clean/profiles/finh_002299.dat};
\addplot[OilBlue, mark=*, mark size=1.1pt, line width=0.8pt] table[x=i, y=ocf]{data/clean/profiles/fin_002299.dat};
\addplot[only marks, OilBlue, mark=*, mark size=1.8pt] table[x=x, y=ocf]{data/clean/profiles/finh_002299.dat};
\legend{货币资金, 有息负债, 经营活动现金流净额}
\end{axis}
\end{tikzpicture}

\captionof{figure}[圣农发展（002299）近年主要财务指标]{圣农发展（002299）近年主要财务指标（左上：营业收入与归母净利润；右上：ROE 与资产负债率；下：货币资金、有息负债与经营活动现金流净额。
2021---2025 年为年报口径，2026H1 为 2026 年半年报/半年末，与其前一年报之间留出间隔、以斜纹或空心点区分；半年报数据未年化）}
\end{center}

\pfl{财务分析} 2026 年 6 月末资产负债率 51.9\%，较 2025 年末 上升 0.3 个百分点（样本中位数 52.2\%，所处三级行业内排第 \zhnumber{4} 位）。 2026 年 6 月末货币资金 11.73 亿元、短期债务（短期借款与一年内到期非流动负债）69.56 亿元、长期债务 7.30 亿元，有息负债合计 76.86 亿元，现金短债比 0.17 倍——缺口明显。 净有息负债 65.13 亿元，相当于其 2026 年上半年营业收入的 60\%。 流动比率 0.67、速动比率 0.36，短期偿付压力较大。 2026 年上半年经营活动现金流净额 9.14 亿元，经营现金流与利润同向为正，内生资金可以支撑运营。 2026 年 6 月末在建工程 4.40 亿元，较上年末增加 0.78 亿元，仍有在建投入。 2026 年 6 月末商誉 10.16 亿元，占归母净资产的 8.8\%，是净资产中最脆弱的一块。

\pfl{重大战略转型} 公司在报告期内的战略动作可归为 4 项：2021-01 2020 年三季度末以 1,244,501,917 股为基数每 10 股派发现金红利 10 元，2021 年 1……；2023-11 并购原嘉吉中国白羽肉鸡项目（太阳谷），从德弘资本手中受让 46\% 股权并受让可转股债权，向上游产能扩张；2025-05-21 完成太阳谷剩余 54\% 股权收购，实现 100\% 控股，公司产能向 10 亿羽/年目标推进；2026-04 2025 年度营业收入 200.94 亿元（同比 +8.12\%）、归母净利润 13.8 亿元（同比 +90.55……。 公告口径上，近三年（2023 年 9 月---2026 年 9 月）公司披露重大重组与并购类公告 3 条、对外投资与转型类公告 3 条、回购与股权激励类公告 31 条，合计公告 377 条。 第一大行业“家禽饲养加工行业”的收入占比由 2021 年年报的 61.7\% 变为 49.8\%（11.9 个百分点），收入结构出现再平衡。 综合判断，公司在报告期内有明确的外延扩张或业务结构调整动作，转型方向与融资安排之间存在直接关联。\footnote{\texttt{https://money.finance.sina.com.cn/corp/view/vCB\_AllBulletinDetail.php}}\footnote{\texttt{http://file.finance.sina.com.cn/211.154.219.97:9494/MRGG/CNSESZ\_STOCK/2021/2021-4/2021-0}}\footnote{\texttt{https://news.qq.com/rain/a/20231117A04IU500}}

\begin{center}\small\setlength{\tabcolsep}{4pt}
\captionof{table}[圣农发展（002299）历史融资与资本运作]{圣农发展（002299）历史融资与资本运作（据公司公告与公开报道整理；金额为披露值，含首次公开发行、再融资、债券、银行借款与重整投资等）}
\begin{tabular}{@{}p{1.5cm}p{3.2cm}p{2.4cm}p{6.3cm}@{}}
\toprule
时间 & 方式 & 规模 & 用途与说明 \\
\midrule
2009-10 & IPO（深交所中小板） & 4,100 万股，发行价 19.75 元/股，募集资金总额 8.0975 亿元、净额 7.6496 亿元 & 优化各生产环节产能（自养自宰模式下的养殖与屠宰项目）；网下配售 820 万股…… \\
2017-10 & 发行股份购买资产（关联交易） & 交易价格 20.20 亿元，发行 128,580,517 股，发行价 15.71 元/股 & 购买圣农实业等 8 名交易对方合计持有的圣农食品 100\% 股权…… \\
2023-11-13 & 对外投资（收购股权及可转股债权） & 圣农发展以 2.4471 亿元受让太阳谷 45.1557\% 股权，廖俊杰先生以 1,625.40 万元受让 3.0000\% 股权；同时受让德弘钰玺持有的太阳谷可转股债权（本金 6,716 万元、438 万元） & 取得原嘉吉中国白羽鸡项目相关主体的股权与债权；公司第六届董事会第二十三次会议审议通过 \\
2025-05-21 & 现金收购股权 & 转让总价款 112,590.00 万元 & 收购德弘钰玺、Ancient Steel (HK) Limited 及廖俊杰先生持有的太…… \\
\bottomrule
\end{tabular}
\end{center}

\pfl{历史融投资情况} 从现金流量表看，2025 年公司取得借款收到的现金 113.31 亿元、偿还债务支付的现金 110.25 亿元、吸收权益性投资 0.56 亿元，筹资活动现金流净额 -11.05 亿元（2023 年为取得借款 77.43 亿元、偿还债务 65.70 亿元）。 2026 年上半年筹资活动现金流净额为 2.88 亿元（取得借款 59.33 亿元、偿还债务 53.29 亿元，未年化），仍处于净融入状态。 2025 年分配股利、利润或偿付利息支付的现金合计 10.70 亿元。 公开信息中与公司直接相关的融资与资本运作包括：IPO（深交所中小板）、发行股份购买资产（关联交易）、对外投资（收购股权及可转股债权）、现金收购股权。 其中规模较大的两笔为：2009-10IPO（深交所中小板）（4,100 万股，发行价 19.75 元/股，募集资金总额 ……）；2025-05-21现金收购股权（转让总价款 112,590.00 万元）。 股本方面，近三年披露股本变动 9 次（股份回购 3 次、激励股份解禁 2 次、其他 1 次）；总股本由 2021 年末的 12.44 亿股变为 12.43 亿股（-0.05\%）。 分红方面，近三年现金分红 6 次、累计每股派现 1.40 元（最近一次 2025-12-31）——在利润承压的阶段仍维持了分红。\footnote{\texttt{http://file.finance.sina.com.cn/211.154.219.97:9494/MRGG/CNSESZ\_STOCK/2016/2016-4/2016-0}}\footnote{\texttt{http://static.cninfo.com.cn/finalpage/2017-08-15/1203800361.PDF}}\footnote{\texttt{http://vip.stock.finance.sina.com.cn/corp/view/vCB\_AllBulletinDetail.php}}

\pfl{未来潜在融资需求分析} 2026 年 6 月末货币资金 11.73 亿元，而短期债务 69.56 亿元（现金短债比 0.17 倍），2026 年上半年经营现金流净额 9.14 亿元。按“短期债务减去账面货币资金”的滚动偿付口径，静态缺口约 57.83 亿元；若再计入上半年亏损的年化额（0.00 亿元），维持一年正常经营所需的外部资金约 57.83---57.83 亿元。 公司 2026 年 6 月末资产负债率 51.9\%、2026 年上半年 ROE 4.4\%（未年化），资产负债表尚有余量，融资需求不是“补血型”的刚性需求，而是围绕并购整合、项目投入与营运资金的主动性需求。 公开披露的下一步融资线索包括：2021-2026 未查得 2018 年以来公司实施定向增发、可转债、配股等股权融资事项，近年扩张（收购太阳谷）以自有资金或自筹资金支付。 资金需求还要叠加在建工程：期末在建工程 4.40 亿元、2026 年上半年购建固定资产等支出 8.64 亿元（未年化），这部分支出在周期底部通常会被主动放缓。\footnote{\texttt{http://notice.10jqka.com.cn/api/pdf/1defcca22c1d799b.pdf}}


\subsubsection{神农集团（605296）}
\label{co:605296}

\pfl{公司基本情况} 公司前身可追溯至 1999 年，云南本土生猪一体化企业，饲料—养殖—屠宰全链条布局，2021 年在沪市主板上市后逆周期扩产、出栏量快速攀升；2026 年上半年猪价深跌致大额亏损，此前两项定增计划先后终止。 公司于 2021-05-28 在上交所首发上市，注册资本 5.26 亿元，总股本 5.25 亿股。 截至 2026-09-24收盘价 30.30 元，流通市值 156.9 亿元，近一年-3.5\%，流通市值在三级行业“生猪养殖”的 11 家公司中排第 \zhnumber{4} 位，近一年涨跌幅高于全样本中位数（-10.1\%）6.6 个百分点。 按 2026 年半年报披露的产品构成，收入占比前三位为生猪养殖分部 90.7\%；饲料加工分部 45.4\%；屠宰及食品加工分部 23.6\%。 与 2021 年年报相比，第一大产品由“饲料”变为“生猪养殖分部”，主业重心已经迁移。 发展过程中的关键节点包括：5 月 28 日在上海证券交易所上市，发行 4,003 万股，发行价 56.08 元/股……（2021）；IPO 合计募集资金约 22 亿元；2024 年 11 月 28 日首发限售股 4.54……（2021）；4 月公告终止 2023 年度非公开发行（原计划最高融资 11.5 亿元），同时启动以简……（2024）。\footnote{\texttt{https://money.finance.sina.com.cn/corp/view/vCB\_AllBulletinDetail.php}}\footnote{\texttt{http://finance.ce.cn/stock/gsgdbd/202411/28/t20241128\_39218362.shtml}}\footnote{\texttt{https://finance.sina.com.cn/stock/observe/2025-03-11/doc-inephsqc9586225.shtml}}

\pfl{历史股价} 以 2021 年 5 月为起点、2026 年 9 月为终点，股价由 56.99 元变为 30.30 元，累计 -46.8\%，区间最高 56.99 元（2021 年 5 月）、最低 20.97 元（2023 年 9 月）；当前价相当于区间最高价的 53.2\%，为区间最低价的 1.44 倍。 月度收益的年化波动率 33.4\%，低于全样本中位数 37.4\%，在三级行业“生猪养殖”的 11 家公司中波动率排第 \zhnumber{9} 位。 从事件看，2023 年 9 月 至 2024 年 5 月的上涨（+93.5\%）与公司2023 年度归属于上市公司股东的净利润 -4.01 亿元，2023 年度不……（2024-04-26）在时间上重合；2021 年 5 月 至 2023 年 9 月的下跌（-63.2\%）与公司2024 年 4 月公告终止该次定增，同时启动简易程序定增方案（2023-06）在时间上重合——股价的拐点多数能在公司自身的资本运作与风险事件上找到对应。\footnote{\texttt{https://money.finance.sina.com.cn/corp/view/vCB\_AllBulletinDetail.php}}\footnote{\texttt{https://finance.sina.com.cn/stock/observe/2025-03-11/doc-inephsqc9586225.shtml}}

\begin{center}
\begin{tikzpicture}
\begin{axis}[
  width=12.6cm, height=3.9cm, scale only axis,
  xmin=2021.255, xmax=2026.888, ymin=19.50, ymax=60.98,
  xtick={2021,2022,2023,2024,2025,2026},
  yearticks,
  yticklabel style={font=\scriptsize},
  ylabel={元/股}, ylabel style={font=\scriptsize},
  grid=major, grid style={LightGray, line width=0.3pt},
  axis line style={InkGray, line width=0.5pt},
  every axis plot/.append style={line width=0.9pt},
]
\addplot[AgriGreen] table[x=x,y=y]{data/clean/profiles/px_605296.dat};
\addplot[SoftRed, dashed, line width=0.7pt] table[x=x,y=y]{data/clean/profiles/pxzz_605296.dat};
\addplot[only marks, mark=*, mark size=1.1pt, SoftRed] table[x=x,y=y]{data/clean/profiles/pxzz_605296.dat};
\end{axis}
\end{tikzpicture}

\captionof{figure}[神农集团（605296）股价与周期骨架]{神农集团（605296）月度收盘价（元/股）与 ZigZag 周期骨架（阈值 25\%，圆点为周期转折点）}
\end{center}

\pfl{过去周期分析} 以 25\% 的 ZigZag 阈值划分，样本区间内股价形成 2 个主升段与 2 个主跌段。 最大主跌段为 2021 年 5 月 至 2023 年 9 月，28 个月-63.2\%；最大主升段出现在 2023 年 9 月 至 2024 年 5 月，8 个月+93.5\%。 最近一次转折出现在 2025 年 8 月（上涨段自 2025 年 1 月起，7 个月+26.3\%），当前处于该段之后的震荡之中。 公司股价较申万农林牧渔指数提前约 32 个月见底，是板块中较早定价周期反转的公司之一。 生猪养殖的定价主线是产能去化与成本曲线的赛跑（第 \ref{sec:hog} 章）：2026 年三季度行业仍处在“产能去化未完成、价格磨底”的阶段，公司 2026 年 6 月末资产负债率 40.7\%、半年 ROE -17.9\%（未年化），在本轮下行中属于随行业同步调整的一类。

\pfl{盈利情况} 2026 年上半年营业收入 24.89 亿元（同比 -11.0\%），归母净利润 -8.09 亿元，由上年同期的盈利转为亏损。 以年报为趋势对照，2023---2025 年营业收入由 38.91 亿元变为 53.52 亿元（两年年均复合增速 +17.3\%），归母净利润由 -4.01 亿元变为 3.39 亿元。 按半年报口径，2026 年上半年的 ROE 为 -17.9\%（未年化）、销售净利率 -32.5\%、毛利率 0.9\%，ROE 在“生猪养殖”11 家公司中排第 \zhnumber{8} 位。 同期全样本半年 ROE 中位数 0.75\%，公司与之相差 18.64 个百分点（弱于中位数公司）。 结合公司披露，逆周期扩产：2024 年出栏 227 万头，2025 年目标出栏 320—350 万头，中期规划 202……；成本战略清晰推进：2023 年生猪养殖完全成本约 16.2 元/公斤，2024 年控制在 14 元/公斤……。\footnote{\texttt{https://finance.sina.com.cn/stock/observe/2025-03-11/doc-inephsqc9586225.shtml}}\footnote{\texttt{https://www.capco.org.cn/hyzl/zthd/shangshigszxd/202505/20250519/j\_202505191552250001747}}

\begin{center}\small\setlength{\tabcolsep}{4pt}
\captionof{table}[神农集团（605296）关键财务指标]{神农集团（605296）关键财务指标（合并报表口径；两个半年报期均未年化；现金短债比 $=$ 货币资金 $\div$ 短期债务）}
\begin{tabular}{@{}lrrrrr@{}}
\toprule
指标 & 2023 年 & 2024 年 & 2025 年 & 2025 年半年报 & 2026 年半年报 \\
\midrule
营业收入（亿元） & 38.91 & 55.84 & 53.52 & 27.98 & 24.89 \\
归母净利润（亿元） & -4.01 & 6.87 & 3.39 & 3.88 & -8.09 \\
毛利率（\%） & 4.4 & 20.8 & 16.3 & 21.9 & 0.9 \\
ROE（\%） & -9.1 & 15.1 & 6.9 & 7.8 & -17.9 \\
资产负债率（\%） & 26.4 & 27.4 & 29.5 & 26.0 & 40.7 \\
经营活动现金流净额（亿元） & 0.19 & 15.34 & 8.10 & 5.57 & -3.75 \\
货币资金（亿元） & 5.44 & 10.42 & 7.15 & 9.39 & 4.50 \\
有息负债合计（亿元） & 7.52 & 7.68 & 7.69 & 7.63 & 13.62 \\
现金短债比（倍） & 1.45 & 2.54 & 3.37 & 3.12 & 0.68 \\
\bottomrule
\end{tabular}
\end{center}

\begin{center}
\begin{tikzpicture}
\begin{groupplot}[
  group style={group size=2 by 1, horizontal sep=0.95cm},
  width=7.2cm, height=4.6cm,
  tick label style={font=\tiny},
  title style={font=\scriptsize\heiti},
  yticklabel style={font=\tiny},
  grid=major, grid style={LightGray, line width=0.3pt},
  axis line style={InkGray, line width=0.5pt},
  enlarge x limits={abs=0.8},
  legend style={font=\scriptsize, at={(0.5,-0.22)}, anchor=north,
                legend columns=2, draw=none, fill=none, column sep=6pt},
]
\nextgroupplot[
  ybar, bar width=4pt,
  ymin=-14.48, ymax=63.52,
  xtick={0,1,2,3,4,5.7},
  xticklabels={2021,2022,2023,2024,2025,2026H1},
  ylabel={亿元},
]
\addplot[fill=AgriGreen!75, draw=AgriGreen, bar shift=-2.6pt] table[x=i, y=rev]{data/clean/profiles/fin_605296.dat};
\addplot[fill=SoftRed!60, draw=SoftRed, bar shift=2.6pt] table[x=i, y=ni]{data/clean/profiles/fin_605296.dat};
\addplot[fill=AgriGreen!30, draw=AgriGreen, pattern=north east lines, bar shift=-2.6pt] table[x=x, y=rev]{data/clean/profiles/finh_605296.dat};
\addplot[fill=SoftRed!20, draw=SoftRed, pattern=north east lines, bar shift=2.6pt] table[x=x, y=ni]{data/clean/profiles/finh_605296.dat};
\legend{营业收入（年/半年）, 归母净利润（年/半年）}
\nextgroupplot[
  ymin=-22.6, ymax=46.6,
  xtick={0,1,2,3,4,5.7},
  xticklabels={2021,2022,2023,2024,2025,2026H1},
  ylabel={\%},
]
\addplot[AgriGold, mark=*, mark size=1.3pt] table[x=i, y=roe]{data/clean/profiles/fin_605296.dat};
\addplot[OilBlue, mark=square*, mark size=1.3pt] table[x=i, y=debt]{data/clean/profiles/fin_605296.dat};
\addplot[only marks, AgriGold, mark=*, mark size=2.0pt] table[x=x, y=roe]{data/clean/profiles/finh_605296.dat};
\addplot[only marks, OilBlue, mark=square*, mark size=2.0pt] table[x=x, y=debt]{data/clean/profiles/finh_605296.dat};
\legend{ROE, 资产负债率}
\end{groupplot}
\end{tikzpicture}
\begin{tikzpicture}
\begin{axis}[
  width=15.2cm, height=3.1cm, scale only axis,
  ymin=-5.66, ymax=17.63,
  xtick={0,1,2,3,4,5.7},
  xticklabels={2021,2022,2023,2024,2025,2026H1},
  tick label style={font=\tiny},
  yticklabel style={font=\tiny},
  ylabel={亿元}, ylabel style={font=\scriptsize},
  ybar, bar width=4pt,
  grid=major, grid style={LightGray, line width=0.3pt},
  axis line style={InkGray, line width=0.5pt},
  enlarge x limits={abs=0.8},
  legend style={font=\scriptsize, at={(0.5,-0.30)}, anchor=north,
                legend columns=3, draw=none, fill=none, column sep=10pt},
]
\addplot[fill=AgriGreen!55, draw=AgriGreen, bar shift=-3.0pt] table[x=i, y=cash]{data/clean/profiles/fin_605296.dat};
\addplot[fill=SoftRed!45, draw=SoftRed, bar shift=3.0pt] table[x=i, y=ibd]{data/clean/profiles/fin_605296.dat};
\addplot[fill=AgriGreen!25, draw=AgriGreen, pattern=north east lines, bar shift=-3.0pt] table[x=x, y=cash]{data/clean/profiles/finh_605296.dat};
\addplot[fill=SoftRed!20, draw=SoftRed, pattern=north east lines, bar shift=3.0pt] table[x=x, y=ibd]{data/clean/profiles/finh_605296.dat};
\addplot[OilBlue, mark=*, mark size=1.1pt, line width=0.8pt] table[x=i, y=ocf]{data/clean/profiles/fin_605296.dat};
\addplot[only marks, OilBlue, mark=*, mark size=1.8pt] table[x=x, y=ocf]{data/clean/profiles/finh_605296.dat};
\legend{货币资金, 有息负债, 经营活动现金流净额}
\end{axis}
\end{tikzpicture}

\captionof{figure}[神农集团（605296）近年主要财务指标]{神农集团（605296）近年主要财务指标（左上：营业收入与归母净利润；右上：ROE 与资产负债率；下：货币资金、有息负债与经营活动现金流净额。
2021---2025 年为年报口径，2026H1 为 2026 年半年报/半年末，与其前一年报之间留出间隔、以斜纹或空心点区分；半年报数据未年化）}
\end{center}

\pfl{财务分析} 2026 年 6 月末资产负债率 40.7\%，较 2025 年末 上升 11.2 个百分点（样本中位数 52.2\%，所处三级行业内排第 \zhnumber{1} 位）。 2026 年 6 月末货币资金 4.50 亿元、短期债务（短期借款与一年内到期非流动负债）6.59 亿元、长期债务 7.03 亿元，有息负债合计 13.62 亿元，现金短债比 0.68 倍——覆盖偏紧。 净有息负债 9.13 亿元，相当于其 2026 年上半年营业收入的 37\%。 流动比率 1.14、速动比率 0.37，短期指标尚可。 2026 年上半年经营活动现金流净额 -3.75 亿元，经营现金流与利润同时为负，现金消耗较快。 2026 年 6 月末在建工程 3.96 亿元，较上年末增加 2.53 亿元，仍有在建投入。

\pfl{重大战略转型} 近三年公司的战略动作集中在以下几个方面：2023-2025 成本战略清晰推进：2023 年生猪养殖完全成本约 16.2 元/公斤，2024 年控制在 14 元/公斤以内，2……；2024-2025 逆周期扩产：2024 年出栏 227 万头，2025 年目标出栏 320—350 万头，中期规划 2027 年出……；2025 2025 年度实现归属于上市公司股东的净利润 338,963,686.41 元，扭亏为盈；2025-11-06 实施 2025 年半年度权益分派，派发现金分红 204,657,892.92 元；2025 年度拟不再进行现金分……。 公告口径上，近三年（2023 年 9 月---2026 年 9 月）公司披露重大重组与并购类公告 1 条、对外投资与转型类公告 5 条、回购与股权激励类公告 78 条，合计公告 550 条。 主营结构的迁移已经落到报表上：2021 年年报的第一大产品为“饲料”（62.9\%），2026 年半年报变为“生猪养殖分部”（90.7\%）。 综合判断，公司在报告期内有明确的外延扩张或业务结构调整动作，转型方向与融资安排之间存在直接关联。\footnote{\texttt{https://money.finance.sina.com.cn/corp/view/vCB\_AllBulletinDetail.php}}\footnote{\texttt{https://www.capco.org.cn/hyzl/zthd/shangshigszxd/202505/20250519/j\_202505191552250001747}}\footnote{\texttt{https://finance.sina.com.cn/stock/observe/2025-03-11/doc-inephsqc9586225.shtml}}

\begin{center}\small\setlength{\tabcolsep}{4pt}
\captionof{table}[神农集团（605296）历史融资与资本运作]{神农集团（605296）历史融资与资本运作（据公司公告与公开报道整理；金额为披露值，含首次公开发行、再融资、债券、银行借款与重整投资等）}
\begin{tabular}{@{}p{1.5cm}p{3.2cm}p{2.4cm}p{6.3cm}@{}}
\toprule
时间 & 方式 & 规模 & 用途与说明 \\
\midrule
2021-05-28 & IPO（上交所主板） & 发行 4,003 万股，发行价 56.08 元/股，募集资金约 22 亿元 & 生猪养殖、饲料加工等项目；证监许可[2021]1314 号核准；主承销商中泰证券 \\
2023-06 & 定向增发（预案，后终止） & 拟募资不超过 11.5 亿元 & 两个饲料加工建设类项目及一个生猪养殖建设类项目；2024 年 4 月公告终止该次定增…… \\
2024-2025 & 以简易程序向特定对象发行股票（后终止） & 拟募资 2.9 亿元（1.37 亿元投红河弥勒五山乡年出栏 24 万头优质仔猪扩繁基地项目、2,500 万元投石林林口铺养殖场生物安全技改、4,333.06 万元投产业链数字化智能化项目、8,437.41 万元补充流动资金） & 产能扩张、技改与补充流动资金；2025 年 3 月披露预案修订稿…… \\
\bottomrule
\end{tabular}
\end{center}

\pfl{历史融投资情况} 从现金流量表看，2025 年公司取得借款收到的现金 3.80 亿元、偿还债务支付的现金 3.94 亿元、吸收权益性投资 1.18 亿元，筹资活动现金流净额 -3.64 亿元（2023 年为取得借款 6.23 亿元、偿还债务 0.32 亿元）。 2026 年上半年筹资活动现金流净额为 5.08 亿元（取得借款 6.59 亿元、偿还债务 0.76 亿元，未年化），仍处于净融入状态。 2025 年分配股利、利润或偿付利息支付的现金合计 2.14 亿元。 公开信息中与公司直接相关的融资与资本运作包括：IPO（上交所主板）、以简易程序向特定对象发行股票（后终止）、定向增发（预案，后终止）。 其中规模较大的两笔为：2024-2025以简易程序向特定对象发行股票（后终止）（拟募资 2.9 亿元（1.37 亿元投红河弥勒五山乡年出栏 ……）；2023-06定向增发（预案，后终止）（拟募资不超过 11.5 亿元）。 股本方面，近三年披露股本变动 10 次（激励股份解禁 3 次、股份回购 2 次、增发新股上市 1 次、股权激励 1 次）；总股本由 2021 年末的 4.00 亿股变为 5.25 亿股（+31.13\%）。 分红方面，近三年现金分红 2 次、累计每股派现 0.64 元（最近一次 2025-06-30）——在利润承压的阶段仍维持了分红。 近三年“再融资”类公告 27 条，涉及定向增发、募集资金使用。\footnote{\texttt{http://finance.ce.cn/stock/gsgdbd/202411/28/t20241128\_39218362.shtml}}\footnote{\texttt{https://finance.sina.com.cn/stock/observe/2025-03-11/doc-inephsqc9586225.shtml}}\footnote{\texttt{https://www.sse.com.cn/disclosure/listedinfo/announcement/c/new/2025-03-08/605296\_202503}}

\pfl{未来潜在融资需求分析} 2026 年 6 月末货币资金 4.50 亿元，而短期债务 6.59 亿元（现金短债比 0.68 倍），2026 年上半年经营现金流净额 -3.75 亿元。按“短期债务减去账面货币资金”的滚动偿付口径，静态缺口约 2.10 亿元；若再计入上半年亏损的年化额（16.18 亿元），维持一年正常经营所需的外部资金约 2.10---18.27 亿元。 公司 2026 年上半年归母净利润为负（-8.09 亿元），但 6 月末资产负债率 40.7\% 尚未进入第一档（亏损且负债率 $>$65\%）的区间，当前融资需求以补充流动资金、支撑低谷期经营性支出为主。 公开披露的下一步融资线索包括：2025-2026 2023 年度定增与 2024 年度简易程序定增均已终止；2025 年半年报分红已实施、2025 年度不再分红，2026……；2026-08 公司管理层在 2026 年半年度投资者关系活动中披露：上半年出栏商品猪 172.4 万头，7 月完全成本 12.0 元/……。 资金需求还要叠加在建工程：期末在建工程 3.96 亿元、2026 年上半年购建固定资产等支出 2.81 亿元（未年化），这部分支出在周期底部通常会被主动放缓。\footnote{\texttt{https://www.nbd.com.cn/articles/2025-07-28/3966061.html}}\footnote{\texttt{https://www.cnfin.com/announ/detail/index.html}}


\subsubsection{立华股份（300761）}
\label{co:300761}

\pfl{公司基本情况} 公司前身可追溯至 1997 年，全国第二大黄羽鸡养殖企业、江苏首家上市畜禽养殖企业，2019 年创业板上市，近年形成“黄羽鸡＋生猪”双轮驱动；2026 年上半年猪价深跌致整体由盈转亏，黄鸡板块盈利形成对冲。 公司于 2019-02-18 在深交所创业板首发上市，注册资本 8.46 亿元，总股本 8.28 亿股。 截至 2026-09-24收盘价 20.43 元，流通市值 132.1 亿元，近一年-10.0\%，流通市值在三级行业“肉鸡养殖”的 8 家公司中排第 \zhnumber{2} 位，近一年涨跌幅高于全样本中位数（-10.1\%）0.0 个百分点。 按 2025 年年报披露的行业构成，收入占比前三位为养殖行业 100.0\%。 发展过程中的关键节点包括：2 月 18 日在深交所创业板上市，发行 4,128 万股（全部为新股），发行价 29.……（2019）；8 月完成向特定对象发行股票，新增股份 8 月 31 日上市，实控人程立力全额认购（2022）；生猪业务放量：2025 年生猪出栏 211.16 万头（同比 +62.88\%），养猪业务……（2024）。\footnote{\texttt{https://money.finance.sina.com.cn/corp/view/vCB\_AllBulletinDetail.php}}\footnote{\texttt{https://vip.stock.finance.sina.com.cn/corp/view/vISSUE\_MarketBulletinDetail.php}}\footnote{\texttt{http://money.finance.sina.com.cn/corp/view/vCB\_AllBulletinDetail.php}}

\pfl{历史股价} 自 2021 年 1 月至 2026 年 9 月，股价由 33.08 元变为 20.43 元，累计 -38.2\%，区间最高 45.81 元（2022 年 8 月）、最低 16.41 元（2026 年 6 月）；当前价相当于区间最高价的 44.6\%，为区间最低价的 1.24 倍。 月度收益的年化波动率 44.1\%，高于全样本中位数 37.4\%，在三级行业“肉鸡养殖”的 8 家公司中波动率排第 \zhnumber{2} 位。 从事件看，2022 年 4 月 至 2022 年 8 月的上涨（+63.9\%）与公司实际控制人程立力全额认购 12.7 亿元定增，募集资金主要投向黄羽鸡一体化养……（2022-08）在时间上重合；2022 年 8 月 至 2023 年 10 月的下跌（-63.0\%）与公司实际控制人程立力全额认购 12.7 亿元定增，募集资金主要投向黄羽鸡一体化养……（2022-08）在时间上重合——股价的拐点多数能在公司自身的资本运作与风险事件上找到对应。\footnote{\texttt{http://www.nbd.com.cn/rss/toutiao/articles/2476589.html}}

\begin{center}
\begin{tikzpicture}
\begin{axis}[
  width=12.6cm, height=3.9cm, scale only axis,
  xmin=2020.922, xmax=2026.888, ymin=15.26, ymax=49.02,
  xtick={2021,2022,2023,2024,2025,2026},
  yearticks,
  yticklabel style={font=\scriptsize},
  ylabel={元/股}, ylabel style={font=\scriptsize},
  grid=major, grid style={LightGray, line width=0.3pt},
  axis line style={InkGray, line width=0.5pt},
  every axis plot/.append style={line width=0.9pt},
]
\addplot[AgriGreen] table[x=x,y=y]{data/clean/profiles/px_300761.dat};
\addplot[SoftRed, dashed, line width=0.7pt] table[x=x,y=y]{data/clean/profiles/pxzz_300761.dat};
\addplot[only marks, mark=*, mark size=1.1pt, SoftRed] table[x=x,y=y]{data/clean/profiles/pxzz_300761.dat};
\end{axis}
\end{tikzpicture}

\captionof{figure}[立华股份（300761）股价与周期骨架]{立华股份（300761）月度收盘价（元/股）与 ZigZag 周期骨架（阈值 25\%，圆点为周期转折点）}
\end{center}

\pfl{过去周期分析} 以 25\% 的 ZigZag 阈值划分，样本区间内股价形成 4 个主升段与 5 个主跌段。 最大主升段出现在 2022 年 4 月 至 2022 年 8 月，4 个月+63.9\%；最大主跌段为 2022 年 8 月 至 2023 年 10 月，14 个月-63.0\%。 最近一次转折出现在 2026 年 6 月（下跌段自 2026 年 4 月起，2 个月-27.9\%），当前处于该段的延续之中。 公司股价与申万农林牧渔指数几乎同步见底，个股并未走出独立行情。 禽链的产能周期只有 6---12 个月、反转快于生猪（第 \ref{sec:other-livestock} 节），2026 年前四个月毛鸡养殖利润同比大幅回升，链条利润的修复先于生猪，公司 2026 年 6 月末资产负债率 43.8\%、半年 ROE -0.8\%（未年化），在本轮下行中属于随行业同步调整的一类。

\pfl{盈利情况} 2026 年上半年营业收入 91.28 亿元（同比 +9.3\%），归母净利润 -0.72 亿元，由上年同期的盈利转为亏损。 以年报为趋势对照，2023---2025 年营业收入由 153.54 亿元变为 187.29 亿元（两年年均复合增速 +10.4\%），归母净利润由 -4.37 亿元变为 5.78 亿元。 按半年报口径，2026 年上半年的 ROE 为 -0.8\%（未年化）、销售净利率 -0.8\%、毛利率 8.0\%，ROE 在“肉鸡养殖”8 家公司中排第 \zhnumber{7} 位。 行业同口径半年 ROE 的中位数为 0.75\%，该公司低于中位数 1.52 个百分点。 结合公司披露，2026 年第二季度单季归母净利润 -2.27 亿元（上年同期 -0.57 亿元），业绩同比下滑主要受畜……；2026 年上半年营业收入 91.28 亿元（同比 +9.27\%，创上市以来同期新高），归母净利润 -7……。\footnote{\texttt{https://pdf.dfcfw.com/pdf/H3\_AP202608241828367091\_1.pdf}}\footnote{\texttt{https://finance.sina.com.cn/roll/2026-08-21/doc-inipazsh1222235.shtml}}

\begin{center}\small\setlength{\tabcolsep}{4pt}
\captionof{table}[立华股份（300761）关键财务指标]{立华股份（300761）关键财务指标（合并报表口径；两个半年报期均未年化；现金短债比 $=$ 货币资金 $\div$ 短期债务）}
\begin{tabular}{@{}lrrrrr@{}}
\toprule
指标 & 2023 年 & 2024 年 & 2025 年 & 2025 年半年报 & 2026 年半年报 \\
\midrule
营业收入（亿元） & 153.54 & 177.25 & 187.29 & 83.53 & 91.28 \\
归母净利润（亿元） & -4.37 & 15.21 & 5.78 & 1.49 & -0.72 \\
毛利率（\%） & 4.8 & 15.9 & 11.0 & 10.0 & 8.0 \\
ROE（\%） & -5.4 & 18.0 & 6.3 & 1.6 & -0.8 \\
资产负债率（\%） & 43.0 & 37.8 & 42.0 & 41.1 & 43.8 \\
经营活动现金流净额（亿元） & 14.50 & 24.08 & 18.48 & 5.59 & 5.21 \\
货币资金（亿元） & 3.66 & 1.83 & 2.64 & 4.61 & 3.67 \\
有息负债合计（亿元） & 28.50 & 22.31 & 32.04 & 29.53 & 35.68 \\
现金短债比（倍） & 0.17 & 0.11 & 0.10 & 0.19 & 0.12 \\
\bottomrule
\end{tabular}
\end{center}

\begin{center}
\begin{tikzpicture}
\begin{groupplot}[
  group style={group size=2 by 1, horizontal sep=0.95cm},
  width=7.2cm, height=4.6cm,
  tick label style={font=\tiny},
  title style={font=\scriptsize\heiti},
  yticklabel style={font=\tiny},
  grid=major, grid style={LightGray, line width=0.3pt},
  axis line style={InkGray, line width=0.5pt},
  enlarge x limits={abs=0.8},
  legend style={font=\scriptsize, at={(0.5,-0.22)}, anchor=north,
                legend columns=2, draw=none, fill=none, column sep=6pt},
]
\nextgroupplot[
  ybar, bar width=4pt,
  ymin=-23.54, ymax=210.29,
  xtick={0,1,2,3,4,5.7},
  xticklabels={2021,2022,2023,2024,2025,2026H1},
  ylabel={亿元},
]
\addplot[fill=AgriGreen!75, draw=AgriGreen, bar shift=-2.6pt] table[x=i, y=rev]{data/clean/profiles/fin_300761.dat};
\addplot[fill=SoftRed!60, draw=SoftRed, bar shift=2.6pt] table[x=i, y=ni]{data/clean/profiles/fin_300761.dat};
\addplot[fill=AgriGreen!30, draw=AgriGreen, pattern=north east lines, bar shift=-2.6pt] table[x=x, y=rev]{data/clean/profiles/finh_300761.dat};
\addplot[fill=SoftRed!20, draw=SoftRed, pattern=north east lines, bar shift=2.6pt] table[x=x, y=ni]{data/clean/profiles/finh_300761.dat};
\legend{营业收入（年/半年）, 归母净利润（年/半年）}
\nextgroupplot[
  ymin=-10.1, ymax=48.8,
  xtick={0,1,2,3,4,5.7},
  xticklabels={2021,2022,2023,2024,2025,2026H1},
  ylabel={\%},
]
\addplot[AgriGold, mark=*, mark size=1.3pt] table[x=i, y=roe]{data/clean/profiles/fin_300761.dat};
\addplot[OilBlue, mark=square*, mark size=1.3pt] table[x=i, y=debt]{data/clean/profiles/fin_300761.dat};
\addplot[only marks, AgriGold, mark=*, mark size=2.0pt] table[x=x, y=roe]{data/clean/profiles/finh_300761.dat};
\addplot[only marks, OilBlue, mark=square*, mark size=2.0pt] table[x=x, y=debt]{data/clean/profiles/finh_300761.dat};
\legend{ROE, 资产负债率}
\end{groupplot}
\end{tikzpicture}
\begin{tikzpicture}
\begin{axis}[
  width=15.2cm, height=3.1cm, scale only axis,
  ymin=-3.57, ymax=39.96,
  xtick={0,1,2,3,4,5.7},
  xticklabels={2021,2022,2023,2024,2025,2026H1},
  tick label style={font=\tiny},
  yticklabel style={font=\tiny},
  ylabel={亿元}, ylabel style={font=\scriptsize},
  ybar, bar width=4pt,
  grid=major, grid style={LightGray, line width=0.3pt},
  axis line style={InkGray, line width=0.5pt},
  enlarge x limits={abs=0.8},
  legend style={font=\scriptsize, at={(0.5,-0.30)}, anchor=north,
                legend columns=3, draw=none, fill=none, column sep=10pt},
]
\addplot[fill=AgriGreen!55, draw=AgriGreen, bar shift=-3.0pt] table[x=i, y=cash]{data/clean/profiles/fin_300761.dat};
\addplot[fill=SoftRed!45, draw=SoftRed, bar shift=3.0pt] table[x=i, y=ibd]{data/clean/profiles/fin_300761.dat};
\addplot[fill=AgriGreen!25, draw=AgriGreen, pattern=north east lines, bar shift=-3.0pt] table[x=x, y=cash]{data/clean/profiles/finh_300761.dat};
\addplot[fill=SoftRed!20, draw=SoftRed, pattern=north east lines, bar shift=3.0pt] table[x=x, y=ibd]{data/clean/profiles/finh_300761.dat};
\addplot[OilBlue, mark=*, mark size=1.1pt, line width=0.8pt] table[x=i, y=ocf]{data/clean/profiles/fin_300761.dat};
\addplot[only marks, OilBlue, mark=*, mark size=1.8pt] table[x=x, y=ocf]{data/clean/profiles/finh_300761.dat};
\legend{货币资金, 有息负债, 经营活动现金流净额}
\end{axis}
\end{tikzpicture}

\captionof{figure}[立华股份（300761）近年主要财务指标]{立华股份（300761）近年主要财务指标（左上：营业收入与归母净利润；右上：ROE 与资产负债率；下：货币资金、有息负债与经营活动现金流净额。
2021---2025 年为年报口径，2026H1 为 2026 年半年报/半年末，与其前一年报之间留出间隔、以斜纹或空心点区分；半年报数据未年化）}
\end{center}

\pfl{财务分析} 2026 年 6 月末资产负债率 43.8\%，较 2025 年末 上升 1.8 个百分点（样本中位数 52.2\%，所处三级行业内排第 \zhnumber{3} 位）。 2026 年 6 月末货币资金 3.67 亿元、短期债务（短期借款与一年内到期非流动负债）29.94 亿元、长期债务 5.74 亿元，有息负债合计 35.68 亿元，现金短债比 0.12 倍——缺口明显。 净有息负债 32.00 亿元，相当于其 2026 年上半年营业收入的 35\%。 流动比率 0.76、速动比率 0.24，短期偿付压力较大。 2026 年上半年经营活动现金流净额 5.21 亿元，净利润为负而经营现金流为正，差额主要来自折旧摊销与营运资本变动，说明亏损尚未完全转化为现金流出。 2026 年 6 月末在建工程 2.79 亿元，较上年末增加 0.93 亿元，仍有在建投入。

\pfl{重大战略转型} 公司在报告期内的战略动作可归为 4 项：2022-08 实际控制人程立力全额认购 12.7 亿元定增，募集资金主要投向黄羽鸡一体化养殖及屠宰、孵化等项目，“all in……；2025 全年营业收入 187.29 亿元（同比 +5.67\%），归母净利润 5.78 亿元（同比 -62.01\%）；肉猪……；2026 食品端延伸与智能化转型提速，鸡猪双主业形成上下半年盈利互补的抗周期结构；2026-05 公司提出养猪成本管控进入行业第一方阵，2025 年肉猪完全成本持续下降，黄鸡业务保持持续盈利。 公告口径上，近三年（2023 年 9 月---2026 年 9 月）公司披露重大重组与并购类公告 0 条、对外投资与转型类公告 2 条、回购与股权激励类公告 47 条，合计公告 453 条。 第一大行业仍是“养殖行业”，收入占比 100.0\%，与 2021 年年报相比没有方向性变化。 综合判断，公司在报告期内有明确的外延扩张或业务结构调整动作，转型方向与融资安排之间存在直接关联。\footnote{\texttt{http://www.nbd.com.cn/rss/toutiao/articles/2476589.html}}\footnote{\texttt{http://money.finance.sina.com.cn/corp/view/vCB\_AllBulletinDetail.php}}\footnote{\texttt{https://stock.stockstar.com/RB2026082600000809.shtml}}

\begin{center}\small\setlength{\tabcolsep}{4pt}
\captionof{table}[立华股份（300761）历史融资与资本运作]{立华股份（300761）历史融资与资本运作（据公司公告与公开报道整理；金额为披露值，含首次公开发行、再融资、债券、银行借款与重整投资等）}
\begin{tabular}{@{}p{1.5cm}p{3.2cm}p{2.4cm}p{6.3cm}@{}}
\toprule
时间 & 方式 & 规模 & 用途与说明 \\
\midrule
2019-02-18 & IPO（深交所创业板） & 4,128 万股，发行价 29.35 元/股，募集资金总额 12.1157 亿元、净额 11.50 亿元 & 安庆立华年出栏 1,750 万羽一体化养鸡建设项目…… \\
2021-09 & 可转换公司债券（方案） & 拟向不特定对象发行可转债募集资金不超过 204,480.94 万元 & 黄羽鸡一体化养殖、屠宰等项目；公司第三届董事会第二次会议通过《关于调整公司向不特定对象发…… \\
2022-08 & 定向增发（向特定对象发行股票） & 55,922,501 股，发行价 22.71 元/股，募集资金总额 12.70 亿元、净额 12.6064 亿元 & 江苏、湖南等地养鸡及屠宰项目、总部基地项目及补充流动资金（项目总投资 27.14 亿元）…… \\
\bottomrule
\end{tabular}
\end{center}

\pfl{历史融投资情况} 从现金流量表看，2025 年公司取得借款收到的现金 34.06 亿元、偿还债务支付的现金 25.05 亿元、吸收权益性投资 1.07 亿元，筹资活动现金流净额 3.94 亿元（2023 年为取得借款 33.53 亿元、偿还债务 28.39 亿元）。 2026 年上半年筹资活动现金流净额为 1.33 亿元（取得借款 20.13 亿元、偿还债务 16.88 亿元，未年化），仍处于净融入状态。 2025 年分配股利、利润或偿付利息支付的现金合计 5.56 亿元。 公司披露的股本与债务融资脉络包括：IPO（深交所创业板）、可转换公司债券（方案）、定向增发（向特定对象发行股票）。 其中规模较大的两笔为：2022-08定向增发（向特定对象发行股票）（55,922,501 股，发行价 22.71 元/股，募集资……）；2021-09可转换公司债券（方案）（拟向不特定对象发行可转债募集资金不超过 204,480.94……）。 股本方面，近三年披露股本变动 5 次（转增 1 次、限售股份上市 1 次）；总股本由 2021 年末的 4.04 亿股变为 8.28 亿股（+104.92\%）。 分红方面，近三年现金分红 5 次、累计每股派现 1.95 元（最近一次 2026-05-25）——在利润承压的阶段仍维持了分红。 近三年“再融资”类公告 3 条，涉及定向增发。\footnote{\texttt{https://money.finance.sina.com.cn/corp/view/vCB\_AllBulletinDetail.php}}\footnote{\texttt{http://static.cninfo.com.cn/finalpage/2021-09-22/1211127566.PDF}}\footnote{\texttt{https://vip.stock.finance.sina.com.cn/corp/view/vISSUE\_MarketBulletinDetail.php}}

\pfl{未来潜在融资需求分析} 2026 年 6 月末货币资金 3.67 亿元，而短期债务 29.94 亿元（现金短债比 0.12 倍），2026 年上半年经营现金流净额 5.21 亿元。按“短期债务减去账面货币资金”的滚动偿付口径，静态缺口约 26.26 亿元；若再计入上半年亏损的年化额（1.44 亿元），维持一年正常经营所需的外部资金约 26.26---27.70 亿元。 公司 2026 年上半年归母净利润为负（-0.72 亿元），但 6 月末资产负债率 43.8\% 尚未进入第一档（亏损且负债率 $>$65\%）的区间，当前融资需求以补充流动资金、支撑低谷期经营性支出为主。 公开披露的下一步融资线索包括：2021-2022 2021 年 9 月调整后的 20.45 亿元可转债方案未查得实施结果，2022 年改以 12.7 亿元定增完成融资（实……；2026-08 2026 年 4 月 22 日投资者关系活动披露：2025 年营收 187.29 亿元、归母净利润 5.78 亿元；20……。 资金需求还要叠加在建工程：期末在建工程 2.79 亿元、2026 年上半年购建固定资产等支出 8.32 亿元（未年化），这部分支出在周期底部通常会被主动放缓。\footnote{\texttt{http://static.cninfo.com.cn/finalpage/2021-09-22/1211127566.PDF}}\footnote{\texttt{https://vip.stock.finance.sina.com.cn/corp/view/vCB\_AllBulletinDetail.php}}


\subsubsection{京基智农（000048）}
\label{co:000048}

\pfl{公司基本情况} 1979 年设立至今，京基智农前身为 1994 年上市的“中国农牧第一股”康达尔，2018—2019 年被京基集团要约收购入主后更名，主业转型为广东、广西区域生猪养殖，房地产转为存量去化业务；2026 年上半年猪价低迷导致大额亏损。 公司于 1994-11-01 在深交所主板首发上市，注册资本 5.30 亿元，总股本 5.31 亿股。 截至 2026-09-24收盘价 18.66 元，流通市值 98.8 亿元，近一年+10.5\%，流通市值在三级行业“生猪养殖”的 11 家公司中排第 \zhnumber{5} 位，近一年涨跌幅高于全样本中位数（-10.1\%）20.6 个百分点。 按 2026 年半年报披露的产品构成，收入占比前三位为猪业产品 71.7\%；饲料产品 14.9\%；商品房 7.1\%。 与 2021 年年报相比，第一大产品由“商品房”变为“猪业产品”，主业重心已经迁移。 发展过程中的关键节点包括：11 月 1 日在深圳证券交易所挂牌上市（股票代码 000048），成为中国首家农牧上市……（1994）；12 月 8 日更名为深圳市中科创业（集团）股份有限公司（后恢复康达尔名称）（1999）；11 月 23 日京基集团有限公司完成对公司的首次要约收购，成为控股股东；2019 年 ……（2018）。\footnote{\texttt{https://finance.sina.com.cn/stock/s/2018-12-14/doc-ihmutuec9248945.shtml}}\footnote{\texttt{http://file.finance.sina.com.cn/211.154.219.97:9494/MRGG/CNSESZ\_STOCK/2013/2013-8/2013-0}}\footnote{\texttt{http://vip.stock.finance.sina.com.cn/corp/view/vCB\_AllBulletinDetail.php}}

\pfl{历史股价} 自 2021 年 1 月至 2026 年 9 月，股价由 22.79 元变为 18.66 元，累计 -18.1\%，区间最高 29.29 元（2021 年 4 月）、最低 12.54 元（2026 年 5 月）；当前价相当于区间最高价的 63.7\%，为区间最低价的 1.49 倍。 月度收益的年化波动率 51.2\%，高于全样本中位数 37.4\%，在三级行业“生猪养殖”的 11 家公司中波动率排第 \zhnumber{1} 位。 从事件看，2025 年 11 月 至 2026 年 3 月的上涨（+75.4\%）与公司董事会审议通过《关于股权投资的议案》相关意向，拟投资 6.5 亿元对江苏汇博……（2025-12-30）在时间上重合；2026 年 5 月 至 2026 年 9 月的上涨（+48.8\%）与公司控股股东京基集团、京基时代分别将其持有的公司股份 921,919 股（占总股……（2026-08-28）在时间上重合——股价的拐点多数能在公司自身的资本运作与风险事件上找到对应。\footnote{\texttt{https://money.finance.sina.com.cn/corp/view/vCB\_AllBulletinDetail.php}}\footnote{\texttt{http://epaper.zqrb.cn/images/2026-09/01/C3/zqrb20260901C3.pdf}}

\begin{center}
\begin{tikzpicture}
\begin{axis}[
  width=12.6cm, height=3.9cm, scale only axis,
  xmin=2020.922, xmax=2026.888, ymin=11.66, ymax=31.34,
  xtick={2021,2022,2023,2024,2025,2026},
  yearticks,
  yticklabel style={font=\scriptsize},
  ylabel={元/股}, ylabel style={font=\scriptsize},
  grid=major, grid style={LightGray, line width=0.3pt},
  axis line style={InkGray, line width=0.5pt},
  every axis plot/.append style={line width=0.9pt},
]
\addplot[AgriGreen] table[x=x,y=y]{data/clean/profiles/px_000048.dat};
\addplot[SoftRed, dashed, line width=0.7pt] table[x=x,y=y]{data/clean/profiles/pxzz_000048.dat};
\addplot[only marks, mark=*, mark size=1.1pt, SoftRed] table[x=x,y=y]{data/clean/profiles/pxzz_000048.dat};
\end{axis}
\end{tikzpicture}

\captionof{figure}[京基智农（000048）股价与周期骨架]{京基智农（000048）月度收盘价（元/股）与 ZigZag 周期骨架（阈值 25\%，圆点为周期转折点）}
\end{center}

\pfl{过去周期分析} 以 25\% 的 ZigZag 阈值划分，样本区间内股价形成 6 个主升段与 5 个主跌段。 最大主升段出现在 2025 年 11 月 至 2026 年 3 月，4 个月+75.4\%；最大主跌段为 2026 年 3 月 至 2026 年 5 月，2 个月-47.0\%。 最近一次转折出现在 2026 年 9 月（上涨段自 2026 年 5 月起，4 个月+48.8\%），当前处于该段之后的震荡之中。 公司股价与申万农林牧渔指数几乎同步见底，个股并未走出独立行情。 生猪养殖的定价主线是产能去化与成本曲线的赛跑（第 \ref{sec:hog} 章）：2026 年三季度行业仍处在“产能去化未完成、价格磨底”的阶段，公司 2026 年 6 月末资产负债率 63.6\%、半年 ROE -12.9\%（未年化），在本轮下行中属于随行业同步调整的一类。

\pfl{盈利情况} 2026 年上半年营业收入 21.30 亿元（同比 -10.7\%），归母净利润 -4.89 亿元，由上年同期的盈利转为亏损。 以年报为趋势对照，2023---2025 年营业收入由 124.17 亿元变为 48.72 亿元（两年年均复合增速 -37.4\%），归母净利润由 17.46 亿元变为 1.53 亿元。 按半年报口径，2026 年上半年的 ROE 为 -12.9\%（未年化）、销售净利率 -23.1\%、毛利率 -8.1\%，ROE 在“生猪养殖”11 家公司中排第 \zhnumber{5} 位。 同期全样本半年 ROE 中位数 0.75\%，公司与之相差 13.63 个百分点（弱于中位数公司）。 结合公司披露，2025 年归母净利润 1.53 亿元，同比下降 78.59\%；扣非归母净利润 1.54 亿元，同比下降……；推进精细化管理与健康养殖，2026 年完全成本目标 12 元/公斤左右；2025 年 1—9 月完全成本……。\footnote{\texttt{https://finance.sina.com.cn/stock/aiassist/yjbg/2026-04-09/doc-inhtxkaq1831111.shtml}}\footnote{\texttt{https://www.cnfin.com/announ/detail/index.html}}

\begin{center}\small\setlength{\tabcolsep}{4pt}
\captionof{table}[京基智农（000048）关键财务指标]{京基智农（000048）关键财务指标（合并报表口径；两个半年报期均未年化；现金短债比 $=$ 货币资金 $\div$ 短期债务）}
\begin{tabular}{@{}lrrrrr@{}}
\toprule
指标 & 2023 年 & 2024 年 & 2025 年 & 2025 年半年报 & 2026 年半年报 \\
\midrule
营业收入（亿元） & 124.17 & 59.60 & 48.72 & 23.87 & 21.30 \\
归母净利润（亿元） & 17.46 & 7.14 & 1.53 & 2.26 & -4.89 \\
毛利率（\%） & 38.7 & 28.0 & 13.8 & 19.3 & -8.1 \\
ROE（\%） & 55.5 & 17.4 & 3.6 & 5.2 & -12.9 \\
资产负债率（\%） & 67.8 & 60.3 & 59.0 & 58.3 & 63.6 \\
经营活动现金流净额（亿元） & -4.15 & 9.97 & 5.97 & 3.11 & 0.41 \\
货币资金（亿元） & 11.07 & 8.74 & 4.97 & 6.93 & 6.08 \\
有息负债合计（亿元） & 26.40 & 21.19 & 17.01 & 18.77 & 19.75 \\
现金短债比（倍） & 1.04 & 0.77 & 0.44 & 0.57 & 0.46 \\
\bottomrule
\end{tabular}
\end{center}

\begin{center}
\begin{tikzpicture}
\begin{groupplot}[
  group style={group size=2 by 1, horizontal sep=0.95cm},
  width=7.2cm, height=4.6cm,
  tick label style={font=\tiny},
  title style={font=\scriptsize\heiti},
  yticklabel style={font=\tiny},
  grid=major, grid style={LightGray, line width=0.3pt},
  axis line style={InkGray, line width=0.5pt},
  enlarge x limits={abs=0.8},
  legend style={font=\scriptsize, at={(0.5,-0.22)}, anchor=north,
                legend columns=2, draw=none, fill=none, column sep=6pt},
]
\nextgroupplot[
  ybar, bar width=4pt,
  ymin=-17.80, ymax=139.66,
  xtick={0,1,2,3,4,5.7},
  xticklabels={2021,2022,2023,2024,2025,2026H1},
  ylabel={亿元},
]
\addplot[fill=AgriGreen!75, draw=AgriGreen, bar shift=-2.6pt] table[x=i, y=rev]{data/clean/profiles/fin_000048.dat};
\addplot[fill=SoftRed!60, draw=SoftRed, bar shift=2.6pt] table[x=i, y=ni]{data/clean/profiles/fin_000048.dat};
\addplot[fill=AgriGreen!30, draw=AgriGreen, pattern=north east lines, bar shift=-2.6pt] table[x=x, y=rev]{data/clean/profiles/finh_000048.dat};
\addplot[fill=SoftRed!20, draw=SoftRed, pattern=north east lines, bar shift=2.6pt] table[x=x, y=ni]{data/clean/profiles/finh_000048.dat};
\legend{营业收入（年/半年）, 归母净利润（年/半年）}
\nextgroupplot[
  ymin=-20.7, ymax=94.6,
  xtick={0,1,2,3,4,5.7},
  xticklabels={2021,2022,2023,2024,2025,2026H1},
  ylabel={\%},
]
\addplot[AgriGold, mark=*, mark size=1.3pt] table[x=i, y=roe]{data/clean/profiles/fin_000048.dat};
\addplot[OilBlue, mark=square*, mark size=1.3pt] table[x=i, y=debt]{data/clean/profiles/fin_000048.dat};
\addplot[only marks, AgriGold, mark=*, mark size=2.0pt] table[x=x, y=roe]{data/clean/profiles/finh_000048.dat};
\addplot[only marks, OilBlue, mark=square*, mark size=2.0pt] table[x=x, y=debt]{data/clean/profiles/finh_000048.dat};
\legend{ROE, 资产负债率}
\end{groupplot}
\end{tikzpicture}
\begin{tikzpicture}
\begin{axis}[
  width=15.2cm, height=3.1cm, scale only axis,
  ymin=-8.09, ymax=40.00,
  xtick={0,1,2,3,4,5.7},
  xticklabels={2021,2022,2023,2024,2025,2026H1},
  tick label style={font=\tiny},
  yticklabel style={font=\tiny},
  ylabel={亿元}, ylabel style={font=\scriptsize},
  ybar, bar width=4pt,
  grid=major, grid style={LightGray, line width=0.3pt},
  axis line style={InkGray, line width=0.5pt},
  enlarge x limits={abs=0.8},
  legend style={font=\scriptsize, at={(0.5,-0.30)}, anchor=north,
                legend columns=3, draw=none, fill=none, column sep=10pt},
]
\addplot[fill=AgriGreen!55, draw=AgriGreen, bar shift=-3.0pt] table[x=i, y=cash]{data/clean/profiles/fin_000048.dat};
\addplot[fill=SoftRed!45, draw=SoftRed, bar shift=3.0pt] table[x=i, y=ibd]{data/clean/profiles/fin_000048.dat};
\addplot[fill=AgriGreen!25, draw=AgriGreen, pattern=north east lines, bar shift=-3.0pt] table[x=x, y=cash]{data/clean/profiles/finh_000048.dat};
\addplot[fill=SoftRed!20, draw=SoftRed, pattern=north east lines, bar shift=3.0pt] table[x=x, y=ibd]{data/clean/profiles/finh_000048.dat};
\addplot[OilBlue, mark=*, mark size=1.1pt, line width=0.8pt] table[x=i, y=ocf]{data/clean/profiles/fin_000048.dat};
\addplot[only marks, OilBlue, mark=*, mark size=1.8pt] table[x=x, y=ocf]{data/clean/profiles/finh_000048.dat};
\legend{货币资金, 有息负债, 经营活动现金流净额}
\end{axis}
\end{tikzpicture}

\captionof{figure}[京基智农（000048）近年主要财务指标]{京基智农（000048）近年主要财务指标（左上：营业收入与归母净利润；右上：ROE 与资产负债率；下：货币资金、有息负债与经营活动现金流净额。
2021---2025 年为年报口径，2026H1 为 2026 年半年报/半年末，与其前一年报之间留出间隔、以斜纹或空心点区分；半年报数据未年化）}
\end{center}

\pfl{财务分析} 2026 年 6 月末资产负债率 63.6\%，较 2025 年末 上升 4.7 个百分点（样本中位数 52.2\%，所处三级行业内排第 \zhnumber{7} 位）。 2026 年 6 月末货币资金 6.08 亿元、短期债务（短期借款与一年内到期非流动负债）13.29 亿元、长期债务 6.46 亿元，有息负债合计 19.75 亿元，现金短债比 0.46 倍——覆盖偏紧。 净有息负债 13.68 亿元，相当于其 2026 年上半年营业收入的 64\%。 流动比率 0.40、速动比率 0.17，短期偿付压力较大。 2026 年上半年经营活动现金流净额 0.41 亿元，净利润为负而经营现金流为正，差额主要来自折旧摊销与营运资本变动，说明亏损尚未完全转化为现金流出。 2026 年 6 月末在建工程 1.32 亿元，较上年末增加 0.05 亿元，仍有在建投入。 公司披露的财务与合规风险包括：控股股东京基集团、京基时代分别将其持有的公司股份 921,919 股（占总股本 0.17……。\footnote{\texttt{http://epaper.zqrb.cn/images/2026-09/01/C3/zqrb20260901C3.pdf}}

\pfl{重大战略转型} 近三年公司的战略动作集中在以下几个方面：2025 推进精细化管理与健康养殖，2026 年完全成本目标 12 元/公斤左右；2025 年 1—9 月完全成本约 13……；2025-12-30 董事会审议通过《关于股权投资的议案》相关意向，拟投资 6.5 亿元对江苏汇博机器人技术股份有限公司增资扩股并取得……；2025-2026 以集中竞价交易方式回购公司股份，截至相关时点累计回购 5,438,000 股（占总股本 1.03\%），最高成交价……；2026-06-29 调整机器人领域投资方案：终止原 6.5 亿元控股型《增资协议》，改为出资 6,000 万元认购江苏汇博新增注册资……。 公告口径上，近三年（2023 年 9 月---2026 年 9 月）公司披露重大重组与并购类公告 1 条、对外投资与转型类公告 0 条、回购与股权激励类公告 61 条，合计公告 460 条。 主营结构的迁移已经落到报表上：2021 年年报的第一大产品为“商品房”（68.8\%），2026 年半年报变为“猪业产品”（71.7\%）。 综合判断，公司在报告期内有明确的外延扩张或业务结构调整动作，转型方向与融资安排之间存在直接关联。\footnote{\texttt{https://disc.static.szse.cn/download/disc/disk02/finalpage/2019-12-05/16f99cde-9710-49c2}}\footnote{\texttt{https://m.21jingji.com/article/20241115/229b667a980d35806b80c2bf7261051a.html}}\footnote{\texttt{https://www.cnfin.com/announ/detail/index.html}}

\pfl{历史融投资情况} 从现金流量表看，2025 年公司取得借款收到的现金 8.22 亿元、偿还债务支付的现金 11.92 亿元、吸收权益性投资 0.07 亿元，筹资活动现金流净额 -8.56 亿元（2023 年为取得借款 13.80 亿元、偿还债务 5.13 亿元）。 2026 年上半年筹资活动现金流净额为 -0.04 亿元（取得借款 9.33 亿元、偿还债务 7.23 亿元，未年化），已转为净流出，处于净还债状态。 2025 年分配股利、利润或偿付利息支付的现金合计 4.57 亿元。 公司披露的股本与债务融资脉络包括：定向增发（预案，发行对象含控股股东）、股权融资（历史合计规模很小）。 其中规模较大的两笔为：2024-11-13定向增发（预案，发行对象含控股股东）（拟募资不超过 19.3 亿元；拟发行股票数量不超过 159,……）；2019-2024股权融资（历史合计规模很小）（上市三十年股权融资总额不过 2.7 亿元）。 股本方面，近三年披露股本变动 4 次（增发新股上市 1 次、股权激励 1 次）；总股本由 2021 年末的 5.23 亿股变为 5.31 亿股（+1.40\%）。 分红方面，近三年现金分红 4 次、累计每股派现 2.11 元（最近一次 2025-06-30）——在利润承压的阶段仍维持了分红。 近三年“再融资”类公告 9 条，涉及定向增发、募集资金使用。\footnote{\texttt{https://m.21jingji.com/article/20241115/229b667a980d35806b80c2bf7261051a.html}}\footnote{\texttt{https://www.stcn.com/article/detail/1411786.html}}

\pfl{未来潜在融资需求分析} 2026 年 6 月末货币资金 6.08 亿元，而短期债务 13.29 亿元（现金短债比 0.46 倍），2026 年上半年经营现金流净额 0.41 亿元。按“短期债务减去账面货币资金”的滚动偿付口径，静态缺口约 7.22 亿元；若再计入上半年亏损的年化额（9.78 亿元），维持一年正常经营所需的外部资金约 7.22---17.00 亿元。 公司 2026 年上半年归母净利润为负（-4.89 亿元），但 6 月末资产负债率 63.6\% 尚未进入第一档（亏损且负债率 $>$65\%）的区间，当前融资需求以补充流动资金、支撑低谷期经营性支出为主。 公开披露的下一步融资线索包括：2024-11 19.3 亿元定增预案为上市以来最大规模股权融资计划（此前三十年股权融资总额不过 2.7 亿元），拟投向云浮、贺州生猪项……；2026-08 2026 年上半年大额亏损、经营活动现金流量净额同比减少 87\%、控股股东股份质押用于补充流动资金，反映资金面压力上升；……。 资金需求还要叠加在建工程：期末在建工程 1.32 亿元、2026 年上半年购建固定资产等支出 0.95 亿元（未年化），这部分支出在周期底部通常会被主动放缓。\footnote{\texttt{https://www.stcn.com/article/detail/1411786.html}}\footnote{\texttt{https://finance.sina.com.cn/stock/zqgd/2026-08-23/doc-iniphwyw3620134.shtml}}


\subsubsection{益生股份（002458）}
\label{co:002458}

\pfl{公司基本情况} 1997 年设立至今，益生股份国内最大的祖代白羽肉种鸡企业，祖代肉种鸡饲养规模连续二十年居全国第一，形成“种鸡＋种猪”双轮驱动；2026 年上半年受益白羽肉鸡行业景气上行与引种受限，净利润同比增长近 49 倍。 公司于 2010-08-10 在深交所主板首发上市，注册资本 14.31 亿元，总股本 11.06 亿股。 截至 2026-09-24收盘价 8.09 元，流通市值 78.2 亿元，近一年-6.9\%，流通市值在三级行业“肉鸡养殖”的 8 家公司中排第 \zhnumber{3} 位，近一年涨跌幅高于全样本中位数（-10.1\%）3.2 个百分点。 按 2026 年半年报披露的行业构成，收入占比前三位为鸡收入 83.8\%；猪收入 13.5\%；乳品收入 1.2\%。 发展过程中的关键节点包括：7 月 28 日发行、8 月 10 日在深交所中小板上市，发行 2,700 万股（网下配……（2010）；招股说明书披露拟募资用于扩大种鸡产能等项目；主承销商安信证券，发行市盈率 32.43 倍（2010）；完成 2022 年度向特定对象发行股票，募集资金总额 11.6 亿元，用于种猪养殖、种鸡……（2023）。\footnote{\texttt{https://vip.stock.finance.sina.com.cn/corp/view/vISSUE\_MarketBulletinDetail.php}}\footnote{\texttt{http://money.finance.sina.com.cn/corp/view/vISSUE\_RaiseExplanationDetail.php}}\footnote{\texttt{https://money.finance.sina.com.cn/corp/view/vCB\_AllBulletinDetail.php}}

\pfl{历史股价} 本报告样本区间（2021 年 1 月 至 2026 年 9 月）内，股价由 10.20 元变为 8.09 元，累计 -20.7\%，区间最高 15.59 元（2023 年 2 月）、最低 7.77 元（2026 年 6 月）；当前价相当于区间最高价的 51.9\%，已接近区间最低价（1.04 倍）。 月度收益的年化波动率 35.2\%，低于全样本中位数 37.4\%，在三级行业“肉鸡养殖”的 8 家公司中波动率排第 \zhnumber{5} 位。 从事件看，2022 年 5 月 至 2023 年 2 月的上涨（+78.0\%）与公司2022 年 9 月 29 日披露预案（拟募资不超过 11.6 亿元）；证监……（2022-09-29）在时间上重合；2023 年 2 月 至 2024 年 8 月的下跌（-48.3\%）与公司11.6 亿元定增完成，募投项目逐步投入使用后种猪产量提升，公司明确朝“种鸡……（2023-12）在时间上重合——股价的拐点多数能在公司自身的资本运作与风险事件上找到对应。\footnote{\texttt{http://money.finance.sina.com.cn/corp/view/vCB\_AllBulletinDetail.php}}\footnote{\texttt{https://vip.stock.finance.sina.com.cn/corp/view/vCB\_AllBulletinDetail.php}}

\begin{center}
\begin{tikzpicture}
\begin{axis}[
  width=12.6cm, height=3.9cm, scale only axis,
  xmin=2020.922, xmax=2026.888, ymin=7.23, ymax=16.68,
  xtick={2021,2022,2023,2024,2025,2026},
  yearticks,
  yticklabel style={font=\scriptsize},
  ylabel={元/股}, ylabel style={font=\scriptsize},
  grid=major, grid style={LightGray, line width=0.3pt},
  axis line style={InkGray, line width=0.5pt},
  every axis plot/.append style={line width=0.9pt},
]
\addplot[AgriGreen] table[x=x,y=y]{data/clean/profiles/px_002458.dat};
\addplot[SoftRed, dashed, line width=0.7pt] table[x=x,y=y]{data/clean/profiles/pxzz_002458.dat};
\addplot[only marks, mark=*, mark size=1.1pt, SoftRed] table[x=x,y=y]{data/clean/profiles/pxzz_002458.dat};
\end{axis}
\end{tikzpicture}

\captionof{figure}[益生股份（002458）股价与周期骨架]{益生股份（002458）月度收盘价（元/股）与 ZigZag 周期骨架（阈值 25\%，圆点为周期转折点）}
\end{center}

\pfl{过去周期分析} 以 25\% 的 ZigZag 阈值划分，样本区间内股价形成 2 个主升段与 3 个主跌段。 最大主升段出现在 2022 年 5 月 至 2023 年 2 月，9 个月+78.0\%；最大主跌段为 2023 年 2 月 至 2024 年 8 月，18 个月-48.3\%。 最近一次转折出现在 2026 年 6 月（下跌段自 2026 年 4 月起，2 个月-30.7\%），当前处于该段的延续之中。 公司股价与申万农林牧渔指数几乎同步见底，个股并未走出独立行情。 禽链的产能周期只有 6---12 个月、反转快于生猪（第 \ref{sec:other-livestock} 节），2026 年前四个月毛鸡养殖利润同比大幅回升，链条利润的修复先于生猪，公司 2026 年上半年实现盈利、半年 ROE 6.9\%（未年化），好于全样本中位数（0.75\%），下行周期对其报表的冲击小于同业。

\pfl{盈利情况} 2026 年上半年营业收入 16.97 亿元（同比 +28.4\%），归母净利润 3.08 亿元，同比 +4897.3\%。 以年报为趋势对照，2023---2025 年营业收入由 32.25 亿元变为 29.52 亿元（两年年均复合增速 -4.3\%），归母净利润由 5.41 亿元变为 1.65 亿元。 按半年报口径，2026 年上半年的 ROE 为 6.9\%（未年化）、销售净利率 18.2\%、毛利率 28.0\%，ROE 在“肉鸡养殖”8 家公司中排第 \zhnumber{1} 位。 同期全样本半年 ROE 中位数 0.75\%，公司与之相差 6.11 个百分点（略好于中位数公司）。

\begin{center}\small\setlength{\tabcolsep}{4pt}
\captionof{table}[益生股份（002458）关键财务指标]{益生股份（002458）关键财务指标（合并报表口径；两个半年报期均未年化；现金短债比 $=$ 货币资金 $\div$ 短期债务）}
\begin{tabular}{@{}lrrrrr@{}}
\toprule
指标 & 2023 年 & 2024 年 & 2025 年 & 2025 年半年报 & 2026 年半年报 \\
\midrule
营业收入（亿元） & 32.25 & 31.36 & 29.52 & 13.21 & 16.97 \\
归母净利润（亿元） & 5.41 & 5.04 & 1.65 & 0.06 & 3.08 \\
毛利率（\%） & 31.6 & 29.4 & 17.8 & 12.1 & 28.0 \\
ROE（\%） & 17.4 & 11.3 & 3.8 & 0.1 & 6.9 \\
资产负债率（\%） & 31.0 & 39.6 & 35.7 & 42.4 & 35.1 \\
经营活动现金流净额（亿元） & 10.04 & 7.53 & 5.62 & 1.43 & 4.61 \\
货币资金（亿元） & 9.09 & 16.27 & 10.01 & 17.65 & 14.79 \\
有息负债合计（亿元） & 10.67 & 20.47 & 17.52 & 24.50 & 18.88 \\
现金短债比（倍） & 1.08 & 0.88 & 0.65 & 0.79 & 0.90 \\
\bottomrule
\end{tabular}
\end{center}

\begin{center}
\begin{tikzpicture}
\begin{groupplot}[
  group style={group size=2 by 1, horizontal sep=0.95cm},
  width=7.2cm, height=4.6cm,
  tick label style={font=\tiny},
  title style={font=\scriptsize\heiti},
  yticklabel style={font=\tiny},
  grid=major, grid style={LightGray, line width=0.3pt},
  axis line style={InkGray, line width=0.5pt},
  enlarge x limits={abs=0.8},
  legend style={font=\scriptsize, at={(0.5,-0.22)}, anchor=north,
                legend columns=2, draw=none, fill=none, column sep=6pt},
]
\nextgroupplot[
  ybar, bar width=4pt,
  ymin=-7.26, ymax=36.56,
  xtick={0,1,2,3,4,5.7},
  xticklabels={2021,2022,2023,2024,2025,2026H1},
  ylabel={亿元},
]
\addplot[fill=AgriGreen!75, draw=AgriGreen, bar shift=-2.6pt] table[x=i, y=rev]{data/clean/profiles/fin_002458.dat};
\addplot[fill=SoftRed!60, draw=SoftRed, bar shift=2.6pt] table[x=i, y=ni]{data/clean/profiles/fin_002458.dat};
\addplot[fill=AgriGreen!30, draw=AgriGreen, pattern=north east lines, bar shift=-2.6pt] table[x=x, y=rev]{data/clean/profiles/finh_002458.dat};
\addplot[fill=SoftRed!20, draw=SoftRed, pattern=north east lines, bar shift=2.6pt] table[x=x, y=ni]{data/clean/profiles/finh_002458.dat};
\legend{营业收入（年/半年）, 归母净利润（年/半年）}
\nextgroupplot[
  ymin=-17.5, ymax=58.4,
  xtick={0,1,2,3,4,5.7},
  xticklabels={2021,2022,2023,2024,2025,2026H1},
  ylabel={\%},
]
\addplot[AgriGold, mark=*, mark size=1.3pt] table[x=i, y=roe]{data/clean/profiles/fin_002458.dat};
\addplot[OilBlue, mark=square*, mark size=1.3pt] table[x=i, y=debt]{data/clean/profiles/fin_002458.dat};
\addplot[only marks, AgriGold, mark=*, mark size=2.0pt] table[x=x, y=roe]{data/clean/profiles/finh_002458.dat};
\addplot[only marks, OilBlue, mark=square*, mark size=2.0pt] table[x=x, y=debt]{data/clean/profiles/finh_002458.dat};
\legend{ROE, 资产负债率}
\end{groupplot}
\end{tikzpicture}
\begin{tikzpicture}
\begin{axis}[
  width=15.2cm, height=3.1cm, scale only axis,
  ymin=-2.05, ymax=22.93,
  xtick={0,1,2,3,4,5.7},
  xticklabels={2021,2022,2023,2024,2025,2026H1},
  tick label style={font=\tiny},
  yticklabel style={font=\tiny},
  ylabel={亿元}, ylabel style={font=\scriptsize},
  ybar, bar width=4pt,
  grid=major, grid style={LightGray, line width=0.3pt},
  axis line style={InkGray, line width=0.5pt},
  enlarge x limits={abs=0.8},
  legend style={font=\scriptsize, at={(0.5,-0.30)}, anchor=north,
                legend columns=3, draw=none, fill=none, column sep=10pt},
]
\addplot[fill=AgriGreen!55, draw=AgriGreen, bar shift=-3.0pt] table[x=i, y=cash]{data/clean/profiles/fin_002458.dat};
\addplot[fill=SoftRed!45, draw=SoftRed, bar shift=3.0pt] table[x=i, y=ibd]{data/clean/profiles/fin_002458.dat};
\addplot[fill=AgriGreen!25, draw=AgriGreen, pattern=north east lines, bar shift=-3.0pt] table[x=x, y=cash]{data/clean/profiles/finh_002458.dat};
\addplot[fill=SoftRed!20, draw=SoftRed, pattern=north east lines, bar shift=3.0pt] table[x=x, y=ibd]{data/clean/profiles/finh_002458.dat};
\addplot[OilBlue, mark=*, mark size=1.1pt, line width=0.8pt] table[x=i, y=ocf]{data/clean/profiles/fin_002458.dat};
\addplot[only marks, OilBlue, mark=*, mark size=1.8pt] table[x=x, y=ocf]{data/clean/profiles/finh_002458.dat};
\legend{货币资金, 有息负债, 经营活动现金流净额}
\end{axis}
\end{tikzpicture}

\captionof{figure}[益生股份（002458）近年主要财务指标]{益生股份（002458）近年主要财务指标（左上：营业收入与归母净利润；右上：ROE 与资产负债率；下：货币资金、有息负债与经营活动现金流净额。
2021---2025 年为年报口径，2026H1 为 2026 年半年报/半年末，与其前一年报之间留出间隔、以斜纹或空心点区分；半年报数据未年化）}
\end{center}

\pfl{财务分析} 2026 年 6 月末资产负债率 35.1\%，较 2025 年末 下降 0.5 个百分点（样本中位数 52.2\%，所处三级行业内排第 \zhnumber{1} 位）。 2026 年 6 月末货币资金 14.79 亿元、短期债务（短期借款与一年内到期非流动负债）16.42 亿元、长期债务 2.46 亿元，有息负债合计 18.88 亿元，现金短债比 0.90 倍——覆盖偏紧。 净有息负债 4.09 亿元，相当于其 2026 年上半年营业收入的 24\%。 流动比率 0.91、速动比率 0.73，短期偿付压力较大。 2026 年上半年经营活动现金流净额 4.61 亿元，经营现金流与利润同向为正，内生资金可以支撑运营。 2026 年 6 月末在建工程 0.66 亿元，较上年末减少 0.33 亿元，在建投入在收缩。 公司披露的财务与合规风险包括：公司收到深交所出具的《关于山东益生种畜禽股份有限公司申请向特定对象发行股票的审核问询函》……。\footnote{\texttt{http://www.stcn.com/article/detail/839042.html}}

\pfl{重大战略转型} 公司在报告期内的战略动作可归为 4 项：2021 公司从丹麦引进 1,300 头纯种母猪，种猪业务已有 4 个种猪场投入使用；2023-12 11.6 亿元定增完成，募投项目逐步投入使用后种猪产量提升，公司明确朝“种鸡和种猪双轮驱动”战略目标推进；2024-03 2023 年度利润分配：以总股本 1,106,412,915 股为基数，每 10 股派发现金红利 2 元（含税）；2026-07 2026 年半年度拟每 10 股派发现金红利 1.5 元（含税）；业绩增长主要受白羽肉鸡苗量价齐升带动。 公告口径上，近三年（2023 年 9 月---2026 年 9 月）公司披露重大重组与并购类公告 0 条、对外投资与转型类公告 2 条、回购与股权激励类公告 51 条，合计公告 451 条。 第一大行业仍是“鸡收入”，收入占比 83.8\%，与 2021 年年报相比没有方向性变化。 综合判断，公司在报告期内有明确的外延扩张或业务结构调整动作，转型方向与融资安排之间存在直接关联。\footnote{\texttt{https://pdf.dfcfw.com/pdf/H3\_AP202203021550115208\_1.pdf}}\footnote{\texttt{https://vip.stock.finance.sina.com.cn/corp/view/vCB\_AllBulletinDetail.php}}

\begin{center}\small\setlength{\tabcolsep}{4pt}
\captionof{table}[益生股份（002458）历史融资与资本运作]{益生股份（002458）历史融资与资本运作（据公司公告与公开报道整理；金额为披露值，含首次公开发行、再融资、债券、银行借款与重整投资等）}
\begin{tabular}{@{}p{1.5cm}p{3.2cm}p{2.4cm}p{6.3cm}@{}}
\toprule
时间 & 方式 & 规模 & 用途与说明 \\
\midrule
2010-08 & IPO（深交所中小板） & 2,700 万股，发行价 24.00 元/股（约 6.48 亿元） & 种鸡养殖及孵化产能扩张等项目；证监许可[2010]752 号核准；主承销商安信证券 \\
2016-01 & 定向增发（非公开发行） & 51,314,641 股，发行价 11.55 元/股，募集资金总额 592,684,103.55 元、净额 582,684,103.55 元 & 种鸡养殖项目等；承销及保荐费用 1,000 万元…… \\
2021-10 & 可转换公司债券（申请未获核准） & 原拟公开募集不超过 10.4 亿元，修订稿缩减为 8.9 亿元 & 4 个种猪养殖项目、3 个饲料生产项目和 2 个种鸡孵化场项目（修订后各砍掉一个种猪…… \\
2022-09-29 & 定向增发（2022 年度向特定对象发行股票） & 113,502,935 股，发行价 10.22 元/股，募集资金总额 11.60 亿元、净额 11.3911 亿元 & 种猪养殖项目（双鸭山祖代种猪场 3.9 亿元…… \\
\bottomrule
\end{tabular}
\end{center}

\pfl{历史融投资情况} 从现金流量表看，2025 年公司取得借款收到的现金 32.68 亿元、偿还债务支付的现金 28.66 亿元、吸收权益性投资 — 亿元，筹资活动现金流净额 -6.30 亿元（2023 年为取得借款 9.28 亿元、偿还债务 17.12 亿元）。 2026 年上半年筹资活动现金流净额为 -0.34 亿元（取得借款 12.30 亿元、偿还债务 9.94 亿元，未年化），已转为净流出，处于净还债状态。 2025 年分配股利、利润或偿付利息支付的现金合计 4.25 亿元。 公开信息中与公司直接相关的融资与资本运作包括：IPO（深交所中小板）、可转换公司债券（申请未获核准）、定向增发（2022 年度向特定对象发行股票）、定向增发（非公开发行）。 其中规模较大的两笔为：2022-09-29定向增发（2022 年度向特定对象发行股票）（113,502,935 股，发行价 10.22 元/股，募集……）；2021-10可转换公司债券（申请未获核准）（原拟公开募集不超过 10.4 亿元，修订稿缩减为 8.9 亿……）。 股本方面，近三年披露股本变动 5 次（增发新股上市 1 次、限售股份上市 1 次）；总股本由 2021 年末的 9.93 亿股变为 11.06 亿股（+11.43\%）。 分红方面，近三年现金分红 6 次、累计每股派现 0.70 元（最近一次 2026-06-30）——在利润承压的阶段仍维持了分红。 近三年“再融资”类公告 10 条，涉及定向增发、募集资金使用。\footnote{\texttt{https://vip.stock.finance.sina.com.cn/corp/view/vISSUE\_MarketBulletinDetail.php}}\footnote{\texttt{https://vip.stock.finance.sina.com.cn/corp/view/vCB\_AllBulletinDetail.php}}\footnote{\texttt{http://www.stcn.com/article/detail/839042.html}}

\pfl{未来潜在融资需求分析} 2026 年 6 月末货币资金 14.79 亿元，而短期债务 16.42 亿元（现金短债比 0.90 倍），2026 年上半年经营现金流净额 4.61 亿元。按“短期债务减去账面货币资金”的滚动偿付口径，静态缺口约 1.63 亿元；若再计入上半年亏损的年化额（0.00 亿元），维持一年正常经营所需的外部资金约 1.63---1.63 亿元。 按第 \ref{sec:financing-need} 节的三档划分，公司属于第三档“扩张型”：2026 年上半年 ROE 6.9\%（未年化，折合约 13.7\%）、6 月末资产负债率 35.1\%，资产负债表仍具备加杠杆空间；融资用途以产能扩张、品类并购与海外布局为主，单笔规模通常在 5---20 亿元量级，会被市场定价为成长而非补血。 公开披露的下一步融资线索包括：2025-2026 公司推进“质量回报双提升”行动方案：2025 年累计引进祖代白羽肉种鸡 26.60 万套，占全国祖代肉种鸡进口总量超 4……；2026-01 2025 年定增募集资金使用情况：累计使用 75,457.92 万元（占 66.24\%），其中 18,500 万元暂时补……。 资金需求还要叠加在建工程：期末在建工程 0.66 亿元、2026 年上半年购建固定资产等支出 1.21 亿元（未年化），这部分支出在周期底部通常会被主动放缓。\footnote{\texttt{http://money.finance.sina.com.cn/corp/view/vCB\_AllBulletinDetail.php}}\footnote{\texttt{http://money.finance.sina.com.cn/corp/view/vCB\_AllBulletinDetail.php}}


\subsubsection{巨星农牧（603477）}
\label{co:603477}

\pfl{公司基本情况} 四川乐山的生猪养殖与饲料企业，前身为皮革企业振静股份（2017 年上交所上市），2020 年发行股份购买巨星农牧 100\% 股权后主业转型为生猪养殖，2023 年 9 月控股股东变更为四川巨星集团、实际控制人变更为唐光跃；2026 年上半年出栏 280.77 万头但受猪价下行拖累由盈转亏。 公司于 2017-12-18 在上交所首发上市，注册资本 5.23 亿元，总股本 5.10 亿股。 截至 2026-09-24收盘价 13.89 元，流通市值 70.8 亿元，近一年-32.9\%，流通市值在三级行业“生猪养殖”的 11 家公司中排第 \zhnumber{6} 位，近一年涨跌幅低于全样本中位数（-10.1\%）22.8 个百分点。 按 2026 年半年报披露的行业构成，收入占比前三位为农牧业务 96.9\%；皮革业务 3.1\%。 与 2021 年年报相比，第一大行业由“养殖行业”变为“农牧业务”，主业重心已经迁移。 发展过程中的关键节点包括：因筹划向巨星农牧股份有限公司全体股东发行股份及支付现金购买其 100\% 股份，公司股票自……（2019）；7 月完成收购巨星农牧有限公司 100\% 股权的资产过户，自 2020 年 7 月 1 ……（2020）；公司名称由“四川振静股份有限公司”变更为“乐山巨星农牧股份有限公司”，A 股证券简称由“……（2021）。\footnote{\texttt{https://money.finance.sina.com.cn/corp/view/vCB\_AllBulletinDetail.php}}\footnote{\texttt{http://static.cninfo.com.cn/finalpage/2022-04-21/1212999930.PDF}}\footnote{\texttt{https://pdf.dfcfw.com/pdf/H2\_AN202203111551964703\_1.pdf}}

\pfl{历史股价} 自 2021 年 1 月至 2026 年 9 月，股价由 10.85 元变为 13.89 元，累计 +28.0\%，区间最高 37.47 元（2023 年 12 月）、最低 9.99 元（2021 年 7 月）；当前价相当于区间最高价的 37.1\%，为区间最低价的 1.39 倍。 月度收益的年化波动率 49.4\%，高于全样本中位数 37.4\%，在三级行业“生猪养殖”的 11 家公司中波动率排第 \zhnumber{2} 位。 从事件看，2022 年 10 月 至 2023 年 12 月的上涨（+95.7\%）与公司四川和邦投资集团有限公司与四川巨星企业集团有限公司签署股份转让协议，以 27……（2023-08-07）在时间上重合——股价的拐点多数能在公司自身的资本运作与风险事件上找到对应。\footnote{\texttt{https://www.cs.com.cn/ssgs/gsxw/202308/t20230809\_6360190.html}}

\begin{center}
\begin{tikzpicture}
\begin{axis}[
  width=12.6cm, height=3.9cm, scale only axis,
  xmin=2020.922, xmax=2026.888, ymin=9.29, ymax=40.09,
  xtick={2021,2022,2023,2024,2025,2026},
  yearticks,
  yticklabel style={font=\scriptsize},
  ylabel={元/股}, ylabel style={font=\scriptsize},
  grid=major, grid style={LightGray, line width=0.3pt},
  axis line style={InkGray, line width=0.5pt},
  every axis plot/.append style={line width=0.9pt},
]
\addplot[AgriGreen] table[x=x,y=y]{data/clean/profiles/px_603477.dat};
\addplot[SoftRed, dashed, line width=0.7pt] table[x=x,y=y]{data/clean/profiles/pxzz_603477.dat};
\addplot[only marks, mark=*, mark size=1.1pt, SoftRed] table[x=x,y=y]{data/clean/profiles/pxzz_603477.dat};
\end{axis}
\end{tikzpicture}

\captionof{figure}[巨星农牧（603477）股价与周期骨架]{巨星农牧（603477）月度收盘价（元/股）与 ZigZag 周期骨架（阈值 25\%，圆点为周期转折点）}
\end{center}

\pfl{过去周期分析} 以 25\% 的 ZigZag 阈值划分，样本区间内股价形成 5 个主升段与 5 个主跌段。 最大主升段出现在 2021 年 7 月 至 2022 年 3 月，8 个月+156.9\%；最大主跌段为 2023 年 12 月 至 2025 年 1 月，13 个月-55.3\%。 最近一次转折出现在 2026 年 9 月（下跌段自 2025 年 5 月起，16 个月-36.4\%），当前处于该段的延续之中。 公司股价较申万农林牧渔指数提前约 58 个月见底，是板块中较早定价周期反转的公司之一。 生猪养殖的定价主线是产能去化与成本曲线的赛跑（第 \ref{sec:hog} 章）：2026 年三季度行业仍处在“产能去化未完成、价格磨底”的阶段，公司 2026 年 6 月末资产负债率 69.4\%，高于全样本中位数 52.2\%，其股价因此更多反映资金面与信用风险，而不是价格弹性本身。

\pfl{盈利情况} 2026 年上半年营业收入 35.57 亿元（同比 -4.3\%），归母净利润 -8.38 亿元，由上年同期的盈利转为亏损。 以年报为趋势对照，2023---2025 年营业收入由 40.41 亿元变为 80.00 亿元（两年年均复合增速 +40.7\%），归母净利润由 -6.45 亿元变为 0.30 亿元。 按半年报口径，2026 年上半年的 ROE 为 -30.3\%（未年化）、销售净利率 -23.6\%、毛利率 -9.7\%，ROE 在“生猪养殖”11 家公司中排第 \zhnumber{9} 位。 同期全样本半年 ROE 中位数 0.75\%，公司与之相差 31.07 个百分点（弱于中位数公司）。 结合公司披露，响应国家生猪产能调控政策主动淘汰低效母猪产能；2026 年上半年商品肥猪完全成本降至 12.6 元/kg……；可转债募投项目德昌巨星生猪繁育一体化项目于 2024 年底完成计划母猪配种并达到满产状态，预计可新增生猪……。\footnote{\texttt{http://stock.finance.sina.com.cn/stock/go.php/vReport\_Show/kind/search/rptid/84262056120}}\footnote{\texttt{https://static.cninfo.com.cn/finalpage/2025-05-22/1223627623.PDF}}

\begin{center}\small\setlength{\tabcolsep}{4pt}
\captionof{table}[巨星农牧（603477）关键财务指标]{巨星农牧（603477）关键财务指标（合并报表口径；两个半年报期均未年化；现金短债比 $=$ 货币资金 $\div$ 短期债务）}
\begin{tabular}{@{}lrrrrr@{}}
\toprule
指标 & 2023 年 & 2024 年 & 2025 年 & 2025 年半年报 & 2026 年半年报 \\
\midrule
营业收入（亿元） & 40.41 & 60.78 & 80.00 & 37.17 & 35.57 \\
归母净利润（亿元） & -6.45 & 5.19 & 0.30 & 1.81 & -8.38 \\
毛利率（\%） & -0.4 & 19.2 & 7.3 & 13.7 & -9.7 \\
ROE（\%） & -19.3 & 15.6 & 0.9 & 5.4 & -30.3 \\
资产负债率（\%） & 59.9 & 61.9 & 67.3 & 64.6 & 69.4 \\
经营活动现金流净额（亿元） & -1.05 & 12.65 & 12.18 & 6.02 & -2.86 \\
货币资金（亿元） & 4.46 & 5.77 & 7.56 & 6.70 & 6.94 \\
有息负债合计（亿元） & 31.42 & 34.47 & 39.84 & 36.21 & 37.45 \\
现金短债比（倍） & 0.29 & 0.30 & 0.36 & 0.33 & 0.34 \\
\bottomrule
\end{tabular}
\end{center}

\begin{center}
\begin{tikzpicture}
\begin{groupplot}[
  group style={group size=2 by 1, horizontal sep=0.95cm},
  width=7.2cm, height=4.6cm,
  tick label style={font=\tiny},
  title style={font=\scriptsize\heiti},
  yticklabel style={font=\tiny},
  grid=major, grid style={LightGray, line width=0.3pt},
  axis line style={InkGray, line width=0.5pt},
  enlarge x limits={abs=0.8},
  legend style={font=\scriptsize, at={(0.5,-0.22)}, anchor=north,
                legend columns=2, draw=none, fill=none, column sep=6pt},
]
\nextgroupplot[
  ybar, bar width=4pt,
  ymin=-17.22, ymax=90.61,
  xtick={0,1,2,3,4,5.7},
  xticklabels={2021,2022,2023,2024,2025,2026H1},
  ylabel={亿元},
]
\addplot[fill=AgriGreen!75, draw=AgriGreen, bar shift=-2.6pt] table[x=i, y=rev]{data/clean/profiles/fin_603477.dat};
\addplot[fill=SoftRed!60, draw=SoftRed, bar shift=2.6pt] table[x=i, y=ni]{data/clean/profiles/fin_603477.dat};
\addplot[fill=AgriGreen!30, draw=AgriGreen, pattern=north east lines, bar shift=-2.6pt] table[x=x, y=rev]{data/clean/profiles/finh_603477.dat};
\addplot[fill=SoftRed!20, draw=SoftRed, pattern=north east lines, bar shift=2.6pt] table[x=x, y=ni]{data/clean/profiles/finh_603477.dat};
\legend{营业收入（年/半年）, 归母净利润（年/半年）}
\nextgroupplot[
  ymin=-38.3, ymax=79.4,
  xtick={0,1,2,3,4,5.7},
  xticklabels={2021,2022,2023,2024,2025,2026H1},
  ylabel={\%},
]
\addplot[AgriGold, mark=*, mark size=1.3pt] table[x=i, y=roe]{data/clean/profiles/fin_603477.dat};
\addplot[OilBlue, mark=square*, mark size=1.3pt] table[x=i, y=debt]{data/clean/profiles/fin_603477.dat};
\addplot[only marks, AgriGold, mark=*, mark size=2.0pt] table[x=x, y=roe]{data/clean/profiles/finh_603477.dat};
\addplot[only marks, OilBlue, mark=square*, mark size=2.0pt] table[x=x, y=debt]{data/clean/profiles/finh_603477.dat};
\legend{ROE, 资产负债率}
\end{groupplot}
\end{tikzpicture}
\begin{tikzpicture}
\begin{axis}[
  width=15.2cm, height=3.1cm, scale only axis,
  ymin=-7.13, ymax=44.97,
  xtick={0,1,2,3,4,5.7},
  xticklabels={2021,2022,2023,2024,2025,2026H1},
  tick label style={font=\tiny},
  yticklabel style={font=\tiny},
  ylabel={亿元}, ylabel style={font=\scriptsize},
  ybar, bar width=4pt,
  grid=major, grid style={LightGray, line width=0.3pt},
  axis line style={InkGray, line width=0.5pt},
  enlarge x limits={abs=0.8},
  legend style={font=\scriptsize, at={(0.5,-0.30)}, anchor=north,
                legend columns=3, draw=none, fill=none, column sep=10pt},
]
\addplot[fill=AgriGreen!55, draw=AgriGreen, bar shift=-3.0pt] table[x=i, y=cash]{data/clean/profiles/fin_603477.dat};
\addplot[fill=SoftRed!45, draw=SoftRed, bar shift=3.0pt] table[x=i, y=ibd]{data/clean/profiles/fin_603477.dat};
\addplot[fill=AgriGreen!25, draw=AgriGreen, pattern=north east lines, bar shift=-3.0pt] table[x=x, y=cash]{data/clean/profiles/finh_603477.dat};
\addplot[fill=SoftRed!20, draw=SoftRed, pattern=north east lines, bar shift=3.0pt] table[x=x, y=ibd]{data/clean/profiles/finh_603477.dat};
\addplot[OilBlue, mark=*, mark size=1.1pt, line width=0.8pt] table[x=i, y=ocf]{data/clean/profiles/fin_603477.dat};
\addplot[only marks, OilBlue, mark=*, mark size=1.8pt] table[x=x, y=ocf]{data/clean/profiles/finh_603477.dat};
\legend{货币资金, 有息负债, 经营活动现金流净额}
\end{axis}
\end{tikzpicture}

\captionof{figure}[巨星农牧（603477）近年主要财务指标]{巨星农牧（603477）近年主要财务指标（左上：营业收入与归母净利润；右上：ROE 与资产负债率；下：货币资金、有息负债与经营活动现金流净额。
2021---2025 年为年报口径，2026H1 为 2026 年半年报/半年末，与其前一年报之间留出间隔、以斜纹或空心点区分；半年报数据未年化）}
\end{center}

\pfl{财务分析} 2026 年 6 月末资产负债率 69.4\%，较 2025 年末 上升 2.1 个百分点（样本中位数 52.2\%，所处三级行业内排第 \zhnumber{8} 位）。 2026 年 6 月末货币资金 6.94 亿元、短期债务（短期借款与一年内到期非流动负债）20.13 亿元、长期债务 17.32 亿元，有息负债合计 37.45 亿元，现金短债比 0.34 倍——覆盖偏紧。 净有息负债 30.51 亿元，相当于其 2026 年上半年营业收入的 86\%。 流动比率 0.80、速动比率 0.23，短期偿付压力较大。 2026 年上半年经营活动现金流净额 -2.86 亿元，经营现金流与利润同时为负，日常运营对债务与股东投入的依赖度较高。 2026 年 6 月末在建工程 4.62 亿元，较上年末增加 0.13 亿元，仍有在建投入。 2026 年 6 月末商誉 5.80 亿元，占归母净资产的 21.1\%，是净资产中最脆弱的一块。 公司披露的财务与合规风险包括：控股股东巨星集团持有公司 15,066.50 万股，其中累计质押 8,981 万股，占其……。\footnote{\texttt{https://static.sse.com.cn/stock/disclosure/announcement/c/202412/603477\_20241205\_R335.pd}}

\pfl{重大战略转型} 公司在报告期内的战略动作可归为 4 项：2020-07 通过发行股份及支付现金收购巨星农牧有限公司 100\% 股权（作价 18.20 亿元），由皮革主业转型为生猪养殖与……；2023-08-07 四川和邦投资集团有限公司与四川巨星企业集团有限公司签署股份转让协议，以 27.67 元/股向巨星集团转让 2,6……；2024-12 可转债募投项目德昌巨星生猪繁育一体化项目于 2024 年底完成计划母猪配种并达到满产状态，预计可新增生猪年出栏量……；2026 响应国家生猪产能调控政策主动淘汰低效母猪产能；2026 年上半年商品肥猪完全成本降至 12.6 元/kg 以内，……。 公告口径上，近三年（2023 年 9 月---2026 年 9 月）公司披露重大重组与并购类公告 0 条、对外投资与转型类公告 0 条、回购与股权激励类公告 33 条，合计公告 587 条。 主营结构的迁移已经落到报表上：2021 年年报的第一大行业为“养殖行业”（83.1\%），2026 年半年报变为“农牧业务”（96.9\%）。 综合判断，公司在报告期内有明确的外延扩张或业务结构调整动作，转型方向与融资安排之间存在直接关联。\footnote{\texttt{http://static.cninfo.com.cn/finalpage/2022-04-21/1212999930.PDF}}\footnote{\texttt{https://www.cs.com.cn/ssgs/gsxw/202308/t20230809\_6360190.html}}\footnote{\texttt{https://static.cninfo.com.cn/finalpage/2025-05-22/1223627623.PDF}}

\begin{center}\small\setlength{\tabcolsep}{4pt}
\captionof{table}[巨星农牧（603477）历史融资与资本运作]{巨星农牧（603477）历史融资与资本运作（据公司公告与公开报道整理；金额为披露值，含首次公开发行、再融资、债券、银行借款与重整投资等）}
\begin{tabular}{@{}p{1.5cm}p{3.2cm}p{2.4cm}p{6.3cm}@{}}
\toprule
时间 & 方式 & 规模 & 用途与说明 \\
\midrule
2017-12 & IPO（上海证券交易所主板，发行主体原为四川振静股份有限公司） & 实际募集资金 3.35 亿元（发行 6,000 万股，发行价 5.58 元/股） & 皮革主业相关项目，其中补充流动资金 7,610.12 万元…… \\
2020-07 & 发行股份及支付现金购买资产（重大资产重组） & 交易作价 18.20 亿元；发行 227,911,629 股、发行价 7.47 元/股，其中股份对价 170,250.00 万元、现金对价 11,750.00 万元 & 收购巨星农牧有限公司 100\% 股权（评估值 182,122.48 万元）…… \\
2021-06-28 & 非公开发行股票（向 8 名特定对象） & 募集资金总额 41,999.9954 万元，净额 40,822.81 万元（38,181,814 股，发行价 11.00 元/股） & 生猪养殖项目（含古蔺巨星皇华种猪场项目 6,000.00 万元）及补流还贷…… \\
2022-04-25 & 公开发行可转换公司债券（债券简称“巨星转债”，代码 113648，期限 6 年） & 面值总额 100,000.00 万元（1,000 万张，每张 100.00 元） & 德昌巨星生猪繁育一体化项目和补充流动资金；经证监许可[2022]663 号文核准…… \\
2026-04 & 向特定对象发行股票（2023 年度定增，6 名发行对象） & 募集资金总额 227,065,985.49 元，净额 218,864,505.63 元（12,908,811 股，发行价 17.59 元/股） & 公猪站建设项目、养殖技术研究基地建设项目、数字智能化建设项目…… \\
\bottomrule
\end{tabular}
\end{center}

\pfl{历史融投资情况} 从现金流量表看，2025 年公司取得借款收到的现金 19.43 亿元、偿还债务支付的现金 15.34 亿元、吸收权益性投资 — 亿元，筹资活动现金流净额 -3.04 亿元（2023 年为取得借款 19.59 亿元、偿还债务 13.44 亿元）。 2026 年上半年筹资活动现金流净额为 2.30 亿元（取得借款 10.46 亿元、偿还债务 13.00 亿元，未年化），仍处于净融入状态。 2025 年分配股利、利润或偿付利息支付的现金合计 1.99 亿元。 公司披露的股本与债务融资脉络包括：IPO（上海证券交易所主板，发行主体原为四川振静股份有限公司）、公开发行可转换公司债券（债券简称“巨星转债”，代码 113648，期限 6 年）、发行股份及支付现金购买资产（重大资产重组）、向特定对象发行股票（2023 年度定增，6 名发行对象）、非公开发行股票（向 8 名特定对象）。 其中规模较大的两笔为：2022-04-25公开发行可转换公司债券（债券简称“巨星转债”，代码 113648，期限 6 年）（面值总额 100,000.00 万元（1,000 万张，每张……）；2017-12IPO（上海证券交易所主板，发行主体原为四川振静股份有限公司）（实际募集资金 3.35 亿元（发行 6,000 万股，发行价……）。 股本方面，近三年披露股本变动 25 次（可转债转股 23 次、限售股份上市 1 次、其他 1 次）；总股本由 2021 年末的 5.06 亿股变为 5.10 亿股（+0.79\%）。 分红方面，近三年现金分红 3 次、累计每股派现 0.22 元（最近一次 2025-12-31）——分红规模与利润的下滑幅度相比仍然有限。 近三年“再融资”类公告 75 条，涉及定向增发、可转债、募集资金使用。\footnote{\texttt{https://basic.10jqka.com.cn/603477/company.html}}\footnote{\texttt{https://vip.stock.finance.sina.com.cn/corp/view/vISSUE\_MarketBulletinDetail.php}}\footnote{\texttt{https://static.sse.com.cn/stock/disclosure/announcement/c/202412/603477\_20241205\_R335.pd}}

\pfl{未来潜在融资需求分析} 2026 年 6 月末货币资金 6.94 亿元，而短期债务 20.13 亿元（现金短债比 0.34 倍），2026 年上半年经营现金流净额 -2.86 亿元。按“短期债务减去账面货币资金”的滚动偿付口径，静态缺口约 13.19 亿元；若再计入上半年亏损的年化额（16.77 亿元），维持一年正常经营所需的外部资金约 13.19---29.96 亿元。 按第 \ref{sec:financing-need} 节的三档划分，公司属于第一档“补血型”：2026 年上半年归母净利润 -8.38 亿元、6 月末资产负债率 69.4\%。 可行的融资路径按现实性排序为：定向增发（需股价配合，且控股股东需放弃优先认购或同步增资）、引入国资或产业投资人（部分公司以控制权让渡为对价）、债务重组或展期（与银行和债券持有人协商）。 公开披露的下一步融资线索包括：2024 公司积极推动定增，拟募集资金总额不超过 8.00 亿元用于生猪养殖产能新建项目、养殖技术研究基地建设项目、数字智能化建设……；2026-02-04 向特定对象发行股票申请获上海证券交易所审核通过，募投项目为公猪站建设项目、养殖技术研究基地建设项目、数字智能化建设项目。 资金需求还要叠加在建工程：期末在建工程 4.62 亿元、2026 年上半年购建固定资产等支出 1.60 亿元（未年化），这部分支出在周期底部通常会被主动放缓。\footnote{\texttt{https://static.cninfo.com.cn/finalpage/2025-05-22/1223627623.PDF}}\footnote{\texttt{https://www.stcn.com/article/detail/3630148.html}}


\subsubsection{罗牛山（000735）}
\label{co:000735}

\pfl{公司基本情况} 公司前身可追溯至 1987 年，海南海口的老牌农业上市公司，1997 年深交所上市的中国首家“菜篮子”股份制上市企业，主营生猪养殖与屠宰、冷链仓储物流、房地产开发销售及职业教育服务四大板块；2026 年上半年出栏生猪 45.62 万头、营收逆势增长，但受猪价深跌拖累归母净亏损 3.10 亿元。 公司于 1997-06-11 在深交所主板首发上市，注册资本 11.52 亿元，总股本 11.52 亿股。 截至 2026-09-24收盘价 5.64 元，流通市值 64.9 亿元，近一年-10.9\%，流通市值在三级行业“生猪养殖”的 11 家公司中排第 \zhnumber{7} 位，近一年涨跌幅低于全样本中位数（-10.1\%）0.8 个百分点。 按 2026 年半年报披露的产品构成，收入占比前三位为生猪 39.3\%；其他 16.5\%；饲料 15.2\%。 与 2021 年年报相比，第一大产品由“肉制品”变为“生猪”，主业重心已经迁移。 发展过程中的关键节点包括：6 月 11 日在深交所上市，发行价 4.35 元/股，是中国首家“菜篮子”股份制上市企……（1997）；实施增发 A 股，实际募集资金净额 4.79 亿元（经证监发行字[2002]107 号文……（2002）；4 月 13 日 2016 年非公开发行新增股份 271,381,578 股在深交所上市……（2016）。\footnote{\texttt{https://q.stock.sohu.com/cn/000735/index.shtml}}\footnote{\texttt{https://basic.10jqka.com.cn/000735/capital.html}}\footnote{\texttt{http://money.finance.sina.com.cn/corp/view/vISSUE\_MarketBulletinDetail.php}}

\pfl{历史股价} 本报告样本区间（2021 年 1 月 至 2026 年 9 月）内，股价由 8.18 元变为 5.64 元，累计 -31.1\%，区间最高 9.04 元（2025 年 12 月）、最低 4.33 元（2024 年 6 月）；当前价相当于区间最高价的 62.4\%，为区间最低价的 1.30 倍。 月度收益的年化波动率 27.3\%，低于全样本中位数 37.4\%，在三级行业“生猪养殖”的 11 家公司中波动率排第 \zhnumber{11} 位。 从事件看，2025 年 12 月 至 2026 年 6 月的下跌（-47.5\%）与公司2026 年上半年受猪价超预期下跌影响，归属于上市公司股东的净利润为 -3.……（2026-06-30）在时间上重合——股价的拐点多数能在公司自身的资本运作与风险事件上找到对应。\footnote{\texttt{https://money.finance.sina.com.cn/corp/view/vCB\_AllBulletinDetail.php}}

\begin{center}
\begin{tikzpicture}
\begin{axis}[
  width=12.6cm, height=3.9cm, scale only axis,
  xmin=2020.922, xmax=2026.888, ymin=4.03, ymax=9.67,
  xtick={2021,2022,2023,2024,2025,2026},
  yearticks,
  yticklabel style={font=\scriptsize},
  ylabel={元/股}, ylabel style={font=\scriptsize},
  grid=major, grid style={LightGray, line width=0.3pt},
  axis line style={InkGray, line width=0.5pt},
  every axis plot/.append style={line width=0.9pt},
]
\addplot[AgriGreen] table[x=x,y=y]{data/clean/profiles/px_000735.dat};
\addplot[SoftRed, dashed, line width=0.7pt] table[x=x,y=y]{data/clean/profiles/pxzz_000735.dat};
\addplot[only marks, mark=*, mark size=1.1pt, SoftRed] table[x=x,y=y]{data/clean/profiles/pxzz_000735.dat};
\end{axis}
\end{tikzpicture}

\captionof{figure}[罗牛山（000735）股价与周期骨架]{罗牛山（000735）月度收盘价（元/股）与 ZigZag 周期骨架（阈值 25\%，圆点为周期转折点）}
\end{center}

\pfl{过去周期分析} 以 25\% 的 ZigZag 阈值划分，样本区间内股价形成 3 个主升段与 3 个主跌段。 最大主升段出现在 2024 年 6 月 至 2025 年 12 月，18 个月+108.8\%；最大主跌段为 2025 年 12 月 至 2026 年 6 月，6 个月-47.5\%。 最近一次转折出现在 2026 年 6 月（下跌段自 2025 年 12 月起，6 个月-47.5\%），当前处于该段的延续之中。 公司股价较申万农林牧渔指数提前约 23 个月见底，是板块中较早定价周期反转的公司之一。 生猪养殖的定价主线是产能去化与成本曲线的赛跑（第 \ref{sec:hog} 章）：2026 年三季度行业仍处在“产能去化未完成、价格磨底”的阶段，公司 2026 年 6 月末资产负债率 52.3\%、半年 ROE -7.4\%（未年化），在本轮下行中属于随行业同步调整的一类。

\pfl{盈利情况} 2026 年上半年营业收入 12.41 亿元（同比 +14.5\%），归母净利润 -3.10 亿元，亏损同比扩大 3.03 亿元。 以年报为趋势对照，2023---2025 年营业收入由 42.42 亿元变为 33.77 亿元（两年年均复合增速 -10.8\%），归母净利润由 -4.88 亿元变为 1.54 亿元。 按半年报口径，2026 年上半年的 ROE 为 -7.4\%（未年化）、销售净利率 -26.2\%、毛利率 0.7\%，ROE 在“生猪养殖”11 家公司中排第 \zhnumber{2} 位。 同期全样本半年 ROE 中位数 0.75\%，公司与之相差 8.19 个百分点（弱于中位数公司）。 结合公司披露，2024 年公司实现营业收入 33.32 亿元，同比下降 21.44\%；实现归属于上市公司股东的净利润 ……；2025 年全年畜牧业收入 15.32 亿元，同比增长 11.94\%；房地产业收入 12.53 亿元，同……。\footnote{\texttt{http://money.finance.sina.com.cn/corp/view/vCB\_AllBulletinDetail.php}}\footnote{\texttt{https://stock.stockstar.com/RB2026082500001421.shtml}}

\begin{center}\small\setlength{\tabcolsep}{4pt}
\captionof{table}[罗牛山（000735）关键财务指标]{罗牛山（000735）关键财务指标（合并报表口径；两个半年报期均未年化；现金短债比 $=$ 货币资金 $\div$ 短期债务）}
\begin{tabular}{@{}lrrrrr@{}}
\toprule
指标 & 2023 年 & 2024 年 & 2025 年 & 2025 年半年报 & 2026 年半年报 \\
\midrule
营业收入（亿元） & 42.42 & 33.32 & 33.77 & 10.84 & 12.41 \\
归母净利润（亿元） & -4.88 & 2.42 & 1.54 & -0.07 & -3.10 \\
毛利率（\%） & 16.1 & 23.8 & 25.6 & 13.7 & 0.7 \\
ROE（\%） & -11.2 & 5.6 & 3.6 & -0.1 & -7.4 \\
资产负债率（\%） & 63.7 & 54.8 & 50.3 & 54.0 & 52.3 \\
经营活动现金流净额（亿元） & 3.75 & 3.71 & -0.93 & -1.27 & -1.38 \\
货币资金（亿元） & 9.53 & 10.95 & 7.59 & 4.24 & 4.11 \\
有息负债合计（亿元） & 29.80 & 16.62 & 17.41 & 14.57 & 19.18 \\
现金短债比（倍） & 1.30 & 1.95 & 0.93 & 0.88 & 0.41 \\
\bottomrule
\end{tabular}
\end{center}

\begin{center}
\begin{tikzpicture}
\begin{groupplot}[
  group style={group size=2 by 1, horizontal sep=0.95cm},
  width=7.2cm, height=4.6cm,
  tick label style={font=\tiny},
  title style={font=\scriptsize\heiti},
  yticklabel style={font=\tiny},
  grid=major, grid style={LightGray, line width=0.3pt},
  axis line style={InkGray, line width=0.5pt},
  enlarge x limits={abs=0.8},
  legend style={font=\scriptsize, at={(0.5,-0.22)}, anchor=north,
                legend columns=2, draw=none, fill=none, column sep=6pt},
]
\nextgroupplot[
  ybar, bar width=4pt,
  ymin=-9.61, ymax=48.09,
  xtick={0,1,2,3,4,5.7},
  xticklabels={2021,2022,2023,2024,2025,2026H1},
  ylabel={亿元},
]
\addplot[fill=AgriGreen!75, draw=AgriGreen, bar shift=-2.6pt] table[x=i, y=rev]{data/clean/profiles/fin_000735.dat};
\addplot[fill=SoftRed!60, draw=SoftRed, bar shift=2.6pt] table[x=i, y=ni]{data/clean/profiles/fin_000735.dat};
\addplot[fill=AgriGreen!30, draw=AgriGreen, pattern=north east lines, bar shift=-2.6pt] table[x=x, y=rev]{data/clean/profiles/finh_000735.dat};
\addplot[fill=SoftRed!20, draw=SoftRed, pattern=north east lines, bar shift=2.6pt] table[x=x, y=ni]{data/clean/profiles/finh_000735.dat};
\legend{营业收入（年/半年）, 归母净利润（年/半年）}
\nextgroupplot[
  ymin=-17.2, ymax=71.2,
  xtick={0,1,2,3,4,5.7},
  xticklabels={2021,2022,2023,2024,2025,2026H1},
  ylabel={\%},
]
\addplot[AgriGold, mark=*, mark size=1.3pt] table[x=i, y=roe]{data/clean/profiles/fin_000735.dat};
\addplot[OilBlue, mark=square*, mark size=1.3pt] table[x=i, y=debt]{data/clean/profiles/fin_000735.dat};
\addplot[only marks, AgriGold, mark=*, mark size=2.0pt] table[x=x, y=roe]{data/clean/profiles/finh_000735.dat};
\addplot[only marks, OilBlue, mark=square*, mark size=2.0pt] table[x=x, y=debt]{data/clean/profiles/finh_000735.dat};
\legend{ROE, 资产负债率}
\end{groupplot}
\end{tikzpicture}
\begin{tikzpicture}
\begin{axis}[
  width=15.2cm, height=3.1cm, scale only axis,
  ymin=-4.49, ymax=33.54,
  xtick={0,1,2,3,4,5.7},
  xticklabels={2021,2022,2023,2024,2025,2026H1},
  tick label style={font=\tiny},
  yticklabel style={font=\tiny},
  ylabel={亿元}, ylabel style={font=\scriptsize},
  ybar, bar width=4pt,
  grid=major, grid style={LightGray, line width=0.3pt},
  axis line style={InkGray, line width=0.5pt},
  enlarge x limits={abs=0.8},
  legend style={font=\scriptsize, at={(0.5,-0.30)}, anchor=north,
                legend columns=3, draw=none, fill=none, column sep=10pt},
]
\addplot[fill=AgriGreen!55, draw=AgriGreen, bar shift=-3.0pt] table[x=i, y=cash]{data/clean/profiles/fin_000735.dat};
\addplot[fill=SoftRed!45, draw=SoftRed, bar shift=3.0pt] table[x=i, y=ibd]{data/clean/profiles/fin_000735.dat};
\addplot[fill=AgriGreen!25, draw=AgriGreen, pattern=north east lines, bar shift=-3.0pt] table[x=x, y=cash]{data/clean/profiles/finh_000735.dat};
\addplot[fill=SoftRed!20, draw=SoftRed, pattern=north east lines, bar shift=3.0pt] table[x=x, y=ibd]{data/clean/profiles/finh_000735.dat};
\addplot[OilBlue, mark=*, mark size=1.1pt, line width=0.8pt] table[x=i, y=ocf]{data/clean/profiles/fin_000735.dat};
\addplot[only marks, OilBlue, mark=*, mark size=1.8pt] table[x=x, y=ocf]{data/clean/profiles/finh_000735.dat};
\legend{货币资金, 有息负债, 经营活动现金流净额}
\end{axis}
\end{tikzpicture}

\captionof{figure}[罗牛山（000735）近年主要财务指标]{罗牛山（000735）近年主要财务指标（左上：营业收入与归母净利润；右上：ROE 与资产负债率；下：货币资金、有息负债与经营活动现金流净额。
2021---2025 年为年报口径，2026H1 为 2026 年半年报/半年末，与其前一年报之间留出间隔、以斜纹或空心点区分；半年报数据未年化）}
\end{center}

\pfl{财务分析} 2026 年 6 月末资产负债率 52.3\%，较 2025 年末 上升 2.0 个百分点（样本中位数 52.2\%，所处三级行业内排第 \zhnumber{2} 位）。 2026 年 6 月末货币资金 4.11 亿元、短期债务（短期借款与一年内到期非流动负债）9.97 亿元、长期债务 9.21 亿元，有息负债合计 19.18 亿元，现金短债比 0.41 倍——覆盖偏紧。 净有息负债 15.07 亿元，相当于其 2026 年上半年营业收入的 121\%。 流动比率 1.00、速动比率 0.29，短期指标尚可。 2026 年上半年经营活动现金流净额 -1.38 亿元，经营现金流与利润同时为负，现金消耗较快。 2026 年 6 月末在建工程 1.69 亿元，较上年末增加 0.06 亿元，仍有在建投入。 公司披露的财务与合规风险包括：控股股东罗牛山集团有限公司持股 197,412,938 股（占 17.14\%），其中质押……。\footnote{\texttt{http://file.finance.sina.com.cn/211.154.219.97:9494/MRGG/CNSESZ\_STOCK/2021/2021-4/2021-0}}

\pfl{重大战略转型} 公司在报告期内的战略动作可归为 4 项：2023-04 公司在投资者关系活动中表示，2020 年定增终止后依靠自身力量建设猪场项目，产能从 50 万头发展到近 130 ……；2024 2024 年利润分配预案：以 1,151,513,578 股为基数，向全体股东每 10 股派发现金红利 0.2 ……；2025 2025 年度拟向全体股东每 10 股派发现金红利 0.35 元（含税）；2026-06-30 房地产业务继续贯彻“聚焦存量去化、均衡现金回笼”的稳健原则，不再盲目追加新增投资；报告期内该板块实现营业收入 1……。 公告口径上，近三年（2023 年 9 月---2026 年 9 月）公司披露重大重组与并购类公告 4 条、对外投资与转型类公告 2 条、回购与股权激励类公告 0 条，合计公告 297 条。 主营结构的迁移已经落到报表上：2021 年年报的第一大产品为“肉制品”（38.6\%），2026 年半年报变为“生猪”（39.3\%）。 综合判断，公司在报告期内有明确的外延扩张或业务结构调整动作，转型方向与融资安排之间存在直接关联。\footnote{\texttt{https://file.finance.qq.com/finance/hs/pdf/2023/04/07/1216347994.PDF}}\footnote{\texttt{http://epaper.stcn.com/pic/202504/19/9d6c1cb55ccc370f3cb0e822d0e14111.pdf}}\footnote{\texttt{https://money.finance.sina.com.cn/corp/view/vCB\_AllBulletinDetail.php}}

\begin{center}\small\setlength{\tabcolsep}{4pt}
\captionof{table}[罗牛山（000735）历史融资与资本运作]{罗牛山（000735）历史融资与资本运作（据公司公告与公开报道整理；金额为披露值，含首次公开发行、再融资、债券、银行借款与重整投资等）}
\begin{tabular}{@{}p{1.5cm}p{3.2cm}p{2.4cm}p{6.3cm}@{}}
\toprule
时间 & 方式 & 规模 & 用途与说明 \\
\midrule
1997-05 & IPO（深圳证券交易所主板） & 实际募集资金净额 1.87 亿元 & 发行价 4.35 元/股，上市日期 1997-06-11 \\
2002-01 & 增发 A 股 & 实际募集资金净额 4.79 亿元 & 募投项目包括海南爱华汽车广场项目（承诺投入 5,000 万元）等…… \\
2016-04-13 & 非公开发行股票（向 6 名特定对象，均以现金认购） & 募集资金总额 1,649,999,994.24 元，净额 1,630,624,839.08 元（271,381,578 股，发行价 6.08 元/股） & “罗牛山十万头现代化猪场项目”“海南农副产品交易配送中心及产业配套项目”和偿还银行贷款…… \\
2020 & 定增（终止） & --- & 原计划扩大生猪养殖规模；2020 年公司为扩大生猪养殖规模启动定增项目…… \\
2023-11-24 & 定增（终止） & 未查得（原 2023 年度向特定对象发行 A 股股票方案） & 公司第十届董事会第十次临时会议及第十届监事会第八次会议审议通过《关于终止公司 2023 …… \\
\bottomrule
\end{tabular}
\end{center}

\pfl{历史融投资情况} 从现金流量表看，2025 年公司取得借款收到的现金 7.40 亿元、偿还债务支付的现金 7.05 亿元、吸收权益性投资 — 亿元，筹资活动现金流净额 -1.49 亿元（2023 年为取得借款 8.72 亿元、偿还债务 5.90 亿元）。 2026 年上半年筹资活动现金流净额为 1.22 亿元（取得借款 5.67 亿元、偿还债务 3.79 亿元，未年化），仍处于净融入状态。 2025 年分配股利、利润或偿付利息支付的现金合计 0.74 亿元。 公司披露的股本与债务融资脉络包括：IPO（深圳证券交易所主板）、增发 A 股、定增（终止）、非公开发行股票（向 6 名特定对象，均以现金认购）。 其中规模较大的两笔为：2002-01增发 A 股（实际募集资金净额 4.79 亿元）；1997-05IPO（深圳证券交易所主板）（实际募集资金净额 1.87 亿元）。 股本方面，近三年披露股本变动 3 次；总股本由 2021 年末的 11.52 亿股变为 11.52 亿股（+0.00\%）。 分红方面，近三年现金分红 4 次、累计每股派现 0.10 元（最近一次 2025-12-31）——在利润承压的阶段仍维持了分红。 近三年“再融资”类公告 1 条，涉及定向增发。\footnote{\texttt{https://basic.10jqka.com.cn/000735/capital.html}}\footnote{\texttt{http://money.finance.sina.com.cn/corp/view/vISSUE\_MarketBulletinDetail.php}}

\pfl{未来潜在融资需求分析} 2026 年 6 月末货币资金 4.11 亿元，而短期债务 9.97 亿元（现金短债比 0.41 倍），2026 年上半年经营现金流净额 -1.38 亿元。按“短期债务减去账面货币资金”的滚动偿付口径，静态缺口约 5.86 亿元；若再计入上半年亏损的年化额（6.19 亿元），维持一年正常经营所需的外部资金约 5.86---12.05 亿元。 公司 2026 年上半年归母净利润为负（-3.10 亿元），但 6 月末资产负债率 52.3\% 尚未进入第一档（亏损且负债率 $>$65\%）的区间，当前融资需求以补充流动资金、支撑低谷期经营性支出为主。 公开披露的下一步融资线索包括：2023-11-24 终止 2023 年度向特定对象发行 A 股股票事项，原扩张产能的股权融资渠道关闭；此后公司主要通过银行借款等方式推进猪场……。 资金需求还要叠加在建工程：期末在建工程 1.69 亿元、2026 年上半年购建固定资产等支出 0.68 亿元（未年化），这部分支出在周期底部通常会被主动放缓。\footnote{\texttt{http://vip.stock.finance.sina.com.cn/corp/view/vCB\_AllBulletinDetail.php}}


\subsubsection{新五丰（600975）}
\label{co:600975}

\pfl{公司基本情况} 湖南省国资委下辖唯一以生猪产业为主业的国有控股上市公司，2004 年上交所上市，2022 年并购天心种业并连续完成多轮再融资，出栏量由 2021 年 44.15 万头增至 2025 年 542.24 万头；2026 年上半年受猪价下行影响营收 28.33 亿元、归母净亏损 7.36 亿元。 公司于 2004-06-09 在上交所首发上市，注册资本 12.50 亿元，总股本 12.61 亿股。 截至 2026-09-24收盘价 5.08 元，流通市值 63.5 亿元，近一年-20.9\%，流通市值在三级行业“生猪养殖”的 11 家公司中排第 \zhnumber{8} 位，近一年涨跌幅低于全样本中位数（-10.1\%）10.8 个百分点。 按 2025 年年报披露的行业构成，收入占比前三位为畜牧业 84.3\%；批发和零售 14.8\%；饲料加工 0.8\%。 与 2021 年年报相比，第一大行业由“批发和零售”变为“畜牧业”，主业重心已经迁移。 发展过程中的关键节点包括：11 月 5 日完成 2021 年度非公开发行股票，总股本由 652,675,584 股……（2021）；实施重大资产重组，通过发行股份及支付现金购买湖南天心种业有限公司 100\% 股权和 20……（2022）；7 月间接控股股东湖南农业发展投资集团有限责任公司成立，由原间接控股股东现代农业集团和湖……（2022）。\footnote{\texttt{https://stockn.xueqiu.com/SH600975/20220427999843.pdf}}\footnote{\texttt{https://pdf.dfcfw.com/pdf/H2\_AN202304241585776835\_1.pdf}}\footnote{\texttt{http://static.cninfo.com.cn/finalpage/2026-04-28/1225213952.PDF}}

\pfl{历史股价} 以 2021 年 1 月为起点、2026 年 9 月为终点，股价由 7.78 元变为 5.08 元，累计 -34.7\%，区间最高 11.00 元（2023 年 4 月）、最低 4.31 元（2026 年 6 月）；当前价相当于区间最高价的 46.2\%，为区间最低价的 1.18 倍。 月度收益的年化波动率 35.7\%，低于全样本中位数 37.4\%，在三级行业“生猪养殖”的 11 家公司中波动率排第 \zhnumber{8} 位。

\begin{center}
\begin{tikzpicture}
\begin{axis}[
  width=12.6cm, height=3.9cm, scale only axis,
  xmin=2020.922, xmax=2026.888, ymin=4.01, ymax=11.77,
  xtick={2021,2022,2023,2024,2025,2026},
  yearticks,
  yticklabel style={font=\scriptsize},
  ylabel={元/股}, ylabel style={font=\scriptsize},
  grid=major, grid style={LightGray, line width=0.3pt},
  axis line style={InkGray, line width=0.5pt},
  every axis plot/.append style={line width=0.9pt},
]
\addplot[AgriGreen] table[x=x,y=y]{data/clean/profiles/px_600975.dat};
\addplot[SoftRed, dashed, line width=0.7pt] table[x=x,y=y]{data/clean/profiles/pxzz_600975.dat};
\addplot[only marks, mark=*, mark size=1.1pt, SoftRed] table[x=x,y=y]{data/clean/profiles/pxzz_600975.dat};
\end{axis}
\end{tikzpicture}

\captionof{figure}[新五丰（600975）股价与周期骨架]{新五丰（600975）月度收盘价（元/股）与 ZigZag 周期骨架（阈值 25\%，圆点为周期转折点）}
\end{center}

\pfl{过去周期分析} 以 25\% 的 ZigZag 阈值划分，样本区间内股价形成 2 个主升段与 3 个主跌段。 最大主升段出现在 2021 年 7 月 至 2022 年 3 月，8 个月+70.7\%；最大主跌段为 2023 年 4 月 至 2026 年 6 月，38 个月-60.8\%。 最近一次转折出现在 2026 年 6 月（下跌段自 2023 年 4 月起，38 个月-60.8\%），当前处于该段的延续之中。 公司股价与申万农林牧渔指数几乎同步见底，个股并未走出独立行情。 生猪养殖的定价主线是产能去化与成本曲线的赛跑（第 \ref{sec:hog} 章）：2026 年三季度行业仍处在“产能去化未完成、价格磨底”的阶段，公司 2026 年 6 月末资产负债率 84.2\%，高于全样本中位数 52.2\%，其股价因此更多反映资金面与信用风险，而不是价格弹性本身。

\pfl{盈利情况} 2026 年上半年营业收入 28.33 亿元（同比 -23.2\%），归母净利润 -7.36 亿元，由上年同期的盈利转为亏损。 以年报为趋势对照，2023---2025 年营业收入由 56.32 亿元变为 72.58 亿元（两年年均复合增速 +13.5\%），归母净利润由 -12.03 亿元变为 -8.76 亿元。 按半年报口径，2026 年上半年的 ROE 为 -41.0\%（未年化）、销售净利率 -26.5\%、毛利率 -5.4\%，ROE 在“生猪养殖”11 家公司中排第 \zhnumber{10} 位。 同期全样本半年 ROE 中位数 0.75\%，公司与之相差 41.75 个百分点（弱于中位数公司）。 结合公司披露，报告期内利润总额 -750,208,689.22 元，业绩下滑主要系生猪销售价格持续低迷；公司援引国家统……；集中资源提升优质产能效率，主动出清高胎龄及低效母猪以优化种群结构；截至 6 月末能繁母猪存栏量由 202……。\footnote{\texttt{https://stock.stockstar.com/RB2026090900020206.shtml}}\footnote{\texttt{https://www.stcn.com/article/detail/4162760.html}}

\begin{center}\small\setlength{\tabcolsep}{4pt}
\captionof{table}[新五丰（600975）关键财务指标]{新五丰（600975）关键财务指标（合并报表口径；两个半年报期均未年化；现金短债比 $=$ 货币资金 $\div$ 短期债务）}
\begin{tabular}{@{}lrrrrr@{}}
\toprule
指标 & 2023 年 & 2024 年 & 2025 年 & 2025 年半年报 & 2026 年半年报 \\
\midrule
营业收入（亿元） & 56.32 & 70.45 & 72.58 & 36.88 & 28.33 \\
归母净利润（亿元） & -12.03 & 0.39 & -8.76 & 0.45 & -7.36 \\
毛利率（\%） & -9.2 & 7.4 & 1.0 & 9.5 & -5.4 \\
ROE（\%） & -42.4 & 1.3 & -33.9 & 1.5 & -41.0 \\
资产负债率（\%） & 75.0 & 74.0 & 78.3 & 72.5 & 84.2 \\
经营活动现金流净额（亿元） & -6.42 & 8.57 & 5.98 & 6.06 & -3.46 \\
货币资金（亿元） & 15.00 & 11.70 & 7.80 & 11.38 & 6.08 \\
有息负债合计（亿元） & 83.49 & 79.19 & 70.34 & 73.86 & 73.49 \\
现金短债比（倍） & 0.64 & 0.66 & 0.43 & 0.76 & 0.50 \\
\bottomrule
\end{tabular}
\end{center}

\begin{center}
\begin{tikzpicture}
\begin{groupplot}[
  group style={group size=2 by 1, horizontal sep=0.95cm},
  width=7.2cm, height=4.6cm,
  tick label style={font=\tiny},
  title style={font=\scriptsize\heiti},
  yticklabel style={font=\tiny},
  grid=major, grid style={LightGray, line width=0.3pt},
  axis line style={InkGray, line width=0.5pt},
  enlarge x limits={abs=0.8},
  legend style={font=\scriptsize, at={(0.5,-0.22)}, anchor=north,
                legend columns=2, draw=none, fill=none, column sep=6pt},
]
\nextgroupplot[
  ybar, bar width=4pt,
  ymin=-20.50, ymax=82.73,
  xtick={0,1,2,3,4,5.7},
  xticklabels={2021,2022,2023,2024,2025,2026H1},
  ylabel={亿元},
]
\addplot[fill=AgriGreen!75, draw=AgriGreen, bar shift=-2.6pt] table[x=i, y=rev]{data/clean/profiles/fin_600975.dat};
\addplot[fill=SoftRed!60, draw=SoftRed, bar shift=2.6pt] table[x=i, y=ni]{data/clean/profiles/fin_600975.dat};
\addplot[fill=AgriGreen!30, draw=AgriGreen, pattern=north east lines, bar shift=-2.6pt] table[x=x, y=rev]{data/clean/profiles/finh_600975.dat};
\addplot[fill=SoftRed!20, draw=SoftRed, pattern=north east lines, bar shift=2.6pt] table[x=x, y=ni]{data/clean/profiles/finh_600975.dat};
\legend{营业收入（年/半年）, 归母净利润（年/半年）}
\nextgroupplot[
  ymin=-52.5, ymax=96.8,
  xtick={0,1,2,3,4,5.7},
  xticklabels={2021,2022,2023,2024,2025,2026H1},
  ylabel={\%},
]
\addplot[AgriGold, mark=*, mark size=1.3pt] table[x=i, y=roe]{data/clean/profiles/fin_600975.dat};
\addplot[OilBlue, mark=square*, mark size=1.3pt] table[x=i, y=debt]{data/clean/profiles/fin_600975.dat};
\addplot[only marks, AgriGold, mark=*, mark size=2.0pt] table[x=x, y=roe]{data/clean/profiles/finh_600975.dat};
\addplot[only marks, OilBlue, mark=square*, mark size=2.0pt] table[x=x, y=debt]{data/clean/profiles/finh_600975.dat};
\legend{ROE, 资产负债率}
\end{groupplot}
\end{tikzpicture}
\begin{tikzpicture}
\begin{axis}[
  width=15.2cm, height=3.1cm, scale only axis,
  ymin=-15.41, ymax=94.28,
  xtick={0,1,2,3,4,5.7},
  xticklabels={2021,2022,2023,2024,2025,2026H1},
  tick label style={font=\tiny},
  yticklabel style={font=\tiny},
  ylabel={亿元}, ylabel style={font=\scriptsize},
  ybar, bar width=4pt,
  grid=major, grid style={LightGray, line width=0.3pt},
  axis line style={InkGray, line width=0.5pt},
  enlarge x limits={abs=0.8},
  legend style={font=\scriptsize, at={(0.5,-0.30)}, anchor=north,
                legend columns=3, draw=none, fill=none, column sep=10pt},
]
\addplot[fill=AgriGreen!55, draw=AgriGreen, bar shift=-3.0pt] table[x=i, y=cash]{data/clean/profiles/fin_600975.dat};
\addplot[fill=SoftRed!45, draw=SoftRed, bar shift=3.0pt] table[x=i, y=ibd]{data/clean/profiles/fin_600975.dat};
\addplot[fill=AgriGreen!25, draw=AgriGreen, pattern=north east lines, bar shift=-3.0pt] table[x=x, y=cash]{data/clean/profiles/finh_600975.dat};
\addplot[fill=SoftRed!20, draw=SoftRed, pattern=north east lines, bar shift=3.0pt] table[x=x, y=ibd]{data/clean/profiles/finh_600975.dat};
\addplot[OilBlue, mark=*, mark size=1.1pt, line width=0.8pt] table[x=i, y=ocf]{data/clean/profiles/fin_600975.dat};
\addplot[only marks, OilBlue, mark=*, mark size=1.8pt] table[x=x, y=ocf]{data/clean/profiles/finh_600975.dat};
\legend{货币资金, 有息负债, 经营活动现金流净额}
\end{axis}
\end{tikzpicture}

\captionof{figure}[新五丰（600975）近年主要财务指标]{新五丰（600975）近年主要财务指标（左上：营业收入与归母净利润；右上：ROE 与资产负债率；下：货币资金、有息负债与经营活动现金流净额。
2021---2025 年为年报口径，2026H1 为 2026 年半年报/半年末，与其前一年报之间留出间隔、以斜纹或空心点区分；半年报数据未年化）}
\end{center}

\pfl{财务分析} 2026 年 6 月末资产负债率 84.2\%，较 2025 年末 上升 5.9 个百分点（样本中位数 52.2\%，所处三级行业内排第 \zhnumber{10} 位）。 2026 年 6 月末货币资金 6.08 亿元、短期债务（短期借款与一年内到期非流动负债）12.10 亿元、长期债务 61.39 亿元，有息负债合计 73.49 亿元，现金短债比 0.50 倍——覆盖偏紧。 净有息负债 67.40 亿元，相当于其 2026 年上半年营业收入的 238\%。 流动比率 1.13、速动比率 0.38，短期指标尚可。 2026 年上半年经营活动现金流净额 -3.46 亿元，经营现金流与利润同时为负，日常运营对债务与股东投入的依赖度较高。

\pfl{重大战略转型} 近三年公司的战略动作集中在以下几个方面：2022 并购天心种业以解决现代农业集团内部同业竞争问题，夯实种猪繁育体系，提升综合竞争、市场拓展与抗风险能力；2023-07 完成募集配套资金 15.5 亿元，为天心种业 5 个募集资金项目提供 7.11 亿元建设资金，进一步夯实种猪繁育……；2026 集中资源提升优质产能效率，主动出清高胎龄及低效母猪以优化种群结构；截至 6 月末能繁母猪存栏量由 2025 年末……。 公告口径上，近三年（2023 年 9 月---2026 年 9 月）公司披露重大重组与并购类公告 3 条、对外投资与转型类公告 2 条、回购与股权激励类公告 3 条，合计公告 425 条。 主营结构的迁移已经落到报表上：2021 年年报的第一大行业为“批发和零售”（55.6\%），2025 年年报变为“畜牧业”（84.3\%）。 综合判断，公司在报告期内有明确的外延扩张或业务结构调整动作，转型方向与融资安排之间存在直接关联。\footnote{\texttt{https://pdf.dfcfw.com/pdf/H3\_AP202305181586740282\_1.pdf}}\footnote{\texttt{https://www.stcn.com/article/detail/1464080.html}}\footnote{\texttt{https://www.stcn.com/article/detail/4162760.html}}

\begin{center}\small\setlength{\tabcolsep}{4pt}
\captionof{table}[新五丰（600975）历史融资与资本运作]{新五丰（600975）历史融资与资本运作（据公司公告与公开报道整理；金额为披露值，含首次公开发行、再融资、债券、银行借款与重整投资等）}
\begin{tabular}{@{}p{1.5cm}p{3.2cm}p{2.4cm}p{6.3cm}@{}}
\toprule
时间 & 方式 & 规模 & 用途与说明 \\
\midrule
2004-06 & IPO（上海证券交易所） & 发行 3,500 万股、发行价 9.20 元/股，预计募集资金净额 30,512.30 万元 & 发行前净资产 14,611.98 万元，发行后净资产增至 45,124.28 万元…… \\
2021-10-27 & 非公开发行股票（向特定对象发行，4 名发行对象） & 募集资金总额 102,999.00 万元（152,365,383 股，发行价 6.76 元/股） & 湖南新五丰存栏 4.32 万头母猪场项目（拟投入募集资金 51,076.40 万元）及补…… \\
2022 & 发行股份及支付现金购买资产（重大资产重组） & 交易作价合计 220,768.24 万元，其中股份对价 195,978.44 万元、现金对价 24,789.80 万元；发行股份价格 7.22 元/股 & 收购天心种业 100\% 股权（含 200 万元国有独享资本公积）…… \\
2023-07 & 发行股份募集配套资金 & 募集资金总额 1,550,579,366.83 元（184,812,797 股，发行价 8.39 元/股） & 部分用于支付本次重组现金对价和补充流动资金…… \\
\bottomrule
\end{tabular}
\end{center}

\pfl{历史融投资情况} 从现金流量表看，2025 年公司取得借款收到的现金 10.18 亿元、偿还债务支付的现金 15.96 亿元、吸收权益性投资 0.15 亿元，筹资活动现金流净额 -8.67 亿元（2023 年为取得借款 29.57 亿元、偿还债务 20.80 亿元）。 2026 年上半年筹资活动现金流净额为 2.15 亿元（取得借款 15.06 亿元、偿还债务 7.90 亿元，未年化），仍处于净融入状态。 2025 年分配股利、利润或偿付利息支付的现金合计 0.54 亿元。 公开信息中与公司直接相关的融资与资本运作包括：IPO（上海证券交易所）、发行股份募集配套资金、发行股份及支付现金购买资产（重大资产重组）、非公开发行股票（向特定对象发行，4 名发行对象）。 其中规模较大的两笔为：2004-06IPO（上海证券交易所）（发行 3,500 万股、发行价 9.20 元/股，预计募集资……）；2022发行股份及支付现金购买资产（重大资产重组）（交易作价合计 220,768.24 万元，其中股份对价 19……）。 股本方面，近三年披露股本变动 6 次（增发新股上市 2 次、限售股份上市 1 次）；总股本由 2021 年末的 8.05 亿股变为 12.61 亿股（+56.67\%）。 分红方面，近三年未实施现金分红，在全部样本中属于分红能力偏弱的一端。 近三年“再融资”类公告 4 条，涉及定向增发、发行股份购买资产。\footnote{\texttt{http://vip.stock.finance.sina.com.cn/corp/view/vISSUE\_RaiseExplanationDetail.php}}\footnote{\texttt{https://www.stcn.com/article/detail/1401109.html}}\footnote{\texttt{http://file.finance.sina.com.cn/211.154.219.97:9494/MRGG/CNSESH\_STOCK/2022/2022-12/2022-}}

\pfl{未来潜在融资需求分析} 2026 年 6 月末货币资金 6.08 亿元，而短期债务 12.10 亿元（现金短债比 0.50 倍），2026 年上半年经营现金流净额 -3.46 亿元。按“短期债务减去账面货币资金”的滚动偿付口径，静态缺口约 6.02 亿元；若再计入上半年亏损的年化额（14.72 亿元），维持一年正常经营所需的外部资金约 6.02---20.74 亿元。 按第 \ref{sec:financing-need} 节的三档划分，公司属于第一档“补血型”：2026 年上半年归母净利润 -7.36 亿元、6 月末资产负债率 84.2\%。负债率已超过 80\%，净资产被亏损侵蚀的程度较深，属于全部样本中资产负债表最紧张的一组。 可行的融资路径按现实性排序为：出售非核心资产与参股股权、定向增发（需股价配合，且控股股东需放弃优先认购或同步增资）、引入国资或产业投资人（部分公司以控制权让渡为对价）。 公开披露的下一步融资线索包括：2025-08-23 第六届董事会第十八次会议审议原“湖南新五丰存栏 4.32 万头母猪场项目”相关事项；该项目投资总额 106,028.24……；2026-01-19 第六届董事会第二十二次会议审议通过《关于使用闲置募集资金临时补充流动资金的议案》，同意使用闲置募集资金不超过 9,000……。\footnote{\texttt{http://epaper.stcn.com/pic/202508/23/de2f70afa08f29e9f46b40cf63b511cf.pdf}}\footnote{\texttt{https://paper.cnstock.com/html/2026-08/31/content\_2263279.htm}}


\subsubsection{华英农业（002321）}
\label{co:002321}

\pfl{公司基本情况} 公司前身可追溯至 1991 年，河南信阳潢川的鸭肉全产业链企业（曾称“世界鸭王”），2009 年深交所中小板上市，2020 年因大额亏损、控股股东非经营性占用资金及无法表示意见的审计报告被实施退市风险警示，2021 年完成破产重整、2022 年摘帽，现聚焦食品与羽绒双主业，2026 年上半年归母净亏损 2,928.25 万元。 公司于 2009-12-16 在深交所主板首发上市，注册资本 21.33 亿元，总股本 21.33 亿股。 截至 2026-09-24收盘价 2.27 元，流通市值 48.2 亿元，近一年-23.8\%，流通市值在三级行业“其他养殖”的 3 家公司中排第 \zhnumber{1} 位，近一年涨跌幅低于全样本中位数（-10.1\%）13.7 个百分点。 按 2026 年半年报披露的行业构成，收入占比前三位为畜牧业 100.0\%。 发展过程中的关键节点包括：12 月 16 日在深交所上市（发行主体后更名为河南华英农业发展股份有限公司），经证监许……（2009）；因 2020 年度财务会计报告被亚太（集团）会计师事务所出具无法表示意见的审计报告，深交……（2021）；4 月 14 日信阳中院作出（2021）豫 15 破 6-4 号《民事裁定书》，裁定确认……（2022）。\footnote{\texttt{https://vip.stock.finance.sina.com.cn/corp/go.php/vISSUE\_NewStock/stockid/002321.phtml}}\footnote{\texttt{http://disc.static.szse.cn/download/disc/disk03/finalpage/2022-06-24/5ac7fee4-7175-4ec4-}}\footnote{\texttt{http://vip.stock.finance.sina.com.cn/corp/view/vCB\_AllBulletinDetail.php}}

\pfl{历史股价} 自 2021 年 1 月至 2026 年 9 月，股价由 3.20 元变为 2.27 元，累计 -29.1\%，区间最高 3.96 元（2021 年 11 月）、最低 1.43 元（2024 年 6 月）；当前价相当于区间最高价的 57.3\%，为区间最低价的 1.59 倍。 月度收益的年化波动率 35.4\%，低于全样本中位数 37.4\%，在三级行业“其他养殖”的 3 家公司中波动率排第 \zhnumber{3} 位。

\begin{center}
\begin{tikzpicture}
\begin{axis}[
  width=12.6cm, height=3.9cm, scale only axis,
  xmin=2020.922, xmax=2026.888, ymin=1.33, ymax=4.24,
  xtick={2021,2022,2023,2024,2025,2026},
  yearticks,
  yticklabel style={font=\scriptsize},
  ylabel={元/股}, ylabel style={font=\scriptsize},
  grid=major, grid style={LightGray, line width=0.3pt},
  axis line style={InkGray, line width=0.5pt},
  every axis plot/.append style={line width=0.9pt},
]
\addplot[AgriGreen] table[x=x,y=y]{data/clean/profiles/px_002321.dat};
\addplot[SoftRed, dashed, line width=0.7pt] table[x=x,y=y]{data/clean/profiles/pxzz_002321.dat};
\addplot[only marks, mark=*, mark size=1.1pt, SoftRed] table[x=x,y=y]{data/clean/profiles/pxzz_002321.dat};
\end{axis}
\end{tikzpicture}

\captionof{figure}[华英农业（002321）股价与周期骨架]{华英农业（002321）月度收盘价（元/股）与 ZigZag 周期骨架（阈值 25\%，圆点为周期转折点）}
\end{center}

\pfl{过去周期分析} 以 25\% 的 ZigZag 阈值划分，样本区间内股价形成 4 个主升段与 4 个主跌段。 最大主升段出现在 2024 年 6 月 至 2025 年 9 月，15 个月+107.7\%；最大主跌段为 2023 年 5 月 至 2024 年 6 月，13 个月-50.0\%。 最近一次转折出现在 2026 年 9 月（上涨段自 2026 年 6 月起，3 个月+35.9\%），当前处于该段之后的震荡之中。 公司股价较申万农林牧渔指数提前约 23 个月见底，是板块中较早定价周期反转的公司之一。 牛、羊、水产的产能周期长于生猪，第 \ref{sec:other-livestock} 章判断肉牛供给的实质性收缩要到 2027 年之后才能兑现，这一时滞决定了个股的股价弹性，公司 2026 年 6 月末资产负债率 66.2\%，高于全样本中位数 52.2\%，其股价因此更多反映资金面与信用风险，而不是价格弹性本身。

\pfl{盈利情况} 2026 年上半年营业收入 23.72 亿元（同比 +14.5\%），归母净利润 -0.29 亿元，亏损同比扩大 0.09 亿元。 以年报为趋势对照，2023---2025 年营业收入由 37.05 亿元变为 54.47 亿元（两年年均复合增速 +21.3\%），归母净利润由 -0.32 亿元变为 -0.60 亿元。 按半年报口径，2026 年上半年的 ROE 为 -3.2\%（未年化）、销售净利率 -0.2\%、毛利率 5.4\%，ROE 在“其他养殖”3 家公司中排第 \zhnumber{3} 位。 同期全样本半年 ROE 中位数 0.75\%，公司与之相差 3.92 个百分点（弱于中位数公司）。 结合公司披露，2026 年上半年营业收入 237,192.89 万元，同比增长 14.52\%；归属于上市公司股东的净利……；以“食品、羽绒双主业”为核心，持续优化产业链协同效率、深化产品结构调整，通过强化成本管控、拓展高附加值产……。\footnote{\texttt{http://vip.stock.finance.sina.com.cn/corp/view/vCB\_AllBulletinDetail.php}}\footnote{\texttt{http://vip.stock.finance.sina.com.cn/corp/view/vCB\_AllBulletinDetail.php}}

\begin{center}\small\setlength{\tabcolsep}{4pt}
\captionof{table}[华英农业（002321）关键财务指标]{华英农业（002321）关键财务指标（合并报表口径；两个半年报期均未年化；现金短债比 $=$ 货币资金 $\div$ 短期债务）}
\begin{tabular}{@{}lrrrrr@{}}
\toprule
指标 & 2023 年 & 2024 年 & 2025 年 & 2025 年半年报 & 2026 年半年报 \\
\midrule
营业收入（亿元） & 37.05 & 47.31 & 54.47 & 20.71 & 23.72 \\
归母净利润（亿元） & -0.32 & 0.32 & -0.60 & -0.21 & -0.29 \\
毛利率（\%） & 9.1 & 7.5 & 5.8 & 6.8 & 5.4 \\
ROE（\%） & -3.2 & 3.3 & -6.4 & -2.2 & -3.2 \\
资产负债率（\%） & 62.1 & 61.1 & 63.4 & 64.8 & 66.2 \\
经营活动现金流净额（亿元） & -1.33 & 2.00 & 3.76 & 1.86 & -2.96 \\
货币资金（亿元） & 2.42 & 2.75 & 4.14 & 3.29 & 1.92 \\
有息负债合计（亿元） & 3.99 & 4.73 & 5.46 & 4.91 & 6.13 \\
现金短债比（倍） & 1.06 & 0.67 & 0.86 & 0.78 & 0.35 \\
\bottomrule
\end{tabular}
\end{center}

\begin{center}
\begin{tikzpicture}
\begin{groupplot}[
  group style={group size=2 by 1, horizontal sep=0.95cm},
  width=7.2cm, height=4.6cm,
  tick label style={font=\tiny},
  title style={font=\scriptsize\heiti},
  yticklabel style={font=\tiny},
  grid=major, grid style={LightGray, line width=0.3pt},
  axis line style={InkGray, line width=0.5pt},
  enlarge x limits={abs=0.8},
  legend style={font=\scriptsize, at={(0.5,-0.22)}, anchor=north,
                legend columns=2, draw=none, fill=none, column sep=6pt},
]
\nextgroupplot[
  ybar, bar width=4pt,
  ymin=-33.37, ymax=64.05,
  xtick={0,1,2,3,4,5.7},
  xticklabels={2021,2022,2023,2024,2025,2026H1},
  ylabel={亿元},
]
\addplot[fill=AgriGreen!75, draw=AgriGreen, bar shift=-2.6pt] table[x=i, y=rev]{data/clean/profiles/fin_002321.dat};
\addplot[fill=SoftRed!60, draw=SoftRed, bar shift=2.6pt] table[x=i, y=ni]{data/clean/profiles/fin_002321.dat};
\addplot[fill=AgriGreen!30, draw=AgriGreen, pattern=north east lines, bar shift=-2.6pt] table[x=x, y=rev]{data/clean/profiles/finh_002321.dat};
\addplot[fill=SoftRed!20, draw=SoftRed, pattern=north east lines, bar shift=2.6pt] table[x=x, y=ni]{data/clean/profiles/finh_002321.dat};
\legend{营业收入（年/半年）, 归母净利润（年/半年）}
\nextgroupplot[
  ymin=-1902.7, ymax=248.5,
  xtick={0,1,2,3,4,5.7},
  xticklabels={2021,2022,2023,2024,2025,2026H1},
  ylabel={\%},
]
\addplot[AgriGold, mark=*, mark size=1.3pt] table[x=i, y=roe]{data/clean/profiles/fin_002321.dat};
\addplot[OilBlue, mark=square*, mark size=1.3pt] table[x=i, y=debt]{data/clean/profiles/fin_002321.dat};
\addplot[only marks, AgriGold, mark=*, mark size=2.0pt] table[x=x, y=roe]{data/clean/profiles/finh_002321.dat};
\addplot[only marks, OilBlue, mark=square*, mark size=2.0pt] table[x=x, y=debt]{data/clean/profiles/finh_002321.dat};
\legend{ROE, 资产负债率}
\end{groupplot}
\end{tikzpicture}
\begin{tikzpicture}
\begin{axis}[
  width=15.2cm, height=3.1cm, scale only axis,
  ymin=-8.61, ymax=11.14,
  xtick={0,1,2,3,4,5.7},
  xticklabels={2021,2022,2023,2024,2025,2026H1},
  tick label style={font=\tiny},
  yticklabel style={font=\tiny},
  ylabel={亿元}, ylabel style={font=\scriptsize},
  ybar, bar width=4pt,
  grid=major, grid style={LightGray, line width=0.3pt},
  axis line style={InkGray, line width=0.5pt},
  enlarge x limits={abs=0.8},
  legend style={font=\scriptsize, at={(0.5,-0.30)}, anchor=north,
                legend columns=3, draw=none, fill=none, column sep=10pt},
]
\addplot[fill=AgriGreen!55, draw=AgriGreen, bar shift=-3.0pt] table[x=i, y=cash]{data/clean/profiles/fin_002321.dat};
\addplot[fill=SoftRed!45, draw=SoftRed, bar shift=3.0pt] table[x=i, y=ibd]{data/clean/profiles/fin_002321.dat};
\addplot[fill=AgriGreen!25, draw=AgriGreen, pattern=north east lines, bar shift=-3.0pt] table[x=x, y=cash]{data/clean/profiles/finh_002321.dat};
\addplot[fill=SoftRed!20, draw=SoftRed, pattern=north east lines, bar shift=3.0pt] table[x=x, y=ibd]{data/clean/profiles/finh_002321.dat};
\addplot[OilBlue, mark=*, mark size=1.1pt, line width=0.8pt] table[x=i, y=ocf]{data/clean/profiles/fin_002321.dat};
\addplot[only marks, OilBlue, mark=*, mark size=1.8pt] table[x=x, y=ocf]{data/clean/profiles/finh_002321.dat};
\legend{货币资金, 有息负债, 经营活动现金流净额}
\end{axis}
\end{tikzpicture}

\captionof{figure}[华英农业（002321）近年主要财务指标]{华英农业（002321）近年主要财务指标（左上：营业收入与归母净利润；右上：ROE 与资产负债率；下：货币资金、有息负债与经营活动现金流净额。
2021---2025 年为年报口径，2026H1 为 2026 年半年报/半年末，与其前一年报之间留出间隔、以斜纹或空心点区分；半年报数据未年化）}
\end{center}

\pfl{财务分析} 2026 年 6 月末资产负债率 66.2\%，较 2025 年末 上升 2.8 个百分点（样本中位数 52.2\%，所处三级行业内排第 \zhnumber{2} 位）。 2026 年 6 月末货币资金 1.92 亿元、短期债务（短期借款与一年内到期非流动负债）5.52 亿元、长期债务 0.61 亿元，有息负债合计 6.13 亿元，现金短债比 0.35 倍——覆盖偏紧。 净有息负债 4.21 亿元，相当于其 2026 年上半年营业收入的 18\%。 流动比率 1.49、速动比率 0.77，短期指标尚可。 2026 年上半年经营活动现金流净额 -2.96 亿元，经营现金流与利润同时为负，日常运营对债务与股东投入的依赖度较高。 2026 年 6 月末商誉 1.12 亿元，占归母净资产的 7.9\%，是净资产中最脆弱的一块。 公司披露的财务与合规风险包括：光大证券对控股股东所持的部分公司股票进行违约处置，以集中竞价方式减持 2,016,400……。\footnote{\texttt{http://vip.stock.finance.sina.com.cn/corp/view/vCB\_AllBulletinDetail.php}}

\pfl{重大战略转型} 从公开披露看，公司的战略推进包括：2022 重整落定后提出冲刺“世界鸭王+世界羽王”双产业目标；2022 公司启动低效资产处置程序，拟处置评估价值为 5.7 亿元的长期股权投资以及评估值为 2.77 亿元的其他应收款资……；2026 以“食品、羽绒双主业”为核心，持续优化产业链协同效率、深化产品结构调整，通过强化成本管控、拓展高附加值产品销路、……。 公告口径上，近三年（2023 年 9 月---2026 年 9 月）公司披露重大重组与并购类公告 0 条、对外投资与转型类公告 0 条、回购与股权激励类公告 8 条，合计公告 339 条。 第一大行业仍是“畜牧业”，收入占比 100.0\%，与 2021 年年报相比没有方向性变化。 综合判断，公司当前的战略重心不是扩张而是化解债务与完成重整，业务层面的调整让位于司法重整程序。\footnote{\texttt{https://m.bjnews.com.cn/detail/1650525774169799.html}}\footnote{\texttt{http://vip.stock.finance.sina.com.cn/corp/view/vCB\_AllBulletinDetail.php}}

\begin{center}\small\setlength{\tabcolsep}{4pt}
\captionof{table}[华英农业（002321）历史融资与资本运作]{华英农业（002321）历史融资与资本运作（据公司公告与公开报道整理；金额为披露值，含首次公开发行、再融资、债券、银行借款与重整投资等）}
\begin{tabular}{@{}p{1.5cm}p{3.2cm}p{2.4cm}p{6.3cm}@{}}
\toprule
时间 & 方式 & 规模 & 用途与说明 \\
\midrule
2009-12 & IPO（深圳证券交易所中小板） & 3,700 万股，发行价 16.98 元/股，实际募集资金 62,826.00 万元 & 依序投向淮阳 4000 万只/年商品肉鸡系列加工生产项目…… \\
2021-12 & 破产重整——重整投资人受让转增股票出资 & 产业投资人鼎新兴华以人民币约 3.62 亿元有条件受让华英农业转增股票约 5.12 亿股，并解决 4,311.25 万元控股股东非经营性资金占用问题 & 受让转增股票、解决控股股东资金占用；2021 年 12 月 23 日公司与鼎新兴华…… \\
2021-12 & 资本公积转增股本（重整出资人权益调整） & 以 2021 年 12 月 31 日总股本 53,429.11 万股为基数，按每 10 股转增 29.92 股实施资本公积转增，共转增 159,859.90 万股，总股本增至 213,289.00 万股；转增股票不向原股东分配 & 其中 531,963,958 股以股抵债形式清偿公司债务…… \\
\bottomrule
\end{tabular}
\end{center}

\pfl{历史融投资情况} 从现金流量表看，2025 年公司取得借款收到的现金 5.61 亿元、偿还债务支付的现金 5.21 亿元、吸收权益性投资 0.00 亿元，筹资活动现金流净额 -2.88 亿元（2023 年为取得借款 3.11 亿元、偿还债务 2.35 亿元）。 2026 年上半年筹资活动现金流净额为 1.97 亿元（取得借款 1.83 亿元、偿还债务 0.92 亿元，未年化），仍处于净融入状态。 2025 年分配股利、利润或偿付利息支付的现金合计 0.82 亿元。 公开信息中与公司直接相关的融资与资本运作包括：IPO（深圳证券交易所中小板）、破产重整——重整投资人受让转增股票出资、资本公积转增股本（重整出资人权益调整）。 其中规模较大的两笔为：2021-12破产重整——重整投资人受让转增股票出资（产业投资人鼎新兴华以人民币约 3.62 亿元有条件受让华英农……）；2009-12IPO（深圳证券交易所中小板）（3,700 万股，发行价 16.98 元/股，实际募集资金 ……）。 股本方面，近三年披露股本变动 3 次；总股本由 2021 年末的 21.33 亿股变为 21.33 亿股（+0.00\%）。 分红方面，近三年未实施现金分红，与重整期间的资金约束一致。 近三年“再融资”类公告 1 条，涉及定向增发。\footnote{\texttt{https://vip.stock.finance.sina.com.cn/corp/view/vISSUE\_RaiseExplanationDetail.php}}\footnote{\texttt{https://www.zhitongcaijing.com/content/detail/628357.html}}\footnote{\texttt{https://q.stock.sohu.com/newpdf/202561987343.pdf}}

\pfl{未来潜在融资需求分析} 2026 年 6 月末货币资金 1.92 亿元，而短期债务 5.52 亿元（现金短债比 0.35 倍），2026 年上半年经营现金流净额 -2.96 亿元。按“短期债务减去账面货币资金”的滚动偿付口径，静态缺口约 3.60 亿元；若再计入上半年亏损的年化额（0.59 亿元），维持一年正常经营所需的外部资金约 3.60---4.18 亿元。 按第 \ref{sec:financing-need} 节的三档划分，公司属于第一档“补血型”：2026 年上半年归母净利润 -0.29 亿元、6 月末资产负债率 66.2\%。 可行的融资路径按现实性排序为：引入国资或产业投资人（部分公司以控制权让渡为对价）、出售非核心资产与参股股权、定向增发（需股价配合，且控股股东需放弃优先认购或同步增资）。 公开披露的下一步融资线索包括：2026-04-10 控股股东鼎新兴华全体合伙人一致同意解散并将所持华英农业股份按各合伙人实缴出资比例通过非交易过户方式分配给全体合伙人；分配……。\footnote{\texttt{http://vip.stock.finance.sina.com.cn/corp/view/vCB\_AllBulletinDetail.php}}


\subsubsection{仙坛股份（002746）}
\label{co:002746}

\pfl{公司基本情况} 山东烟台白羽肉鸡养殖屠宰一体化龙头企业，2001 年成立、2015 年 2 月深交所上市，商品代肉鸡以“公司+农户”委托养殖为主，2020 年定增逾 10 亿元投建诸城年产 1.2 亿羽肉鸡产业生态项目；2026 年上半年营业收入 29.95 亿元、归母净利润 1.80 亿元，同比增长约三成。 公司于 2015-02-16 在深交所主板首发上市，注册资本 8.61 亿元，总股本 8.61 亿股。 截至 2026-09-24收盘价 6.19 元，流通市值 43.9 亿元，近一年-1.3\%，流通市值在三级行业“肉鸡养殖”的 8 家公司中排第 \zhnumber{4} 位，近一年涨跌幅高于全样本中位数（-10.1\%）8.8 个百分点。 按 2026 年半年报披露的行业构成，收入占比前三位为肉鸡养殖及屠宰加工 100.0\%。 发展过程中的关键节点包括：1 月 25 日由山东仙坛集团有限公司整体变更为山东仙坛股份有限公司，发起人为王寿纯、曲……（2011）；2 月 16 日在深圳证券交易所上市，经证监许可[2015]189 号文核准公开发行 3……（2015）。\footnote{\texttt{https://pdf.dfcfw.com/pdf/H2\_AN201502050008310568\_1.pdf}}\footnote{\texttt{https://vip.stock.finance.sina.com.cn/corp/view/vISSUE\_MarketBulletinDetail.php}}

\pfl{历史股价} 本报告样本区间（2021 年 1 月 至 2026 年 9 月）内，股价由 10.04 元变为 6.19 元，累计 -38.3\%，区间最高 10.98 元（2021 年 3 月）、最低 5.10 元（2026 年 6 月）；当前价相当于区间最高价的 56.4\%，为区间最低价的 1.21 倍。 月度收益的年化波动率 37.6\%，高于全样本中位数 37.4\%，在三级行业“肉鸡养殖”的 8 家公司中波动率排第 \zhnumber{4} 位。 从事件看，2021 年 12 月 至 2026 年 6 月的下跌（-51.4\%）与公司诸城年产 1.2 亿羽肉鸡产业生态项目一期配套的屠宰加工厂、饲料厂和商品鸡养……（2026-06-30）在时间上重合——股价的拐点多数能在公司自身的资本运作与风险事件上找到对应。\footnote{\texttt{https://stock.stockstar.com/RB2026082700056657.shtml}}

\begin{center}
\begin{tikzpicture}
\begin{axis}[
  width=12.6cm, height=3.9cm, scale only axis,
  xmin=2020.922, xmax=2026.888, ymin=4.74, ymax=11.75,
  xtick={2021,2022,2023,2024,2025,2026},
  yearticks,
  yticklabel style={font=\scriptsize},
  ylabel={元/股}, ylabel style={font=\scriptsize},
  grid=major, grid style={LightGray, line width=0.3pt},
  axis line style={InkGray, line width=0.5pt},
  every axis plot/.append style={line width=0.9pt},
]
\addplot[AgriGreen] table[x=x,y=y]{data/clean/profiles/px_002746.dat};
\addplot[SoftRed, dashed, line width=0.7pt] table[x=x,y=y]{data/clean/profiles/pxzz_002746.dat};
\addplot[only marks, mark=*, mark size=1.1pt, SoftRed] table[x=x,y=y]{data/clean/profiles/pxzz_002746.dat};
\end{axis}
\end{tikzpicture}

\captionof{figure}[仙坛股份（002746）股价与周期骨架]{仙坛股份（002746）月度收盘价（元/股）与 ZigZag 周期骨架（阈值 25\%，圆点为周期转折点）}
\end{center}

\pfl{过去周期分析} 以 25\% 的 ZigZag 阈值划分，样本区间内股价形成 2 个主升段与 2 个主跌段。 最大主跌段为 2021 年 12 月 至 2026 年 6 月，54 个月-51.4\%；最大主升段出现在 2021 年 11 月 至 2021 年 12 月，1 个月+57.7\%。 最近一次转折出现在 2026 年 6 月（下跌段自 2021 年 12 月起，54 个月-51.4\%），当前处于该段的延续之中。 公司股价与申万农林牧渔指数几乎同步见底，个股并未走出独立行情。 禽链的产能周期只有 6---12 个月、反转快于生猪（第 \ref{sec:other-livestock} 节），2026 年前四个月毛鸡养殖利润同比大幅回升，链条利润的修复先于生猪，公司 2026 年上半年实现盈利、半年 ROE 3.9\%（未年化），好于全样本中位数（0.75\%），下行周期对其报表的冲击小于同业。

\pfl{盈利情况} 2026 年上半年营业收入 29.95 亿元（同比 +17.9\%），归母净利润 1.80 亿元，同比 +32.0\%。 以年报为趋势对照，2023---2025 年营业收入由 57.81 亿元变为 57.72 亿元（两年年均复合增速 -0.1\%），归母净利润由 2.26 亿元变为 2.43 亿元。 按半年报口径，2026 年上半年的 ROE 为 3.9\%（未年化）、销售净利率 6.2\%、毛利率 8.7\%，ROE 在“肉鸡养殖”8 家公司中排第 \zhnumber{3} 位。 同期全样本半年 ROE 中位数 0.75\%，公司与之相差 3.12 个百分点（略好于中位数公司）。 结合公司披露，2026 年上半年经营活动产生的现金流量净额仅 6,006.74 万元，同比大幅下降 82.17\%；短期……；媒体报道公司 2020 年发起的定增募投项目“年产 1.2 亿羽肉鸡产业生态项目”近几年持续亏损，交付时……。\footnote{\texttt{https://stock.stockstar.com/RB2026082700056657.shtml}}\footnote{\texttt{https://news.qq.com/rain/a/20260521A05RGM00}}

\begin{center}\small\setlength{\tabcolsep}{4pt}
\captionof{table}[仙坛股份（002746）关键财务指标]{仙坛股份（002746）关键财务指标（合并报表口径；两个半年报期均未年化；现金短债比 $=$ 货币资金 $\div$ 短期债务）}
\begin{tabular}{@{}lrrrrr@{}}
\toprule
指标 & 2023 年 & 2024 年 & 2025 年 & 2025 年半年报 & 2026 年半年报 \\
\midrule
营业收入（亿元） & 57.81 & 52.91 & 57.72 & 25.40 & 29.95 \\
归母净利润（亿元） & 2.26 & 2.06 & 2.43 & 1.37 & 1.80 \\
毛利率（\%） & 3.8 & 5.2 & 6.5 & 7.9 & 8.7 \\
ROE（\%） & 5.0 & 4.5 & 5.3 & 2.9 & 3.9 \\
资产负债率（\%） & 31.4 & 38.6 & 41.8 & 41.8 & 43.0 \\
经营活动现金流净额（亿元） & 3.30 & 1.55 & 6.79 & 3.37 & 0.60 \\
货币资金（亿元） & 25.15 & 39.01 & 36.01 & 39.16 & 36.20 \\
有息负债合计（亿元） & 10.55 & 18.60 & 23.83 & 22.08 & 25.92 \\
现金短债比（倍） & 2.77 & 2.38 & 1.61 & 1.90 & 1.43 \\
\bottomrule
\end{tabular}
\end{center}

\begin{center}
\begin{tikzpicture}
\begin{groupplot}[
  group style={group size=2 by 1, horizontal sep=0.95cm},
  width=7.2cm, height=4.6cm,
  tick label style={font=\tiny},
  title style={font=\scriptsize\heiti},
  yticklabel style={font=\tiny},
  grid=major, grid style={LightGray, line width=0.3pt},
  axis line style={InkGray, line width=0.5pt},
  enlarge x limits={abs=0.8},
  legend style={font=\scriptsize, at={(0.5,-0.22)}, anchor=north,
                legend columns=2, draw=none, fill=none, column sep=6pt},
]
\nextgroupplot[
  ybar, bar width=4pt,
  ymin=-5.78, ymax=64.75,
  xtick={0,1,2,3,4,5.7},
  xticklabels={2021,2022,2023,2024,2025,2026H1},
  ylabel={亿元},
]
\addplot[fill=AgriGreen!75, draw=AgriGreen, bar shift=-2.6pt] table[x=i, y=rev]{data/clean/profiles/fin_002746.dat};
\addplot[fill=SoftRed!60, draw=SoftRed, bar shift=2.6pt] table[x=i, y=ni]{data/clean/profiles/fin_002746.dat};
\addplot[fill=AgriGreen!30, draw=AgriGreen, pattern=north east lines, bar shift=-2.6pt] table[x=x, y=rev]{data/clean/profiles/finh_002746.dat};
\addplot[fill=SoftRed!20, draw=SoftRed, pattern=north east lines, bar shift=2.6pt] table[x=x, y=ni]{data/clean/profiles/finh_002746.dat};
\legend{营业收入（年/半年）, 归母净利润（年/半年）}
\nextgroupplot[
  ymin=-3.4, ymax=47.3,
  xtick={0,1,2,3,4,5.7},
  xticklabels={2021,2022,2023,2024,2025,2026H1},
  ylabel={\%},
]
\addplot[AgriGold, mark=*, mark size=1.3pt] table[x=i, y=roe]{data/clean/profiles/fin_002746.dat};
\addplot[OilBlue, mark=square*, mark size=1.3pt] table[x=i, y=debt]{data/clean/profiles/fin_002746.dat};
\addplot[only marks, AgriGold, mark=*, mark size=2.0pt] table[x=x, y=roe]{data/clean/profiles/finh_002746.dat};
\addplot[only marks, OilBlue, mark=square*, mark size=2.0pt] table[x=x, y=debt]{data/clean/profiles/finh_002746.dat};
\legend{ROE, 资产负债率}
\end{groupplot}
\end{tikzpicture}
\begin{tikzpicture}
\begin{axis}[
  width=15.2cm, height=3.1cm, scale only axis,
  ymin=-4.87, ymax=43.80,
  xtick={0,1,2,3,4,5.7},
  xticklabels={2021,2022,2023,2024,2025,2026H1},
  tick label style={font=\tiny},
  yticklabel style={font=\tiny},
  ylabel={亿元}, ylabel style={font=\scriptsize},
  ybar, bar width=4pt,
  grid=major, grid style={LightGray, line width=0.3pt},
  axis line style={InkGray, line width=0.5pt},
  enlarge x limits={abs=0.8},
  legend style={font=\scriptsize, at={(0.5,-0.30)}, anchor=north,
                legend columns=3, draw=none, fill=none, column sep=10pt},
]
\addplot[fill=AgriGreen!55, draw=AgriGreen, bar shift=-3.0pt] table[x=i, y=cash]{data/clean/profiles/fin_002746.dat};
\addplot[fill=SoftRed!45, draw=SoftRed, bar shift=3.0pt] table[x=i, y=ibd]{data/clean/profiles/fin_002746.dat};
\addplot[fill=AgriGreen!25, draw=AgriGreen, pattern=north east lines, bar shift=-3.0pt] table[x=x, y=cash]{data/clean/profiles/finh_002746.dat};
\addplot[fill=SoftRed!20, draw=SoftRed, pattern=north east lines, bar shift=3.0pt] table[x=x, y=ibd]{data/clean/profiles/finh_002746.dat};
\addplot[OilBlue, mark=*, mark size=1.1pt, line width=0.8pt] table[x=i, y=ocf]{data/clean/profiles/fin_002746.dat};
\addplot[only marks, OilBlue, mark=*, mark size=1.8pt] table[x=x, y=ocf]{data/clean/profiles/finh_002746.dat};
\legend{货币资金, 有息负债, 经营活动现金流净额}
\end{axis}
\end{tikzpicture}

\captionof{figure}[仙坛股份（002746）近年主要财务指标]{仙坛股份（002746）近年主要财务指标（左上：营业收入与归母净利润；右上：ROE 与资产负债率；下：货币资金、有息负债与经营活动现金流净额。
2021---2025 年为年报口径，2026H1 为 2026 年半年报/半年末，与其前一年报之间留出间隔、以斜纹或空心点区分；半年报数据未年化）}
\end{center}

\pfl{财务分析} 2026 年 6 月末资产负债率 43.0\%，较 2025 年末 上升 1.2 个百分点（样本中位数 52.2\%，所处三级行业内排第 \zhnumber{2} 位）。 2026 年 6 月末货币资金 36.20 亿元、短期债务（短期借款与一年内到期非流动负债）25.39 亿元、长期债务 0.53 亿元，有息负债合计 25.92 亿元，现金短债比 1.43 倍——覆盖充分。 净有息负债 -10.28 亿元，相当于其 2026 年上半年营业收入的 -34\%。 流动比率 1.67、速动比率 1.41，短期指标尚可。 2026 年上半年经营活动现金流净额 0.60 亿元，经营现金流与利润同向为正，内生资金可以支撑运营。

\pfl{重大战略转型} 从公开披露看，公司的战略推进包括：2020-12 定增 10.45 亿元投向诸城年产 1.2 亿羽肉鸡产业生态项目，被定位为“乡村振兴齐鲁样板工程”，建成后公司肉……；2025 公司 2025 年年度报告称诸城项目截至目前已有 36 个商品鸡养殖场投入生产，其余 2 个预计年内相继建设投产……；2026-06-30 诸城年产 1.2 亿羽肉鸡产业生态项目一期配套的屠宰加工厂、饲料厂和商品鸡养殖场已全部投产，二期屠宰加工厂已正式……；2026-07-28 全资子公司拟使用 2.40 亿元募集资金对全资孙公司山东仙润鸿食品有限公司进行增资，并分期提供总额 1.00 亿……。 公告口径上，近三年（2023 年 9 月---2026 年 9 月）公司披露重大重组与并购类公告 2 条、对外投资与转型类公告 1 条、回购与股权激励类公告 0 条，合计公告 327 条。 第一大行业仍是“肉鸡养殖及屠宰加工”，收入占比 100.0\%，与 2021 年年报相比没有方向性变化。 综合判断，公司在报告期内有明确的外延扩张或业务结构调整动作，转型方向与融资安排之间存在直接关联。\footnote{\texttt{https://m.21jingji.com/article/20220731/herald/a09065d6214763dbcf1f8a0c16532cbf\_ths.html}}\footnote{\texttt{http://money.finance.sina.com.cn/corp/view/vCB\_AllBulletinDetail.php}}\footnote{\texttt{https://stock.stockstar.com/RB2026082700056657.shtml}}

\begin{center}\small\setlength{\tabcolsep}{4pt}
\captionof{table}[仙坛股份（002746）历史融资与资本运作]{仙坛股份（002746）历史融资与资本运作（据公司公告与公开报道整理；金额为披露值，含首次公开发行、再融资、债券、银行借款与重整投资等）}
\begin{tabular}{@{}p{1.5cm}p{3.2cm}p{2.4cm}p{6.3cm}@{}}
\toprule
时间 & 方式 & 规模 & 用途与说明 \\
\midrule
2015-02 & IPO（深圳证券交易所中小板） & 3,985 万股，发行价 5.28 元/股，募集资金 2.1 亿元 & 经证监许可[2015]189 号文核准；主承销商民生证券…… \\
2015-06 & 非公开发行股票（预案，后经董事会、股东大会调整） & 拟发行不超过 2,619 万股，发行价格不低于 31.98 元/股，募集资金不超过 8.4 亿元 & 熟食品加工项目、商品鸡饲养立体养殖技术改造项目及补充流动资金…… \\
2020-12 & 非公开发行股票 & 发行 1.10 亿股，募集资金总额 10.45 亿元、净额 10.31 亿元 & 年产 1.2 亿羽肉鸡产业生态项目，该项目包括 7 个商品鸡场…… \\
\bottomrule
\end{tabular}
\end{center}

\pfl{历史融投资情况} 从现金流量表看，2025 年公司取得借款收到的现金 22.71 亿元、偿还债务支付的现金 17.50 亿元、吸收权益性投资 — 亿元，筹资活动现金流净额 -1.32 亿元（2023 年为取得借款 9.38 亿元、偿还债务 4.72 亿元）。 2026 年上半年筹资活动现金流净额为 0.94 亿元（取得借款 20.65 亿元、偿还债务 18.55 亿元，未年化），仍处于净融入状态。 2025 年分配股利、利润或偿付利息支付的现金合计 3.31 亿元。 公司披露的股本与债务融资脉络包括：IPO（深圳证券交易所中小板）、非公开发行股票、非公开发行股票（预案，后经董事会、股东大会调整）。 其中规模较大的两笔为：2015-02IPO（深圳证券交易所中小板）（3,985 万股，发行价 5.28 元/股，募集资金 2.1……）；2020-12非公开发行股票（发行 1.10 亿股，募集资金总额 10.45 亿元、净额 ……）。 股本方面，近三年披露股本变动 3 次；总股本由 2021 年末的 8.61 亿股变为 8.61 亿股（+0.00\%）。 分红方面，近三年现金分红 7 次、累计每股派现 0.95 元（最近一次 2026-06-30）——在利润承压的阶段仍维持了分红。\footnote{\texttt{http://finance.china.com.cn/stock/ssgs/20150623/3190867.shtml}}\footnote{\texttt{https://pdf.dfcfw.com/pdf/H3\_AP202110191523732220\_1.pdf}}

\pfl{未来潜在融资需求分析} 2026 年 6 月末货币资金 36.20 亿元，而短期债务 25.39 亿元（现金短债比 1.43 倍），2026 年上半年经营现金流净额 0.60 亿元。账面货币资金可以覆盖短期债务，缺口不体现在静态偿付层面。 公司 2026 年 6 月末资产负债率 43.0\%、2026 年上半年 ROE 3.9\%（未年化），资产负债表尚有余量，融资需求不是“补血型”的刚性需求，而是围绕并购整合、项目投入与营运资金的主动性需求。 公开披露的下一步融资线索包括：2026-07-28 以募集资金 2.40 亿元增资孙公司并向其提供 1.00 亿元无息借款，用于推进诸城年产 1.2 亿羽肉鸡产业生态项目，……。\footnote{\texttt{https://www.mrjjxw.com/articles/2026-07-28/4524327.html}}


\subsubsection{益客食品（301116）}
\label{co:301116}

\pfl{公司基本情况} 公司前身可追溯至 2008 年，江苏肉禽（肉鸭、肉鸡）屠宰与加工企业，2022 年 1 月创业板上市，形成肉鸡、肉鸭双产业链协同，产品覆盖鸭产品、鸭苗、鸡产品、鸡苗、饲料、调理品、熟食与羽绒；控股股东益客农牧持股 70.11\%，2026 年上半年营业收入 99.02 亿元、归母净利润 3,370.62 万元实现扭亏。 公司于 2022-01-18 在深交所创业板首发上市，注册资本 4.49 亿元，总股本 4.49 亿股。 截至 2026-09-24收盘价 9.12 元，流通市值 40.9 亿元，近一年-15.6\%，流通市值在三级行业“肉鸡养殖”的 8 家公司中排第 \zhnumber{5} 位，近一年涨跌幅低于全样本中位数（-10.1\%）5.5 个百分点。 按 2025 年年报披露的行业构成，收入占比前三位为屠宰 70.5\%；饲料 15.0\%；羽绒 5.4\%。 与 2021 年年报相比，第一大行业由“屠宰板块”变为“屠宰”，主业重心已经迁移。 发展过程中的关键节点包括：1 月 18 日在深圳证券交易所创业板上市，经证监许可[2021]3800 号批复核准公……（2022）；4 月 7 日披露向控股股东益客农牧发行股票预案（拟募资不超 4.5 亿元、价格 12.……（2023）；2 月 18 日第三届董事会第九次会议审议通过回购公司股份事项；截至 2024 年半年报……（2024）。\footnote{\texttt{http://money.finance.sina.com.cn/corp/view/vISSUE\_MarketBulletinDetail.php}}\footnote{\texttt{https://stcn.com/article/detail/1317323.html}}\footnote{\texttt{http://static.cninfo.com.cn/finalpage/2024-08-27/1220984063.PDF}}

\pfl{历史股价} 本报告样本区间（2022 年 1 月 至 2026 年 9 月）内，股价由 22.05 元变为 9.12 元，累计 -58.6\%，区间最高 22.05 元（2022 年 1 月）、最低 8.56 元（2026 年 6 月）；当前价相当于区间最高价的 41.4\%，为区间最低价的 1.07 倍。 月度收益的年化波动率 34.8\%，低于全样本中位数 37.4\%，在三级行业“肉鸡养殖”的 8 家公司中波动率排第 \zhnumber{6} 位。 从事件看，2022 年 1 月 至 2024 年 6 月的下跌（-59.5\%）与公司2024 年 1 月获中国证监会同意注册批复；2024 年 9 月 11 日……（2023-04-07）在时间上重合；2024 年 6 月 至 2024 年 11 月的上涨（+53.3\%）与公司终止向控股股东益客农牧发行股票（原拟募资不超 4.5 亿元补充流动资金）事项……（2024-09-11）在时间上重合——股价的拐点多数能在公司自身的资本运作与风险事件上找到对应。\footnote{\texttt{https://stcn.com/article/detail/1317323.html}}

\begin{center}
\begin{tikzpicture}
\begin{axis}[
  width=12.6cm, height=3.9cm, scale only axis,
  xmin=2021.922, xmax=2026.888, ymin=7.96, ymax=23.59,
  xtick={2022,2023,2024,2025,2026},
  yearticks,
  yticklabel style={font=\scriptsize},
  ylabel={元/股}, ylabel style={font=\scriptsize},
  grid=major, grid style={LightGray, line width=0.3pt},
  axis line style={InkGray, line width=0.5pt},
  every axis plot/.append style={line width=0.9pt},
]
\addplot[AgriGreen] table[x=x,y=y]{data/clean/profiles/px_301116.dat};
\addplot[SoftRed, dashed, line width=0.7pt] table[x=x,y=y]{data/clean/profiles/pxzz_301116.dat};
\addplot[only marks, mark=*, mark size=1.1pt, SoftRed] table[x=x,y=y]{data/clean/profiles/pxzz_301116.dat};
\end{axis}
\end{tikzpicture}

\captionof{figure}[益客食品（301116）股价与周期骨架]{益客食品（301116）月度收盘价（元/股）与 ZigZag 周期骨架（阈值 25\%，圆点为周期转折点）}
\end{center}

\pfl{过去周期分析} 以 25\% 的 ZigZag 阈值划分，样本区间内股价形成 1 个主升段与 2 个主跌段。 最大主升段出现在 2024 年 6 月 至 2024 年 11 月，5 个月+53.3\%；最大主跌段为 2022 年 1 月 至 2024 年 6 月，29 个月-59.5\%。 最近一次转折出现在 2026 年 6 月（下跌段自 2024 年 11 月起，19 个月-37.5\%），当前处于该段的延续之中。 公司股价与申万农林牧渔指数几乎同步见底，个股并未走出独立行情。 禽链的产能周期只有 6---12 个月、反转快于生猪（第 \ref{sec:other-livestock} 节），2026 年前四个月毛鸡养殖利润同比大幅回升，链条利润的修复先于生猪，公司 2026 年 6 月末资产负债率 68.9\%，高于全样本中位数 52.2\%，其股价因此更多反映资金面与信用风险，而不是价格弹性本身。

\pfl{盈利情况} 2026 年上半年营业收入 99.02 亿元（同比 +13.2\%），归母净利润 0.34 亿元，实现扭亏。 以年报为趋势对照，2023---2025 年营业收入由 218.88 亿元变为 189.62 亿元（两年年均复合增速 -6.9\%），归母净利润由 -1.71 亿元变为 -2.80 亿元。 按半年报口径，2026 年上半年的 ROE 为 2.2\%（未年化）、销售净利率 0.4\%、毛利率 4.0\%，ROE 在“肉鸡养殖”8 家公司中排第 \zhnumber{4} 位。 同期全样本半年 ROE 中位数 0.75\%，公司与之相差 1.43 个百分点（略好于中位数公司）。 结合公司披露，2025 年公司主要产品销量同比变化不大，营业收入同比下降约 9\%，下降主要系主要产品销售价格下降所致；……；2025 年肉鸭产业链屠宰、鸭苗等业务环节市场需求不及预期，呈现周期性、暂时性产能相对过剩，产品价格降幅……。\footnote{\texttt{https://static.cninfo.com.cn/finalpage/2026-04-24/1225160345.PDF}}

\begin{center}\small\setlength{\tabcolsep}{4pt}
\captionof{table}[益客食品（301116）关键财务指标]{益客食品（301116）关键财务指标（合并报表口径；两个半年报期均未年化；现金短债比 $=$ 货币资金 $\div$ 短期债务）}
\begin{tabular}{@{}lrrrrr@{}}
\toprule
指标 & 2023 年 & 2024 年 & 2025 年 & 2025 年半年报 & 2026 年半年报 \\
\midrule
营业收入（亿元） & 218.88 & 208.37 & 189.62 & 87.47 & 99.02 \\
归母净利润（亿元） & -1.71 & 1.00 & -2.80 & -1.81 & 0.34 \\
毛利率（\%） & 2.2 & 4.4 & 2.4 & 1.9 & 4.0 \\
ROE（\%） & -9.1 & 5.7 & -17.2 & -10.8 & 2.2 \\
资产负债率（\%） & 61.3 & 61.3 & 65.2 & 68.2 & 68.9 \\
经营活动现金流净额（亿元） & 6.33 & 9.42 & 3.36 & -2.75 & -0.70 \\
货币资金（亿元） & 4.88 & 2.54 & 2.69 & 3.98 & 4.84 \\
有息负债合计（亿元） & 19.44 & 16.33 & 19.04 & 24.45 & 23.69 \\
现金短债比（倍） & 0.41 & 0.31 & 0.19 & 0.25 & 0.28 \\
\bottomrule
\end{tabular}
\end{center}

\begin{center}
\begin{tikzpicture}
\begin{groupplot}[
  group style={group size=2 by 1, horizontal sep=0.95cm},
  width=7.2cm, height=4.6cm,
  tick label style={font=\tiny},
  title style={font=\scriptsize\heiti},
  yticklabel style={font=\tiny},
  grid=major, grid style={LightGray, line width=0.3pt},
  axis line style={InkGray, line width=0.5pt},
  enlarge x limits={abs=0.8},
  legend style={font=\scriptsize, at={(0.5,-0.22)}, anchor=north,
                legend columns=2, draw=none, fill=none, column sep=6pt},
]
\nextgroupplot[
  ybar, bar width=4pt,
  ymin=-24.97, ymax=245.48,
  xtick={0,1,2,3,4,5.7},
  xticklabels={2021,2022,2023,2024,2025,2026H1},
  ylabel={亿元},
]
\addplot[fill=AgriGreen!75, draw=AgriGreen, bar shift=-2.6pt] table[x=i, y=rev]{data/clean/profiles/fin_301116.dat};
\addplot[fill=SoftRed!60, draw=SoftRed, bar shift=2.6pt] table[x=i, y=ni]{data/clean/profiles/fin_301116.dat};
\addplot[fill=AgriGreen!30, draw=AgriGreen, pattern=north east lines, bar shift=-2.6pt] table[x=x, y=rev]{data/clean/profiles/finh_301116.dat};
\addplot[fill=SoftRed!20, draw=SoftRed, pattern=north east lines, bar shift=2.6pt] table[x=x, y=ni]{data/clean/profiles/finh_301116.dat};
\legend{营业收入（年/半年）, 归母净利润（年/半年）}
\nextgroupplot[
  ymin=-24.1, ymax=77.5,
  xtick={0,1,2,3,4,5.7},
  xticklabels={2021,2022,2023,2024,2025,2026H1},
  ylabel={\%},
]
\addplot[AgriGold, mark=*, mark size=1.3pt] table[x=i, y=roe]{data/clean/profiles/fin_301116.dat};
\addplot[OilBlue, mark=square*, mark size=1.3pt] table[x=i, y=debt]{data/clean/profiles/fin_301116.dat};
\addplot[only marks, AgriGold, mark=*, mark size=2.0pt] table[x=x, y=roe]{data/clean/profiles/finh_301116.dat};
\addplot[only marks, OilBlue, mark=square*, mark size=2.0pt] table[x=x, y=debt]{data/clean/profiles/finh_301116.dat};
\legend{ROE, 资产负债率}
\end{groupplot}
\end{tikzpicture}
\begin{tikzpicture}
\begin{axis}[
  width=15.2cm, height=3.1cm, scale only axis,
  ymin=-4.19, ymax=26.73,
  xtick={0,1,2,3,4,5.7},
  xticklabels={2021,2022,2023,2024,2025,2026H1},
  tick label style={font=\tiny},
  yticklabel style={font=\tiny},
  ylabel={亿元}, ylabel style={font=\scriptsize},
  ybar, bar width=4pt,
  grid=major, grid style={LightGray, line width=0.3pt},
  axis line style={InkGray, line width=0.5pt},
  enlarge x limits={abs=0.8},
  legend style={font=\scriptsize, at={(0.5,-0.30)}, anchor=north,
                legend columns=3, draw=none, fill=none, column sep=10pt},
]
\addplot[fill=AgriGreen!55, draw=AgriGreen, bar shift=-3.0pt] table[x=i, y=cash]{data/clean/profiles/fin_301116.dat};
\addplot[fill=SoftRed!45, draw=SoftRed, bar shift=3.0pt] table[x=i, y=ibd]{data/clean/profiles/fin_301116.dat};
\addplot[fill=AgriGreen!25, draw=AgriGreen, pattern=north east lines, bar shift=-3.0pt] table[x=x, y=cash]{data/clean/profiles/finh_301116.dat};
\addplot[fill=SoftRed!20, draw=SoftRed, pattern=north east lines, bar shift=3.0pt] table[x=x, y=ibd]{data/clean/profiles/finh_301116.dat};
\addplot[OilBlue, mark=*, mark size=1.1pt, line width=0.8pt] table[x=i, y=ocf]{data/clean/profiles/fin_301116.dat};
\addplot[only marks, OilBlue, mark=*, mark size=1.8pt] table[x=x, y=ocf]{data/clean/profiles/finh_301116.dat};
\legend{货币资金, 有息负债, 经营活动现金流净额}
\end{axis}
\end{tikzpicture}

\captionof{figure}[益客食品（301116）近年主要财务指标]{益客食品（301116）近年主要财务指标（左上：营业收入与归母净利润；右上：ROE 与资产负债率；下：货币资金、有息负债与经营活动现金流净额。
2021---2025 年为年报口径，2026H1 为 2026 年半年报/半年末，与其前一年报之间留出间隔、以斜纹或空心点区分；半年报数据未年化）}
\end{center}

\pfl{财务分析} 2026 年 6 月末资产负债率 68.9\%，较 2025 年末 上升 3.6 个百分点（样本中位数 52.2\%，所处三级行业内排第 \zhnumber{8} 位）。 2026 年 6 月末货币资金 4.84 亿元、短期债务（短期借款与一年内到期非流动负债）17.56 亿元、长期债务 6.14 亿元，有息负债合计 23.69 亿元，现金短债比 0.28 倍——缺口明显。 净有息负债 18.86 亿元，相当于其 2026 年上半年营业收入的 19\%。 流动比率 0.74、速动比率 0.31，短期偿付压力较大。 2026 年上半年经营活动现金流净额 -0.70 亿元，净利润为正但经营现金流为负，利润的现金含量偏低。

\pfl{重大战略转型} 公司在报告期内的战略动作可归为 3 项：2024-02 公告控股股东、实际控制人及家属、董事为公司及子公司银行综合授信提供关联担保，额度最高 28 亿、担保总额度不超过……；2026-06-10 拟募资不超 9 亿元用于益客食品科技谷项目、肉禽产业链综合研发项目、数智化升级建设项目与补充流动资金，控股股东拟……；2026-09 公司称 2026 年上半年业绩扭亏主要系鸡产品、鸭产品、饲料、鸡苗、调理品、熟食、羽绒等主要产品同比毛利升高；同……。 公告口径上，近三年（2023 年 9 月---2026 年 9 月）公司披露重大重组与并购类公告 0 条、对外投资与转型类公告 14 条、回购与股权激励类公告 23 条，合计公告 478 条。 主营结构的迁移已经落到报表上：2021 年年报的第一大行业为“屠宰板块”（76.6\%），2025 年年报变为“屠宰”（70.5\%）。 综合判断，公司在报告期内有明确的外延扩张或业务结构调整动作，转型方向与融资安排之间存在直接关联。\footnote{\texttt{http://money.finance.sina.com.cn/corp/view/vCB\_AllBulletinDetail.php}}\footnote{\texttt{https://www.stcn.com/article/detail/3954365.html}}\footnote{\texttt{https://www.lanjinger.com/d/1787734135191661490}}

\begin{center}\small\setlength{\tabcolsep}{4pt}
\captionof{table}[益客食品（301116）历史融资与资本运作]{益客食品（301116）历史融资与资本运作（据公司公告与公开报道整理；金额为披露值，含首次公开发行、再融资、债券、银行借款与重整投资等）}
\begin{tabular}{@{}p{1.5cm}p{3.2cm}p{2.4cm}p{6.3cm}@{}}
\toprule
时间 & 方式 & 规模 & 用途与说明 \\
\midrule
2022-01 & IPO（深圳证券交易所创业板） & 发行 44,897,960 股、发行价 11.40 元/股，募集资金总额 51,183.67 万元，扣除发行费用 6,761.34 万元（不含增值税）后净额 44,422.33 万元 & 肉鸭屠宰线建设、禽肉熟食扩建、供应链数字化、专用饲料研发及产业化项目及补充流动资金…… \\
2023-04-07 & 向特定对象发行股票（向控股股东益客农牧，后终止） & 拟募集资金不超过 4.5 亿元（发行不超过 3,955.69 万股，发行价格 12.64 元/股，后调整为不超过 45,000.00 万元） & 扣除发行费用后全部用于补充流动资金；2024 年 1 月获中国证监会同意注册批复…… \\
2026-06-10 & 向特定对象发行股票（预案，含控股股东认购） & 拟募集资金不超过 9 亿元（发行股票数量不超过 1.35 亿股），其中控股股东益客农牧拟以现金认购不低于 5,000 万元且不高于 1.5 亿元 & 益客食品科技谷项目、肉禽产业链综合研发项目、数智化升级建设项目、补充流动资金项目…… \\
\bottomrule
\end{tabular}
\end{center}

\pfl{历史融投资情况} 从现金流量表看，2025 年公司取得借款收到的现金 18.21 亿元、偿还债务支付的现金 14.96 亿元、吸收权益性投资 0.03 亿元，筹资活动现金流净额 1.52 亿元（2023 年为取得借款 12.26 亿元、偿还债务 10.89 亿元）。 2026 年上半年筹资活动现金流净额为 2.35 亿元（取得借款 12.07 亿元、偿还债务 7.76 亿元，未年化），仍处于净融入状态。 2025 年分配股利、利润或偿付利息支付的现金合计 1.37 亿元。 公开信息中与公司直接相关的融资与资本运作包括：IPO（深圳证券交易所创业板）、向特定对象发行股票（向控股股东益客农牧，后终止）、向特定对象发行股票（预案，含控股股东认购）。 其中规模较大的两笔为：2022-01IPO（深圳证券交易所创业板）（发行 44,897,960 股、发行价 11.40 元/股，……）；2026-06-10向特定对象发行股票（预案，含控股股东认购）（拟募集资金不超过 9 亿元（发行股票数量不超过 1.35 亿……）。 股本方面，近三年披露股本变动 5 次（限售股份上市 1 次、其他 1 次）；总股本由 2021 年末的 4.04 亿股变为 4.49 亿股（+11.11\%）。 分红方面，近三年现金分红 3 次、累计每股派现 0.20 元（最近一次 2024-12-31）——在利润承压的阶段仍维持了分红。 近三年“再融资”类公告 48 条，涉及定向增发、募集资金使用。\footnote{\texttt{http://file.finance.sina.com.cn/211.154.219.97:9494/MRGG/CNSESZ\_STOCK/2024/2024-4/2024-0}}\footnote{\texttt{https://stcn.com/article/detail/1317323.html}}\footnote{\texttt{https://www.stcn.com/article/detail/3954365.html}}

\pfl{未来潜在融资需求分析} 2026 年 6 月末货币资金 4.84 亿元，而短期债务 17.56 亿元（现金短债比 0.28 倍），2026 年上半年经营现金流净额 -0.70 亿元。按“短期债务减去账面货币资金”的滚动偿付口径，静态缺口约 12.72 亿元；若再计入上半年亏损的年化额（0.00 亿元），维持一年正常经营所需的外部资金约 12.72---12.72 亿元。 公司 2026 年 6 月末资产负债率 68.9\%、2026 年上半年 ROE 2.2\%（未年化），资产负债表尚有余量，融资需求不是“补血型”的刚性需求，而是围绕并购整合、项目投入与营运资金的主动性需求。 公开披露的下一步融资线索包括：2024-09-11 终止向控股股东益客农牧发行股票（原拟募资不超 4.5 亿元补充流动资金）事项，公司称不会对正常生产经营与持续稳定发展造成……；2026-06-10 拟向包括控股股东益客农牧在内的不超过 35 名特定对象发行股票募资不超过 9 亿元，为 2024 年终止前次定增后的新一……；2026-08-21 2026 年度向特定对象发行 A 股股票申请获深圳证券交易所受理（募集说明书申报稿同日披露）。\footnote{\texttt{https://stcn.com/article/detail/1317323.html}}\footnote{\texttt{https://www.stcn.com/article/detail/3954365.html}}\footnote{\texttt{https://www.cninfo.com.cn/new/disclosure/stock}}


\subsubsection{天邦食品（002124）}
\label{co:002124}

\pfl{公司基本情况} 浙江余姚起家的生猪养殖与食品加工企业（曾称“浙江猪王”），1996 年成立、2007 年深交所上市，2021 年起相继出售疫苗、饲料与史记生物种猪业务以聚焦生猪养殖，2024 年 8 月进入预重整程序，2026 年 9 月与产业投资人解除重整投资协议并重新公开招募投资人，截至 2026 年 6 月末归母净资产 8.23 亿元。 公司于 2007-04-03 在深交所主板首发上市，注册资本 22.22 亿元，总股本 22.22 亿股。 截至 2026-09-24收盘价 2.07 元，流通市值 40.8 亿元，近一年-31.0\%，流通市值在三级行业“生猪养殖”的 11 家公司中排第 \zhnumber{9} 位，近一年涨跌幅低于全样本中位数（-10.1\%）20.9 个百分点。 按 2026 年半年报披露的行业构成，收入占比前三位为养殖 65.4\%；食品 31.4\%；饲料产品 3.2\%。 与 2021 年年报相比，第一大行业由“生猪养殖”变为“养殖”，主业重心已经迁移。 发展过程中的关键节点包括：2 月经证监发行字[2007]48 号文核准公开发行人民币普通股（A 股）1,850.0……（2007）；公司从 CG 美国公司引进世界一流种猪及美式养殖模式，高起点进入生猪养殖行业（2013）；上半年出售生物制品疫苗业务，下半年出售水产饲料业务和部分猪用饲料业务，主营业务聚焦为生猪……（2021）。\footnote{\texttt{https://vip.stock.finance.sina.com.cn/corp/go.php/vISSUE\_NewStock/stockid/002124.phtml}}\footnote{\texttt{http://file.finance.sina.com.cn/211.154.219.97:9494/MRGG/CNSESZ\_STOCK/2021/2021-4/2021-0}}\footnote{\texttt{https://pdf.dfcfw.com/pdf/H2\_AN202306121590862005\_1.pdf}}

\pfl{历史股价} 本报告样本区间（2021 年 1 月 至 2026 年 9 月）内，股价由 14.27 元变为 2.07 元，累计 -85.5\%，区间最高 17.55 元（2021 年 2 月）、最低 2.07 元（2026 年 9 月）；当前价相当于区间最高价的 11.8\%，已接近区间最低价（1.00 倍）。 月度收益的年化波动率 45.3\%，高于全样本中位数 37.4\%，在三级行业“生猪养殖”的 11 家公司中波动率排第 \zhnumber{3} 位。 从事件看，2022 年 3 月 至 2024 年 6 月的下跌（-76.1\%）与公司申请重整及预重整事项经公司 2024 年第二次临时股东大会审议通过（2024-04-08）在时间上重合——股价的拐点多数能在公司自身的资本运作与风险事件上找到对应。\footnote{\texttt{https://static.cninfo.com.cn/finalpage/2026-08-25/1225495960.PDF}}

\begin{center}
\begin{tikzpicture}
\begin{axis}[
  width=12.6cm, height=3.9cm, scale only axis,
  xmin=2020.922, xmax=2026.888, ymin=1.93, ymax=18.78,
  xtick={2021,2022,2023,2024,2025,2026},
  yearticks,
  yticklabel style={font=\scriptsize},
  ylabel={元/股}, ylabel style={font=\scriptsize},
  grid=major, grid style={LightGray, line width=0.3pt},
  axis line style={InkGray, line width=0.5pt},
  every axis plot/.append style={line width=0.9pt},
]
\addplot[AgriGreen] table[x=x,y=y]{data/clean/profiles/px_002124.dat};
\addplot[SoftRed, dashed, line width=0.7pt] table[x=x,y=y]{data/clean/profiles/pxzz_002124.dat};
\addplot[only marks, mark=*, mark size=1.1pt, SoftRed] table[x=x,y=y]{data/clean/profiles/pxzz_002124.dat};
\end{axis}
\end{tikzpicture}

\captionof{figure}[天邦食品（002124）股价与周期骨架]{天邦食品（002124）月度收盘价（元/股）与 ZigZag 周期骨架（阈值 25\%，圆点为周期转折点）}
\end{center}

\pfl{过去周期分析} 以 25\% 的 ZigZag 阈值划分，样本区间内股价形成 2 个主升段与 3 个主跌段。 最大主升段出现在 2021 年 8 月 至 2022 年 3 月，7 个月+59.3\%；最大主跌段为 2022 年 3 月 至 2024 年 6 月，27 个月-76.1\%。 最近一次转折出现在 2026 年 9 月（下跌段自 2024 年 11 月起，22 个月-36.9\%），当前处于该段的延续之中。 与申万农林牧渔指数的最近一次底部相比，该公司见底晚约 4 个月，在本轮下行中属于调整更慢的一类。 生猪养殖的定价主线是产能去化与成本曲线的赛跑（第 \ref{sec:hog} 章）：2026 年三季度行业仍处在“产能去化未完成、价格磨底”的阶段，公司 2026 年 6 月末资产负债率 91.2\%，高于全样本中位数 52.2\%，其股价因此更多反映资金面与信用风险，而不是价格弹性本身。

\pfl{盈利情况} 2026 年上半年营业收入 36.56 亿元（同比 -22.1\%），归母净利润 -15.43 亿元，由上年同期的盈利转为亏损。 以年报为趋势对照，2023---2025 年营业收入由 102.32 亿元变为 86.89 亿元（两年年均复合增速 -7.8\%），归母净利润由 -28.83 亿元变为 -13.09 亿元。 按半年报口径，2026 年上半年的 ROE 为 -96.8\%（未年化）、销售净利率 -42.2\%、毛利率 3.0\%，ROE 在“生猪养殖”11 家公司中排第 \zhnumber{11} 位。 行业同口径半年 ROE 的中位数为 0.75\%，该公司低于中位数 97.51 个百分点。 结合公司披露，第九届董事会第十一次（临时）会议审议同意公司与产业投资人签署的重整投资协议已解除，同意解除与南宁漓源签署……；公司继续公开招募重整投资人，招募目的为引入具备产业协同与资金实力的重整投资人为公司提供产业及资金等方面支……。\footnote{\texttt{https://finance.sina.com.cn/roll/2026-09-22/doc-inisrnpr7417227.shtml}}\footnote{\texttt{http://money.finance.sina.com.cn/corp/view/vCB\_AllBulletinDetail.php}}

\begin{center}\small\setlength{\tabcolsep}{4pt}
\captionof{table}[天邦食品（002124）关键财务指标]{天邦食品（002124）关键财务指标（合并报表口径；两个半年报期均未年化；现金短债比 $=$ 货币资金 $\div$ 短期债务）}
\begin{tabular}{@{}lrrrrr@{}}
\toprule
指标 & 2023 年 & 2024 年 & 2025 年 & 2025 年半年报 & 2026 年半年报 \\
\midrule
营业收入（亿元） & 102.32 & 96.44 & 86.89 & 46.93 & 36.56 \\
归母净利润（亿元） & -28.83 & 14.59 & -13.09 & 3.52 & -15.43 \\
毛利率（\%） & -8.9 & 16.5 & 8.6 & 17.4 & 3.0 \\
ROE（\%） & -116.3 & 49.5 & -43.4 & 9.1 & -96.8 \\
资产负债率（\%） & 86.7 & 72.6 & 79.8 & 69.8 & 91.2 \\
经营活动现金流净额（亿元） & 2.64 & -9.64 & 9.11 & 4.89 & 2.71 \\
货币资金（亿元） & 12.54 & 1.50 & 2.25 & 2.34 & 1.55 \\
有息负债合计（亿元） & 61.93 & 52.86 & 49.26 & 51.10 & 48.40 \\
现金短债比（倍） & 0.30 & 0.04 & 0.06 & 0.07 & 0.04 \\
\bottomrule
\end{tabular}
\end{center}

\begin{center}
\begin{tikzpicture}
\begin{groupplot}[
  group style={group size=2 by 1, horizontal sep=0.95cm},
  width=7.2cm, height=4.6cm,
  tick label style={font=\tiny},
  title style={font=\scriptsize\heiti},
  yticklabel style={font=\tiny},
  grid=major, grid style={LightGray, line width=0.3pt},
  axis line style={InkGray, line width=0.5pt},
  enlarge x limits={abs=0.8},
  legend style={font=\scriptsize, at={(0.5,-0.22)}, anchor=north,
                legend columns=2, draw=none, fill=none, column sep=6pt},
]
\nextgroupplot[
  ybar, bar width=4pt,
  ymin=-59.59, ymax=123.03,
  xtick={0,1,2,3,4,5.7},
  xticklabels={2021,2022,2023,2024,2025,2026H1},
  ylabel={亿元},
]
\addplot[fill=AgriGreen!75, draw=AgriGreen, bar shift=-2.6pt] table[x=i, y=rev]{data/clean/profiles/fin_002124.dat};
\addplot[fill=SoftRed!60, draw=SoftRed, bar shift=2.6pt] table[x=i, y=ni]{data/clean/profiles/fin_002124.dat};
\addplot[fill=AgriGreen!30, draw=AgriGreen, pattern=north east lines, bar shift=-2.6pt] table[x=x, y=rev]{data/clean/profiles/finh_002124.dat};
\addplot[fill=SoftRed!20, draw=SoftRed, pattern=north east lines, bar shift=2.6pt] table[x=x, y=ni]{data/clean/profiles/finh_002124.dat};
\legend{营业收入（年/半年）, 归母净利润（年/半年）}
\nextgroupplot[
  ymin=-132.9, ymax=112.0,
  xtick={0,1,2,3,4,5.7},
  xticklabels={2021,2022,2023,2024,2025,2026H1},
  ylabel={\%},
]
\addplot[AgriGold, mark=*, mark size=1.3pt] table[x=i, y=roe]{data/clean/profiles/fin_002124.dat};
\addplot[OilBlue, mark=square*, mark size=1.3pt] table[x=i, y=debt]{data/clean/profiles/fin_002124.dat};
\addplot[only marks, AgriGold, mark=*, mark size=2.0pt] table[x=x, y=roe]{data/clean/profiles/finh_002124.dat};
\addplot[only marks, OilBlue, mark=square*, mark size=2.0pt] table[x=x, y=debt]{data/clean/profiles/finh_002124.dat};
\legend{ROE, 资产负债率}
\end{groupplot}
\end{tikzpicture}
\begin{tikzpicture}
\begin{axis}[
  width=15.2cm, height=3.1cm, scale only axis,
  ymin=-18.56, ymax=90.22,
  xtick={0,1,2,3,4,5.7},
  xticklabels={2021,2022,2023,2024,2025,2026H1},
  tick label style={font=\tiny},
  yticklabel style={font=\tiny},
  ylabel={亿元}, ylabel style={font=\scriptsize},
  ybar, bar width=4pt,
  grid=major, grid style={LightGray, line width=0.3pt},
  axis line style={InkGray, line width=0.5pt},
  enlarge x limits={abs=0.8},
  legend style={font=\scriptsize, at={(0.5,-0.30)}, anchor=north,
                legend columns=3, draw=none, fill=none, column sep=10pt},
]
\addplot[fill=AgriGreen!55, draw=AgriGreen, bar shift=-3.0pt] table[x=i, y=cash]{data/clean/profiles/fin_002124.dat};
\addplot[fill=SoftRed!45, draw=SoftRed, bar shift=3.0pt] table[x=i, y=ibd]{data/clean/profiles/fin_002124.dat};
\addplot[fill=AgriGreen!25, draw=AgriGreen, pattern=north east lines, bar shift=-3.0pt] table[x=x, y=cash]{data/clean/profiles/finh_002124.dat};
\addplot[fill=SoftRed!20, draw=SoftRed, pattern=north east lines, bar shift=3.0pt] table[x=x, y=ibd]{data/clean/profiles/finh_002124.dat};
\addplot[OilBlue, mark=*, mark size=1.1pt, line width=0.8pt] table[x=i, y=ocf]{data/clean/profiles/fin_002124.dat};
\addplot[only marks, OilBlue, mark=*, mark size=1.8pt] table[x=x, y=ocf]{data/clean/profiles/finh_002124.dat};
\legend{货币资金, 有息负债, 经营活动现金流净额}
\end{axis}
\end{tikzpicture}

\captionof{figure}[天邦食品（002124）近年主要财务指标]{天邦食品（002124）近年主要财务指标（左上：营业收入与归母净利润；右上：ROE 与资产负债率；下：货币资金、有息负债与经营活动现金流净额。
2021---2025 年为年报口径，2026H1 为 2026 年半年报/半年末，与其前一年报之间留出间隔、以斜纹或空心点区分；半年报数据未年化）}
\end{center}

\pfl{财务分析} 2026 年 6 月末资产负债率 91.2\%，较 2025 年末 上升 11.4 个百分点（样本中位数 52.2\%，所处三级行业内排第 \zhnumber{11} 位）。 2026 年 6 月末货币资金 1.55 亿元、短期债务（短期借款与一年内到期非流动负债）35.01 亿元、长期债务 13.38 亿元，有息负债合计 48.40 亿元，现金短债比 0.04 倍——缺口明显。 净有息负债 46.85 亿元，相当于其 2026 年上半年营业收入的 128\%。 流动比率 0.18、速动比率 0.05，短期偿付压力较大。 2026 年上半年经营活动现金流净额 2.71 亿元，净利润为负而经营现金流为正，差额主要来自折旧摊销与营运资本变动，说明亏损尚未完全转化为现金流出。 2026 年 6 月末在建工程 7.73 亿元，较上年末减少 0.21 亿元，在建投入在收缩。 公司披露的财务与合规风险包括：公司提示：如法院依法裁定受理公司重整申请，公司股票将被实施退市风险警示（*ST）；如重整……；公司收到（2025）浙 02 破申 7 号之二《决定书》，宁波中院准许延期申请，决定延长……。\footnote{\texttt{http://money.finance.sina.com.cn/corp/view/vCB\_AllBulletinDetail.php}}\footnote{\texttt{https://static.cninfo.com.cn/finalpage/2026-08-25/1225495960.PDF}}

\pfl{重大战略转型} 从公开披露看，公司的战略推进包括：2022-06 以 10.2 亿元转让史记生物技术（南京）有限公司 51\% 股权给三亚史记生物科技有限公司，由 100\% 控股变……；2023-12 2023 年度向特定对象发行股票募集资金将剩余的募集资金用途全部变更为新募投项目中的“天邦股份数智化猪场升级项目……；2026-06-30 公司持续推进生物安全升级、创新管理等各项降本增效措施，2026 年上半年断奶仔猪成本降至 278 元/头，同比下……；2026-09-22 公司继续公开招募重整投资人，招募目的为引入具备产业协同与资金实力的重整投资人为公司提供产业及资金等方面支持，化解……。 公告口径上，近三年（2023 年 9 月---2026 年 9 月）公司披露重大重组与并购类公告 4 条、对外投资与转型类公告 5 条、回购与股权激励类公告 0 条，合计公告 492 条。 主营结构的迁移已经落到报表上：2021 年年报的第一大行业为“生猪养殖”（60.4\%），2026 年半年报变为“养殖”（65.4\%）。 综合判断，公司当前的战略重心不是扩张而是化解债务与完成重整，业务层面的调整让位于司法重整程序。\footnote{\texttt{https://www.tongwei.cn/uploads/2022/fdd55691ffc603c5d3e7ddf292a7ca12.pdf}}\footnote{\texttt{http://static.cninfo.com.cn/finalpage/2021-09-22/1211121269.PDF}}\footnote{\texttt{http://file.finance.sina.com.cn/211.154.219.97:9494/MRGG/CNSESZ\_STOCK/2023/2023-4/2023-0}}

\begin{center}\small\setlength{\tabcolsep}{4pt}
\captionof{table}[天邦食品（002124）历史融资与资本运作]{天邦食品（002124）历史融资与资本运作（据公司公告与公开报道整理；金额为披露值，含首次公开发行、再融资、债券、银行借款与重整投资等）}
\begin{tabular}{@{}p{1.5cm}p{3.2cm}p{2.4cm}p{6.3cm}@{}}
\toprule
时间 & 方式 & 规模 & 用途与说明 \\
\midrule
2007-04 & IPO（深圳证券交易所中小板） & 实际发行 1,850.00 万股，发行价 10.25 元/股；实际募集资金合计 18,962.50 万元，发行费用合计 1,811.37 万元，募集资金净额 17,151.13 万元 & 主承销商平安证券有限责任公司，承销方式余额包销…… \\
2015-12 & 非公开发行股票 & 发行 80,789,941 股，发行价 5.57 元/股，募集资金总额 450,000,000.00 元，扣除保荐及承销费用后净额 441,000,000.00 元 & 经证监许可[2015]2417 号文核准；公司于 2015 年 12 月在宁波市市场监督…… \\
2017-03-16 & 非公开发行股票（向特定对象） & 发行 136,405,529 股，发行价 10.85 元/股，募集资金总额 1,479,999,989.65 元，净额 1,464,199,989.75 元 & 募投项目，后多次变更用途（累计变更用途的募集资金总额占比曾达 87.55\%）…… \\
2020-08-21 & 非公开发行股票（向 7 家投资者） & 发行 154,080,921 股，发行价 17.30 元/股，募集资金总额 2,665,599,933.30 元，扣除发行费用 31,301,599.27 元后实际募集资金净额 2,634,298,334.03 元 & 经证监许可〔2020〕249 号文核准；募集资金于 2020 年 8 月 10 日到账…… \\
2021-09 & 出售饲料资产并开展战略合作（向通威股份） & 通威股份以 12.71 亿元收购天邦股份旗下水产饲料全部资产及猪料部分资产 & 回笼资金、聚焦生猪养殖与食品加工；2021 年 9 月 20 日公司第七届董事会第四十五…… \\
2022-06-30 & 转让子公司股权（资产出售） & 以 102,000 万元转让史记生物技术（南京）有限公司 51\% 股权（评估值 101,888.09 万元），产生税前资产处置收益 14.52 亿元 & 改善盈利及资产负债状况、聚焦生猪养殖业务；受让方为三亚史记生物科技有限公司（史记生物管理…… \\
2023-12-29 & 向特定对象发行股票（定增） & 发行 382,741,206 股，发行价 3.13 元/股，募集资金总额 1,197,979,974.78 元，净额 1,186,554,139.48 元 & 数智化猪场升级项目及补充流动资金（拟投入数智化猪场升级项目 200,000.00 万元…… \\
2025-05-08 & 重整投资（预重整投资人出资） & 产业投资人厦门建发物产有限公司、南宁漓源粮油饲料有限公司拟以 1.85 元/股价格分别取得 2 亿股天邦食品转增股份 & 用于天邦食品重整；协议签订前两家公司未持有天邦食品股份；重整后两家各持 2 亿股…… \\
2025-05-08 & 重整投资协议（产业投资人） & 1.85 元/股（不得低于前 20 个交易日股票交易均价 3.22 元/股的相应比例） & 受让转增股份、参与重整投资；经投资人评审委员会综合评定…… \\
\bottomrule
\end{tabular}
\end{center}

\pfl{历史融投资情况} 从现金流量表看，2025 年公司取得借款收到的现金 24.15 亿元、偿还债务支付的现金 24.60 亿元、吸收权益性投资 0.03 亿元，筹资活动现金流净额 -5.64 亿元（2023 年为取得借款 36.06 亿元、偿还债务 35.15 亿元）。 2026 年上半年筹资活动现金流净额为 -2.88 亿元（取得借款 6.87 亿元、偿还债务 7.19 亿元，未年化），已转为净流出，处于净还债状态。 2025 年分配股利、利润或偿付利息支付的现金合计 1.31 亿元。 公司披露的股本与债务融资脉络包括：IPO（深圳证券交易所中小板）、出售饲料资产并开展战略合作（向通威股份）、向特定对象发行股票（定增）、转让子公司股权（资产出售）、重整投资协议（产业投资人）、重整投资（预重整投资人出资）。 其中规模较大的两笔为：2025-05-08重整投资（预重整投资人出资）（产业投资人厦门建发物产有限公司、南宁漓源粮油饲料有限公司拟以……）；2007-04IPO（深圳证券交易所中小板）（实际发行 1,850.00 万股，发行价 10.25 元/股……）。 股本方面，近三年披露股本变动 5 次（增发新股上市 1 次、限售股份上市 1 次）；总股本由 2021 年末的 18.39 亿股变为 22.22 亿股（+20.81\%）。 分红方面，近三年未实施现金分红，与重整期间的资金约束一致。 近三年“再融资”类公告 29 条，涉及定向增发、募集资金使用。\footnote{\texttt{https://vip.stock.finance.sina.com.cn/corp/go.php/vISSUE\_NewStock/stockid/002124.phtml}}\footnote{\texttt{http://money.finance.sina.com.cn/corp/view/vCB\_AllBulletinDetail.php}}\footnote{\texttt{http://vip.stock.finance.sina.com.cn/corp/view/vISSUE\_MarketBulletinDetail.php}}

\pfl{未来潜在融资需求分析} 2026 年 6 月末货币资金 1.55 亿元，而短期债务 35.01 亿元（现金短债比 0.04 倍），2026 年上半年经营现金流净额 2.71 亿元。按“短期债务减去账面货币资金”的滚动偿付口径，静态缺口约 33.46 亿元；若再计入上半年亏损的年化额（30.86 亿元），维持一年正常经营所需的外部资金约 33.46---64.32 亿元。 按第 \ref{sec:financing-need} 节的三档划分，公司属于第一档“补血型”：2026 年上半年归母净利润 -15.43 亿元、6 月末资产负债率 91.2\%。负债率已超过 80\%，净资产被亏损侵蚀的程度较深，属于全部样本中资产负债表最紧张的一组。 可行的融资路径按现实性排序为：引入国资或产业投资人（部分公司以控制权让渡为对价）、出售非核心资产与参股股权、定向增发（需股价配合，且控股股东需放弃优先认购或同步增资）。 公开披露的下一步融资线索包括：2024-09-30 披露《关于公开招募重整投资人的公告》（公告编号：2024-106），2024 年 10 月 22 日披露《关于公开招募重……；2026-06-30 2026 年上半年公司总计计提利息费用约 1.53 亿元，因未满产或空置猪场产生的费用约 1.38 亿元，两项费用合计减……；2026-09-22 公司继续公开招募重整投资人，意向重整投资人报名截止时间为 2026 年 10 月 15 日 18:00；预重整管理人有权……。 资金需求还要叠加在建工程：期末在建工程 7.73 亿元、2026 年上半年购建固定资产等支出 0.50 亿元（未年化），这部分支出在周期底部通常会被主动放缓。\footnote{\texttt{http://static.cninfo.com.cn/finalpage/2025-11-07/1224791120.PDF}}\footnote{\texttt{https://disc.static.szse.cn/download/disc/disk03/finalpage/2026-07-14/c8de8806-f3dc-4b87}}\footnote{\texttt{https://finance.sina.cn/stock/ssgs/2026-09-27/detail-inithkhf4545398.d.html}}


\subsubsection{东瑞股份（001201）}
\label{co:001201}

\pfl{公司基本情况} 广东河源的供港澳活大猪供应商与粤港澳大湾区“菜篮子”重要生产基地，2002 年成立、2021 年 4 月深交所主板上市（发行价 63.38 元、募资总额 20.07 亿元），形成饲料、育种、扩繁、养殖、屠宰冷链一体化生猪全产业链，2026 年上半年出栏生猪 88.60 万头但归母净亏损 4.25 亿元。 公司于 2021-04-28 在深交所主板首发上市，注册资本 2.58 亿元，总股本 2.58 亿股。 截至 2026-09-24收盘价 14.66 元，流通市值 31.2 亿元，近一年-10.6\%，流通市值在三级行业“生猪养殖”的 11 家公司中排第 \zhnumber{10} 位，近一年涨跌幅低于全样本中位数（-10.1\%）0.5 个百分点。 按 2026 年半年报披露的行业构成，收入占比前三位为养殖行业 95.3\%；饲料加工业 2.7\%；屠宰、肉食业务 1.7\%。 发展过程中的关键节点包括：公司取得供港澳地区生猪资格；此后成为内地供港活大猪前三大供应商之一和粤港澳大湾区“菜篮子……（2003）；8 月公司正式更名为“东瑞食品集团有限公司”（2013）；12 月 2 日由东瑞有限整体变更设立东瑞食品集团股份有限公司，注册资本 9,500.0……（2016）。\footnote{\texttt{https://finance.sina.cn/fund/jjgdxw/2021-06-24/detail-ikqcfnca3041841.d.html}}\footnote{\texttt{http://money.finance.sina.com.cn/corp/view/vISSUE\_RaiseExplanationDetail.php}}

\pfl{历史股价} 本报告样本区间（2021 年 4 月 至 2026 年 9 月）内，股价由 73.93 元变为 14.66 元，累计 -80.2\%，区间最高 73.93 元（2021 年 4 月）、最低 10.86 元（2026 年 6 月）；当前价相当于区间最高价的 19.8\%，为区间最低价的 1.35 倍。 月度收益的年化波动率 37.2\%，低于全样本中位数 37.4\%，在三级行业“生猪养殖”的 11 家公司中波动率排第 \zhnumber{6} 位。 从事件看，2022 年 3 月 至 2024 年 2 月的下跌（-59.2\%）与公司经证监许可〔2023〕1377 号核准；新增股份上市时间 2023 年 12……（2023-12-28）在时间上重合；2024 年 5 月 至 2026 年 6 月的下跌（-54.4\%）与公司2026 年上半年营业收入 10.39 亿元，同比减少 5.09\%；归属于上……（2026-06-30）在时间上重合——股价的拐点多数能在公司自身的资本运作与风险事件上找到对应。\footnote{\texttt{https://www.100est.com/res/financial-report/r2023/SZ001201\_202404291631627113.pdf}}\footnote{\texttt{http://vip.stock.finance.sina.com.cn/corp/view/vCB\_AllBulletinDetail.php}}

\begin{center}
\begin{tikzpicture}
\begin{axis}[
  width=12.6cm, height=3.9cm, scale only axis,
  xmin=2021.172, xmax=2026.888, ymin=10.10, ymax=79.11,
  xtick={2021,2022,2023,2024,2025,2026},
  yearticks,
  yticklabel style={font=\scriptsize},
  ylabel={元/股}, ylabel style={font=\scriptsize},
  grid=major, grid style={LightGray, line width=0.3pt},
  axis line style={InkGray, line width=0.5pt},
  every axis plot/.append style={line width=0.9pt},
]
\addplot[AgriGreen] table[x=x,y=y]{data/clean/profiles/px_001201.dat};
\addplot[SoftRed, dashed, line width=0.7pt] table[x=x,y=y]{data/clean/profiles/pxzz_001201.dat};
\addplot[only marks, mark=*, mark size=1.1pt, SoftRed] table[x=x,y=y]{data/clean/profiles/pxzz_001201.dat};
\end{axis}
\end{tikzpicture}

\captionof{figure}[东瑞股份（001201）股价与周期骨架]{东瑞股份（001201）月度收盘价（元/股）与 ZigZag 周期骨架（阈值 25\%，圆点为周期转折点）}
\end{center}

\pfl{过去周期分析} 以 25\% 的 ZigZag 阈值划分，样本区间内股价形成 3 个主升段与 3 个主跌段。 最大主升段出现在 2026 年 6 月 至 2026 年 9 月，3 个月+35.0\%；最大主跌段为 2022 年 3 月 至 2024 年 2 月，23 个月-59.2\%。 最近一次转折出现在 2026 年 9 月（上涨段自 2026 年 6 月起，3 个月+35.0\%），当前处于该段之后的震荡之中。 公司股价与申万农林牧渔指数几乎同步见底，个股并未走出独立行情。 生猪养殖的定价主线是产能去化与成本曲线的赛跑（第 \ref{sec:hog} 章）：2026 年三季度行业仍处在“产能去化未完成、价格磨底”的阶段，公司 2026 年 6 月末资产负债率 52.9\%、半年 ROE -14.3\%（未年化），在本轮下行中属于随行业同步调整的一类。

\pfl{盈利情况} 2026 年上半年营业收入 10.39 亿元（同比 -5.1\%），归母净利润 -4.25 亿元，由上年同期的盈利转为亏损。 以年报为趋势对照，2023---2025 年营业收入由 10.37 亿元变为 21.35 亿元（两年年均复合增速 +43.5\%），归母净利润由 -5.16 亿元变为 -1.61 亿元。 按半年报口径，2026 年上半年的 ROE 为 -14.3\%（未年化）、销售净利率 -40.9\%、毛利率 -6.6\%，ROE 在“生猪养殖”11 家公司中排第 \zhnumber{6} 位。 同期全样本半年 ROE 中位数 0.75\%，公司与之相差 15.07 个百分点（弱于中位数公司）。 结合公司披露，2026 年半年度业绩预告预计净利润为负值：归属于上市公司股东的净利润亏损 38,000 万元–46,0……。\footnote{\texttt{https://disc.static.szse.cn/download/disc/disk03/finalpage/2026-07-15/549a2007-84ba-44ca}}

\begin{center}\small\setlength{\tabcolsep}{4pt}
\captionof{table}[东瑞股份（001201）关键财务指标]{东瑞股份（001201）关键财务指标（合并报表口径；两个半年报期均未年化；现金短债比 $=$ 货币资金 $\div$ 短期债务）}
\begin{tabular}{@{}lrrrrr@{}}
\toprule
指标 & 2023 年 & 2024 年 & 2025 年 & 2025 年半年报 & 2026 年半年报 \\
\midrule
营业收入（亿元） & 10.37 & 15.20 & 21.35 & 10.95 & 10.39 \\
归母净利润（亿元） & -5.16 & -0.66 & -1.61 & 0.03 & -4.25 \\
毛利率（\%） & -1.4 & 21.3 & 12.3 & 17.3 & -6.6 \\
ROE（\%） & -18.1 & -1.9 & -4.9 & 0.1 & -14.3 \\
资产负债率（\%） & 42.5 & 44.1 & 46.7 & 45.0 & 52.9 \\
经营活动现金流净额（亿元） & -1.78 & 0.58 & 1.74 & 1.23 & -1.38 \\
货币资金（亿元） & 11.75 & 5.07 & 4.58 & 5.45 & 4.55 \\
有息负债合计（亿元） & 19.85 & 20.52 & 21.32 & 21.88 & 24.27 \\
现金短债比（倍） & 1.26 & 0.47 & 0.34 & 0.47 & 0.29 \\
\bottomrule
\end{tabular}
\end{center}

\begin{center}
\begin{tikzpicture}
\begin{groupplot}[
  group style={group size=2 by 1, horizontal sep=0.95cm},
  width=7.2cm, height=4.6cm,
  tick label style={font=\tiny},
  title style={font=\scriptsize\heiti},
  yticklabel style={font=\tiny},
  grid=major, grid style={LightGray, line width=0.3pt},
  axis line style={InkGray, line width=0.5pt},
  enlarge x limits={abs=0.8},
  legend style={font=\scriptsize, at={(0.5,-0.22)}, anchor=north,
                legend columns=2, draw=none, fill=none, column sep=6pt},
]
\nextgroupplot[
  ybar, bar width=4pt,
  ymin=-7.81, ymax=24.53,
  xtick={0,1,2,3,4,5.7},
  xticklabels={2021,2022,2023,2024,2025,2026H1},
  ylabel={亿元},
]
\addplot[fill=AgriGreen!75, draw=AgriGreen, bar shift=-2.6pt] table[x=i, y=rev]{data/clean/profiles/fin_001201.dat};
\addplot[fill=SoftRed!60, draw=SoftRed, bar shift=2.6pt] table[x=i, y=ni]{data/clean/profiles/fin_001201.dat};
\addplot[fill=AgriGreen!30, draw=AgriGreen, pattern=north east lines, bar shift=-2.6pt] table[x=x, y=rev]{data/clean/profiles/finh_001201.dat};
\addplot[fill=SoftRed!20, draw=SoftRed, pattern=north east lines, bar shift=2.6pt] table[x=x, y=ni]{data/clean/profiles/finh_001201.dat};
\legend{营业收入（年/半年）, 归母净利润（年/半年）}
\nextgroupplot[
  ymin=-23.8, ymax=60.0,
  xtick={0,1,2,3,4,5.7},
  xticklabels={2021,2022,2023,2024,2025,2026H1},
  ylabel={\%},
]
\addplot[AgriGold, mark=*, mark size=1.3pt] table[x=i, y=roe]{data/clean/profiles/fin_001201.dat};
\addplot[OilBlue, mark=square*, mark size=1.3pt] table[x=i, y=debt]{data/clean/profiles/fin_001201.dat};
\addplot[only marks, AgriGold, mark=*, mark size=2.0pt] table[x=x, y=roe]{data/clean/profiles/finh_001201.dat};
\addplot[only marks, OilBlue, mark=square*, mark size=2.0pt] table[x=x, y=debt]{data/clean/profiles/finh_001201.dat};
\legend{ROE, 资产负债率}
\end{groupplot}
\end{tikzpicture}
\begin{tikzpicture}
\begin{axis}[
  width=15.2cm, height=3.1cm, scale only axis,
  ymin=-4.39, ymax=27.40,
  xtick={0,1,2,3,4,5.7},
  xticklabels={2021,2022,2023,2024,2025,2026H1},
  tick label style={font=\tiny},
  yticklabel style={font=\tiny},
  ylabel={亿元}, ylabel style={font=\scriptsize},
  ybar, bar width=4pt,
  grid=major, grid style={LightGray, line width=0.3pt},
  axis line style={InkGray, line width=0.5pt},
  enlarge x limits={abs=0.8},
  legend style={font=\scriptsize, at={(0.5,-0.30)}, anchor=north,
                legend columns=3, draw=none, fill=none, column sep=10pt},
]
\addplot[fill=AgriGreen!55, draw=AgriGreen, bar shift=-3.0pt] table[x=i, y=cash]{data/clean/profiles/fin_001201.dat};
\addplot[fill=SoftRed!45, draw=SoftRed, bar shift=3.0pt] table[x=i, y=ibd]{data/clean/profiles/fin_001201.dat};
\addplot[fill=AgriGreen!25, draw=AgriGreen, pattern=north east lines, bar shift=-3.0pt] table[x=x, y=cash]{data/clean/profiles/finh_001201.dat};
\addplot[fill=SoftRed!20, draw=SoftRed, pattern=north east lines, bar shift=3.0pt] table[x=x, y=ibd]{data/clean/profiles/finh_001201.dat};
\addplot[OilBlue, mark=*, mark size=1.1pt, line width=0.8pt] table[x=i, y=ocf]{data/clean/profiles/fin_001201.dat};
\addplot[only marks, OilBlue, mark=*, mark size=1.8pt] table[x=x, y=ocf]{data/clean/profiles/finh_001201.dat};
\legend{货币资金, 有息负债, 经营活动现金流净额}
\end{axis}
\end{tikzpicture}

\captionof{figure}[东瑞股份（001201）近年主要财务指标]{东瑞股份（001201）近年主要财务指标（左上：营业收入与归母净利润；右上：ROE 与资产负债率；下：货币资金、有息负债与经营活动现金流净额。
2021---2025 年为年报口径，2026H1 为 2026 年半年报/半年末，与其前一年报之间留出间隔、以斜纹或空心点区分；半年报数据未年化）}
\end{center}

\pfl{财务分析} 2026 年 6 月末资产负债率 52.9\%，较 2025 年末 上升 6.2 个百分点（样本中位数 52.2\%，所处三级行业内排第 \zhnumber{3} 位）。 2026 年 6 月末货币资金 4.55 亿元、短期债务（短期借款与一年内到期非流动负债）15.69 亿元、长期债务 8.58 亿元，有息负债合计 24.27 亿元，现金短债比 0.29 倍——缺口明显。 净有息负债 19.72 亿元，相当于其 2026 年上半年营业收入的 190\%。 流动比率 0.67、速动比率 0.26，短期偿付压力较大。 2026 年上半年经营活动现金流净额 -1.38 亿元，经营现金流与利润同时为负，现金消耗较快。 2026 年 6 月末在建工程 0.68 亿元，较上年末增加 0.27 亿元，仍有在建投入。

\pfl{重大战略转型} 公司在报告期内的战略动作可归为 3 项：2024-12-06 终止“惠州东瑞多层楼房生猪养殖项目”并将该项目募集资金及利息、理财收益永久补充流动资金，相关议案经第三届董事会第……；2026-06-30 受猪价下行影响业绩承压后，公司经营重心转向提质增效；2026-06-30 公司坚持以生猪全产业链布局和供港澳活大猪业务为特色，2025 年生猪养殖行业收入占比达 96.33\%。 公告口径上，近三年（2023 年 9 月---2026 年 9 月）公司披露重大重组与并购类公告 0 条、对外投资与转型类公告 2 条、回购与股权激励类公告 18 条，合计公告 411 条。 第一大行业仍是“养殖行业”，收入占比 95.3\%，与 2021 年年报相比没有方向性变化。 综合判断，公司在报告期内有明确的外延扩张或业务结构调整动作，转型方向与融资安排之间存在直接关联。\footnote{\texttt{https://money.finance.sina.com.cn/corp/view/vCB\_AllBulletinDetail.php}}\footnote{\texttt{https://stock.stockstar.com/RB2026090700018540.shtml}}

\begin{center}\small\setlength{\tabcolsep}{4pt}
\captionof{table}[东瑞股份（001201）历史融资与资本运作]{东瑞股份（001201）历史融资与资本运作（据公司公告与公开报道整理；金额为披露值，含首次公开发行、再融资、债券、银行借款与重整投资等）}
\begin{tabular}{@{}p{1.5cm}p{3.2cm}p{2.4cm}p{6.3cm}@{}}
\toprule
时间 & 方式 & 规模 & 用途与说明 \\
\midrule
2021-04 & IPO（深圳证券交易所主板） & 公开发行 3,167 万股、发行价 63.38 元/股，募集资金总额 20.07 亿元，扣除发行费用后净额 18.10 亿元 & 紫金东瑞农业发展有限公司富竹生态养殖项目、连平东瑞农牧发展有限公司传统养殖改高床养殖节能…… \\
2021-05-21 & 募集资金置换自筹资金 & 以募集资金置换预先投入募投项目的自筹资金 43,962.75 万元 & 置换预先投入募投项目的自筹资金；截至 2021 年 12 月 31 日…… \\
2023-12-28 & 向特定对象发行股票（2023 年度定增） & 发行 44,978,401 股，募集资金总额 92,475.59 万元（含发行费用 1,374.39 万元），净额 91,101.20 万元 & 东源东瑞农牧发展有限公司东源县船塘现代农业综合体项目（黄沙基地）及补充公司流动资金项目（…… \\
\bottomrule
\end{tabular}
\end{center}

\pfl{历史融投资情况} 从现金流量表看，2025 年公司取得借款收到的现金 12.66 亿元、偿还债务支付的现金 12.29 亿元、吸收权益性投资 — 亿元，筹资活动现金流净额 -0.46 亿元（2023 年为取得借款 16.35 亿元、偿还债务 6.26 亿元）。 2026 年上半年筹资活动现金流净额为 2.97 亿元（取得借款 8.78 亿元、偿还债务 5.83 亿元，未年化），仍处于净融入状态。 2025 年分配股利、利润或偿付利息支付的现金合计 0.67 亿元。 公开信息中与公司直接相关的融资与资本运作包括：IPO（深圳证券交易所主板）、募集资金置换自筹资金、向特定对象发行股票（2023 年度定增）。 其中规模较大的两笔为：2023-12-28向特定对象发行股票（2023 年度定增）（发行 44,978,401 股，募集资金总额 92,475.……）；2021-05-21募集资金置换自筹资金（以募集资金置换预先投入募投项目的自筹资金 43,962.75……）。 股本方面，近三年披露股本变动 6 次（限售股份上市 2 次、增发新股上市 1 次）；总股本由 2021 年末的 1.77 亿股变为 2.58 亿股（+45.36\%）。 分红方面，近三年未实施现金分红，在全部样本中属于分红能力偏弱的一端。 近三年“再融资”类公告 43 条，涉及定向增发、募集资金使用。\footnote{\texttt{https://finance.sina.com.cn/jjxw/2026-09-24/doc-inisxnuv5181817.shtml}}\footnote{\texttt{https://disc.static.szse.cn/disc/disk03/finalpage/2022-04-28/4b56bb3a-bdfa-410c-a2f5-e5e}}\footnote{\texttt{https://www.100est.com/res/financial-report/r2023/SZ001201\_202404291631627113.pdf}}

\pfl{未来潜在融资需求分析} 2026 年 6 月末货币资金 4.55 亿元，而短期债务 15.69 亿元（现金短债比 0.29 倍），2026 年上半年经营现金流净额 -1.38 亿元。按“短期债务减去账面货币资金”的滚动偿付口径，静态缺口约 11.14 亿元；若再计入上半年亏损的年化额（8.50 亿元），维持一年正常经营所需的外部资金约 11.14---19.64 亿元。 公司 2026 年上半年归母净利润为负（-4.25 亿元），但 6 月末资产负债率 52.9\% 尚未进入第一档（亏损且负债率 $>$65\%）的区间，当前融资需求以补充流动资金、支撑低谷期经营性支出为主。 公开披露的下一步融资线索包括：2024-12-31 IPO 与 2023 年度定增的募集资金均已全部使用完毕、募投项目全部结项、专户全部注销；截至 2024 年末无存续的募……。 资金需求还要叠加在建工程：期末在建工程 0.68 亿元、2026 年上半年购建固定资产等支出 0.89 亿元（未年化），这部分支出在周期底部通常会被主动放缓。\footnote{\texttt{https://static.cninfo.com.cn/finalpage/2025-04-26/1223312424.PDF}}


\subsubsection{天山生物（300313）}
\label{co:300313}

\pfl{公司基本情况} 天山生物成立于 2003 年，新疆昌吉的肉牛全产业链畜牧生物企业，以牛品种改良、冻精与活畜养殖起家；因23.73亿元跨界并购大象广告遭合同诈骗巨亏，2025年4月被实施退市风险警示，2026年6月摘帽并恢复“天山生物”简称。 公司于 2012-04-25 在深交所创业板首发上市，注册资本 2.40 亿元，总股本 3.13 亿股。 截至 2026-09-24收盘价 12.37 元，流通市值 24.4 亿元，近一年+46.9\%，流通市值在三级行业“其他养殖”的 3 家公司中排第 \zhnumber{2} 位，近一年涨跌幅高于全样本中位数（-10.1\%）57.0 个百分点。 按 2026 年半年报披露的产品构成，收入占比前三位为肉类 60.1\%；饲料 13.7\%；活畜 11.6\%。 与 2021 年年报相比，第一大产品由“活畜”变为“肉类”，主业重心已经迁移。 发展过程中的关键节点包括：经证监许可〔2012〕380号核准首次公开发行2,273万股，发行价13.00元/股，2……（2012）；向呼图壁县天山农业发展有限公司、天山农牧业发展有限公司等发行股份购买资产并募集配套资金（……（2016）；2017年9月与陈德宏、华融渝稳、华融天泽等36名交易对方签署《发行股份及支付现金购买资……（2017）。\footnote{\texttt{http://static.cninfo.com.cn/finalpage/2016-03-26/1202085151.PDF}}\footnote{\texttt{http://static.cninfo.com.cn/finalpage/2018-01-03/1204301717.PDF}}\footnote{\texttt{http://static.cninfo.com.cn/finalpage/2024-03-30/1219474293.PDF}}

\pfl{历史股价} 本报告样本区间（2021 年 1 月 至 2026 年 9 月）内，股价由 14.95 元变为 12.37 元，累计 -17.3\%，区间最高 15.12 元（2021 年 2 月）、最低 4.76 元（2022 年 4 月）；当前价相当于区间最高价的 81.8\%，为区间最低价的 2.60 倍。 月度收益的年化波动率 70.5\%，高于全样本中位数 37.4\%，在三级行业“其他养殖”的 3 家公司中波动率排第 \zhnumber{1} 位。 从事件看，2023 年 5 月 至 2023 年 10 月的上涨（+162.6\%）与公司控股股东湖州皓辉与中电投新农创科技有限公司签署控制权转让框架协议，拟协议转让……（2023-10）在时间上重合；2025 年 4 月 至 2026 年 9 月的上涨（+150.4\%）与公司披露《2026年向特定对象发行A股股票募集说明书（申报稿）》，募集资金拟全部……（2026-08-26）在时间上重合——股价的拐点多数能在公司自身的资本运作与风险事件上找到对应。\footnote{\texttt{https://m.21jingji.com/article/20231010/herald/9acdfefbf4ff77f27434167e56bd3617.html}}\footnote{\texttt{https://money.finance.sina.com.cn/corp/view/vCB\_AllBulletinDetail.php}}

\begin{center}
\begin{tikzpicture}
\begin{axis}[
  width=12.6cm, height=3.9cm, scale only axis,
  xmin=2020.922, xmax=2026.888, ymin=4.43, ymax=16.18,
  xtick={2021,2022,2023,2024,2025,2026},
  yearticks,
  yticklabel style={font=\scriptsize},
  ylabel={元/股}, ylabel style={font=\scriptsize},
  grid=major, grid style={LightGray, line width=0.3pt},
  axis line style={InkGray, line width=0.5pt},
  every axis plot/.append style={line width=0.9pt},
]
\addplot[AgriGreen] table[x=x,y=y]{data/clean/profiles/px_300313.dat};
\addplot[SoftRed, dashed, line width=0.7pt] table[x=x,y=y]{data/clean/profiles/pxzz_300313.dat};
\addplot[only marks, mark=*, mark size=1.1pt, SoftRed] table[x=x,y=y]{data/clean/profiles/pxzz_300313.dat};
\end{axis}
\end{tikzpicture}

\captionof{figure}[天山生物（300313）股价与周期骨架]{天山生物（300313）月度收盘价（元/股）与 ZigZag 周期骨架（阈值 25\%，圆点为周期转折点）}
\end{center}

\pfl{过去周期分析} 以 25\% 的 ZigZag 阈值划分，样本区间内股价形成 4 个主升段与 4 个主跌段。 最大主升段出现在 2023 年 5 月 至 2023 年 10 月，5 个月+162.6\%；最大主跌段为 2023 年 10 月 至 2025 年 4 月，18 个月-61.4\%。 最近一次转折出现在 2026 年 9 月（上涨段自 2025 年 4 月起，17 个月+150.4\%），当前处于该段之后的震荡之中。 公司股价较申万农林牧渔指数提前约 49 个月见底，是板块中较早定价周期反转的公司之一。 牛、羊、水产的产能周期长于生猪，第 \ref{sec:other-livestock} 章判断肉牛供给的实质性收缩要到 2027 年之后才能兑现，这一时滞决定了个股的股价弹性，公司 2026 年 6 月末资产负债率 82.0\%，高于全样本中位数 52.2\%，其股价因此更多反映资金面与信用风险，而不是价格弹性本身。

\pfl{盈利情况} 2026 年上半年营业收入 2.20 亿元（同比 +427.7\%），归母净利润 0.04 亿元，实现扭亏。 以年报为趋势对照，2023---2025 年营业收入由 1.37 亿元变为 2.00 亿元（两年年均复合增速 +20.7\%），归母净利润由 -0.22 亿元变为 -0.27 亿元。 按半年报口径，2026 年上半年的 ROE 为 12.1\%（未年化）、销售净利率 2.4\%、毛利率 9.1\%，ROE 在“其他养殖”3 家公司中排第 \zhnumber{1} 位。 行业同口径半年 ROE 的中位数为 0.75\%，该公司高于中位数 11.34 个百分点。 结合公司披露，因大象广告事件，2018年对大象广告长期股权投资计提资产减值准备17.96亿元，并对公司为大象广告担保形……。\footnote{\texttt{https://news.stcn.com/sd/202110/t20211011\_3745305.html}}

\begin{center}\small\setlength{\tabcolsep}{4pt}
\captionof{table}[天山生物（300313）关键财务指标]{天山生物（300313）关键财务指标（合并报表口径；两个半年报期均未年化；现金短债比 $=$ 货币资金 $\div$ 短期债务）}
\begin{tabular}{@{}lrrrrr@{}}
\toprule
指标 & 2023 年 & 2024 年 & 2025 年 & 2025 年半年报 & 2026 年半年报 \\
\midrule
营业收入（亿元） & 1.37 & 1.38 & 2.00 & 0.42 & 2.20 \\
归母净利润（亿元） & -0.22 & -0.66 & -0.27 & -0.08 & 0.04 \\
毛利率（\%） & 17.0 & -5.2 & 7.4 & 18.4 & 9.1 \\
ROE（\%） & -30.1 & 0.0 & --- & 0.0 & 12.1 \\
资产负债率（\%） & 77.2 & 97.7 & 84.1 & 100.8 & 82.0 \\
经营活动现金流净额（亿元） & -0.14 & 0.14 & 0.02 & -0.03 & -0.11 \\
货币资金（亿元） & 0.17 & 0.25 & 0.39 & 0.19 & 0.26 \\
有息负债合计（亿元） & 0.86 & 0.84 & 0.38 & 0.83 & 0.93 \\
现金短债比（倍） & 0.20 & 0.48 & 18.90 & 0.38 & 0.45 \\
\bottomrule
\end{tabular}
\end{center}

\begin{center}
\begin{tikzpicture}
\begin{groupplot}[
  group style={group size=2 by 1, horizontal sep=0.95cm},
  width=7.2cm, height=4.6cm,
  tick label style={font=\tiny},
  title style={font=\scriptsize\heiti},
  yticklabel style={font=\tiny},
  grid=major, grid style={LightGray, line width=0.3pt},
  axis line style={InkGray, line width=0.5pt},
  enlarge x limits={abs=0.8},
  legend style={font=\scriptsize, at={(0.5,-0.22)}, anchor=north,
                legend columns=2, draw=none, fill=none, column sep=6pt},
]
\nextgroupplot[
  ybar, bar width=4pt,
  ymin=-0.95, ymax=2.54,
  xtick={0,1,2,3,4,5.7},
  xticklabels={2021,2022,2023,2024,2025,2026H1},
  ylabel={亿元},
]
\addplot[fill=AgriGreen!75, draw=AgriGreen, bar shift=-2.6pt] table[x=i, y=rev]{data/clean/profiles/fin_300313.dat};
\addplot[fill=SoftRed!60, draw=SoftRed, bar shift=2.6pt] table[x=i, y=ni]{data/clean/profiles/fin_300313.dat};
\addplot[fill=AgriGreen!30, draw=AgriGreen, pattern=north east lines, bar shift=-2.6pt] table[x=x, y=rev]{data/clean/profiles/finh_300313.dat};
\addplot[fill=SoftRed!20, draw=SoftRed, pattern=north east lines, bar shift=2.6pt] table[x=x, y=ni]{data/clean/profiles/finh_300313.dat};
\legend{营业收入（年/半年）, 归母净利润（年/半年）}
\nextgroupplot[
  ymin=-44.3, ymax=110.9,
  xtick={0,1,2,3,4,5.7},
  xticklabels={2021,2022,2023,2024,2025,2026H1},
  ylabel={\%},
]
\addplot[AgriGold, mark=*, mark size=1.3pt] table[x=i, y=roe]{data/clean/profiles/fin_300313.dat};
\addplot[OilBlue, mark=square*, mark size=1.3pt] table[x=i, y=debt]{data/clean/profiles/fin_300313.dat};
\addplot[only marks, AgriGold, mark=*, mark size=2.0pt] table[x=x, y=roe]{data/clean/profiles/finh_300313.dat};
\addplot[only marks, OilBlue, mark=square*, mark size=2.0pt] table[x=x, y=debt]{data/clean/profiles/finh_300313.dat};
\legend{ROE, 资产负债率}
\end{groupplot}
\end{tikzpicture}
\begin{tikzpicture}
\begin{axis}[
  width=15.2cm, height=3.1cm, scale only axis,
  ymin=-0.53, ymax=1.67,
  xtick={0,1,2,3,4,5.7},
  xticklabels={2021,2022,2023,2024,2025,2026H1},
  tick label style={font=\tiny},
  yticklabel style={font=\tiny},
  ylabel={亿元}, ylabel style={font=\scriptsize},
  ybar, bar width=4pt,
  grid=major, grid style={LightGray, line width=0.3pt},
  axis line style={InkGray, line width=0.5pt},
  enlarge x limits={abs=0.8},
  legend style={font=\scriptsize, at={(0.5,-0.30)}, anchor=north,
                legend columns=3, draw=none, fill=none, column sep=10pt},
]
\addplot[fill=AgriGreen!55, draw=AgriGreen, bar shift=-3.0pt] table[x=i, y=cash]{data/clean/profiles/fin_300313.dat};
\addplot[fill=SoftRed!45, draw=SoftRed, bar shift=3.0pt] table[x=i, y=ibd]{data/clean/profiles/fin_300313.dat};
\addplot[fill=AgriGreen!25, draw=AgriGreen, pattern=north east lines, bar shift=-3.0pt] table[x=x, y=cash]{data/clean/profiles/finh_300313.dat};
\addplot[fill=SoftRed!20, draw=SoftRed, pattern=north east lines, bar shift=3.0pt] table[x=x, y=ibd]{data/clean/profiles/finh_300313.dat};
\addplot[OilBlue, mark=*, mark size=1.1pt, line width=0.8pt] table[x=i, y=ocf]{data/clean/profiles/fin_300313.dat};
\addplot[only marks, OilBlue, mark=*, mark size=1.8pt] table[x=x, y=ocf]{data/clean/profiles/finh_300313.dat};
\legend{货币资金, 有息负债, 经营活动现金流净额}
\end{axis}
\end{tikzpicture}

\captionof{figure}[天山生物（300313）近年主要财务指标]{天山生物（300313）近年主要财务指标（左上：营业收入与归母净利润；右上：ROE 与资产负债率；下：货币资金、有息负债与经营活动现金流净额。
2021---2025 年为年报口径，2026H1 为 2026 年半年报/半年末，与其前一年报之间留出间隔、以斜纹或空心点区分；半年报数据未年化）}
\end{center}

\pfl{财务分析} 2026 年 6 月末资产负债率 82.0\%，较 2025 年末 下降 2.0 个百分点（样本中位数 52.2\%，所处三级行业内排第 \zhnumber{3} 位）。 2026 年 6 月末货币资金 0.26 亿元、短期债务（短期借款与一年内到期非流动负债）0.58 亿元、长期债务 0.35 亿元，有息负债合计 0.93 亿元，现金短债比 0.45 倍——覆盖偏紧。 净有息负债 0.66 亿元，相当于其 2026 年上半年营业收入的 30\%。 流动比率 1.01、速动比率 0.60，短期指标尚可。 2026 年上半年经营活动现金流净额 -0.11 亿元，净利润为正但经营现金流为负，利润的现金含量偏低。 公司披露的财务与合规风险包括：向深圳证券交易所申请撤销对公司股票交易实施的退市风险警示及其他风险警示；2024年末归属于上市公司股东的净资产为-377.93万元，扣除后营业收入9,228.1……。\footnote{\texttt{http://static.cninfo.com.cn/finalpage/2026-04-25/1225191616.PDF}}\footnote{\texttt{http://static.cninfo.com.cn/finalpage/2025-04-29/1223376987.PDF}}

\pfl{重大战略转型} 从公开披露看，公司的战略推进包括：2023-10 控股股东湖州皓辉与中电投新农创科技有限公司签署控制权转让框架协议，拟协议转让40,718,359股（占13.01……；2025-04 公司提出“大肉牛战略”：上游强化肉牛品种改良与冻精、胚胎等遗传物质业务，下游启动屠宰、流通环节及牛肉深加工与品牌……；2025-04 公司称将适时启动资本运作项目以改善资本结构，并通过大股东借款协议展期、银行短期借款置换为中长期等方式优化债务结构；2025-09 厦门舍德受让湖州皓辉所持5,386.13万股（占总股本17.21\%）及湖州中植对公司享有的债权（本金7,649万……。 公告口径上，近三年（2023 年 9 月---2026 年 9 月）公司披露重大重组与并购类公告 2 条、对外投资与转型类公告 3 条、回购与股权激励类公告 2 条，合计公告 406 条。 主营结构的迁移已经落到报表上：2021 年年报的第一大产品为“活畜”（63.8\%），2026 年半年报变为“肉类”（60.1\%）。 综合判断，公司当前的战略重心不是扩张而是化解债务与完成重整，业务层面的调整让位于司法重整程序。\footnote{\texttt{https://m.21jingji.com/article/20231010/herald/9acdfefbf4ff77f27434167e56bd3617.html}}\footnote{\texttt{http://static.cninfo.com.cn/finalpage/2025-04-29/1223376987.PDF}}

\begin{center}\small\setlength{\tabcolsep}{4pt}
\captionof{table}[天山生物（300313）历史融资与资本运作]{天山生物（300313）历史融资与资本运作（据公司公告与公开报道整理；金额为披露值，含首次公开发行、再融资、债券、银行借款与重整投资等）}
\begin{tabular}{@{}p{1.5cm}p{3.2cm}p{2.4cm}p{6.3cm}@{}}
\toprule
时间 & 方式 & 规模 & 用途与说明 \\
\midrule
2012-04 & IPO（深交所创业板） & 发行2,273万股，发行价13.00元/股，募集资金总额29,549.00万元，募集资金净额26,099.82万元 & 募投项目建设；2012年年度报告披露口径 \\
2016-06 & 发行股份购买资产并募集配套资金（证监许可〔2015〕2393号） & 募集配套资金3,900万元，扣除发行费用671.35万元后净额3,228.65万元，2016年6月21日全部到位；新增股份2016年7月12日上市 & 配套募集资金用途相关项目建设；配套资金发行2,401,477股，发行价16.24元/股 \\
2017-09 & 发行股份及支付现金购买资产（收购大象广告96.21\%股权） & 交易作价237,261.42万元，其中现金对价57,696.41万元、股份对价179,565.04万元 & 跨界并购户外广告资产；该次交易的会计处理于2023年年报中作追溯调整 \\
2026-03 & 向特定对象发行A股股票（预案） & 发行数量不超过5,900万股，发行价格5.91元/股，募集资金总额不超过34,869.00万元 & 扣除发行费用后拟全部用于补充流动资金；发行对象为公司实际控制人陈明艺（不超过900万股）…… \\
\bottomrule
\end{tabular}
\end{center}

\pfl{历史融投资情况} 从现金流量表看，2025 年公司取得借款收到的现金 — 亿元、偿还债务支付的现金 0.50 亿元、吸收权益性投资 0.02 亿元，筹资活动现金流净额 0.17 亿元（2023 年为取得借款 1.07 亿元、偿还债务 1.18 亿元）。 2026 年上半年筹资活动现金流净额为 0.01 亿元（取得借款 — 亿元、偿还债务 0.01 亿元，未年化），仍处于净融入状态。 2025 年分配股利、利润或偿付利息支付的现金合计 0.05 亿元。 公开信息中与公司直接相关的融资与资本运作包括：IPO（深交所创业板）、发行股份及支付现金购买资产（收购大象广告96.21\%股权）、发行股份购买资产并募集配套资金（证监许可〔2015〕2393号）、向特定对象发行A股股票（预案）。 其中规模较大的两笔为：2016-06发行股份购买资产并募集配套资金（证监许可〔2015〕2393号）（募集配套资金3,900万元，扣除发行费用671.35万元后净……）；2012-04IPO（深交所创业板）（发行2,273万股，发行价13.00元/股，募集资金总额29……）。 股本方面，近三年披露股本变动 3 次；总股本由 2021 年末的 3.13 亿股变为 3.13 亿股（+0.00\%）。 分红方面，近三年未实施现金分红，与重整期间的资金约束一致。 近三年“再融资”类公告 28 条，涉及定向增发、发行股份购买资产、募集资金使用。\footnote{\texttt{https://money.finance.sina.com.cn/corp/view/vCB\_AllBulletinDetail.php}}\footnote{\texttt{http://file.finance.sina.com.cn/211.154.219.97:9494/MRGG/CNSESZ\_STOCK/2019/2019-8/2019-0}}\footnote{\texttt{http://static.cninfo.com.cn/finalpage/2024-03-30/1219474293.PDF}}

\pfl{未来潜在融资需求分析} 2026 年 6 月末货币资金 0.26 亿元，而短期债务 0.58 亿元（现金短债比 0.45 倍），2026 年上半年经营现金流净额 -0.11 亿元。账面货币资金可以覆盖短期债务，缺口不体现在静态偿付层面。 公司 2026 年 6 月末资产负债率 82.0\%、2026 年上半年 ROE 12.1\%（未年化），资产负债表尚有余量，融资需求不是“补血型”的刚性需求，而是围绕并购整合、项目投入与营运资金的主动性需求。 公开披露的下一步融资线索包括：2026-03-27 第六届董事会第三次会议审议通过向特定对象发行A股股票议案，拟募集资金不超过34,869.00万元全部用于补充流动资金；2026-06-18 董事会审议通过《2026年度向特定对象发行A股股票预案（修订稿）》等议案（此前2026年5月22日2025年度股东会已审……；2026-08-26 披露《2026年向特定对象发行A股股票募集说明书（申报稿）》，募集资金拟全部用于补充流动资金，不涉及投资项目报批事项。\footnote{\texttt{http://static.cninfo.com.cn/finalpage/2025-04-29/1223376987.PDF}}\footnote{\texttt{https://www.cnfin.com/announ/detail/index.html}}\footnote{\texttt{https://money.finance.sina.com.cn/corp/view/vCB\_AllBulletinDetail.php}}


\subsubsection{民和股份（002234）}
\label{co:002234}

\pfl{公司基本情况} 1997 年设立至今，民和股份山东蓬莱的白羽肉鸡一体化企业，以父母代肉种鸡饲养和商品代鸡苗销售为核心，兼营鸡肉制品与有机废弃物资源化利用；白羽鸡行业已连续亏损多年，公司2025年、2026年上半年持续亏损，母公司存在累计未弥补亏损。 公司于 2008-05-16 在深交所主板首发上市，注册资本 3.49 亿元，总股本 3.49 亿股。 截至 2026-09-24收盘价 7.75 元，流通市值 23.4 亿元，近一年-10.0\%，流通市值在三级行业“肉鸡养殖”的 8 家公司中排第 \zhnumber{6} 位，近一年涨跌幅高于全样本中位数（-10.1\%）0.1 个百分点。 按 2026 年半年报披露的产品构成，收入占比前三位为销售商品 99.7\%；让渡资产使用权 0.3\%。 与 2021 年年报相比，第一大产品由“雏鸡”变为“销售商品”，主业重心已经迁移。 发展过程中的关键节点包括：2000年12月1日整体变更为山东民和牧业股份有限公司，变更后股本总额4,025万元（2000）；2006年12月25日经第三次临时股东大会决议，以2005年末未分配利润向全体股东每10……（2006）；经证监发行字〔2008〕547号核准首次公开发行2,700万股，发行价10.61元/股，……（2008）。\footnote{\texttt{https://baike.baidu.com/item/\%E6\%B0\%91\%E5\%92\%8C\%E8\%82\%A1\%E4\%BB\%BD/12650038}}\footnote{\texttt{https://vip.stock.finance.sina.com.cn/corp/view/vISSUE\_MarketBulletinDetail.php}}

\pfl{历史股价} 本报告样本区间（2021 年 1 月 至 2026 年 9 月）内，股价由 13.34 元变为 7.75 元，累计 -41.9\%，区间最高 22.65 元（2023 年 1 月）、最低 6.41 元（2026 年 6 月）；当前价相当于区间最高价的 34.2\%，为区间最低价的 1.21 倍。 月度收益的年化波动率 33.0\%，低于全样本中位数 37.4\%，在三级行业“肉鸡养殖”的 8 家公司中波动率排第 \zhnumber{7} 位。 从事件看，2021 年 1 月 至 2023 年 1 月的上涨（+69.8\%）与公司23名发行对象以现金认购，新增股份2021年4月23日上市；资金于2021年……（2021-04）在时间上重合——股价的拐点多数能在公司自身的资本运作与风险事件上找到对应。\footnote{\texttt{https://vip.stock.finance.sina.com.cn/corp/view/vISSUE\_MarketBulletinDetail.php}}

\begin{center}
\begin{tikzpicture}
\begin{axis}[
  width=12.6cm, height=3.9cm, scale only axis,
  xmin=2020.922, xmax=2026.888, ymin=5.96, ymax=24.24,
  xtick={2021,2022,2023,2024,2025,2026},
  yearticks,
  yticklabel style={font=\scriptsize},
  ylabel={元/股}, ylabel style={font=\scriptsize},
  grid=major, grid style={LightGray, line width=0.3pt},
  axis line style={InkGray, line width=0.5pt},
  every axis plot/.append style={line width=0.9pt},
]
\addplot[AgriGreen] table[x=x,y=y]{data/clean/profiles/px_002234.dat};
\addplot[SoftRed, dashed, line width=0.7pt] table[x=x,y=y]{data/clean/profiles/pxzz_002234.dat};
\addplot[only marks, mark=*, mark size=1.1pt, SoftRed] table[x=x,y=y]{data/clean/profiles/pxzz_002234.dat};
\end{axis}
\end{tikzpicture}

\captionof{figure}[民和股份（002234）股价与周期骨架]{民和股份（002234）月度收盘价（元/股）与 ZigZag 周期骨架（阈值 25\%，圆点为周期转折点）}
\end{center}

\pfl{过去周期分析} 以 25\% 的 ZigZag 阈值划分，样本区间内股价形成 1 个主升段与 1 个主跌段。 最大主升段出现在 2021 年 1 月 至 2023 年 1 月，24 个月+69.8\%；最大主跌段为 2023 年 1 月 至 2026 年 6 月，41 个月-71.7\%。 最近一次转折出现在 2026 年 6 月（下跌段自 2023 年 1 月起，41 个月-71.7\%），当前处于该段的延续之中。 公司股价与申万农林牧渔指数几乎同步见底，个股并未走出独立行情。 禽链的产能周期只有 6---12 个月、反转快于生猪（第 \ref{sec:other-livestock} 节），2026 年前四个月毛鸡养殖利润同比大幅回升，链条利润的修复先于生猪，公司 2026 年 6 月末资产负债率 53.9\%、半年 ROE -5.5\%（未年化），在本轮下行中属于随行业同步调整的一类。

\pfl{盈利情况} 2026 年上半年营业收入 10.27 亿元（同比 +4.2\%），归母净利润 -0.95 亿元，亏损同比收窄 1.19 亿元。 以年报为趋势对照，2023---2025 年营业收入由 20.74 亿元变为 21.62 亿元（两年年均复合增速 +2.1\%），归母净利润由 -3.86 亿元变为 -2.69 亿元。 按半年报口径，2026 年上半年的 ROE 为 -5.5\%（未年化）、销售净利率 -9.7\%、毛利率 5.0\%，ROE 在“肉鸡养殖”8 家公司中排第 \zhnumber{8} 位。 同期全样本半年 ROE 中位数 0.75\%，公司与之相差 6.28 个百分点（弱于中位数公司）。 结合公司披露，2026年上半年营业收入10.27亿元（同比+4.2\%），归母净利润-9,524万元，扣非净利润-9,9……；母公司财务报表存在累计未弥补亏损239,331,714.91元，公司不满足实施现金分红的前提条件。\footnote{\texttt{https://finance.sina.com.cn/stock/zqgd/2026-08-21/doc-inipakut1652161.shtml}}\footnote{\texttt{http://money.finance.sina.com.cn/corp/view/vCB\_AllBulletinDetail.php}}

\begin{center}\small\setlength{\tabcolsep}{4pt}
\captionof{table}[民和股份（002234）关键财务指标]{民和股份（002234）关键财务指标（合并报表口径；两个半年报期均未年化；现金短债比 $=$ 货币资金 $\div$ 短期债务）}
\begin{tabular}{@{}lrrrrr@{}}
\toprule
指标 & 2023 年 & 2024 年 & 2025 年 & 2025 年半年报 & 2026 年半年报 \\
\midrule
营业收入（亿元） & 20.74 & 21.60 & 21.62 & 9.86 & 10.27 \\
归母净利润（亿元） & -3.86 & -2.49 & -2.69 & -2.14 & -0.95 \\
毛利率（\%） & 3.3 & 8.1 & 4.8 & -0.3 & 5.0 \\
ROE（\%） & -15.6 & -11.5 & -14.1 & -11.1 & -5.5 \\
资产负债率（\%） & 39.4 & 47.3 & 51.8 & 52.0 & 53.9 \\
经营活动现金流净额（亿元） & -0.65 & -0.56 & 1.40 & 0.55 & 0.58 \\
货币资金（亿元） & 5.88 & 6.22 & 5.27 & 6.14 & 5.14 \\
有息负债合计（亿元） & 11.25 & 14.80 & 14.45 & 14.83 & 13.65 \\
现金短债比（倍） & 0.53 & 0.44 & 0.40 & 0.45 & 0.42 \\
\bottomrule
\end{tabular}
\end{center}

\begin{center}
\begin{tikzpicture}
\begin{groupplot}[
  group style={group size=2 by 1, horizontal sep=0.95cm},
  width=7.2cm, height=4.6cm,
  tick label style={font=\tiny},
  title style={font=\scriptsize\heiti},
  yticklabel style={font=\tiny},
  grid=major, grid style={LightGray, line width=0.3pt},
  axis line style={InkGray, line width=0.5pt},
  enlarge x limits={abs=0.8},
  legend style={font=\scriptsize, at={(0.5,-0.22)}, anchor=north,
                legend columns=2, draw=none, fill=none, column sep=6pt},
]
\nextgroupplot[
  ybar, bar width=4pt,
  ymin=-7.13, ymax=24.75,
  xtick={0,1,2,3,4,5.7},
  xticklabels={2021,2022,2023,2024,2025,2026H1},
  ylabel={亿元},
]
\addplot[fill=AgriGreen!75, draw=AgriGreen, bar shift=-2.6pt] table[x=i, y=rev]{data/clean/profiles/fin_002234.dat};
\addplot[fill=SoftRed!60, draw=SoftRed, bar shift=2.6pt] table[x=i, y=ni]{data/clean/profiles/fin_002234.dat};
\addplot[fill=AgriGreen!30, draw=AgriGreen, pattern=north east lines, bar shift=-2.6pt] table[x=x, y=rev]{data/clean/profiles/finh_002234.dat};
\addplot[fill=SoftRed!20, draw=SoftRed, pattern=north east lines, bar shift=2.6pt] table[x=x, y=ni]{data/clean/profiles/finh_002234.dat};
\legend{营业收入（年/半年）, 归母净利润（年/半年）}
\nextgroupplot[
  ymin=-21.2, ymax=60.9,
  xtick={0,1,2,3,4,5.7},
  xticklabels={2021,2022,2023,2024,2025,2026H1},
  ylabel={\%},
]
\addplot[AgriGold, mark=*, mark size=1.3pt] table[x=i, y=roe]{data/clean/profiles/fin_002234.dat};
\addplot[OilBlue, mark=square*, mark size=1.3pt] table[x=i, y=debt]{data/clean/profiles/fin_002234.dat};
\addplot[only marks, AgriGold, mark=*, mark size=2.0pt] table[x=x, y=roe]{data/clean/profiles/finh_002234.dat};
\addplot[only marks, OilBlue, mark=square*, mark size=2.0pt] table[x=x, y=debt]{data/clean/profiles/finh_002234.dat};
\legend{ROE, 资产负债率}
\end{groupplot}
\end{tikzpicture}
\begin{tikzpicture}
\begin{axis}[
  width=15.2cm, height=3.1cm, scale only axis,
  ymin=-4.66, ymax=16.92,
  xtick={0,1,2,3,4,5.7},
  xticklabels={2021,2022,2023,2024,2025,2026H1},
  tick label style={font=\tiny},
  yticklabel style={font=\tiny},
  ylabel={亿元}, ylabel style={font=\scriptsize},
  ybar, bar width=4pt,
  grid=major, grid style={LightGray, line width=0.3pt},
  axis line style={InkGray, line width=0.5pt},
  enlarge x limits={abs=0.8},
  legend style={font=\scriptsize, at={(0.5,-0.30)}, anchor=north,
                legend columns=3, draw=none, fill=none, column sep=10pt},
]
\addplot[fill=AgriGreen!55, draw=AgriGreen, bar shift=-3.0pt] table[x=i, y=cash]{data/clean/profiles/fin_002234.dat};
\addplot[fill=SoftRed!45, draw=SoftRed, bar shift=3.0pt] table[x=i, y=ibd]{data/clean/profiles/fin_002234.dat};
\addplot[fill=AgriGreen!25, draw=AgriGreen, pattern=north east lines, bar shift=-3.0pt] table[x=x, y=cash]{data/clean/profiles/finh_002234.dat};
\addplot[fill=SoftRed!20, draw=SoftRed, pattern=north east lines, bar shift=3.0pt] table[x=x, y=ibd]{data/clean/profiles/finh_002234.dat};
\addplot[OilBlue, mark=*, mark size=1.1pt, line width=0.8pt] table[x=i, y=ocf]{data/clean/profiles/fin_002234.dat};
\addplot[only marks, OilBlue, mark=*, mark size=1.8pt] table[x=x, y=ocf]{data/clean/profiles/finh_002234.dat};
\legend{货币资金, 有息负债, 经营活动现金流净额}
\end{axis}
\end{tikzpicture}

\captionof{figure}[民和股份（002234）近年主要财务指标]{民和股份（002234）近年主要财务指标（左上：营业收入与归母净利润；右上：ROE 与资产负债率；下：货币资金、有息负债与经营活动现金流净额。
2021---2025 年为年报口径，2026H1 为 2026 年半年报/半年末，与其前一年报之间留出间隔、以斜纹或空心点区分；半年报数据未年化）}
\end{center}

\pfl{财务分析} 2026 年 6 月末资产负债率 53.9\%，较 2025 年末 上升 2.2 个百分点（样本中位数 52.2\%，所处三级行业内排第 \zhnumber{6} 位）。 2026 年 6 月末货币资金 5.14 亿元、短期债务（短期借款与一年内到期非流动负债）12.38 亿元、长期债务 1.28 亿元，有息负债合计 13.65 亿元，现金短债比 0.42 倍——覆盖偏紧。 净有息负债 8.51 亿元，相当于其 2026 年上半年营业收入的 83\%。 流动比率 0.76、速动比率 0.40，短期偿付压力较大。 2026 年上半年经营活动现金流净额 0.58 亿元，净利润为负而经营现金流为正，差额主要来自折旧摊销与营运资本变动，说明亏损尚未完全转化为现金流出。 2026 年 6 月末在建工程 1.39 亿元，较上年末减少 0.80 亿元，在建投入在收缩。 公司披露的财务与合规风险包括：股票交易实行退市风险警示特别处理，股票简称由“民和股份”变更为“*ST民和”，日涨跌幅限……；2013年全年实现营业收入996,313,767.09元、较上年下降16.36\%，当年继……。\footnote{\texttt{https://companies.caixin.com/2014-04-20/100667679.html}}\footnote{\texttt{http://static.cninfo.com.cn/finalpage/2014-04-19/63870251.PDF}}

\pfl{重大战略转型} 近三年公司的战略动作集中在以下几个方面：2015 子公司沼气提纯压缩项目完工转固并开始计提折旧，同时投入水溶肥业务广告、人力与研发费用（公司将其列为当期亏损原因之……；2025 公司已形成以父母代肉种鸡饲养、商品代肉雏鸡生产销售，肉鸡饲养、屠宰加工，鸡肉熟食和调理品生产销售为核心，有机废弃……。 公告口径上，近三年（2023 年 9 月---2026 年 9 月）公司披露重大重组与并购类公告 0 条、对外投资与转型类公告 0 条、回购与股权激励类公告 0 条，合计公告 261 条。 主营结构的迁移已经落到报表上：2021 年年报的第一大产品为“雏鸡”（52.7\%），2026 年半年报变为“销售商品”（99.7\%）。 综合判断，公司当前的战略重心不是扩张而是化解债务与完成重整，业务层面的调整让位于司法重整程序。\footnote{\texttt{https://stock.stockstar.com/SS2015071500000270.shtml}}\footnote{\texttt{https://vip.stock.finance.sina.com.cn/corp/view/vCB\_AllBulletinDetail.php}}

\begin{center}\small\setlength{\tabcolsep}{4pt}
\captionof{table}[民和股份（002234）历史融资与资本运作]{民和股份（002234）历史融资与资本运作（据公司公告与公开报道整理；金额为披露值，含首次公开发行、再融资、债券、银行借款与重整投资等）}
\begin{tabular}{@{}p{1.5cm}p{3.2cm}p{2.4cm}p{6.3cm}@{}}
\toprule
时间 & 方式 & 规模 & 用途与说明 \\
\midrule
2008-05 & IPO（深交所中小板） & 发行2,700万股（网下配售540万股、网上定价发行2,160万股），发行价10.61元/股，扣除发行费用后实际募集资金净额268,692,500.00元 & 募投项目建设（含粪污处理大型沼气工程及资源化利用项目）…… \\
2011-11 & 非公开发行（A股） & 发行43,523,316股，发行价19.30元/股，募集资金总额839,999,998.80元，募集资金净额812,299,998.80元 & 募投项目建设；10名发行对象认购股份限售期12个月，新增股份2011年11月15日上市 \\
2021-04 & 非公开发行（证监许可〔2020〕3346号） & 发行46,913,580股，发行价12.15元/股，募集资金总额569,999,997.00元，募集资金净额560,469,997.04元 & 募投项目建设；23名发行对象以现金认购，新增股份2021年4月23日上市…… \\
\bottomrule
\end{tabular}
\end{center}

\pfl{历史融投资情况} 从现金流量表看，2025 年公司取得借款收到的现金 17.12 亿元、偿还债务支付的现金 14.93 亿元、吸收权益性投资 — 亿元，筹资活动现金流净额 -0.18 亿元（2023 年为取得借款 12.96 亿元、偿还债务 11.98 亿元）。 2026 年上半年筹资活动现金流净额为 0.34 亿元（取得借款 8.75 亿元、偿还债务 8.02 亿元，未年化），仍处于净融入状态。 2025 年分配股利、利润或偿付利息支付的现金合计 0.44 亿元。 公司披露的股本与债务融资脉络包括：IPO（深交所中小板）、非公开发行（A股）、非公开发行（证监许可〔2020〕3346号）。 股本方面，近三年披露股本变动 3 次；总股本由 2021 年末的 3.49 亿股变为 3.49 亿股（+0.00\%）。 分红方面，近三年未实施现金分红，与重整期间的资金约束一致。 近三年“再融资”类公告 1 条，涉及定向增发。\footnote{\texttt{https://vip.stock.finance.sina.com.cn/corp/view/vISSUE\_MarketBulletinDetail.php}}\footnote{\texttt{https://vip.stock.finance.sina.com.cn/corp/view/vISSUE\_MarketBulletinDetail.php}}\footnote{\texttt{https://vip.stock.finance.sina.com.cn/corp/view/vISSUE\_MarketBulletinDetail.php}}

\pfl{未来潜在融资需求分析} 2026 年 6 月末货币资金 5.14 亿元，而短期债务 12.38 亿元（现金短债比 0.42 倍），2026 年上半年经营现金流净额 0.58 亿元。按“短期债务减去账面货币资金”的滚动偿付口径，静态缺口约 7.24 亿元；若再计入上半年亏损的年化额（1.90 亿元），维持一年正常经营所需的外部资金约 7.24---9.14 亿元。 公司 2026 年上半年归母净利润为负（-0.95 亿元），但 6 月末资产负债率 53.9\% 尚未进入第一档（亏损且负债率 $>$65\%）的区间，当前融资需求以补充流动资金、支撑低谷期经营性支出为主。 资金需求还要叠加在建工程：期末在建工程 1.39 亿元、2026 年上半年购建固定资产等支出 0.43 亿元（未年化），这部分支出在周期底部通常会被主动放缓。


\subsubsection{*ST福成（600965）}
\label{co:600965}

\pfl{公司基本情况} *ST福成成立于 1998 年，河北三河的牛羊肉屠宰加工与餐饮企业，2015年发行股份收购大股东旗下灵山宝塔陵园，成为A股唯一拥有殡葬业务的公司；2026年5月因2025年财务报表被出具无法表示意见，被实施退市风险警示并叠加其他风险警示，简称变更为“*ST福成”。 公司于 2004-07-13 在上交所首发上市，注册资本 8.19 亿元，总股本 8.19 亿股。 截至 2026-09-24收盘价 2.69 元，流通市值 22.0 亿元，近一年-45.0\%，流通市值在三级行业“其他养殖”的 3 家公司中排第 \zhnumber{3} 位，近一年涨跌幅低于全样本中位数（-10.1\%）34.9 个百分点。 发展过程中的关键节点包括：2001年2月28日在廊坊市工商行政管理局注册登记，三河五丰福成食品有限公司变更为河北福……（2001）；2004年6月21日经中国证监会批准首次向社会公众发行人民币普通股8,000万股，200……（2004）；2013年11月26日取得证监许可〔2013〕1478号批复，向三河福生投资有限公司等发……（2013）。\footnote{\texttt{http://www.aastocks.com/tc/cnhk/analysis/company-fundamental/company-profile}}\footnote{\texttt{http://vip.stock.finance.sina.com.cn/corp/go.php/vCI\_CorpRule/stockid/600965.phtml}}\footnote{\texttt{https://qxb-pdf-osscache.qixin.com/AnBaseinfo/21ad91804d0da4a66f49059ee313c3c7.pdf}}

\pfl{历史股价} 本报告样本区间（2021 年 1 月 至 2026 年 9 月）内，股价由 4.82 元变为 2.69 元，累计 -44.2\%，区间最高 8.29 元（2022 年 12 月）、最低 2.55 元（2026 年 6 月）；当前价相当于区间最高价的 32.4\%，为区间最低价的 1.05 倍。 月度收益的年化波动率 39.2\%，高于全样本中位数 37.4\%，在三级行业“其他养殖”的 3 家公司中波动率排第 \zhnumber{2} 位。 从事件看，2021 年 1 月 至 2022 年 12 月的上涨（+72.0\%）与公司公告李福成收到河北省衡水市冀州区人民法院《刑事判决书（2020）冀1181刑……（2022-01-28）在时间上重合——股价的拐点多数能在公司自身的资本运作与风险事件上找到对应。\footnote{\texttt{https://m.21jingji.com/article/20220128/herald/0efe4101e40806d4544d164bd206782e\_zaker.ht}}

\begin{center}
\begin{tikzpicture}
\begin{axis}[
  width=12.6cm, height=3.9cm, scale only axis,
  xmin=2020.922, xmax=2026.888, ymin=2.37, ymax=8.87,
  xtick={2021,2022,2023,2024,2025,2026},
  yearticks,
  yticklabel style={font=\scriptsize},
  ylabel={元/股}, ylabel style={font=\scriptsize},
  grid=major, grid style={LightGray, line width=0.3pt},
  axis line style={InkGray, line width=0.5pt},
  every axis plot/.append style={line width=0.9pt},
]
\addplot[AgriGreen] table[x=x,y=y]{data/clean/profiles/px_600965.dat};
\addplot[SoftRed, dashed, line width=0.7pt] table[x=x,y=y]{data/clean/profiles/pxzz_600965.dat};
\addplot[only marks, mark=*, mark size=1.1pt, SoftRed] table[x=x,y=y]{data/clean/profiles/pxzz_600965.dat};
\end{axis}
\end{tikzpicture}

\captionof{figure}[*ST福成（600965）股价与周期骨架]{*ST福成（600965）月度收盘价（元/股）与 ZigZag 周期骨架（阈值 25\%，圆点为周期转折点）}
\end{center}

\pfl{过去周期分析} 以 25\% 的 ZigZag 阈值划分，样本区间内股价形成 2 个主升段与 2 个主跌段。 最大主跌段为 2026 年 1 月 至 2026 年 6 月，5 个月-60.1\%；最大主升段出现在 2024 年 8 月 至 2026 年 1 月，17 个月+87.9\%。 最近一次转折出现在 2026 年 6 月（下跌段自 2026 年 1 月起，5 个月-60.1\%），当前处于该段的延续之中。 公司股价与申万农林牧渔指数几乎同步见底，个股并未走出独立行情。 牛、羊、水产的产能周期长于生猪，第 \ref{sec:other-livestock} 章判断肉牛供给的实质性收缩要到 2027 年之后才能兑现，这一时滞决定了个股的股价弹性，公司 2026 年上半年实现盈利、半年 ROE 0.8\%（未年化），好于全样本中位数（0.75\%），下行周期对其报表的冲击小于同业。

\pfl{盈利情况} 2026 年上半年营业收入 5.69 亿元（同比 +6.9\%），归母净利润 0.18 亿元，同比 -48.9\%。 以年报为趋势对照，2023---2025 年营业收入由 10.46 亿元变为 11.63 亿元（两年年均复合增速 +5.4\%），归母净利润由 0.99 亿元变为 0.57 亿元。 按半年报口径，2026 年上半年的 ROE 为 0.8\%（未年化）、销售净利率 3.4\%、毛利率 19.7\%，ROE 在“其他养殖”3 家公司中排第 \zhnumber{2} 位。 同期全样本半年 ROE 中位数 0.75\%，公司与之相差 0.08 个百分点（略好于中位数公司）。 结合公司披露，审计机构披露的无法表示意见涉及事项：2022年至2023年支付刘家河养牛场、兴隆庄养牛场、西吴养牛场三处……；湘潭市中级人民法院作出二审最终裁定，驳回上诉、维持原判，该案尚在执行中；截至2025年12月31日陵园公……。\footnote{\texttt{https://money.finance.sina.com.cn/corp/view/vCB\_AllBulletinDetail.php}}

\begin{center}\small\setlength{\tabcolsep}{4pt}
\captionof{table}[*ST福成（600965）关键财务指标]{*ST福成（600965）关键财务指标（合并报表口径；两个半年报期均未年化；现金短债比 $=$ 货币资金 $\div$ 短期债务）}
\begin{tabular}{@{}lrrrrr@{}}
\toprule
指标 & 2023 年 & 2024 年 & 2025 年 & 2025 年半年报 & 2026 年半年报 \\
\midrule
营业收入（亿元） & 10.46 & 10.22 & 11.63 & 5.32 & 5.69 \\
归母净利润（亿元） & 0.99 & 0.52 & 0.57 & 0.36 & 0.18 \\
毛利率（\%） & 37.7 & 25.4 & 22.0 & 25.7 & 19.7 \\
ROE（\%） & 4.7 & 2.4 & 2.6 & 1.6 & 0.8 \\
资产负债率（\%） & 15.4 & 16.1 & 16.5 & 16.3 & 16.7 \\
经营活动现金流净额（亿元） & 0.35 & 1.04 & 2.14 & 0.37 & -0.24 \\
货币资金（亿元） & 1.37 & 0.97 & 2.38 & 0.58 & 2.03 \\
有息负债合计（亿元） & 0.73 & 1.18 & 1.30 & 1.17 & 0.98 \\
现金短债比（倍） & 14.20 & 1.12 & 33.53 & 0.66 & 43.48 \\
\bottomrule
\end{tabular}
\end{center}

\begin{center}
\begin{tikzpicture}
\begin{groupplot}[
  group style={group size=2 by 1, horizontal sep=0.95cm},
  width=7.2cm, height=4.6cm,
  tick label style={font=\tiny},
  title style={font=\scriptsize\heiti},
  yticklabel style={font=\tiny},
  grid=major, grid style={LightGray, line width=0.3pt},
  axis line style={InkGray, line width=0.5pt},
  enlarge x limits={abs=0.8},
  legend style={font=\scriptsize, at={(0.5,-0.22)}, anchor=north,
                legend columns=2, draw=none, fill=none, column sep=6pt},
]
\nextgroupplot[
  ybar, bar width=4pt,
  ymin=-1.27, ymax=14.21,
  xtick={0,1,2,3,4,5.7},
  xticklabels={2021,2022,2023,2024,2025,2026H1},
  ylabel={亿元},
]
\addplot[fill=AgriGreen!75, draw=AgriGreen, bar shift=-2.6pt] table[x=i, y=rev]{data/clean/profiles/fin_600965.dat};
\addplot[fill=SoftRed!60, draw=SoftRed, bar shift=2.6pt] table[x=i, y=ni]{data/clean/profiles/fin_600965.dat};
\addplot[fill=AgriGreen!30, draw=AgriGreen, pattern=north east lines, bar shift=-2.6pt] table[x=x, y=rev]{data/clean/profiles/finh_600965.dat};
\addplot[fill=SoftRed!20, draw=SoftRed, pattern=north east lines, bar shift=2.6pt] table[x=x, y=ni]{data/clean/profiles/finh_600965.dat};
\legend{营业收入（年/半年）, 归母净利润（年/半年）}
\nextgroupplot[
  ymin=-1.6, ymax=21.3,
  xtick={0,1,2,3,4,5.7},
  xticklabels={2021,2022,2023,2024,2025,2026H1},
  ylabel={\%},
]
\addplot[AgriGold, mark=*, mark size=1.3pt] table[x=i, y=roe]{data/clean/profiles/fin_600965.dat};
\addplot[OilBlue, mark=square*, mark size=1.3pt] table[x=i, y=debt]{data/clean/profiles/fin_600965.dat};
\addplot[only marks, AgriGold, mark=*, mark size=2.0pt] table[x=x, y=roe]{data/clean/profiles/finh_600965.dat};
\addplot[only marks, OilBlue, mark=square*, mark size=2.0pt] table[x=x, y=debt]{data/clean/profiles/finh_600965.dat};
\legend{ROE, 资产负债率}
\end{groupplot}
\end{tikzpicture}
\begin{tikzpicture}
\begin{axis}[
  width=15.2cm, height=3.1cm, scale only axis,
  ymin=-1.26, ymax=11.25,
  xtick={0,1,2,3,4,5.7},
  xticklabels={2021,2022,2023,2024,2025,2026H1},
  tick label style={font=\tiny},
  yticklabel style={font=\tiny},
  ylabel={亿元}, ylabel style={font=\scriptsize},
  ybar, bar width=4pt,
  grid=major, grid style={LightGray, line width=0.3pt},
  axis line style={InkGray, line width=0.5pt},
  enlarge x limits={abs=0.8},
  legend style={font=\scriptsize, at={(0.5,-0.30)}, anchor=north,
                legend columns=3, draw=none, fill=none, column sep=10pt},
]
\addplot[fill=AgriGreen!55, draw=AgriGreen, bar shift=-3.0pt] table[x=i, y=cash]{data/clean/profiles/fin_600965.dat};
\addplot[fill=SoftRed!45, draw=SoftRed, bar shift=3.0pt] table[x=i, y=ibd]{data/clean/profiles/fin_600965.dat};
\addplot[fill=AgriGreen!25, draw=AgriGreen, pattern=north east lines, bar shift=-3.0pt] table[x=x, y=cash]{data/clean/profiles/finh_600965.dat};
\addplot[fill=SoftRed!20, draw=SoftRed, pattern=north east lines, bar shift=3.0pt] table[x=x, y=ibd]{data/clean/profiles/finh_600965.dat};
\addplot[OilBlue, mark=*, mark size=1.1pt, line width=0.8pt] table[x=i, y=ocf]{data/clean/profiles/fin_600965.dat};
\addplot[only marks, OilBlue, mark=*, mark size=1.8pt] table[x=x, y=ocf]{data/clean/profiles/finh_600965.dat};
\legend{货币资金, 有息负债, 经营活动现金流净额}
\end{axis}
\end{tikzpicture}

\captionof{figure}[*ST福成（600965）近年主要财务指标]{*ST福成（600965）近年主要财务指标（左上：营业收入与归母净利润；右上：ROE 与资产负债率；下：货币资金、有息负债与经营活动现金流净额。
2021---2025 年为年报口径，2026H1 为 2026 年半年报/半年末，与其前一年报之间留出间隔、以斜纹或空心点区分；半年报数据未年化）}
\end{center}

\pfl{财务分析} 2026 年 6 月末资产负债率 16.7\%，较 2025 年末 上升 0.2 个百分点（样本中位数 52.2\%，所处三级行业内排第 \zhnumber{1} 位）。 2026 年 6 月末货币资金 2.03 亿元、短期债务（短期借款与一年内到期非流动负债）0.05 亿元、长期债务 0.93 亿元，有息负债合计 0.98 亿元，现金短债比 43.48 倍——覆盖充分。 净有息负债 -1.05 亿元，相当于其 2026 年上半年营业收入的 -18\%。 流动比率 4.82、速动比率 1.15，短期指标尚可。 2026 年上半年经营活动现金流净额 -0.24 亿元，净利润为正但经营现金流为负，利润的现金含量偏低。 公司披露的财务与合规风险包括：因触及《上海证券交易所股票上市规则》第9.3.2条第一款第（三）项及第9.8.1条第一款……；公告李福成收到河北省衡水市冀州区人民法院《刑事判决书（2020）冀1181刑初94号》：……。\footnote{\texttt{https://money.finance.sina.com.cn/corp/view/vCB\_AllBulletinDetail.php}}\footnote{\texttt{https://m.21jingji.com/article/20220128/herald/0efe4101e40806d4544d164bd206782e\_zaker.ht}}

\pfl{重大战略转型} 公司在报告期内的战略动作可归为 3 项：2014-2015 以15亿元交易作价向控股股东福成投资集团有限公司发行股份收购三河灵山宝塔陵园100\%股权，跨界进入殡葬服务行业，……；2015-2018 先后向控股股东福成集团收购宝塔陵园100\%股权、通过并购基金持有杭州钱江陵园60\%股权，并从第三方收购天德陵园6……；2026-04-30 针对审计意见涉及的固定资产与收入问题，公司提出加强项目建设监督控制、派专人整理工程项目档案、加强财务与运营环节信……。 公告口径上，近三年（2023 年 9 月---2026 年 9 月）公司披露重大重组与并购类公告 0 条、对外投资与转型类公告 1 条、回购与股权激励类公告 0 条，合计公告 221 条。 综合判断，公司当前的战略重心不是扩张而是化解债务与完成重整，业务层面的调整让位于司法重整程序。\footnote{\texttt{http://epaper.stcn.com/att/202604/30/ZQ30B823-TT\_eBook.pdf}}\footnote{\texttt{https://finance.sina.cn/chanjing/gsxw/2021-06-10/detail-ikqciyzi8885678.d.html}}\footnote{\texttt{https://money.finance.sina.com.cn/corp/view/vCB\_AllBulletinDetail.php}}

\begin{center}\small\setlength{\tabcolsep}{4pt}
\captionof{table}[*ST福成（600965）历史融资与资本运作]{*ST福成（600965）历史融资与资本运作（据公司公告与公开报道整理；金额为披露值，含首次公开发行、再融资、债券、银行借款与重整投资等）}
\begin{tabular}{@{}p{1.5cm}p{3.2cm}p{2.4cm}p{6.3cm}@{}}
\toprule
时间 & 方式 & 规模 & 用途与说明 \\
\midrule
2004-07 & IPO（上海证券交易所） & 首次公开发行人民币普通股8,000万股，发行价3.15元/股 & 2004年6月21日经中国证监会批准，2004年7月13日在上交所上市 \\
2013-12 & 发行股份购买资产（证监许可〔2013〕1478号） & 向三河福生投资有限公司等发行股份，购买福成肥牛餐饮管理有限公司100\%股权、三河市福成都市食品有限公司100\%股权 & 收购大股东旗下餐饮与食品资产；2013年12月4日标的资产过户完成…… \\
2015-07 & 发行股份购买资产（证监许可〔2015〕298号） & 向控股股东福成投资集团有限公司发行290,697,674股，交易作价15亿元；重组预案披露发行价格5.23元/股 & 收购三河灵山宝塔陵园有限公司100\%股权，进入殡葬服务行业…… \\
\bottomrule
\end{tabular}
\end{center}

\pfl{历史融投资情况} 从现金流量表看，2025 年公司取得借款收到的现金 1.00 亿元、偿还债务支付的现金 0.80 亿元、吸收权益性投资 — 亿元，筹资活动现金流净额 -0.25 亿元（2023 年为取得借款 — 亿元、偿还债务 — 亿元）。 2026 年上半年筹资活动现金流净额为 -0.00 亿元（取得借款 0.33 亿元、偿还债务 0.30 亿元，未年化），已转为净流出，处于净还债状态。 2025 年分配股利、利润或偿付利息支付的现金合计 0.34 亿元。 公司披露的股本与债务融资脉络包括：IPO（上海证券交易所）、发行股份购买资产（证监许可〔2013〕1478号）、发行股份购买资产（证监许可〔2015〕298号）。 其中规模较大的两笔为：2015-07发行股份购买资产（证监许可〔2015〕298号）（向控股股东福成投资集团有限公司发行290,697,674股，……）。 股本方面，近三年披露股本变动 3 次；总股本由 2021 年末的 8.19 亿股变为 8.19 亿股（+0.00\%）。 分红方面，近三年现金分红 4 次、累计每股派现 0.26 元（最近一次 2025-04-24）——分红规模与利润的下滑幅度相比仍然有限。\footnote{\texttt{http://vip.stock.finance.sina.com.cn/corp/go.php/vCI\_CorpRule/stockid/600965.phtml}}\footnote{\texttt{https://qxb-pdf-osscache.qixin.com/AnBaseinfo/21ad91804d0da4a66f49059ee313c3c7.pdf}}\footnote{\texttt{https://www.charltonslaw.com/newsletters/china-news-alerts/cn/2014/423/16.pdf}}

\pfl{未来潜在融资需求分析} 2026 年 6 月末货币资金 2.03 亿元，而短期债务 0.05 亿元（现金短债比 43.48 倍），2026 年上半年经营现金流净额 -0.24 亿元。账面货币资金可以覆盖短期债务，缺口不体现在静态偿付层面。 公司 2026 年 6 月末资产负债率 16.7\%、2026 年上半年 ROE 0.8\%（未年化），资产负债表稳健、盈利水平平淡，暂不存在刚性融资需求；潜在的融资需求以技改扩产、并购或被并购为主。


\subsubsection{晓鸣股份（300967）}
\label{co:300967}

\pfl{公司基本情况} 晓鸣股份成立于 1992 年，宁夏银川的蛋鸡制种龙头，集祖代、父母代蛋种鸡饲养与父母代、商品代雏鸡销售于一体，商品代雏鸡国内市占率约25\%；2026年上半年受蛋鸡周期底部拖累，营业收入同比下降约24\%、净利润同比下降99\%。 公司于 2021-04-13 在深交所创业板首发上市，注册资本 1.88 亿元，总股本 1.89 亿股。 截至 2026-09-24收盘价 16.15 元，流通市值 20.2 亿元，近一年-32.0\%，流通市值在三级行业“肉鸡养殖”的 8 家公司中排第 \zhnumber{7} 位，近一年涨跌幅低于全样本中位数（-10.1\%）21.9 个百分点。 按 2025 年年报披露的行业构成，收入占比前三位为养殖行业 100.0\%；其他(补充) 0.0\%。 发展过程中的关键节点包括：宁夏晓鸣农牧股份有限公司于2011年7月5日成立，由魏晓明、杜建峰、石玉鑫、朱万前、吴忠……（2011）；经证监许可〔2021〕613号同意注册，首次公开发行4,700万股人民币普通股，发行价4……（2021）；经证监许可〔2023〕253号同意注册，向不特定对象发行可转换公司债券“晓鸣转债”，20……（2023）。\footnote{\texttt{http://file.finance.sina.com.cn/211.154.219.97:9494/MRGG/CNSESZ\_STOCK/2022/2022-8/2022-0}}\footnote{\texttt{https://vip.stock.finance.sina.com.cn/corp/view/vISSUE\_MarketBulletinDetail.php}}\footnote{\texttt{https://finance.sina.com.cn/stock/s/2026-07-08/doc-inihattw8228676.shtml}}

\pfl{历史股价} 自 2021 年 4 月至 2026 年 9 月，股价由 18.12 元变为 16.15 元，累计 -10.9\%，区间最高 27.90 元（2021 年 10 月）、最低 9.07 元（2024 年 6 月）；当前价相当于区间最高价的 57.9\%，为区间最低价的 1.78 倍。 月度收益的年化波动率 47.0\%，高于全样本中位数 37.4\%，在三级行业“肉鸡养殖”的 8 家公司中波动率排第 \zhnumber{1} 位。

\begin{center}
\begin{tikzpicture}
\begin{axis}[
  width=12.6cm, height=3.9cm, scale only axis,
  xmin=2021.172, xmax=2026.888, ymin=8.44, ymax=29.85,
  xtick={2021,2022,2023,2024,2025,2026},
  yearticks,
  yticklabel style={font=\scriptsize},
  ylabel={元/股}, ylabel style={font=\scriptsize},
  grid=major, grid style={LightGray, line width=0.3pt},
  axis line style={InkGray, line width=0.5pt},
  every axis plot/.append style={line width=0.9pt},
]
\addplot[AgriGreen] table[x=x,y=y]{data/clean/profiles/px_300967.dat};
\addplot[SoftRed, dashed, line width=0.7pt] table[x=x,y=y]{data/clean/profiles/pxzz_300967.dat};
\addplot[only marks, mark=*, mark size=1.1pt, SoftRed] table[x=x,y=y]{data/clean/profiles/pxzz_300967.dat};
\end{axis}
\end{tikzpicture}

\captionof{figure}[晓鸣股份（300967）股价与周期骨架]{晓鸣股份（300967）月度收盘价（元/股）与 ZigZag 周期骨架（阈值 25\%，圆点为周期转折点）}
\end{center}

\pfl{过去周期分析} 以 25\% 的 ZigZag 阈值划分，样本区间内股价形成 3 个主升段与 4 个主跌段。 最大主升段出现在 2024 年 6 月 至 2025 年 8 月，14 个月+169.0\%；最大主跌段为 2023 年 3 月 至 2024 年 6 月，15 个月-52.8\%。 最近一次转折出现在 2026 年 6 月（下跌段自 2025 年 8 月起，10 个月-41.1\%），当前处于该段的延续之中。 公司股价较申万农林牧渔指数提前约 23 个月见底，是板块中较早定价周期反转的公司之一。 禽链的产能周期只有 6---12 个月、反转快于生猪（第 \ref{sec:other-livestock} 节），2026 年前四个月毛鸡养殖利润同比大幅回升，链条利润的修复先于生猪，公司 2026 年 6 月末资产负债率 52.3\%、半年 ROE 0.2\%（未年化），在本轮下行中属于随行业同步调整的一类。

\pfl{盈利情况} 2026 年上半年营业收入 5.68 亿元（同比 -24.5\%），归母净利润 0.02 亿元，同比 -99.0\%。 以年报为趋势对照，2023---2025 年营业收入由 8.31 亿元变为 12.21 亿元（两年年均复合增速 +21.2\%），归母净利润由 -1.55 亿元变为 0.85 亿元。 按半年报口径，2026 年上半年的 ROE 为 0.2\%（未年化）、销售净利率 0.3\%、毛利率 12.9\%，ROE 在“肉鸡养殖”8 家公司中排第 \zhnumber{5} 位。 同期全样本半年 ROE 中位数 0.75\%，公司与之相差 0.51 个百分点（弱于中位数公司）。 结合公司披露，公司上半年经营现金流下降，公司解释主要为本期净利润减少所致；一季度行业周期底部拖累导致上半年业绩大幅下滑；第五届董事会第二十六次会议及2026年第二次临时股东会审议通过向下修正“晓鸣转债”转股价格。\footnote{\texttt{https://finance.sina.com.cn/jjxw/2026-08-17/doc-ininrqsk9081861.shtml}}\footnote{\texttt{https://vip.stock.finance.sina.com.cn/corp/go.php/vCB\_AllBulletin/stockid/300967.phtml}}

\begin{center}\small\setlength{\tabcolsep}{4pt}
\captionof{table}[晓鸣股份（300967）关键财务指标]{晓鸣股份（300967）关键财务指标（合并报表口径；两个半年报期均未年化；现金短债比 $=$ 货币资金 $\div$ 短期债务）}
\begin{tabular}{@{}lrrrrr@{}}
\toprule
指标 & 2023 年 & 2024 年 & 2025 年 & 2025 年半年报 & 2026 年半年报 \\
\midrule
营业收入（亿元） & 8.31 & 9.71 & 12.21 & 7.52 & 5.68 \\
归母净利润（亿元） & -1.55 & 0.45 & 0.85 & 1.85 & 0.02 \\
毛利率（\%） & 0.9 & 16.5 & 22.0 & 34.0 & 12.9 \\
ROE（\%） & -20.0 & 5.9 & 10.8 & 21.3 & 0.2 \\
资产负债率（\%） & 54.4 & 53.6 & 48.9 & 45.4 & 52.3 \\
经营活动现金流净额（亿元） & 0.61 & 3.06 & 2.40 & 2.24 & 1.40 \\
货币资金（亿元） & 0.71 & 1.80 & 1.18 & 1.96 & 1.43 \\
有息负债合计（亿元） & 7.08 & 7.18 & 6.06 & 6.04 & 6.43 \\
现金短债比（倍） & 0.79 & 0.99 & 1.44 & 2.62 & 1.49 \\
\bottomrule
\end{tabular}
\end{center}

\begin{center}
\begin{tikzpicture}
\begin{groupplot}[
  group style={group size=2 by 1, horizontal sep=0.95cm},
  width=7.2cm, height=4.6cm,
  tick label style={font=\tiny},
  title style={font=\scriptsize\heiti},
  yticklabel style={font=\tiny},
  grid=major, grid style={LightGray, line width=0.3pt},
  axis line style={InkGray, line width=0.5pt},
  enlarge x limits={abs=0.8},
  legend style={font=\scriptsize, at={(0.5,-0.22)}, anchor=north,
                legend columns=2, draw=none, fill=none, column sep=6pt},
]
\nextgroupplot[
  ybar, bar width=4pt,
  ymin=-2.92, ymax=13.86,
  xtick={0,1,2,3,4,5.7},
  xticklabels={2021,2022,2023,2024,2025,2026H1},
  ylabel={亿元},
]
\addplot[fill=AgriGreen!75, draw=AgriGreen, bar shift=-2.6pt] table[x=i, y=rev]{data/clean/profiles/fin_300967.dat};
\addplot[fill=SoftRed!60, draw=SoftRed, bar shift=2.6pt] table[x=i, y=ni]{data/clean/profiles/fin_300967.dat};
\addplot[fill=AgriGreen!30, draw=AgriGreen, pattern=north east lines, bar shift=-2.6pt] table[x=x, y=rev]{data/clean/profiles/finh_300967.dat};
\addplot[fill=SoftRed!20, draw=SoftRed, pattern=north east lines, bar shift=2.6pt] table[x=x, y=ni]{data/clean/profiles/finh_300967.dat};
\legend{营业收入（年/半年）, 归母净利润（年/半年）}
\nextgroupplot[
  ymin=-25.9, ymax=61.8,
  xtick={0,1,2,3,4,5.7},
  xticklabels={2021,2022,2023,2024,2025,2026H1},
  ylabel={\%},
]
\addplot[AgriGold, mark=*, mark size=1.3pt] table[x=i, y=roe]{data/clean/profiles/fin_300967.dat};
\addplot[OilBlue, mark=square*, mark size=1.3pt] table[x=i, y=debt]{data/clean/profiles/fin_300967.dat};
\addplot[only marks, AgriGold, mark=*, mark size=2.0pt] table[x=x, y=roe]{data/clean/profiles/finh_300967.dat};
\addplot[only marks, OilBlue, mark=square*, mark size=2.0pt] table[x=x, y=debt]{data/clean/profiles/finh_300967.dat};
\legend{ROE, 资产负债率}
\end{groupplot}
\end{tikzpicture}
\begin{tikzpicture}
\begin{axis}[
  width=15.2cm, height=3.1cm, scale only axis,
  ymin=-0.72, ymax=8.04,
  xtick={0,1,2,3,4,5.7},
  xticklabels={2021,2022,2023,2024,2025,2026H1},
  tick label style={font=\tiny},
  yticklabel style={font=\tiny},
  ylabel={亿元}, ylabel style={font=\scriptsize},
  ybar, bar width=4pt,
  grid=major, grid style={LightGray, line width=0.3pt},
  axis line style={InkGray, line width=0.5pt},
  enlarge x limits={abs=0.8},
  legend style={font=\scriptsize, at={(0.5,-0.30)}, anchor=north,
                legend columns=3, draw=none, fill=none, column sep=10pt},
]
\addplot[fill=AgriGreen!55, draw=AgriGreen, bar shift=-3.0pt] table[x=i, y=cash]{data/clean/profiles/fin_300967.dat};
\addplot[fill=SoftRed!45, draw=SoftRed, bar shift=3.0pt] table[x=i, y=ibd]{data/clean/profiles/fin_300967.dat};
\addplot[fill=AgriGreen!25, draw=AgriGreen, pattern=north east lines, bar shift=-3.0pt] table[x=x, y=cash]{data/clean/profiles/finh_300967.dat};
\addplot[fill=SoftRed!20, draw=SoftRed, pattern=north east lines, bar shift=3.0pt] table[x=x, y=ibd]{data/clean/profiles/finh_300967.dat};
\addplot[OilBlue, mark=*, mark size=1.1pt, line width=0.8pt] table[x=i, y=ocf]{data/clean/profiles/fin_300967.dat};
\addplot[only marks, OilBlue, mark=*, mark size=1.8pt] table[x=x, y=ocf]{data/clean/profiles/finh_300967.dat};
\legend{货币资金, 有息负债, 经营活动现金流净额}
\end{axis}
\end{tikzpicture}

\captionof{figure}[晓鸣股份（300967）近年主要财务指标]{晓鸣股份（300967）近年主要财务指标（左上：营业收入与归母净利润；右上：ROE 与资产负债率；下：货币资金、有息负债与经营活动现金流净额。
2021---2025 年为年报口径，2026H1 为 2026 年半年报/半年末，与其前一年报之间留出间隔、以斜纹或空心点区分；半年报数据未年化）}
\end{center}

\pfl{财务分析} 2026 年 6 月末资产负债率 52.3\%，较 2025 年末 上升 3.4 个百分点（样本中位数 52.2\%，所处三级行业内排第 \zhnumber{5} 位）。 2026 年 6 月末货币资金 1.43 亿元、短期债务（短期借款与一年内到期非流动负债）0.96 亿元、长期债务 5.47 亿元，有息负债合计 6.43 亿元，现金短债比 1.49 倍——覆盖充分。 净有息负债 5.00 亿元，相当于其 2026 年上半年营业收入的 88\%。 流动比率 1.08、速动比率 0.69，短期指标尚可。 2026 年上半年经营活动现金流净额 1.40 亿元，经营现金流与利润同向为正，内生资金可以支撑运营。

\pfl{重大战略转型} 从公开披露看，公司的战略推进包括：2022 将募投项目“阿拉善种鸡养殖场建设项目”的部分募集资金用途变更为“红寺堡智慧农业产业示范园父母代种业基地项目（一期……；2025 公司提出打造中国蛋种鸡产业生物安全第一品牌，通过创新养殖模式、技术研发与管理提升推进蛋鸡制种业务；2026-08 二季度随鸡蛋行情上行实现环比扭亏，公司同步加码种源储备、推进全产业链布局，合同负债、鸡苗售价等先行指标释放行业回……。 公告口径上，近三年（2023 年 9 月---2026 年 9 月）公司披露重大重组与并购类公告 0 条、对外投资与转型类公告 3 条、回购与股权激励类公告 26 条，合计公告 454 条。 第一大行业仍是“养殖行业”，收入占比 100.0\%，与 2021 年年报相比没有方向性变化。 综合判断，公司在报告期内有明确的外延扩张或业务结构调整动作，转型方向与融资安排之间存在直接关联。\footnote{\texttt{https://disc.static.szse.cn/download/disc/disk03/finalpage/2023-04-25/c48322a3-9ac6-4c29}}\footnote{\texttt{http://jrjgj.nx.gov.cn/djgz/spxw/202110/t20211026\_3105010.html}}\footnote{\texttt{https://finance.sina.com.cn/wm/2026-08-18/doc-ininttrk8594671.shtml}}

\begin{center}\small\setlength{\tabcolsep}{4pt}
\captionof{table}[晓鸣股份（300967）历史融资与资本运作]{晓鸣股份（300967）历史融资与资本运作（据公司公告与公开报道整理；金额为披露值，含首次公开发行、再融资、债券、银行借款与重整投资等）}
\begin{tabular}{@{}p{1.5cm}p{3.2cm}p{2.4cm}p{6.3cm}@{}}
\toprule
时间 & 方式 & 规模 & 用途与说明 \\
\midrule
2021-04 & IPO（深交所创业板，证监许可〔2021〕613号） & 发行4,700万股，募集资金总额213,380,000.00元，扣除发行费用66,077,866.74元后实际募集资金净额147,302,133.26元 & 募投项目建设（后部分变更为红寺堡智慧农业产业示范园父母代种业基地项目）…… \\
2023-04 & 向不特定对象发行可转换公司债券（证监许可〔2023〕253号） & 发行3,290,000张，每张面值100元，募集资金总额32,900万元（3.29亿元）；初始转股价格19.43元/股，转股期2023年10月12日至2029年4月5日 & 募投项目建设（截至2025年底募集资金已使用完毕，募投项目均达到预定可使用状态）…… \\
2024-12 & 股份回购（用于员工持股计划或股权激励） & 回购总金额不低于2,000.00万元（含）且不超过4,000.00万元（含），回购价格不超过19.74元/股 & 员工持股计划或股权激励；2024年12月20日第五届董事会第十七次会议…… \\
\bottomrule
\end{tabular}
\end{center}

\pfl{历史融投资情况} 从现金流量表看，2025 年公司取得借款收到的现金 1.08 亿元、偿还债务支付的现金 2.30 亿元、吸收权益性投资 — 亿元，筹资活动现金流净额 -2.03 亿元（2023 年为取得借款 1.07 亿元、偿还债务 0.94 亿元）。 2026 年上半年筹资活动现金流净额为 -0.29 亿元（取得借款 0.59 亿元、偿还债务 0.32 亿元，未年化），已转为净流出，处于净还债状态。 2025 年分配股利、利润或偿付利息支付的现金合计 0.49 亿元。 公司披露的股本与债务融资脉络包括：IPO（深交所创业板，证监许可〔2021〕613号）、向不特定对象发行可转换公司债券（证监许可〔2023〕253号）、股份回购（用于员工持股计划或股权激励）。 其中规模较大的两笔为：2023-04向不特定对象发行可转换公司债券（证监许可〔2023〕253号）（发行3,290,000张，每张面值100元，募集资金总额32……）；2024-12股份回购（用于员工持股计划或股权激励）（回购总金额不低于2,000.00万元（含）且不超过4,000……）。 股本方面，近三年披露股本变动 10 次（可转债转股 4 次、激励股份解禁 1 次、股份回购 1 次、限售股份上市 1 次）；总股本由 2021 年末的 1.88 亿股变为 1.89 亿股（+1.06\%）。 分红方面，近三年现金分红 2 次、累计每股派现 0.45 元（最近一次 2025-12-31）——在利润承压的阶段仍维持了分红。 近三年“再融资”类公告 33 条，涉及定向增发、可转债、募集资金使用。\footnote{\texttt{https://vip.stock.finance.sina.com.cn/corp/view/vCB\_AllBulletinDetail.php}}\footnote{\texttt{https://finance.sina.com.cn/stock/s/2026-07-08/doc-inihattw8228676.shtml}}\footnote{\texttt{https://money.finance.sina.com.cn/corp/view/vCB\_AllBulletinDetail.php}}

\pfl{未来潜在融资需求分析} 2026 年 6 月末货币资金 1.43 亿元，而短期债务 0.96 亿元（现金短债比 1.49 倍），2026 年上半年经营现金流净额 1.40 亿元。账面货币资金可以覆盖短期债务，缺口不体现在静态偿付层面。 公司 2026 年 6 月末资产负债率 52.3\%、2026 年上半年 ROE 0.2\%（未年化），资产负债表尚有余量，融资需求不是“补血型”的刚性需求，而是围绕并购整合、项目投入与营运资金的主动性需求。 公开披露的下一步融资线索包括：2024-12 公司以自有资金和/或自筹资金实施2,000万元至4,000万元股份回购，资金用于员工持股计划或股权激励；2026-03-31 “晓鸣转债”债券余额3.29亿元、到期兑付日2029年4月5日，未转股比例99.91\%，公司面临转债到期兑付压力；2026-07-30 董事会与股东会审议通过向下修正“晓鸣转债”转股价格，以推动转债转股。\footnote{\texttt{https://money.finance.sina.com.cn/corp/view/vCB\_AllBulletinDetail.php}}\footnote{\texttt{https://vip.stock.finance.sina.com.cn/corp/view/vCB\_AllBulletinDetail.php}}\footnote{\texttt{https://vip.stock.finance.sina.com.cn/corp/go.php/vCB\_AllBulletin/stockid/300967.phtml}}


\subsubsection{天域生物（603717）}
\label{co:603717}

\pfl{公司基本情况} 天域生物成立于 2000 年，重庆起家、总部在上海的园林生态工程企业，2019年起跨界生猪养殖、2021年新增光伏新能源并于2024年更名“天域生物”；2020—2024年连续五年亏损，控股股东及主要股东股权高比例质押并被司法冻结，实控人因违规减持被证监会处罚。 公司于 2017-03-27 在上交所首发上市，注册资本 2.90 亿元，总股本 2.90 亿股。 截至 2026-09-24收盘价 5.93 元，流通市值 17.2 亿元，近一年-33.2\%，流通市值在三级行业“生猪养殖”的 11 家公司中排第 \zhnumber{11} 位，近一年涨跌幅低于全样本中位数（-10.1\%）23.1 个百分点。 按 2026 年半年报披露的产品构成，收入占比前三位为生态农牧食品业务 90.6\%；生态环境业务 5.8\%；生态能源业务 3.8\%。 与 2021 年年报相比，第一大产品由“园林生态工程”变为“生态农牧食品业务”，主业重心已经迁移。 发展过程中的关键节点包括：2013年11月30日根据公司股东会决议改制成为股份有限公司（2013年12月增资至注册……（2013）；2014年2月13日股东大会决议，公司名称由“重庆天域园林股份有限公司”更名为“天域园林……（2014）；2017年3月在中国境内公开发行A股股票，2017年3月27日在上海证券交易所主板上市（……（2017）。\footnote{\texttt{https://money.finance.sina.com.cn/corp/view/vCB\_AllBulletinDetail.php}}\footnote{\texttt{http://basic.10jqka.com.cn/603717/company.html}}

\pfl{历史股价} 自 2021 年 1 月至 2026 年 9 月，股价由 8.91 元变为 5.93 元，累计 -33.4\%，区间最高 15.17 元（2022 年 2 月）、最低 5.43 元（2024 年 8 月）；当前价相当于区间最高价的 39.1\%，为区间最低价的 1.09 倍。 月度收益的年化波动率 42.3\%，高于全样本中位数 37.4\%，在三级行业“生猪养殖”的 11 家公司中波动率排第 \zhnumber{4} 位。

\begin{center}
\begin{tikzpicture}
\begin{axis}[
  width=12.6cm, height=3.9cm, scale only axis,
  xmin=2020.922, xmax=2026.888, ymin=5.05, ymax=16.23,
  xtick={2021,2022,2023,2024,2025,2026},
  yearticks,
  yticklabel style={font=\scriptsize},
  ylabel={元/股}, ylabel style={font=\scriptsize},
  grid=major, grid style={LightGray, line width=0.3pt},
  axis line style={InkGray, line width=0.5pt},
  every axis plot/.append style={line width=0.9pt},
]
\addplot[AgriGreen] table[x=x,y=y]{data/clean/profiles/px_603717.dat};
\addplot[SoftRed, dashed, line width=0.7pt] table[x=x,y=y]{data/clean/profiles/pxzz_603717.dat};
\addplot[only marks, mark=*, mark size=1.1pt, SoftRed] table[x=x,y=y]{data/clean/profiles/pxzz_603717.dat};
\end{axis}
\end{tikzpicture}

\captionof{figure}[天域生物（603717）股价与周期骨架]{天域生物（603717）月度收盘价（元/股）与 ZigZag 周期骨架（阈值 25\%，圆点为周期转折点）}
\end{center}

\pfl{过去周期分析} 以 25\% 的 ZigZag 阈值划分，样本区间内股价形成 3 个主升段与 3 个主跌段。 最大主跌段为 2022 年 2 月 至 2022 年 9 月，7 个月-55.7\%；最大主升段出现在 2021 年 1 月 至 2022 年 2 月，13 个月+70.3\%。 最近一次转折出现在 2026 年 7 月（下跌段自 2025 年 10 月起，9 个月-39.2\%），当前处于该段的延续之中。 公司股价较申万农林牧渔指数提前约 21 个月见底，是板块中较早定价周期反转的公司之一。 生猪养殖的定价主线是产能去化与成本曲线的赛跑（第 \ref{sec:hog} 章）：2026 年三季度行业仍处在“产能去化未完成、价格磨底”的阶段，公司 2026 年 6 月末资产负债率 81.0\%，高于全样本中位数 52.2\%，其股价因此更多反映资金面与信用风险，而不是价格弹性本身。

\pfl{盈利情况} 2026 年上半年营业收入 3.51 亿元（同比 -15.0\%），归母净利润 -0.54 亿元，由上年同期的盈利转为亏损。 以年报为趋势对照，2023---2025 年营业收入由 6.70 亿元变为 7.26 亿元（两年年均复合增速 +4.0\%），归母净利润由 -4.62 亿元变为 -1.07 亿元。 按半年报口径，2026 年上半年的 ROE 为 -14.6\%（未年化）、销售净利率 -14.7\%、毛利率 -1.0\%，ROE 在“生猪养殖”11 家公司中排第 \zhnumber{7} 位。 行业同口径半年 ROE 的中位数为 0.75\%，该公司低于中位数 15.35 个百分点。 结合公司披露，披露《关于累计涉及诉讼、仲裁事项的公告》，其中“损害公司债权人利益责任纠纷”一案涉案金额759.08万元……；截至2024年12月31日合并财务报表累计未分配利润为-728,796,592.68元，实收股本290,……。\footnote{\texttt{https://epaper.stcn.com/pic/202603/13/c9b8a982c44b507f828b8a2b4a2b1ac9.pdf}}\footnote{\texttt{http://epaper.zqrb.cn/html/2025-04/16/content\_1136306.htm}}

\begin{center}\small\setlength{\tabcolsep}{4pt}
\captionof{table}[天域生物（603717）关键财务指标]{天域生物（603717）关键财务指标（合并报表口径；两个半年报期均未年化；现金短债比 $=$ 货币资金 $\div$ 短期债务）}
\begin{tabular}{@{}lrrrrr@{}}
\toprule
指标 & 2023 年 & 2024 年 & 2025 年 & 2025 年半年报 & 2026 年半年报 \\
\midrule
营业收入（亿元） & 6.70 & 8.02 & 7.26 & 4.12 & 3.51 \\
归母净利润（亿元） & -4.62 & -1.07 & -1.07 & 0.11 & -0.54 \\
毛利率（\%） & -34.4 & 10.8 & -0.1 & 14.6 & -1.0 \\
ROE（\%） & -52.7 & -18.2 & -23.0 & 2.1 & -14.6 \\
资产负债率（\%） & 74.5 & 78.0 & 79.5 & 75.6 & 81.0 \\
经营活动现金流净额（亿元） & -1.76 & 1.35 & 2.58 & 0.90 & -0.05 \\
货币资金（亿元） & 1.52 & 1.78 & 1.53 & 1.58 & 1.35 \\
有息负债合计（亿元） & 11.80 & 11.44 & 10.41 & 11.09 & 10.49 \\
现金短债比（倍） & 0.56 & 0.64 & 0.54 & 0.56 & 0.49 \\
\bottomrule
\end{tabular}
\end{center}

\begin{center}
\begin{tikzpicture}
\begin{groupplot}[
  group style={group size=2 by 1, horizontal sep=0.95cm},
  width=7.2cm, height=4.6cm,
  tick label style={font=\tiny},
  title style={font=\scriptsize\heiti},
  yticklabel style={font=\tiny},
  grid=major, grid style={LightGray, line width=0.3pt},
  axis line style={InkGray, line width=0.5pt},
  enlarge x limits={abs=0.8},
  legend style={font=\scriptsize, at={(0.5,-0.22)}, anchor=north,
                legend columns=2, draw=none, fill=none, column sep=6pt},
]
\nextgroupplot[
  ybar, bar width=4pt,
  ymin=-6.02, ymax=11.14,
  xtick={0,1,2,3,4,5.7},
  xticklabels={2021,2022,2023,2024,2025,2026H1},
  ylabel={亿元},
]
\addplot[fill=AgriGreen!75, draw=AgriGreen, bar shift=-2.6pt] table[x=i, y=rev]{data/clean/profiles/fin_603717.dat};
\addplot[fill=SoftRed!60, draw=SoftRed, bar shift=2.6pt] table[x=i, y=ni]{data/clean/profiles/fin_603717.dat};
\addplot[fill=AgriGreen!30, draw=AgriGreen, pattern=north east lines, bar shift=-2.6pt] table[x=x, y=rev]{data/clean/profiles/finh_603717.dat};
\addplot[fill=SoftRed!20, draw=SoftRed, pattern=north east lines, bar shift=2.6pt] table[x=x, y=ni]{data/clean/profiles/finh_603717.dat};
\legend{营业收入（年/半年）, 归母净利润（年/半年）}
\nextgroupplot[
  ymin=-63.4, ymax=94.4,
  xtick={0,1,2,3,4,5.7},
  xticklabels={2021,2022,2023,2024,2025,2026H1},
  ylabel={\%},
]
\addplot[AgriGold, mark=*, mark size=1.3pt] table[x=i, y=roe]{data/clean/profiles/fin_603717.dat};
\addplot[OilBlue, mark=square*, mark size=1.3pt] table[x=i, y=debt]{data/clean/profiles/fin_603717.dat};
\addplot[only marks, AgriGold, mark=*, mark size=2.0pt] table[x=x, y=roe]{data/clean/profiles/finh_603717.dat};
\addplot[only marks, OilBlue, mark=square*, mark size=2.0pt] table[x=x, y=debt]{data/clean/profiles/finh_603717.dat};
\legend{ROE, 资产负债率}
\end{groupplot}
\end{tikzpicture}
\begin{tikzpicture}
\begin{axis}[
  width=15.2cm, height=3.1cm, scale only axis,
  ymin=-3.64, ymax=13.48,
  xtick={0,1,2,3,4,5.7},
  xticklabels={2021,2022,2023,2024,2025,2026H1},
  tick label style={font=\tiny},
  yticklabel style={font=\tiny},
  ylabel={亿元}, ylabel style={font=\scriptsize},
  ybar, bar width=4pt,
  grid=major, grid style={LightGray, line width=0.3pt},
  axis line style={InkGray, line width=0.5pt},
  enlarge x limits={abs=0.8},
  legend style={font=\scriptsize, at={(0.5,-0.30)}, anchor=north,
                legend columns=3, draw=none, fill=none, column sep=10pt},
]
\addplot[fill=AgriGreen!55, draw=AgriGreen, bar shift=-3.0pt] table[x=i, y=cash]{data/clean/profiles/fin_603717.dat};
\addplot[fill=SoftRed!45, draw=SoftRed, bar shift=3.0pt] table[x=i, y=ibd]{data/clean/profiles/fin_603717.dat};
\addplot[fill=AgriGreen!25, draw=AgriGreen, pattern=north east lines, bar shift=-3.0pt] table[x=x, y=cash]{data/clean/profiles/finh_603717.dat};
\addplot[fill=SoftRed!20, draw=SoftRed, pattern=north east lines, bar shift=3.0pt] table[x=x, y=ibd]{data/clean/profiles/finh_603717.dat};
\addplot[OilBlue, mark=*, mark size=1.1pt, line width=0.8pt] table[x=i, y=ocf]{data/clean/profiles/fin_603717.dat};
\addplot[only marks, OilBlue, mark=*, mark size=1.8pt] table[x=x, y=ocf]{data/clean/profiles/finh_603717.dat};
\legend{货币资金, 有息负债, 经营活动现金流净额}
\end{axis}
\end{tikzpicture}

\captionof{figure}[天域生物（603717）近年主要财务指标]{天域生物（603717）近年主要财务指标（左上：营业收入与归母净利润；右上：ROE 与资产负债率；下：货币资金、有息负债与经营活动现金流净额。
2021---2025 年为年报口径，2026H1 为 2026 年半年报/半年末，与其前一年报之间留出间隔、以斜纹或空心点区分；半年报数据未年化）}
\end{center}

\pfl{财务分析} 2026 年 6 月末资产负债率 81.0\%，较 2025 年末 上升 1.5 个百分点（样本中位数 52.2\%，所处三级行业内排第 \zhnumber{9} 位）。 2026 年 6 月末货币资金 1.35 亿元、短期债务（短期借款与一年内到期非流动负债）2.75 亿元、长期债务 7.73 亿元，有息负债合计 10.49 亿元，现金短债比 0.49 倍——覆盖偏紧。 净有息负债 9.13 亿元，相当于其 2026 年上半年营业收入的 260\%。 流动比率 0.90、速动比率 0.72，短期偿付压力较大。 2026 年上半年经营活动现金流净额 -0.05 亿元，经营现金流与利润同时为负，日常运营对债务与股东投入的依赖度较高。 公司披露的财务与合规风险包括：控股股东、实际控制人罗卫国所持部分股份被司法冻结3,129,282股、司法标记6,900……；中国证监会重庆监管局作出《行政处罚决定书》〔2026〕6号：罗卫国为实现减持目的通过基金……。\footnote{\texttt{https://www.stcn.com/article/detail/3727459.html}}\footnote{\texttt{https://www.csrc.gov.cn/chongqing/c104814/c7653993/content.shtml}}

\pfl{重大战略转型} 从公开披露看，公司的战略推进包括：2021 公司重金押向分布式光伏发电业务（光伏新能源为2021年末新设板块），2022年、2023年前三季度该板块毛利仅为……；2024 2024年生态农牧食品业务板块在生猪价格同比上涨叠加饲料成本下降的背景下实现扭亏为盈；公司继续缩减生态环境业务规……；2025 2025年计划退休的史东伟解除与罗卫国的一致行动关系，明确放弃公司控制权，其持有的3,233.88万股（占总股本……；2025-2026 终止生态环境类募投项目“天长市龙岗红色古镇文化旅游景区工程总承包项目”，将剩余募集资金永久补充流动资金。 公告口径上，近三年（2023 年 9 月---2026 年 9 月）公司披露重大重组与并购类公告 1 条、对外投资与转型类公告 2 条、回购与股权激励类公告 0 条，合计公告 535 条。 主营结构的迁移已经落到报表上：2021 年年报的第一大产品为“园林生态工程”（92.6\%），2026 年半年报变为“生态农牧食品业务”（90.6\%）。 综合判断，公司在报告期内有明确的外延扩张或业务结构调整动作，转型方向与融资安排之间存在直接关联。\footnote{\texttt{https://www.stcn.com/article/detail/1090436.html}}\footnote{\texttt{http://m.epaper.zqrb.cn/html/2021-01/06/content\_691630.htm}}

\begin{center}\small\setlength{\tabcolsep}{4pt}
\captionof{table}[天域生物（603717）历史融资与资本运作]{天域生物（603717）历史融资与资本运作（据公司公告与公开报道整理；金额为披露值，含首次公开发行、再融资、债券、银行借款与重整投资等）}
\begin{tabular}{@{}p{1.5cm}p{3.2cm}p{2.4cm}p{6.3cm}@{}}
\toprule
时间 & 方式 & 规模 & 用途与说明 \\
\midrule
2017-03 & IPO（上海证券交易所主板） & 发行不超过4,317.79万股，发行价14.63元/股，预计募集资金总额63,169.27万元、预计净额57,302.87万元；实际筹资净额57,018.61万元 & 募投项目建设；发行前（2016年12月31日）净资产60,157.32万元…… \\
2021 & 非公开发行股票 & 筹资净额39,386.55万元 & 募投项目建设；历次筹资情况表披露口径（中德证券关于公司2022年度向特定对象发行股票的核…… \\
2026-01 & 终止募投项目并将剩余募集资金永久补充流动资金 & 项目“天长市龙岗红色古镇文化旅游景区工程总承包项目”募集资金计划投资总额275,705,879.66元，累计投入216,661,029.30元 & 剩余募集资金永久补充流动资金；2025年12月19日第五届董事会第五次会议…… \\
2026-02 & 撤回向特定对象发行股票申请 & 原拟募集资金总额不超过46,523.53万元（曾为不超过52,673.53万元、发行不超过80,417,610股） & 原拟用于补充流动资金和偿还债务；2026年2月11日公司与保荐人国海证券向上交所提交撤回…… \\
\bottomrule
\end{tabular}
\end{center}

\pfl{历史融投资情况} 从现金流量表看，2025 年公司取得借款收到的现金 1.16 亿元、偿还债务支付的现金 1.73 亿元、吸收权益性投资 0.01 亿元，筹资活动现金流净额 -1.45 亿元（2023 年为取得借款 3.28 亿元、偿还债务 2.24 亿元）。 2026 年上半年筹资活动现金流净额为 -0.16 亿元（取得借款 0.40 亿元、偿还债务 0.32 亿元，未年化），已转为净流出，处于净还债状态。 2025 年分配股利、利润或偿付利息支付的现金合计 0.50 亿元。 公司披露的股本与债务融资脉络包括：IPO（上海证券交易所主板）、撤回向特定对象发行股票申请、终止募投项目并将剩余募集资金永久补充流动资金、非公开发行股票。 其中规模较大的两笔为：2017-03IPO（上海证券交易所主板）（发行不超过4,317.79万股，发行价14.63元/股，预计……）；2026-02撤回向特定对象发行股票申请（原拟募集资金总额不超过46,523.53万元（曾为不超过52……）。 股本方面，近三年披露股本变动 3 次；总股本由 2021 年末的 2.90 亿股变为 2.90 亿股（+0.00\%）。 分红方面，近三年未实施现金分红，在全部样本中属于分红能力偏弱的一端。 近三年“再融资”类公告 30 条，涉及定向增发、募集资金使用。\footnote{\texttt{http://money.finance.sina.com.cn/corp/view/vCB\_AllBulletinDetail.php}}\footnote{\texttt{https://static.sse.com.cn/stock/disclosure/announcement/c/202302/603717\_20230228\_3R28.pd}}\footnote{\texttt{https://money.finance.sina.com.cn/corp/view/vCB\_AllBulletinDetail.php}}

\pfl{未来潜在融资需求分析} 2026 年 6 月末货币资金 1.35 亿元，而短期债务 2.75 亿元（现金短债比 0.49 倍），2026 年上半年经营现金流净额 -0.05 亿元。按“短期债务减去账面货币资金”的滚动偿付口径，静态缺口约 1.40 亿元；若再计入上半年亏损的年化额（1.09 亿元），维持一年正常经营所需的外部资金约 1.40---2.49 亿元。 按第 \ref{sec:financing-need} 节的三档划分，公司属于第一档“补血型”：2026 年上半年归母净利润 -0.54 亿元、6 月末资产负债率 81.0\%。负债率已超过 80\%，净资产被亏损侵蚀的程度较深，属于全部样本中资产负债表最紧张的一组。 可行的融资路径按现实性排序为：引入国资或产业投资人（部分公司以控制权让渡为对价）、债务重组或展期（与银行和债券持有人协商）、定向增发（需股价配合，且控股股东需放弃优先认购或同步增资）。 公开披露的下一步融资线索包括：2025-12 公司推出向特定对象发行A股股票方案，拟募集资金总额不超过46,523.53万元，扣除发行费用后拟全部用于补充流动资金和偿……；2026-04-30 2025年年度股东会审议向控股股东罗卫国申请不超过1亿元借款额度、向持股5\%以上股东史东伟申请不超过1亿元借款额度（借款……；2026-05 股东罗卫国和史东伟合计拟提供最高1.5亿元借款，公司无需提供任何抵押、质押、保证等担保。\footnote{\texttt{https://static.sse.com.cn/stock/disclosure/announcement/c/202512/603717\_20251211\_9BAF.pd}}\footnote{\texttt{http://epaper.stcn.com/att/202604/30/ZQ30B631-SS\_eBook.pdf}}\footnote{\texttt{https://finance.sina.com.cn/wm/2026-05-28/doc-inhzmtwq8586724.shtml}}


\subsubsection{湘佳股份（002982）}
\label{co:002982}

\pfl{公司基本情况} 2003 年设立至今，湘佳股份湖南石门的黄羽肉鸡一体化企业，集饲料、种禽、养殖、屠宰、冷链与冰鲜禽肉销售于一体，被称为“中国生鲜家禽第一股”，并延伸至生猪养殖；2025年归母净利润3,294.67万元，2026年上半年由盈转亏。 公司于 2020-04-24 在深交所主板首发上市，注册资本 2.03 亿元，总股本 1.43 亿股。 截至 2026-09-24收盘价 12.42 元，流通市值 16.2 亿元，近一年-16.0\%，流通市值在三级行业“肉鸡养殖”的 8 家公司中排第 \zhnumber{8} 位，近一年涨跌幅低于全样本中位数（-10.1\%）5.9 个百分点。 按 2026 年半年报披露的行业构成，收入占比前三位为养殖业 100.0\%。 与 2021 年年报相比，第一大行业由“养殖行业”变为“养殖业”，主业重心已经迁移。 发展过程中的关键节点包括：以2012年3月31日为基准日整体变更为湖南湘佳牧业股份有限公司，2012年6月18日在……（2012）；2020年4月15日开启申购，2020年4月24日在深圳证券交易所上市（中小板），发行价……（2020）；经证监许可〔2021〕1814号核准，公开发行可转换公司债券“湘佳转债”640万张，募集……（2022）。\footnote{\texttt{https://disc.static.szse.cn/disc/disk03/finalpage/2026-04-21/a3b01e22-7b3a-4536-ab72-a86}}\footnote{\texttt{https://dfjrjgj.hunan.gov.cn/dfjrjgj/tslm\_71665/jrzx/202004/t20200416\_11955062.html}}\footnote{\texttt{https://www.cnfin.com/announ/detail/index.html}}

\pfl{历史股价} 以 2021 年 1 月为起点、2026 年 9 月为终点，股价由 57.70 元变为 12.42 元，累计 -78.5\%，区间最高 62.98 元（2021 年 2 月）、最低 10.55 元（2026 年 6 月）；当前价相当于区间最高价的 19.7\%，为区间最低价的 1.18 倍。 月度收益的年化波动率 38.6\%，高于全样本中位数 37.4\%，在三级行业“肉鸡养殖”的 8 家公司中波动率排第 \zhnumber{3} 位。

\begin{center}
\begin{tikzpicture}
\begin{axis}[
  width=12.6cm, height=3.9cm, scale only axis,
  xmin=2020.922, xmax=2026.888, ymin=9.81, ymax=67.39,
  xtick={2021,2022,2023,2024,2025,2026},
  yearticks,
  yticklabel style={font=\scriptsize},
  ylabel={元/股}, ylabel style={font=\scriptsize},
  grid=major, grid style={LightGray, line width=0.3pt},
  axis line style={InkGray, line width=0.5pt},
  every axis plot/.append style={line width=0.9pt},
]
\addplot[AgriGreen] table[x=x,y=y]{data/clean/profiles/px_002982.dat};
\addplot[SoftRed, dashed, line width=0.7pt] table[x=x,y=y]{data/clean/profiles/pxzz_002982.dat};
\addplot[only marks, mark=*, mark size=1.1pt, SoftRed] table[x=x,y=y]{data/clean/profiles/pxzz_002982.dat};
\end{axis}
\end{tikzpicture}

\captionof{figure}[湘佳股份（002982）股价与周期骨架]{湘佳股份（002982）月度收盘价（元/股）与 ZigZag 周期骨架（阈值 25\%，圆点为周期转折点）}
\end{center}

\pfl{过去周期分析} 以 25\% 的 ZigZag 阈值划分，样本区间内股价形成 3 个主升段与 4 个主跌段。 最大主跌段为 2023 年 2 月 至 2024 年 6 月，16 个月-68.0\%；最大主升段出现在 2022 年 4 月 至 2023 年 2 月，10 个月+40.2\%。 最近一次转折出现在 2026 年 6 月（下跌段自 2025 年 4 月起，14 个月-48.4\%），当前处于该段的延续之中。 公司股价与申万农林牧渔指数几乎同步见底，个股并未走出独立行情。 禽链的产能周期只有 6---12 个月、反转快于生猪（第 \ref{sec:other-livestock} 节），2026 年前四个月毛鸡养殖利润同比大幅回升，链条利润的修复先于生猪，公司 2026 年 6 月末资产负债率 60.9\%、半年 ROE -0.7\%（未年化），在本轮下行中属于随行业同步调整的一类。

\pfl{盈利情况} 2026 年上半年营业收入 23.16 亿元（同比 +8.6\%），归母净利润 -0.12 亿元，由上年同期的盈利转为亏损。 以年报为趋势对照，2023---2025 年营业收入由 38.90 亿元变为 44.71 亿元（两年年均复合增速 +7.2\%），归母净利润由 -1.47 亿元变为 0.33 亿元。 按半年报口径，2026 年上半年的 ROE 为 -0.7\%（未年化）、销售净利率 -0.6\%、毛利率 12.9\%，ROE 在“肉鸡养殖”8 家公司中排第 \zhnumber{6} 位。 行业同口径半年 ROE 的中位数为 0.75\%，该公司低于中位数 1.45 个百分点。 结合公司披露，2026年上半年黄羽肉鸡行业行情低迷、价格承压，生猪毛利率同比骤降近26个百分点，公司由盈转亏；2026年半年度报告列示的主要风险包括：禽类发生疫病的风险、业绩波动风险、进入新行业的风险、原材料价格波……。\footnote{\texttt{https://stock.stockstar.com/RB2026082400002252.shtml}}\footnote{\texttt{https://money.finance.sina.com.cn/corp/view/vCB\_AllBulletinDetail.php}}

\begin{center}\small\setlength{\tabcolsep}{4pt}
\captionof{table}[湘佳股份（002982）关键财务指标]{湘佳股份（002982）关键财务指标（合并报表口径；两个半年报期均未年化；现金短债比 $=$ 货币资金 $\div$ 短期债务）}
\begin{tabular}{@{}lrrrrr@{}}
\toprule
指标 & 2023 年 & 2024 年 & 2025 年 & 2025 年半年报 & 2026 年半年报 \\
\midrule
营业收入（亿元） & 38.90 & 42.04 & 44.71 & 21.34 & 23.16 \\
归母净利润（亿元） & -1.47 & 0.98 & 0.33 & 0.06 & -0.12 \\
毛利率（\%） & 12.6 & 18.0 & 14.3 & 12.6 & 12.9 \\
ROE（\%） & -8.4 & 5.8 & 1.9 & 0.3 & -0.7 \\
资产负债率（\%） & 55.0 & 56.6 & 60.2 & 58.8 & 60.9 \\
经营活动现金流净额（亿元） & 1.62 & 3.80 & 3.77 & 0.56 & 0.84 \\
货币资金（亿元） & 3.42 & 3.85 & 4.39 & 3.28 & 3.94 \\
有息负债合计（亿元） & 13.87 & 15.95 & 18.12 & 18.06 & 19.28 \\
现金短债比（倍） & 1.06 & 1.04 & 0.93 & 0.87 & 0.98 \\
\bottomrule
\end{tabular}
\end{center}

\begin{center}
\begin{tikzpicture}
\begin{groupplot}[
  group style={group size=2 by 1, horizontal sep=0.95cm},
  width=7.2cm, height=4.6cm,
  tick label style={font=\tiny},
  title style={font=\scriptsize\heiti},
  yticklabel style={font=\tiny},
  grid=major, grid style={LightGray, line width=0.3pt},
  axis line style={InkGray, line width=0.5pt},
  enlarge x limits={abs=0.8},
  legend style={font=\scriptsize, at={(0.5,-0.22)}, anchor=north,
                legend columns=2, draw=none, fill=none, column sep=6pt},
]
\nextgroupplot[
  ybar, bar width=4pt,
  ymin=-6.09, ymax=50.25,
  xtick={0,1,2,3,4,5.7},
  xticklabels={2021,2022,2023,2024,2025,2026H1},
  ylabel={亿元},
]
\addplot[fill=AgriGreen!75, draw=AgriGreen, bar shift=-2.6pt] table[x=i, y=rev]{data/clean/profiles/fin_002982.dat};
\addplot[fill=SoftRed!60, draw=SoftRed, bar shift=2.6pt] table[x=i, y=ni]{data/clean/profiles/fin_002982.dat};
\addplot[fill=AgriGreen!30, draw=AgriGreen, pattern=north east lines, bar shift=-2.6pt] table[x=x, y=rev]{data/clean/profiles/finh_002982.dat};
\addplot[fill=SoftRed!20, draw=SoftRed, pattern=north east lines, bar shift=2.6pt] table[x=x, y=ni]{data/clean/profiles/finh_002982.dat};
\legend{营业收入（年/半年）, 归母净利润（年/半年）}
\nextgroupplot[
  ymin=-14.0, ymax=67.9,
  xtick={0,1,2,3,4,5.7},
  xticklabels={2021,2022,2023,2024,2025,2026H1},
  ylabel={\%},
]
\addplot[AgriGold, mark=*, mark size=1.3pt] table[x=i, y=roe]{data/clean/profiles/fin_002982.dat};
\addplot[OilBlue, mark=square*, mark size=1.3pt] table[x=i, y=debt]{data/clean/profiles/fin_002982.dat};
\addplot[only marks, AgriGold, mark=*, mark size=2.0pt] table[x=x, y=roe]{data/clean/profiles/finh_002982.dat};
\addplot[only marks, OilBlue, mark=square*, mark size=2.0pt] table[x=x, y=debt]{data/clean/profiles/finh_002982.dat};
\legend{ROE, 资产负债率}
\end{groupplot}
\end{tikzpicture}
\begin{tikzpicture}
\begin{axis}[
  width=15.2cm, height=3.1cm, scale only axis,
  ymin=-1.93, ymax=21.59,
  xtick={0,1,2,3,4,5.7},
  xticklabels={2021,2022,2023,2024,2025,2026H1},
  tick label style={font=\tiny},
  yticklabel style={font=\tiny},
  ylabel={亿元}, ylabel style={font=\scriptsize},
  ybar, bar width=4pt,
  grid=major, grid style={LightGray, line width=0.3pt},
  axis line style={InkGray, line width=0.5pt},
  enlarge x limits={abs=0.8},
  legend style={font=\scriptsize, at={(0.5,-0.30)}, anchor=north,
                legend columns=3, draw=none, fill=none, column sep=10pt},
]
\addplot[fill=AgriGreen!55, draw=AgriGreen, bar shift=-3.0pt] table[x=i, y=cash]{data/clean/profiles/fin_002982.dat};
\addplot[fill=SoftRed!45, draw=SoftRed, bar shift=3.0pt] table[x=i, y=ibd]{data/clean/profiles/fin_002982.dat};
\addplot[fill=AgriGreen!25, draw=AgriGreen, pattern=north east lines, bar shift=-3.0pt] table[x=x, y=cash]{data/clean/profiles/finh_002982.dat};
\addplot[fill=SoftRed!20, draw=SoftRed, pattern=north east lines, bar shift=3.0pt] table[x=x, y=ibd]{data/clean/profiles/finh_002982.dat};
\addplot[OilBlue, mark=*, mark size=1.1pt, line width=0.8pt] table[x=i, y=ocf]{data/clean/profiles/fin_002982.dat};
\addplot[only marks, OilBlue, mark=*, mark size=1.8pt] table[x=x, y=ocf]{data/clean/profiles/finh_002982.dat};
\legend{货币资金, 有息负债, 经营活动现金流净额}
\end{axis}
\end{tikzpicture}

\captionof{figure}[湘佳股份（002982）近年主要财务指标]{湘佳股份（002982）近年主要财务指标（左上：营业收入与归母净利润；右上：ROE 与资产负债率；下：货币资金、有息负债与经营活动现金流净额。
2021---2025 年为年报口径，2026H1 为 2026 年半年报/半年末，与其前一年报之间留出间隔、以斜纹或空心点区分；半年报数据未年化）}
\end{center}

\pfl{财务分析} 2026 年 6 月末资产负债率 60.9\%，较 2025 年末 上升 0.8 个百分点（样本中位数 52.2\%，所处三级行业内排第 \zhnumber{7} 位）。 2026 年 6 月末货币资金 3.94 亿元、短期债务（短期借款与一年内到期非流动负债）4.02 亿元、长期债务 15.26 亿元，有息负债合计 19.28 亿元，现金短债比 0.98 倍——覆盖偏紧。 净有息负债 15.34 亿元，相当于其 2026 年上半年营业收入的 66\%。 流动比率 1.35、速动比率 0.80，短期指标尚可。 2026 年上半年经营活动现金流净额 0.84 亿元，净利润为负而经营现金流为正，差额主要来自折旧摊销与营运资本变动，说明亏损尚未完全转化为现金流出。 2026 年 6 月末在建工程 2.35 亿元，较上年末减少 0.04 亿元，在建投入在收缩。 公司披露的财务与合规风险包括：股东股权质押情况：质押总比例2.27\%，质押总股数462.04万股，质押总笔数3笔。\footnote{\texttt{https://data.eastmoney.com/notices/detail/002982/AN202607311827527019.html}}

\pfl{重大战略转型} 近三年公司的战略动作集中在以下几个方面：2020 IPO募资中部分用于屠宰产能扩张：子公司岳阳湘佳新建年屠宰3,000万羽优质鸡加工厂项目（一期），建设工期2年，……；2020-2022 公司布局“肉鸡＋种猪”双轮驱动：子公司湖南泰淼鲜丰食品有限公司拟建设1万头种猪养殖项目、20万头商品猪养殖场项目……；2026 将“1350万羽优质鸡标准化养殖基地建设项目”中“夹山镇浮坪村鸡场”的实施地点变更为广西壮族自治区浦北县石埇镇，……。 公告口径上，近三年（2023 年 9 月---2026 年 9 月）公司披露重大重组与并购类公告 0 条、对外投资与转型类公告 5 条、回购与股权激励类公告 46 条，合计公告 482 条。 主营结构的迁移已经落到报表上：2021 年年报的第一大行业为“养殖行业”（100.0\%），2026 年半年报变为“养殖业”（100.0\%）。 综合判断，公司在报告期内有明确的外延扩张或业务结构调整动作，转型方向与融资安排之间存在直接关联。\footnote{\texttt{https://pdf.dfcfw.com/pdf/H3\_AP202006151385129176\_1.pdf}}\footnote{\texttt{http://m.eeo.com.cn/2022/0430/532774.shtml}}\footnote{\texttt{https://www.cnfin.com/announ/detail/index.html}}

\begin{center}\small\setlength{\tabcolsep}{4pt}
\captionof{table}[湘佳股份（002982）历史融资与资本运作]{湘佳股份（002982）历史融资与资本运作（据公司公告与公开报道整理；金额为披露值，含首次公开发行、再融资、债券、银行借款与重整投资等）}
\begin{tabular}{@{}p{1.5cm}p{3.2cm}p{2.4cm}p{6.3cm}@{}}
\toprule
时间 & 方式 & 规模 & 用途与说明 \\
\midrule
2020-04 & IPO（深交所中小板，证监许可〔2020〕41号） & 发行25,630,000股，发行价29.63元/股，募集资金总额759,416,900.00元，扣除发行费用后实际募集资金净额645,635,700.00元 & 1,250万羽优质鸡标准化养殖基地建设项目（计划26,980.00万元）…… \\
2022-04 & 公开发行可转换公司债券“湘佳转债”（证监许可〔2021〕1814号） & 发行640.00万张，每张面值100.00元，按面值发行，发行总额64,000.00万元；扣除承销及保荐费720.00万元后实际到账63,280.00万元，扣除其他发行费用233.65万元后募集资金净额63,046.35万元 & 1,350万羽优质鸡标准化养殖基地建设项目、1万头种猪养殖基地项目和补充流动资金…… \\
2024年度 & 利润分配（现金分红＋资本公积转增股本） & 每10股派发现金红利1元（含税）、以资本公积金每10股转增4股 & 股东回报；2025年度分配预案为每10股派1元（含税） \\
\bottomrule
\end{tabular}
\end{center}

\pfl{历史融投资情况} 从现金流量表看，2025 年公司取得借款收到的现金 5.91 亿元、偿还债务支付的现金 4.16 亿元、吸收权益性投资 — 亿元，筹资活动现金流净额 0.70 亿元（2023 年为取得借款 3.84 亿元、偿还债务 2.44 亿元）。 2026 年上半年筹资活动现金流净额为 0.63 亿元（取得借款 4.07 亿元、偿还债务 3.27 亿元，未年化），仍处于净融入状态。 2025 年分配股利、利润或偿付利息支付的现金合计 0.84 亿元。 公司披露的股本与债务融资脉络包括：IPO（深交所中小板，证监许可〔2020〕41号）、公开发行可转换公司债券“湘佳转债”（证监许可〔2021〕1814号）、利润分配（现金分红＋资本公积转增股本）。 其中规模较大的两笔为：2022-04公开发行可转换公司债券“湘佳转债”（证监许可〔2021〕1814号）（发行640.00万张，每张面值100.00元，按面值发行，发……）。 股本方面，近三年披露股本变动 14 次（可转债转股 8 次、其他 3 次、限售股份上市 1 次、转增 1 次）；总股本由 2021 年末的 1.02 亿股变为 1.43 亿股（+40.00\%）。 分红方面，近三年现金分红 5 次、累计每股派现 0.60 元（最近一次 2025-12-31）——在利润承压的阶段仍维持了分红。 近三年“再融资”类公告 15 条，涉及可转债。\footnote{\texttt{http://static.cninfo.com.cn/finalpage/2024-04-10/1219550513.PDF}}\footnote{\texttt{https://www.cnfin.com/announ/detail/index.html}}\footnote{\texttt{https://vip.stock.finance.sina.com.cn/corp/view/vCB\_AllBulletinDetail.php}}

\pfl{未来潜在融资需求分析} 2026 年 6 月末货币资金 3.94 亿元，而短期债务 4.02 亿元（现金短债比 0.98 倍），2026 年上半年经营现金流净额 0.84 亿元。账面货币资金可以覆盖短期债务，缺口不体现在静态偿付层面。 公司 2026 年上半年归母净利润为负（-0.12 亿元），但 6 月末资产负债率 60.9\% 尚未进入第一档（亏损且负债率 $>$65\%）的区间，当前融资需求以补充流动资金、支撑低谷期经营性支出为主。 公开披露的下一步融资线索包括：2025-03 “湘佳转债”发行规模6.40亿元、债券余额6.40亿元，到期兑付日2028年4月19日；初始转股价42.56元/股，20……；2026-07-24 2026年第一次临时股东会决议公告（公司治理/融资相关议案审议）；2026-08-01 公司公告向银行申请综合授信额度，以满足公司日常生产经营的资金需求。 资金需求还要叠加在建工程：期末在建工程 2.35 亿元、2026 年上半年购建固定资产等支出 1.36 亿元（未年化），这部分支出在周期底部通常会被主动放缓。\footnote{\texttt{http://notice.10jqka.com.cn/api/pdf/73f2163d6a7bb45a.pdf}}\footnote{\texttt{https://data.eastmoney.com/notices/detail/002982/AN202607231827294437.html}}\footnote{\texttt{https://data.eastmoney.com/notices/detail/002982/AN202607311827527019.html}}


